\documentclass[11pt,a4paper,oneside]{book}
% ---------------------------------------------------------------------------
% preamble.tex -- packages, theorem environments, boxes and macros shared by
% all lectures.  Only packages that ship with the local TinyTeX are used.
% ---------------------------------------------------------------------------

% ---- fonts and encoding ---------------------------------------------------
\usepackage[T1]{fontenc}
\usepackage[utf8]{inputenc}
\usepackage{lmodern}                   % Latin Modern text, sans, typewriter and math
\linespread{1.04}
\setlength{\emergencystretch}{2em}
% Latin Modern has no italic small caps: fall back to upright small caps silently
\input{t1lmr.fd}   % declare the family now so that the shape below can be added
\DeclareFontShape{T1}{lmr}{m}{scit}{<->ssub * lmr/m/sc}{}

% ---- mathematics ----------------------------------------------------------
\usepackage{amsmath,amssymb,amsthm,mathtools,bm}

% ---- layout ---------------------------------------------------------------
\usepackage[a4paper,left=3.1cm,right=3.1cm,top=3cm,bottom=3.2cm]{geometry}
\usepackage{booktabs,array,enumitem,float}
\usepackage{xcolor}
\usepackage{titlesec}
\setlist{itemsep=2pt,topsep=4pt,parsep=0pt}

\definecolor{accent}{HTML}{1B4F8A}     % deep blue
\definecolor{warn}{HTML}{B5470B}       % burnt orange
\definecolor{claim}{HTML}{8E1B3A}      % wine red, for unverified claims
\definecolor{softgray}{HTML}{666666}

% ---- headings -------------------------------------------------------------
\renewcommand{\chaptername}{Lecture}
\titleformat{\chapter}[display]
  {\sffamily\bfseries}
  {\color{accent}\large\MakeUppercase{\chaptertitlename}\ \thechapter}
  {4pt}
  {\raggedright\Huge}
  [\vspace{5pt}{\color{accent}\titlerule[1.2pt]}]
\titlespacing*{\chapter}{0pt}{-14pt}{22pt}
\titleformat{\section}
  {\sffamily\Large\bfseries\color{accent}}{\thesection}{0.7em}{}
\titleformat{\subsection}
  {\sffamily\large\bfseries\color{accent}}{\thesubsection}{0.7em}{}
\titleformat{\subsubsection}[runin]
  {\sffamily\bfseries}{}{0pt}{}[.]
\titlespacing*{\section}{0pt}{3.2ex plus 1ex minus .3ex}{1.4ex plus .2ex}
\titlespacing*{\subsection}{0pt}{2.6ex plus .8ex minus .3ex}{1ex plus .2ex}
\titlespacing*{\subsubsection}{0pt}{1.6ex plus .5ex}{0.8em}
\setcounter{secnumdepth}{2}
\setcounter{tocdepth}{1}
\pagestyle{plain}

% ---- theorem environments (one shared counter, numbered within a lecture) -
\theoremstyle{plain}
\newtheorem{theorem}{Theorem}[chapter]
\newtheorem{lemma}[theorem]{Lemma}
\newtheorem{proposition}[theorem]{Proposition}
\newtheorem{corollary}[theorem]{Corollary}
\newtheorem{conjecture}[theorem]{Conjecture}
\theoremstyle{definition}
\newtheorem{definition}[theorem]{Definition}
\newtheorem{example}[theorem]{Example}
\theoremstyle{remark}
\newtheorem*{remark}{Remark}
\newtheorem*{observation}{Observation}

\theoremstyle{definition}
\newtheorem{exercise}{Exercise}[chapter]

% ---- boxes ----------------------------------------------------------------
\usepackage{tikz}
\usetikzlibrary{arrows.meta,positioning,calc,shapes.geometric,fit,backgrounds,decorations.pathreplacing}
\usepackage[skins,breakable]{tcolorbox}

\tcbset{
  notebox/.style={
    enhanced, breakable, frame hidden, boxrule=0pt, arc=1.5pt,
    left=10pt, right=8pt, top=6pt, bottom=6pt,
    fonttitle=\sffamily\bfseries\small, coltitle=black, attach boxed title to top left={yshift=-2mm,xshift=0mm},
    before skip=10pt, after skip=10pt
  }
}
\newtcolorbox{bigpicture}[2][]{notebox, #1, colback=accent!6, colframe=accent,
  borderline west={2.6pt}{0pt}{accent}, title={\color{accent}#2}, before upper={\small\noindent}}
\newtcolorbox{pitfall}[2][]{notebox, #1, colback=warn!7, colframe=warn,
  borderline west={2.6pt}{0pt}{warn}, title={\color{warn}Pitfall: #2}}
\newtcolorbox{statusbox}[1]{notebox, colback=claim!6, colframe=claim,
  borderline west={2.6pt}{0pt}{claim}, title={\color{claim}Status note: #1}}
\newtcolorbox{goalbox}[1]{notebox, colback=black!4, colframe=black!50,
  borderline west={2.6pt}{0pt}{black!50}, title={#1}}

% ---- bibliography ---------------------------------------------------------
\usepackage[round,authoryear]{natbib}
\setlength{\bibsep}{3pt}

% ---- hyperlinks (load last) -----------------------------------------------
\usepackage[colorlinks=true,linkcolor=accent,citecolor=accent!70!black,urlcolor=accent,
            pdfusetitle,bookmarksnumbered]{hyperref}

% ---------------------------------------------------------------------------
% macros
% ---------------------------------------------------------------------------
\newcommand{\NP}{\mathsf{NP}}
\newcommand{\Ptime}{\mathsf{P}}
\newcommand{\coNP}{\mathsf{coNP}}
\newcommand{\PCP}{\mathsf{PCP}}
\newcommand{\BPP}{\mathsf{BPP}}
\newcommand{\poly}{\operatorname{poly}}
\newcommand{\polylog}{\operatorname{polylog}}
\newcommand{\eps}{\varepsilon}
\newcommand{\N}{\mathbb{N}}
\newcommand{\R}{\mathbb{R}}
\newcommand{\Z}{\mathbb{Z}}
\newcommand{\F}{\mathbb{F}}
\newcommand{\bits}{\{0,1\}}
\DeclareMathOperator*{\E}{\mathbb{E}}
\DeclareMathOperator{\OPT}{OPT}
\DeclareMathOperator{\val}{val}
\DeclareMathOperator{\Acc}{Acc}
\newcommand{\Sig}{\Sigma}
\newcommand{\UG}{\textsc{UG}}
\newcommand{\LC}{\textsc{LabelCover}}
\newcommand{\sat}[1]{\textsc{#1}}
\newcommand{\ie}{i.e.\ }
\newcommand{\eg}{e.g.\ }
\newcommand{\defn}[1]{\emph{#1}}

% Gap_{c,s}-Problem
\newcommand{\Gap}[3]{\textsc{Gap}_{#1,#2}\text{-}\textsc{#3}}

% a "star" difficulty marker for exercises
\newcommand{\stars}[1]{\textsf{\small\color{softgray}[#1]}}


\title{From the PCP Theorem to Unique Games\\[4pt]
  \Large Proofs, Algorithms and Geometry of Hardness of Approximation}
\author{Lecture notes}
\date{Draft of \today}

\begin{document}
\frontmatter

% ---------------------------------------------------------------- title page
\begin{titlepage}
  \centering
  \vspace*{3.2cm}
  {\sffamily\color{accent}\large\bfseries LECTURE NOTES\par}
  \vspace{1.2cm}
  {\sffamily\bfseries\fontsize{30}{36}\selectfont From the PCP Theorem\\to Unique Games\par}
  \vspace{0.8cm}
  {\color{accent}\rule{0.55\textwidth}{1.2pt}\par}
  \vspace{0.8cm}
  {\sffamily\Large Proofs, algorithms and geometry\\of hardness of approximation\par}
  \vfill
  {\sffamily\small\color{softgray}
   Sixteen lectures from \emph{proof checking} to the \emph{Unique Games Conjecture}\\[2pt]
   Draft of \today\par}
  \vspace{1.5cm}
\end{titlepage}

% ---------------------------------------------------------------- about
\chapter*{About these notes}
\addcontentsline{toc}{chapter}{About these notes}

These notes retrace one of the most fruitful lines of research in theoretical computer
science: how the question \emph{``is there a proof of this statement?''} turned into the
question \emph{``how well can this optimization problem be approximated?''}, and how the
second question, in turn, led to a single conjecture---the \emph{Unique Games
Conjecture} (UGC)---that seems to explain the answer for a startlingly large family of
problems.  The plan is one lecture per paper-sized idea: the language of games and the PCP theorem
(Lectures 1--4), amplification by parallel repetition and by Dinur's gap amplification (5--6),
Håstad's optimal inapproximability (7), the UGC and its consequences in analysis, geometry and
semidefinite programming (8--12), the $2$-to-$2$ Games theorem (13--15), and a critical reading
of the 2026 claim of a full proof of the UGC (16).

\paragraph{Prerequisites.}
Undergraduate algorithms and the basics of $\NP$-completeness.  Everything else
(approximation ratios, gap problems, games, \dots) is defined when first needed.

\paragraph{How to read the boxes.}
\begin{itemize}
  \item \textbf{Blue boxes} summarize the big picture or a key idea; if you are
        short on time, read the blue boxes and the statements of the theorems.
  \item \textbf{Orange boxes} flag common misconceptions.
  \item \textbf{Wine-red boxes} state the \emph{epistemic status} of a result.  Whenever
        a claim has not been independently verified, we say so there and nowhere do we
        treat it as established.
\end{itemize}

\paragraph{How results are labelled.}
A result with an attribution in its heading, \eg ``Theorem (PCP theorem;
Arora--Safra, Arora--Lund--Motwani--Sudan--Szegedy)'', is quoted from the literature; we
give a proof only if one follows the statement, and otherwise point to where the proof
can be found.  Statements that
depend on an unproven conjecture carry that conjecture in their name, \eg
``Theorem (assuming the UGC)''.

\paragraph{Exercises.}
Each lecture ends with exercises marked \stars{$\star$} (routine),
\stars{$\star\star$} (needs an idea) or \stars{$\star\star\star$} (substantial).
Several exercises come with a hint.

\tableofcontents

% ---------------------------------------------------------------- lectures
\mainmatter
\chapter{Why PCP and Unique Games Matter}
\label{lec:why}

{\sffamily\large\color{softgray}From checking proofs to approximating optima\par}
\vspace{0.5em}

\noindent
Color every vertex of a graph red or blue, and call an edge \emph{cut} if its endpoints
get different colors.  Flipping a fair coin for each vertex cuts every edge with
probability~$\tfrac12$, so some coloring cuts at least half of all edges.  In 1994
Goemans and Williamson \citep{GW95} showed how to do much better: using semidefinite
programming and a random hyperplane, one can find in polynomial time a cut that contains
at least $0.8785\ldots$ times as many edges as a \emph{maximum} cut.  More than thirty
years later that constant has not moved.  Is it a limit of our imagination, or a limit of
mathematics?

Here is a remarkable partial answer.  The best \emph{unconditional} hardness result says
that, unless $\Ptime=\NP$, no polynomial-time algorithm beats $16/17\approx 0.941$.  But
if the \emph{Unique Games Conjecture} (UGC) is true, then no polynomial-time algorithm can
beat $0.8785\ldots+\eps$ for any $\eps>0$ \citep{KKMO07,MOO10}: the strange constant defined
by a calculus minimization is \emph{exactly} the truth.  The conjecture was formulated by
Khot in 2002 \citep{Khot02}.  As we write this (October 2026) it is in an unusual
position: a manuscript dated 23~September~2026 claims a proof (Section~\ref{sec:map}).  We
treat the conjecture as a conjecture throughout, and say so whenever a result depends on it.

Every step of this story---what ``hard to approximate'' means, why we can prove it at all,
and why the conjecture takes such a peculiar form---traces back to a single idea:

\begin{quote}\itshape
To prove that an optimization problem is hard to approximate, we need proofs that can be
checked by reading only a few bits.
\end{quote}

\noindent
This lecture builds the vocabulary (gap problems, probabilistically checkable proofs,
two-prover games, Label Cover, Unique Games), proves the handful of short results that
explain \emph{why} the vocabulary looks the way it does, and draws the map on which the
remaining fifteen lectures sit.

\begin{goalbox}{After this lecture you should be able to\ldots}
\begin{enumerate}[leftmargin=*,label=(\arabic*)]
  \item define a gap problem, and explain how $\NP$-hardness of a gap problem rules out
        approximation algorithms (Lemma~\ref{lem:gap-approx});
  \item explain why the classical reductions of the 1970s cannot produce a gap, and what the
        PCP theorem adds (Section~\ref{sec:nogap});
  \item translate between PCP verifiers, CSPs, two-prover games and Label Cover, and prove
        the two forms of the PCP theorem equivalent (Proposition~\ref{prop:pcp-equiv});
  \item state the Unique Games Conjecture with its quantifiers in the right order, and say
        why a Unique Game with perfect completeness is easy;
  \item locate each later lecture on a map of the field (Section~\ref{sec:map}).
\end{enumerate}
\end{goalbox}


%% ===========================================================================
\section{A scoreboard of approximability}
\label{sec:scoreboard}
%% ===========================================================================

Let us begin with the answer key.  Table~\ref{tab:scoreboard} lists four problems, the best
approximation ratio known for each, and what is known about the \emph{limit} of
approximability.

\begin{table}[t]
\centering\small
\renewcommand{\arraystretch}{1.35}
\begin{tabular}{@{}>{\raggedright\arraybackslash}p{2.75cm}
                   >{\raggedright\arraybackslash}p{2.95cm}
                   >{\raggedright\arraybackslash}p{3.85cm}
                   >{\raggedright\arraybackslash}p{3.45cm}@{}}
\toprule
Problem & Best algorithm & $\NP$-hard beyond & If the UGC holds, hard beyond\\
\midrule
\sat{Max-E3-Lin}$(\F_2)$ & $\tfrac12$ (random) & $\tfrac12+\eps$ \citep{Hastad01} & (already tight)\\
\sat{Max-E3-SAT} & $\tfrac78$ (random) & $\tfrac78+\eps$ \citep{Hastad01} & (already tight)\\
\midrule
\sat{Max-Cut} & $\alpha_{\mathrm{GW}}\approx 0.8786$ \citep{GW95} &
   $\tfrac{16}{17}+\eps\approx 0.941$ \citep{Hastad01,TSSW00} &
   $\alpha_{\mathrm{GW}}+\eps$ \citep{KKMO07,MOO10}\\[3pt]
\sat{Vertex Cover} & $2$ (maximal matching) &
   $\sqrt2-\eps\approx 1.414$ \citep{DKKMS18,KMS18} &
   $2-\eps$ \citep{KhotRegev08}\\
\bottomrule
\end{tabular}
\caption{Approximability thresholds for four basic problems.  Ratios are
$\le 1$ for maximization and $\ge 1$ for minimization problems (Definition~\ref{def:approx}).
``$\NP$-hard beyond $\alpha$'' means that no polynomial-time algorithm achieves ratio better
than $\alpha$ unless $\Ptime=\NP$.  The first two rows are unconditionally tight; for the last
two, the UGC would close the gap between algorithm and hardness.}
\label{tab:scoreboard}
\end{table}

The first two rows are the great success of the PCP method, and they are astonishing.
\sat{Max-E3-Lin}$(\F_2)$ asks for an assignment satisfying as many as possible of a given
list of linear equations $x_i+x_j+x_k=b \pmod 2$.  A uniformly random assignment satisfies
half of the equations, whatever the system is.  H\aa stad proved that nothing polynomial-time
does better, even if you are promised that a $(1-\eps)$-fraction of the equations can be
satisfied simultaneously: \emph{guessing is optimal}.  The same is true for \sat{Max-E3-SAT},
where a random assignment satisfies a $\tfrac78$-fraction of the clauses.

The last two rows are where the story stops being a triumph.  For \sat{Max-Cut} and
\sat{Vertex Cover} the best algorithms (a clever semidefinite program; a trivial maximal matching)
leave a visible gap to the best hardness proofs.  Assuming the UGC, both gaps close
\emph{simultaneously} and \emph{exactly}.  In Lecture~11 we shall see that this is not a
coincidence: Raghavendra \citep{Raghavendra08} proved that, assuming the UGC, for \emph{every}
constraint satisfaction problem the optimal approximation ratio is the integrality gap of
one generic semidefinite program.  In the other direction, the constant in the \sat{Max-Cut} row is
\begin{equation}\label{eq:alphaGW}
  \alpha_{\mathrm{GW}} \;=\; \min_{0<\theta\le\pi}\;\frac{2}{\pi}\cdot\frac{\theta}{1-\cos\theta}\;\approx\;0.87856,
\end{equation}
and the hardness proof of Khot, Kindler, Mossel and O'Donnell reproduces precisely this
minimization, with $\theta$ the angle between two random-hyperplane vectors.
We will see in Lecture~9 how the same expression emerges from a Fourier-analytic argument
about Gaussian noise.

\begin{bigpicture}[unbreakable]{The question behind the whole course}
For which problems can an efficient algorithm do better than the trivial one, by how much, and
how do we \emph{prove} that nothing does better?  The PCP theorem answers this for problems
like \sat{Max-E3-SAT}; the Unique Games Conjecture (still unproven) would answer it for a much
larger class.
\end{bigpicture}


%% ===========================================================================
\section{Approximation algorithms in one page}
\label{sec:approx}
%% ===========================================================================

All problems in this lecture are constraint satisfaction problems or close relatives.  An
\defn{instance} of a \sat{Max-CSP} consists of variables $x_1,\dots,x_n$ with values in a
finite domain and a list of $m$ \defn{constraints}, each depending on a few variables.  The
\defn{value} of an assignment is the \emph{fraction} of constraints it satisfies, and
$\OPT(I)$ is the maximum value over all assignments.  Examples:
\sat{Max-3SAT} (clauses with at most three literals),
\sat{Max-E3-Lin}$(\F_2)$ (equations $x_i+x_j+x_k=b$),
\sat{Max-Cut} (constraints $x_u\ne x_v$, one per edge).

\begin{definition}[Approximation ratio]\label{def:approx}
A polynomial-time algorithm $\mathcal A$ is an \defn{$\alpha$-approximation} for a
maximization problem ($0<\alpha\le1$) if on every instance $I$ it outputs a solution of value
at least $\alpha\cdot\OPT(I)$.  For a minimization problem $\mathcal A$ is a
\defn{$\rho$-approximation} ($\rho\ge1$) if it outputs a solution of value at most
$\rho\cdot\OPT(I)$.  The \defn{approximation threshold} of a problem is the best constant
$\alpha$ (or $\rho$) for which an approximation algorithm exists.
\end{definition}

\begin{example}[Four one-paragraph algorithms]\label{ex:random}
\leavevmode
\begin{itemize}
  \item \emph{\sat{Max-Cut}, ratio $\tfrac12$.}  Put each vertex on a random side.  Every edge is
        cut with probability $\tfrac12$, so the expected number of cut edges is $m/2\ge\OPT/2$.
  \item \emph{\sat{Max-E3-Lin}$(\F_2)$, ratio $\tfrac12$.}  For a random assignment and any
        fixed equation $x_i+x_j+x_k=b$, condition on $x_i,x_j$: exactly one value of $x_k$
        satisfies the equation, so it holds with probability~$\tfrac12$.
  \item \emph{\sat{Max-E3-SAT}, ratio $\tfrac78$} (every clause has three literals on three
        distinct variables).  A clause is violated by exactly one of the eight settings of its
        variables, hence is satisfied with probability $\tfrac78$.
  \item \emph{\sat{Vertex Cover}, ratio $2$.}  Let $M$ be a maximal matching.  Its $2|M|$
        endpoints form a vertex cover (an uncovered edge could be added to $M$), while any
        cover must contain one endpoint of each of the $|M|$ disjoint edges of $M$, so
        $\OPT\ge|M|$.
\end{itemize}
The three randomized algorithms can be derandomized by the method of conditional
expectations (Exercise~\ref{ex:derand}).
\end{example}

So an algorithm that does \emph{nothing intelligent} already achieves a constant ratio.  The
whole subject starts from the question whether intelligence can ever pay off---and for
\sat{Max-E3-SAT} and \sat{Max-E3-Lin} it turns out that it cannot.  The tool that proves this is the
subject of the next three sections.


%% ===========================================================================
\section{Gap problems: the language of hardness}
\label{sec:gap}
%% ===========================================================================

How does one prove that an \emph{approximation} problem is $\NP$-hard?  Decision problems
are what the theory of $\NP$-completeness knows how to handle, so we turn the approximation
problem into a decision problem---a \emph{promise} problem.

\begin{definition}[Gap problem]\label{def:gap}
Let $\Pi$ be a \sat{Max-CSP} and $0\le s<c\le1$.  The \defn{gap problem}
$\Gap{c}{s}{$\Pi$}$ is the following promise problem.  Given an instance $I$ of $\Pi$ such that
either
\[
  \text{(YES)}\ \ \OPT(I)\ge c \qquad\text{or}\qquad \text{(NO)}\ \ \OPT(I)\le s,
\]
decide which of the two cases holds.  On instances outside both cases, any answer is acceptable.
We call $c$ the \defn{completeness} and $s$ the \defn{soundness}; the gap problem has
\defn{perfect completeness} if $c=1$.  The problem is \defn{$\NP$-hard} if there is a
polynomial-time map $f$ from $3$-CNF formulas to instances of $\Pi$ with
$\varphi$ satisfiable $\Rightarrow f(\varphi)$ is a YES instance, and $\varphi$ unsatisfiable
$\Rightarrow f(\varphi)$ is a NO instance.
\end{definition}

\begin{lemma}[Gap hardness implies inapproximability]\label{lem:gap-approx}
If $\Gap{c}{s}{$\Pi$}$ is $\NP$-hard, then for every $\alpha>s/c$ there is no polynomial-time
$\alpha$-approximation algorithm for $\Pi$ unless $\Ptime=\NP$.
\end{lemma}
\begin{proof}
Suppose $\mathcal A$ is a polynomial-time $\alpha$-approximation with $\alpha>s/c$.  We use it
to decide $3$SAT.  Given $\varphi$, compute $I=f(\varphi)$, run $\mathcal A$ on $I$, evaluate the
value $v$ of the solution it returns (this is easy), and accept iff $v>s$.
If $\varphi$ is satisfiable then $\OPT(I)\ge c$, so $v\ge\alpha c>s$.
If $\varphi$ is unsatisfiable then $v\le\OPT(I)\le s$.
Hence this is a polynomial-time algorithm for $3$SAT.
\end{proof}

\noindent Conversely, an approximation algorithm with ratio below $s/c$ would solve the gap
problem---so the gap problem \emph{is} the decision core of approximation.  In this language, the
rows of Table~\ref{tab:scoreboard} read as follows:
\begin{align*}
  &\text{H\aa stad:} && \Gap{1-\eps}{\tfrac12+\eps}{Max-E3-Lin}\ \text{ and }\ \Gap{1}{\tfrac78+\eps}{Max-E3-SAT}\ \text{ are $\NP$-hard};\\
  &\text{UGC:} && \Gap{c}{\alpha_{\mathrm{GW}}\, c+\eps}{Max-Cut}\ \text{ is $\NP$-hard (for suitable $c$).}
\end{align*}

\begin{remark}[Why completeness $1-\eps$ rather than $1$]
Perfect completeness is not always attainable.  A system of linear equations over $\F_2$ whose
equations can \emph{all} be satisfied is solved by Gaussian elimination, so
$\Gap{1}{s}{Max-E3-Lin}\in\Ptime$ for every $s<1$ (Exercise~\ref{ex:gauss}).  Håstad's theorem has to
take $c=1-\eps$ for exactly this reason, and we will see the same phenomenon, for a similar
reason, in the Unique Games Conjecture (Section~\ref{sec:ug}).
\end{remark}


%% ===========================================================================
\section{Why ordinary reductions cannot create a gap}
\label{sec:nogap}
%% ===========================================================================

The NP-completeness theory of the 1970s \citep{Cook71,Levin73,Karp72} is a theory of
\emph{exact} problems, and its reductions are, in a precise sense, too gentle for
approximation.  This section explains why, and in doing so isolates what the PCP theorem
must supply.

\subsection{Cook--Levin yields a gap of one clause}

Recall how $3$SAT is shown $\NP$-hard.  Let $L\in\NP$ with verifier $V$ and let $x$ be an
input.  The Cook--Levin reduction outputs a formula $\varphi_x$ that encodes the computation
tableau of $V$ on $(x,w)$, where the witness $w$ is left as free variables.  The clauses say that
(a)~the first row of the tableau contains $x$, (b)~every $2\times3$ window of the tableau is
consistent with the transition function of $V$, and (c)~the last row contains the accepting state.

\begin{observation}\label{obs:cl}
For \emph{every} $x$---also for $x\notin L$---there is an assignment to $\varphi_x$ that
violates only the clause in~(c).  Hence $\OPT(\varphi_x)\ge 1-1/m$, where $m$ is the number
of clauses.
\end{observation}
\begin{proof}[Proof sketch]
Choose any witness $w$ and write down the true computation of $V$ on $(x,w)$ as a tableau.
Constraints (a) and (b) hold by construction; (c) fails if $V$ rejects $(x,w)$, but it is a single
clause.
\end{proof}

So the classical reduction maps YES instances to formulas with $\OPT=1$ and NO instances to
formulas with $\OPT=1-\tfrac1m$ exactly.  It proves that $\Gap{1}{1-1/m}{3SAT}$
is $\NP$-hard, which is nothing but the statement that maximizing the number of satisfied
clauses \emph{exactly} is $\NP$-hard.  The gap is a single clause out of $m$, and by
Lemma~\ref{lem:gap-approx} it rules out only approximation ratios above $1-\tfrac1m$, which
is no restriction at all on an algorithm.

\subsection{Karp reductions transport gaps, they do not create them}

Suppose now that someone \emph{gives us} a gap, and we want to carry it to another problem.
The standard reduction from $3$SAT to \sat{Independent Set} does this perfectly.

\begin{proposition}[The triangle reduction]\label{prop:triangle}
Let $\varphi$ be a $3$-CNF with $m$ clauses, and define the graph $G_\varphi$ as follows:
there is a vertex for every pair (clause $C_j$, literal $\ell\in C_j$); the three vertices of each
clause form a triangle; and two vertices $(j,\ell)$, $(j',\ell')$ are adjacent when
$\ell'=\neg\ell$.  Then the size of a maximum independent set of $G_\varphi$ equals the
maximum number of clauses of $\varphi$ that can be simultaneously satisfied.
\end{proposition}
\begin{proof}
($\ge$) Given an assignment satisfying $t$ clauses, pick for each satisfied clause one
satisfied literal.  The $t$ chosen vertices lie in different triangles, and no two of them carry
complementary literals, since both are true under the same assignment.  So they are independent.
($\le$) Let $S$ be an independent set.  It contains at most one vertex per triangle, and the
literals it contains are pairwise non-complementary, so they can all be set to true by a single
assignment, which then satisfies at least $|S|$ clauses.
\end{proof}

Consequently $\Gap{c}{s}{3SAT}$ reduces to the problem of distinguishing independent sets of
size $\ge cm$ from independent sets of size $\le sm$ in a graph with $3m$ vertices:
the gap is carried over \emph{unchanged}.  The same is true for almost all the textbook
reductions among $\NP$-complete problems.  They are \emph{local gadget constructions} that
preserve the quantity ``how many constraints can be satisfied''; so they preserve a gap if
there is one at the start, and they have none to preserve if there is none.

\begin{figure}[t]
\centering
\begin{tikzpicture}[x=1cm,y=1cm,every node/.style={font=\small}]
  % ---------- panel (a)
  \node[anchor=west,font=\sffamily\small\bfseries,color=accent] at (0,4.5) {(a)\; Cook--Levin / Karp reduction};
  \draw[-{Stealth[length=2mm]}] (0,3.4) -- (12.6,3.4);
  \foreach \x/\l in {0/{$\tfrac12$}, 12/{$1$}} \draw (\x,3.33)--(\x,3.47) node[below=6pt]{\l};
  \draw[line width=3pt,warn] (0,3.4) -- (11.55,3.4);
  \fill[claim] (12,3.4) circle (3pt);
  \node[warn,above] at (5.5,3.5) {NO: $\OPT\le 1-\tfrac1m$};
  \node[claim,above] at (12,3.5) {YES};
  \draw[decorate,decoration={brace,amplitude=2.5pt,mirror}] (11.55,2.55) -- (12,2.55)
        node[midway,below=3pt,font=\footnotesize]{gap $=\tfrac1m$};
  % ---------- panel (b)
  \node[anchor=west,font=\sffamily\small\bfseries,color=accent] at (0,0.95) {(b)\; After the PCP theorem};
  \draw[-{Stealth[length=2mm]}] (0,-0.15) -- (12.6,-0.15);
  \foreach \x/\l in {0/{$\tfrac12$}, 12/{$1$}} \draw (\x,-0.22)--(\x,-0.08) node[below=6pt]{\l};
  \draw[line width=3pt,warn] (0,-0.15) -- (9.3,-0.15);
  \fill[claim] (12,-0.15) circle (3pt);
  \node[warn,above] at (4.6,-0.05) {NO: $\OPT\le 1-\varepsilon_0$};
  \node[claim,above] at (12,-0.05) {YES};
  \draw[decorate,decoration={brace,amplitude=3pt,mirror}] (9.3,-1.0) -- (12,-1.0)
        node[midway,below=4pt,font=\footnotesize]{gap $=\varepsilon_0$, a constant};
\end{tikzpicture}
\caption{Where the optimum $\OPT(\varphi)$ lives for YES and NO instances of a reduction to
\sat{Max-3SAT} (schematic, not to scale).  The classical reductions leave a single clause between the two
cases; the PCP theorem moves the NO instances away by a constant fraction.}
\label{fig:gap}
\end{figure}

\begin{bigpicture}{Gap creation versus gap preservation}
Hardness of approximation has two ingredients.  \textbf{(i) Gap creation}: one reduction from
$\NP$ (\eg from $3$SAT) to a constraint satisfaction problem that turns ``satisfiable'' vs.\
``unsatisfiable'' into ``everything satisfiable'' vs.\ ``at most a $(1-\eps_0)$ fraction
satisfiable.''  This is the PCP theorem (Section~\ref{sec:pcp}), and it does not follow from
anything in the $\NP$-completeness toolbox.  \textbf{(ii) Gap preservation or
amplification}: a long chain of reductions that carry the gap to other problems, sharpening its
parameters on the way.  Lectures 3, 4 and~6 are about~(i); Lectures 5 and~7 are about~(ii); and Lectures 8--16 explore what
happens when ordinary gap amplification runs out of steam, which is where the Unique Games
Conjecture comes in.
\end{bigpicture}


%% ===========================================================================
\section{The PCP theorem: two faces}
\label{sec:pcp}
%% ===========================================================================

\subsection{Proofs that can be spot-checked}

A classical $\NP$ verifier reads its entire witness.  Suppose instead we insist that the
verifier be \emph{randomized} and may read only a constant number of bits of a (very long)
proof $\pi$.  Of course it must then sometimes accept wrong proofs; the demand is that it
reject every proof of a false statement with probability bounded away from $0$.

\begin{definition}[Probabilistically checkable proofs]\label{def:pcp}
A \defn{probabilistic verifier} $V$ with randomness $r(n)$ and query complexity $q(n)$ is a
polynomial-time algorithm that, on input $x$ of length $n$ and oracle access to a proof string
$\pi$, draws a random string $\rho\in\{0,1\}^{r(n)}$, \emph{non-adaptively} computes $q(n)$
positions $i_1(x,\rho),\dots,i_q(x,\rho)$, reads $\pi_{i_1},\dots,\pi_{i_q}$, and accepts or rejects
according to a predicate $P_{x,\rho}\colon\{0,1\}^q\to\{0,1\}$.
Let $0\le s<c\le1$.  A language $L$ belongs to $\PCP_{c,s}[r,q]$ if there is such a
verifier with
\begin{align*}
  x\in L &\ \Longrightarrow\ \exists\pi:\ \Pr_\rho[V^\pi(x)\text{ accepts}]\ge c &&\text{(completeness)},\\
  x\notin L &\ \Longrightarrow\ \forall\pi:\ \Pr_\rho[V^\pi(x)\text{ accepts}]\le s &&\text{(soundness)}.
\end{align*}
We write $\PCP[r,q]=\PCP_{1,1/2}[r,q]$.
\end{definition}

Two extreme cases show that the trade-off is not vacuous.  With $r=0$ and $q=n^{O(1)}$ the verifier
is deterministic and reads the whole witness: $\PCP[0,\poly(n)]=\NP$.  With $q=0$ the proof is
useless, and we get randomized polynomial time.  The question is how much \emph{randomness}
(which buys coverage of the proof) and how many \emph{queries} are needed.  Note that
$\PCP[O(\log n),O(1)]\subseteq\NP$ trivially: there are only $2^{O(\log n)}=\poly(n)$ random
strings, and together they query only polynomially many positions of $\pi$, so a nondeterministic
machine can guess these bits and check all random strings.  The surprise is the other direction.

\begin{figure}[t]
\centering
\begin{tikzpicture}[x=1cm,y=1cm,every node/.style={font=\small}]
  % proof string
  \foreach \i in {0,...,15} {\draw[fill=white] (\i*0.68,0) rectangle +(0.68,0.6);}
  \foreach \i in {2,7,12} {\draw[fill=accent!25] (\i*0.68,0) rectangle +(0.68,0.6);}
  \node[anchor=east] at (-0.15,0.3) {proof $\pi$};
  \node at (5.44,-0.55) {$\pi\in\{0,1\}^N$,\quad $N=\poly(n)$};
  % verifier
  \node[draw,rounded corners=3pt,fill=black!5,minimum width=3.4cm,minimum height=0.9cm] (V) at (5.44,2.8) {Verifier $V(x)$};
  \node[anchor=east,align=right] at (3.5,2.8) {random string\\ $\rho\in\{0,1\}^{O(\log n)}$};
  \draw[->] (3.55,2.8) -- (V.west);
  \foreach \i in {2,7,12} \draw[-{Stealth[length=1.8mm]},accent,thick] (V.south) to[out=-90,in=90] (\i*0.68+0.34,0.62);
  \node[anchor=west,align=left] at (7.6,2.8) {reads $\pi_{i_1},\pi_{i_2},\pi_{i_3}$ \;($q=O(1)$)\\
       accepts iff $P_{x,\rho}(\pi_{i_1},\pi_{i_2},\pi_{i_3})=1$};
\end{tikzpicture}
\caption{A PCP verifier.  For each of polynomially many random strings $\rho$ it inspects a constant
number of proof bits and applies a local test.  A correct proof passes every test; every proof of a
false statement fails at least half of them.}
\label{fig:pcp}
\end{figure}

\begin{theorem}[PCP theorem, proof-checking form; Arora--Safra, Arora--Lund--Motwani--Sudan--Szegedy]\label{thm:pcp1}
$\NP=\PCP[O(\log n),O(1)]$.
\end{theorem}

\begin{theorem}[PCP theorem, hardness form]\label{thm:pcp2}
There is a constant $\eps_0>0$ such that $\Gap{1}{1-\eps_0}{3SAT}$ is $\NP$-hard.
\end{theorem}

\noindent
Here $\Gap{1}{1-\eps_0}{3SAT}$ asks to distinguish satisfiable $3$-CNF formulas from formulas
in which every assignment violates an $\eps_0$-fraction of the clauses.  Compare with
Figure~\ref{fig:gap}: this is precisely the picture in panel~(b).  The two theorems look very
different---one about proofs, one about optimization---but they say the same thing.

\begin{proposition}[Equivalence of the two forms]\label{prop:pcp-equiv}
Theorem~\ref{thm:pcp1} holds if and only if Theorem~\ref{thm:pcp2} holds.
\end{proposition}

\begin{proof}
\emph{Hardness form $\Rightarrow$ proof-checking form.}  Let $L\in\NP$.  Compose the Cook--Levin
reduction with the reduction witnessing the $\NP$-hardness of $\Gap{1}{1-\eps_0}{3SAT}$ to obtain a
polynomial-time map $x\mapsto\varphi_x$ such that $x\in L$ implies that $\varphi_x$ is
satisfiable, and $x\notin L$ implies that every assignment violates at least $\eps_0 m$ clauses,
where $m=|\varphi_x|$.  The verifier expects as proof a satisfying assignment $\pi$ of $\varphi_x$.
It picks a clause uniformly at random ($\log m=O(\log n)$ random bits), reads the (at most) three
variables occurring in it, and accepts iff the clause is satisfied.  If $x\in L$, an honest
proof is accepted with probability~$1$.  If $x\notin L$ then \emph{whatever} $\pi$ is, it is an
assignment and thus violates an $\eps_0$-fraction of the clauses, so $V$ accepts with probability
at most $1-\eps_0$.  Now repeat the test $k=\lceil\ln2/\eps_0\rceil$ times with
independent clauses, accepting iff all $k$ clauses are satisfied.  The proof $\pi$ is fixed once
and for all, so the $k$ tests are independent given $\pi$ and the acceptance probability is at most
$(1-\eps_0)^k\le e^{-\eps_0k}\le\tfrac12$.  The verifier uses $k\log m=O(\log n)$ random bits and
$3k=O(1)$ queries.

\emph{Proof-checking form $\Rightarrow$ hardness form.}  Let $V$ be a verifier for $\sat{3SAT}\in\NP$
with $r=O(\log n)$ random bits, $q=O(1)$ queries, perfect completeness and soundness~$\tfrac12$.
(We may assume $q\ge3$.)  Given a formula $\psi$, we build a $3$-CNF $\varphi$ as follows.
Since $2^r=\poly(n)$, we can enumerate every random string $\rho$ and compute the query positions
$i_1(\rho),\dots,i_q(\rho)$ and the predicate $P_\rho$.  Introduce one Boolean variable $y_i$ for every
position $i$ of the proof.  For each $\rho$ we want to enforce $P_\rho(y_{i_1},\dots,y_{i_q})=1$.  This
constraint is a CNF with at most $2^q$ clauses of width $q$: one clause for each $a\in\{0,1\}^q$ with
$P_\rho(a)=0$, namely the clause that is false exactly on the input~$a$.  Finally, each clause
$(\ell_1\vee\dots\vee\ell_q)$ of width $q$ is converted into $q-2$ clauses of width~$3$ with fresh
variables $z_1,\dots,z_{q-3}$:
\[
  (\ell_1\vee\ell_2\vee z_1)\wedge(\neg z_1\vee\ell_3\vee z_2)\wedge\dots\wedge(\neg z_{q-3}\vee\ell_{q-1}\vee\ell_q).
\]
If some $\ell_j$ is true, one can choose the $z$'s to satisfy all $q-2$ clauses; if all $\ell_j$ are
false then $z_1$ must be true, hence $z_2$, \dots, hence $z_{q-3}$, and the last clause is violated, so
\emph{at least one} of the $q-2$ clauses fails for every choice of $z$.
Let $\varphi_\rho$ be the resulting $3$-CNF; it has at most $K:=(q-2)2^q$ clauses, and let $\varphi=\bigwedge_\rho\varphi_\rho$,
which has $m\le K\cdot 2^r$ clauses.

\emph{Completeness.}  If $\psi$ is satisfiable, there is a proof $\pi$ accepted for every $\rho$.  Setting $y=\pi$ and
choosing the $z$'s suitably satisfies all of $\varphi$.

\emph{Soundness.}  If $\psi$ is unsatisfiable, every proof $\pi$ is rejected for at least half of
the random strings.  Given any assignment to $\varphi$, read off $\pi=y$.  For every $\rho$ on which $V^\pi$
rejects, the assignment violates \emph{at least one} clause of $\varphi_\rho$.  There are at least $2^{r}/2$
such $\rho$, so at least $2^r/2$ clauses are violated out of $m\le K2^r$, \ie a fraction
$\ge 1/(2K)$.  Thus $\Gap{1}{1-\eps_0}{3SAT}$ is $\NP$-hard with $\eps_0=1/(2K)$, a constant.
\end{proof}

\begin{remark}[A dictionary]
The proof is our first instance of a translation that we shall use constantly:
\emph{a verifier is a constraint satisfaction problem}.  Random strings $\rho$ correspond to
constraints; positions of the proof correspond to variables; the predicate $P_\rho$ is the
constraint; the number $q$ of queries is the arity; soundness $s$ is the bound on the optimum
in the NO case.  We extend the dictionary to games in Section~\ref{sec:games}
(Table~\ref{tab:dict}).
\end{remark}

\begin{pitfall}{a proof is not a witness}
In Definition~\ref{def:pcp} the soundness condition quantifies over \emph{every} string~$\pi$,
including strings that have nothing to do with any computation.  A PCP is useful precisely
because its verifier cannot be fooled by an adversarially chosen $\pi$ that is merely
\emph{locally plausible}.  In the first half of the proof above this is why we can read off
an assignment from $\pi$ and apply the NO-case promise to it.
\end{pitfall}

The PCP theorem is a deep result: the original proofs by \citet{AroraSafra98} and
\citet{ALMSS98} are algebraic (low-degree polynomials and recursive proof composition), and \citet{Dinur07}
later found a purely combinatorial proof by gap amplification on expander graphs.  We take
Theorem~\ref{thm:pcp1} for granted in this lecture, return to its proofs in
Lectures~3, 4 and~6, and ask here what it gives us.

\begin{remark}[Two-sided history]
Theorem~\ref{thm:pcp2} reads as a statement about \emph{approximation}, yet the road to it was
paved by results about \emph{interactive proofs}: multi-prover interactive proofs
\citep{BGKW88}, $\mathsf{NEXP}=\mathsf{MIP}$ \citep{BFL91}, and the discovery by
\citet{FGLSS96} that a PCP verifier can be turned into a graph whose largest clique encodes the
maximum acceptance probability (Exercise~\ref{ex:fglss}), and hence that PCPs imply
hardness of approximating \sat{Max-Clique}.  Section~\ref{sec:games} follows the thread through
the multi-prover systems.
\end{remark}


%% ===========================================================================
\section{Two provers, Label Cover, and parallel repetition}
\label{sec:games}
%% ===========================================================================

The gap in the PCP theorem is a small constant: soundness $1-\eps_0$.  To prove hardness
results with sharp constants like $\tfrac78$ and $\tfrac12$ we need soundness $s\to0$, and
we need it while keeping the \emph{number of queries} under control.  The cleanest setting for
this is the following.

\subsection{Games, labelings, projections}

\begin{definition}[Two-prover one-round game / Label Cover]\label{def:game}
A \defn{two-prover one-round game} $G$ consists of a bipartite graph $(U,W,E)$ with finite
label sets $\Sigma_U,\Sigma_W$ and, for each edge $e=(u,w)$, a \defn{constraint}
$C_e\subseteq\Sigma_U\times\Sigma_W$.  The referee draws an edge $e=(u,w)$ uniformly at random,
sends $u$ to Alice and $w$ to Bob, receives answers $a\in\Sigma_U$ and $b\in\Sigma_W$ from
the provers (who cannot communicate), and accepts iff $(a,b)\in C_e$.  The \defn{value} of $G$,
\[
  \val(G)\;=\;\max_{\ell_U\colon U\to\Sigma_U,\ \ell_W\colon W\to\Sigma_W}\ \Pr_{e=(u,w)\sim E}\big[(\ell_U(u),\ell_W(w))\in C_e\big],
\]
is the maximum over \defn{labelings}, \ie over deterministic strategies of the two provers.
Finding an optimal labeling is the optimization problem \defn{Label Cover}.  The game is a
\defn{projection game} if every $C_e$ is the graph of a function $\pi_e\colon\Sigma_U\to\Sigma_W$:
$C_e=\{(a,\pi_e(a))\}$.  It is a \defn{Unique Game} if moreover $|\Sigma_U|=|\Sigma_W|=k$
and every $\pi_e$ is a bijection.
\end{definition}

Restricting to deterministic strategies loses nothing.

\begin{lemma}[Determinism suffices]\label{lem:determ}
The maximum acceptance probability over all (randomized, possibly correlated through shared
randomness) strategies of Alice and Bob equals $\val(G)$.
\end{lemma}
\begin{proof}
A randomized strategy with shared random string $r$ is, for each fixed $r$, a deterministic
strategy.  Its acceptance probability is the average over $r$ of the acceptance probabilities of
the deterministic strategies, hence at most the maximum.
\end{proof}

\begin{lemma}[Random labeling]\label{lem:random-label}
In a projection game, $\val(G)\ge 1/|\Sigma_W|$.  In particular, for every Unique Game with $k$ labels
$\val(G)\ge1/k$.
\end{lemma}
\begin{proof}
Fix any labels for the $u\in U$ and give every $w\in W$ an independent uniformly random
label.  For each edge $e=(u,w)$ the label $\pi_e(\ell_U(u))$ is fixed, so $w$ matches it with
probability exactly $1/|\Sigma_W|$.  By linearity of expectation some labeling achieves at least that.
\end{proof}

\noindent
Lemma~\ref{lem:random-label} will matter later: no hardness theorem can have soundness $\delta$ with
fewer than $1/\delta$ labels.

\begin{example}[The clause--variable game]\label{ex:cv}
Let $\varphi$ be a $3$-CNF in which every clause has three distinct variables.  Alice plays the
\emph{clause prover}, Bob the \emph{variable prover}.  The referee picks a uniformly random clause
$C$ and a uniformly random variable $x$ in it.  Alice receives $C$ and must answer with a
satisfying assignment of its three variables ($7$ possibilities); Bob receives $x$ and answers
with a bit.  They win iff Bob's bit equals Alice's value for $x$.  This is a projection game with
$|\Sigma_U|=7$ and $|\Sigma_W|=2$ (Figure~\ref{fig:cv}); the projection $\pi_{(C,x)}$ sends an
assignment of $C$ to its restriction to~$x$.
\end{example}

\begin{figure}[t]
\centering
\begin{tikzpicture}[x=1cm,y=1cm,every node/.style={font=\small},
   cl/.style={draw,rounded corners=2pt,minimum width=4.1cm,minimum height=0.62cm,fill=accent!8},
   vr/.style={draw,circle,minimum size=0.75cm,fill=warn!10}]
  \node[font=\sffamily\small\bfseries,color=accent] at (0,3.3) {Alice: clause prover};
  \node[font=\sffamily\small\bfseries,color=warn]   at (8,3.3) {Bob: variable prover};
  \node[cl] (c1) at (0,2.2) {$C_1=(x_1\vee x_2\vee\neg x_3)$};
  \node[cl] (c2) at (0,0.9) {$C_2=(\neg x_1\vee x_3\vee x_4)$};
  \node[cl] (c3) at (0,-0.4) {$C_3=(\neg x_2\vee\neg x_3\vee\neg x_4)$};
  \node[vr] (x1) at (8,2.6) {$x_1$};
  \node[vr] (x2) at (8,1.5) {$x_2$};
  \node[vr] (x3) at (8,0.4) {$x_3$};
  \node[vr] (x4) at (8,-0.7) {$x_4$};
  \draw[black!55] (c1.east)--(x1) (c1.east)--(x2) (c1.east)--(x3)
        (c2.east)--(x1) (c2.east)--(x3) (c2.east)--(x4)
        (c3.east)--(x2) (c3.east)--(x3) (c3.east)--(x4);
  \draw[accent,very thick] (c2.east)--(x3);
  \node[font=\footnotesize,align=center,color=accent] at (0,-1.45)
     {answer: one of the $7$ satisfying\\ assignments of the clause};
  \node[font=\footnotesize,align=center,color=warn] at (8,-1.45)
     {answer: a bit\\ (value of the variable)};
\end{tikzpicture}
\caption{The clause--variable game of a $3$-CNF.  An edge is a pair (clause, variable occurring in it).
The highlighted edge $(C_2,x_3)$ is accepted iff Alice's assignment to $C_2$ and Bob's bit agree on~$x_3$,
\ie iff the projection $\pi_{(C_2,x_3)}$ of Alice's label to the coordinate $x_3$ equals Bob's label.}
\label{fig:cv}
\end{figure}

\begin{proposition}[Soundness of the clause--variable game]\label{prop:cv}
If $\varphi$ is satisfiable then $\val(G_\varphi)=1$.  If every assignment violates at least an
$\eps$-fraction of the clauses then $\val(G_\varphi)\le1-\eps/3$.
\end{proposition}
\begin{proof}
If $\varphi$ is satisfiable, let both provers answer according to a satisfying assignment.
Conversely, fix deterministic strategies; Bob's answers define an assignment $b$ to the variables.
By assumption $b$ violates at least an $\eps$-fraction of the clauses.  Consider a clause $C$ violated by $b$.
Alice's answer $a_C$ is a satisfying assignment of $C$ (otherwise she loses outright), so $a_C\ne b|_C$,
and therefore $a_C$ disagrees with $b$ on at least one of the three variables of $C$.  The referee asks
exactly that variable with probability at least $\tfrac13$, and rejects.  Hence the rejection probability is
at least $\eps\cdot\tfrac13$.
\end{proof}

\noindent
Combined with Theorem~\ref{thm:pcp2}, Proposition~\ref{prop:cv} shows that it is $\NP$-hard to
distinguish $\val(G)=1$ from $\val(G)\le1-\eps_0/3$ for projection games over alphabets of
size~$7$ and $2$.  We have now met the ``game'' column of the dictionary.

\begin{table}[t]
\centering\small
\renewcommand{\arraystretch}{1.3}
\begin{tabular}{@{}llll@{}}
\toprule
PCP verifier & CSP & Two-prover game & Label Cover\\
\midrule
random string $\rho$ & constraint & question pair $(u,w)$ & edge $e$\\
proof position & variable & (a prover's answer table) & vertex\\
$q$ queries & arity $q$ & $2$ provers & arity $2$\\
predicate $P_\rho$ & constraint predicate & referee's test & relation $C_e$\\
alphabet $\{0,1\}$ & domain & answer set & label set $\Sigma$\\
soundness $s$ & $\OPT\le s$ & $\val(G)\le s$ & $\val\le s$\\
\bottomrule
\end{tabular}
\caption{The dictionary.  The same object is a verifier, a constraint satisfaction problem, a game, and a
Label Cover instance.}
\label{tab:dict}
\end{table}

\subsection{Parallel repetition}

To drive the soundness from $1-\eps_0/3$ to a small $\delta$ we want to repeat.
For a one-oracle verifier, as we saw in the proof of Proposition~\ref{prop:pcp-equiv}, independent
repetition multiplies the acceptance probabilities.  For a game we are not free to repeat
sequentially: that would require several rounds of communication, which we want to avoid
(otherwise the provers could coordinate).  Instead one repeats \emph{in parallel}.

\begin{definition}[Parallel repetition]
The \defn{$k$-fold parallel repetition} $G^{\otimes k}$ of a game $G$: the referee draws $k$ independent
edges $e_i=(u_i,w_i)$, sends $(u_1,\dots,u_k)$ to Alice and $(w_1,\dots,w_k)$ to Bob, receives
$a\in\Sigma_U^k$ and $b\in\Sigma_W^k$, and accepts iff $(a_i,b_i)\in C_{e_i}$ for every $i$.
\end{definition}

Alice and Bob can play each coordinate independently with an optimal strategy of $G$, so
$\val(G^{\otimes k})\ge\val(G)^k$.  Is this an equality?  It is not, and the failure is
dramatic enough to deserve a warning.

\begin{pitfall}{parallel repetition is not independent repetition}
Let $\ell\ge3$ be odd.  In the \defn{odd cycle game} the questions are the vertices of an
$\ell$-cycle.  With probability $\tfrac12$ the referee sends the \emph{same} vertex to both provers and
demands the same color ($0$ or~$1$); with probability $\tfrac12$ it sends the two \emph{endpoints
of a random edge}, in random order, and demands different colors.  The provers are trying to convince
the referee that an odd cycle is $2$-colorable.  Its value is $1-\tfrac1{2\ell}$ (Exercise~\ref{ex:oddcycle});
this is a Unique Game with two labels.  Naively, $k$ repetitions should yield value about
$(1-\tfrac1{2\ell})^k\approx e^{-k/2\ell}$, which for $k=\ell^2/100$ is exponentially small in $\ell$.
In fact \citet{Raz11} showed that for $k\le\ell^2$ there is a strategy winning $G^{\otimes k}$ with
probability $1-O(\sqrt k/\ell)$, so that even for $k=c\,\ell^2$ with a small constant~$c$ the
value of $G^{\otimes k}$ is still at least~$\tfrac12$ (this is essentially tight by \citealp{FKO07}).
Exercise~\ref{ex:fortnow} gives a much smaller example in which already the $2$-fold game beats $\val^2$.
Examples of this kind go back to Fortnow and Feige around 1990, and for several years it was a
famous open problem whether repetition drives the value to $0$ at all.
\end{pitfall}

The qualitative fact that repetition \emph{does} drive the value to $0$ is the theorem of Raz.

\begin{theorem}[Parallel repetition; Raz 1998, Holenstein 2009, Rao 2011]\label{thm:raz}
There is an absolute constant $c>0$ such that for every two-prover one-round game $G$ with
$\val(G)\le1-\eta$ and answer sets of size at most~$\sigma$,
\[
  \val(G^{\otimes k})\;\le\;\bigl(1-\eta^{c}\bigr)^{\Omega(k/\log\sigma)}.
\]
For projection games, Rao showed $\val(G^{\otimes k})\le(1-\Omega(\eta^2))^k$, independently of $\sigma$.
\end{theorem}

\noindent
Raz's original paper has $c=32$ \citep{Raz98}; Holenstein simplified the proof and improved this to $c=3$
\citep{Holenstein09}.  (We prove none of this until Lecture~5.)  Putting the PCP theorem, the
clause--variable game and parallel repetition together, we obtain the result that all of
Lectures 7--16 take as their starting point.

\begin{theorem}[Label Cover theorem]\label{thm:lc}
For every $\delta>0$ there is a $\sigma=\poly(1/\delta)$ and a reduction, running in time
$n^{O(\log(1/\delta))}$ on formulas of size $n$, from $3$SAT to projection games with label sets
of size at most $\sigma$, such that
\[
  \varphi\text{ satisfiable}\ \Rightarrow\ \val(G_\varphi)=1,\qquad
  \varphi\text{ unsatisfiable}\ \Rightarrow\ \val(G_\varphi)\le\delta.
\]
In other words $\Gap{1}{\delta}{LabelCover}_{\sigma}$ is $\NP$-hard.
\end{theorem}
\begin{proof}
Apply Theorem~\ref{thm:pcp2} to obtain $\varphi'$ with the $\eps_0$-gap, then form the
clause--variable game $G=G_{\varphi'}$; by Proposition~\ref{prop:cv} we have $\val(G)=1$ or $\val(G)\le1-\eta$
with $\eta=\eps_0/3$.  Take $G^{\otimes k}$ with $k=O(\log(1/\delta)/\eta^2)=O(\log(1/\delta))$.  Repetition
of a projection game is again a projection game (the projection is $\pi_{e_1}\times\dots\times\pi_{e_k}$).
If $\val(G)=1$ then $\val(G^{\otimes k})=1$.  Otherwise Theorem~\ref{thm:raz} gives $\val(G^{\otimes k})\le(1-\Omega(\eta^2))^k\le\delta$.
The new label sets have sizes $7^k$ and $2^k$, which is $\poly(1/\delta)$, and $G^{\otimes k}$ has at most $(3m)^k$ edges.
\end{proof}

\noindent
Two remarks about the price we paid.  First, the alphabet size $\sigma$ is polynomial in $1/\delta$,
and by Lemma~\ref{lem:random-label} it \emph{must} be at least $1/\delta$.  Second, the projection games we
have produced are $d$-to-$1$ with large $d$: each label of Bob has about $\sigma/2^k$ preimages.
Both facts will be the stage on which Unique Games enters.


%% ===========================================================================
\section{Unique Games and the Unique Games Conjecture}
\label{sec:ug}
%% ===========================================================================

\subsection{Definition and a first observation}

We now restrict the projection constraints as much as possible.  It is convenient (and customary)
to define Unique Games on a general graph.

\begin{definition}[Unique Games]\label{def:ug}
A \defn{Unique Games instance} $\mathcal U=(G,[k],\{\pi_e\})$ consists of a graph $G=(V,E)$, a label
set $[k]$, and for every edge $e=(u,v)$ a \emph{bijection} $\pi_e\colon[k]\to[k]$.  A labeling
$\ell\colon V\to[k]$ \defn{satisfies} $e$ if $\pi_e(\ell(u))=\ell(v)$, and
$\val(\mathcal U)=\max_\ell\Pr_{e\sim E}[\ell\text{ satisfies }e]$.  Since $\pi_e$ is a bijection,
$\pi_e(\ell(u))=\ell(v)\iff\ell(u)=\pi_e^{-1}(\ell(v))$, so the orientation of an edge is immaterial.
\end{definition}

\begin{example}\label{ex:ug-examples}
\sat{Max-Cut} is a Unique Game with $k=2$ and $\pi_e=$ the transposition $0\leftrightarrow1$ for every edge.
More generally, a system of equations $x_u-x_v=c_e\pmod k$ is a Unique Game with $\pi_e(i)=i+c_e$: two-variable
linear equations over $\Z_k$.  By contrast, \sat{Max-3-Cut} ($x_u\ne x_v$ with $k=3$) is \emph{not} a Unique Game:
given $x_u$, the constraint allows two values of $x_v$, not one (Exercise~\ref{ex:not-ug}).
\end{example}

The word ``unique'' means that, given the label of one endpoint, the label of the other endpoint is forced.
That makes perfect completeness trivial.

\begin{proposition}[Unique Games with perfect completeness are easy]\label{prop:ug-easy}
Given a Unique Games instance with $k$ labels, one can decide in time $O(k\,(|V|+|E|))$ whether $\val(\mathcal U)=1$, and find a
perfect labeling if there is one.
\end{proposition}
\begin{proof}
Treat each connected component separately.  Pick a root $r$ and try each of the $k$ labels $a$ for it.  Propagate
along a BFS tree: if $u$ has label $\ell(u)$ and $(u,v)$ is a tree edge then set $\ell(v)=\pi_e(\ell(u))$
(or $\pi_e^{-1}(\ell(u))$ if the edge is traversed against its orientation).  Check all edges.
If a perfect labeling $\ell^\ast$ exists, then with $a=\ell^\ast(r)$ the propagation reconstructs $\ell^\ast$ by induction along
the tree, because bijectivity makes each step forced.  So some $a$ passes the check.
\end{proof}

\begin{remark}
The argument breaks down for projection games, where the label of $u$ determines that of $w$ but
not vice versa: propagating from $w$ to $u$ requires choosing among the $d$ preimages, and the
branching is exponential.  Indeed, deciding $\val=1$ for projection games is $\NP$-hard (Proposition~\ref{prop:cv}).
Likewise, as in Section~\ref{sec:gap}, if a hard Unique Games gap problem exists, it must have
completeness $1-\eps<1$: an approximately satisfiable instance may have a bad edge on every
propagation path, and the instance has to be built so that no local repair works.
\end{remark}

\subsection{The conjecture}

\begin{conjecture}[Unique Games Conjecture; Khot 2002]\label{conj:ugc}
For every $\eps,\delta>0$ there is a $k=k(\eps,\delta)$ such that $\Gap{1-\eps}{\delta}{UG}_k$ is $\NP$-hard:
it is $\NP$-hard to distinguish Unique Games instances over $k$ labels with $\val\ge1-\eps$ from those with
$\val\le\delta$.
\end{conjecture}

\noindent
Every word of this statement carries weight.

\begin{enumerate}[label=(\roman*)]
  \item \emph{Quantifiers.}  First $\eps$ and $\delta$ are chosen, \emph{then} the alphabet~$k$.  The
        alphabet is a constant but it grows without bound as $\eps,\delta\to0$.  It must:
        by Lemma~\ref{lem:random-label}, $k\ge1/\delta$.
  \item \emph{A family of problems.}  The UGC is not a statement about one problem, but about
        a family of promise problems parametrised by $(\eps,\delta)$.  (Hence ``UG is $\NP$-hard to
        approximate'' is a misleading summary.)
  \item \emph{Completeness $1-\eps$.}  By Proposition~\ref{prop:ug-easy} we cannot ask for completeness~$1$.
  \item \emph{Soundness $\delta\to 0$.}  This is the opposite extreme to the PCP theorem, which only
        gets soundness $1-\eps_0$: the conjecture says that even ``almost satisfiable'' and ``almost
        completely unsatisfiable'' are indistinguishable.
  \item \emph{Reductions.}  ``$\NP$-hard'' means via polynomial-time Karp reductions from $3$SAT.  Some
        applications are known to go through under weaker hypotheses, \eg randomized reductions.
\end{enumerate}

\begin{remark}
There are many variants of the conjecture: with bipartite graphs, with linear constraints
$\pi_e(i)=i+c_e$ (a special case, believed hard too), with ``strong'' versions in which
soundness is measured on vertex sets rather than edges.  Their relationships are surveyed by
\citet{Khot10} and \citet{Trevisan12}.  The manuscript discussed in Section~\ref{sec:map} addresses the
usual edge-value form of Conjecture~\ref{conj:ugc} with translation constraints on a bipartite graph.
\end{remark}

\subsection{Why Unique Games?}

If we can already prove $\Gap{1}{\delta}{LabelCover}$ is $\NP$-hard for every $\delta$ (Theorem~\ref{thm:lc}), why would
anyone want a new conjecture?  Because Label Cover is the wrong tool for exactly the problems in the last two
rows of Table~\ref{tab:scoreboard}.

\begin{figure}[t]
\centering
\begin{tikzpicture}[x=1cm,y=1cm,every node/.style={font=\small},
   la/.style={circle,fill=accent,inner sep=1.7pt}, lb/.style={circle,fill=warn,inner sep=1.7pt}]
  % left panel: 2-to-1
  \node[font=\sffamily\small\bfseries] at (1.4,0.55) {projection ($2$-to-$1$)};
  \foreach \i in {1,...,6} \node[la] (a\i) at (0,-\i*0.5) {};
  \foreach \j in {1,2,3} \node[lb] (b\j) at (2.8,-\j+0.25) {};
  \draw[->,thick] (a1)--(b1); \draw[->,thick] (a2)--(b1);
  \draw[->,thick] (a3)--(b2); \draw[->,thick] (a4)--(b2);
  \draw[->,thick] (a5)--(b3); \draw[->,thick] (a6)--(b3);
  \node at (0,-3.5) {$\Sigma_U$}; \node at (2.8,-3.5) {$\Sigma_W$};
  % right panel: bijection
  \begin{scope}[xshift=7.2cm]
  \node[font=\sffamily\small\bfseries] at (1.4,0.55) {unique (bijection)};
  \foreach \i in {1,...,4} \node[la] (c\i) at (0,-\i*0.75+0.25) {};
  \foreach \j in {1,...,4} \node[lb] (d\j) at (2.8,-\j*0.75+0.25) {};
  \draw[->,thick] (c1)--(d3); \draw[->,thick] (c2)--(d1);
  \draw[->,thick] (c3)--(d4); \draw[->,thick] (c4)--(d2);
  \node at (0,-3.5) {$[k]$}; \node at (2.8,-3.5) {$[k]$};
  \end{scope}
\end{tikzpicture}
\caption{A constraint $\pi_e$ between the label sets of two vertices.  In a projection game (left) a label
of $w$ has several preimages; in a Unique Game (right) the constraint is a perfect matching.}
\label{fig:uniq}
\end{figure}

\begin{bigpicture}{The two barriers behind Khot's conjecture}
\textbf{(1) The two-query barrier.}  H\aa stad's optimal results for \sat{Max-E3-SAT} and \sat{Max-E3-Lin} come
from PCPs that read \emph{three} bits.  A problem such as \sat{Max-Cut}, \sat{Vertex Cover} or
\sat{Min-2SAT-Deletion} has constraints on \emph{two} variables, so a PCP proving its hardness
must use only \emph{two} queries---and two-query tests are intrinsically much weaker.
Label Cover is the natural two-query starting point, with alphabet as long as we please.
\textbf{(2) The many-to-one barrier.}  Roughly speaking, a hardness reduction from Label Cover replaces every vertex
by a code (the \emph{long code}) and every constraint $\pi_e$ by a local test between two codewords;
the analysis of a cheating prover must \emph{decode} the code tables back into a labeling of the
Label Cover instance.  When $\pi_e$ is $d$-to-$1$ with large $d$, the decoding can only recover the label of $w$
up to a list of $d$ candidates for $u$, and the loss in soundness is too large to reach
the thresholds of Table~\ref{tab:scoreboard}.  When $\pi_e$ is a bijection there is no such
loss.  Khot's insight \citep{Khot02} is that one can hope that a \emph{unique} Label Cover is just as hard as general Label Cover,
once completeness is relaxed to $1-\eps$, and that this hypothesis---the UGC---would then imply tight
hardness for two-variable problems.  (We make all of this precise in Lectures 8 and 9.)
\end{bigpicture}

\subsection{Evidence for and against}

Nobody knows whether the UGC is true.  The state of the evidence, up to this
writing, is a mixture of three things.

\begin{itemize}
  \item \emph{Algorithms.}  If $\val\ge1-\eps$, one can find in polynomial time a labeling of value $1-O(\sqrt{\eps\log k})$
        \citep{CMM06}, and in time $\exp(n^{\eps^{\Omega(1)}})$ a labeling of value $1-\eps^{\Omega(1)}$ \citep{ABS15}.
        No algorithm contradicts the conjecture, but the subexponential algorithm shows that Unique Games
        is not as hard as $3$SAT is believed to be (Lecture~12).
  \item \emph{Integrality gaps.}  The semidefinite program that powers the algorithms above cannot
        refute the conjecture: it has integrality gaps \citep{KhotVishnoi15} (Lecture~10).
  \item \emph{Partial hardness.}  The $2$-to-$2$ Games Conjecture, a weakened version of the UGC in
        which each constraint is $2$-to-$2$ instead of a bijection, is a \emph{theorem}
        \citep{KMS17,DKKMS18,KMS18}; see \citet{Khot19} for an exposition.  It already implies
        $\NP$-hardness of approximating \sat{Vertex Cover} within $\sqrt2-\eps$, the best-known (Lectures~13--15).
\end{itemize}


%% ===========================================================================
\section{The map of the territory}
\label{sec:map}
%% ===========================================================================

Figure~\ref{fig:map} places the main results of this course on a single page; Table~\ref{tab:timeline}
gives the chronology.  Read an arrow as ``is used to prove'' (or ``builds on'').  One item on the map
needs a health warning.

\begin{figure}[p]
\centering
\makebox[\textwidth][c]{%
\hyphenpenalty=10000 \exhyphenpenalty=10000
\begin{tikzpicture}[x=1cm,y=1cm,
   every node/.style={font=\footnotesize},
   box/.style={draw,rounded corners=3pt,align=center,inner sep=5pt,fill=accent!7,
               text width=4.3cm,minimum height=1.7cm},
   side/.style={box,text width=4.5cm},
   boxw/.style={box,fill=warn!10,draw=warn,thick},
   boxc/.style={side,fill=claim!7,draw=claim},
   arr/.style={-{Stealth[length=2.2mm]},thick,accent!80!black},
   darr/.style={{Stealth[length=2.2mm]}-{Stealth[length=2.2mm]},thick,dashed,warn},
   carr/.style={-{Stealth[length=2.2mm]},thick,dashed,claim}]
  % spine
  \node[box] (cook) at (0,0)    {\textbf{$3$SAT is $\NP$-complete}\\ Cook, Levin, Karp 1971--72};
  \node[box] (pcp)  at (0,-3.0) {\textbf{PCP theorem}\\ $\Gap{1}{1-\eps_0}{3SAT}$ is $\NP$-hard\\
       Arora--Safra, ALMSS 1992;\\ Dinur 2006\\ \emph{Lectures 3, 4, 6}};
  \node[box] (lc)   at (0,-6.0) {\textbf{Label Cover with}\\ \textbf{tiny soundness}\\
       Raz 1995: parallel repetition\\ \emph{Lectures 2, 5}};
  \node[boxw] (ugc) at (0,-9.0) {\textbf{Unique Games}\\ \textbf{Conjecture}\\ Khot 2002\\ \emph{Lecture 8}};
  \node[box,minimum height=3.0cm] (cons) at (0,-12.6) {\textbf{If the UGC holds}\\
       \sat{Max-Cut} is tight \emph{(L9)}\\ \sat{Vertex Cover}: $2-\eps$ \emph{(L9)}\\
       integrality gaps \emph{(L10)}\\ every CSP \emph{(L11)}};
  % left column
  \node[side] (has) at (-6.0,-9.0)  {\textbf{Optimal $3$-bit PCPs}\\ H\aa stad 1997\\
       \sat{Max-E3-SAT}, \sat{Max-E3-Lin}\\ \emph{Lecture 7}};
  \node[side] (two) at (-6.0,-12.6) {\textbf{$2$-to-$2$ Games theorem}\\ reduction from H\aa stad's $3$-Lin\\
       a ``half-way'' UGC\\ \sat{Vertex Cover}: $\sqrt2-\eps$\\ \emph{Lectures 13--15}};
  % right column
  \node[boxc] (oai) at (6.0,-6.0)  {\textbf{Claimed proof of the UGC}\\ OpenAI manuscript, 23 Sep.\ 2026\\
       \emph{not independently verified}\\ \emph{Lecture 16}};
  \node[side] (abs) at (6.0,-9.0)  {\textbf{Subexponential algorithm}\\ \textbf{for Unique Games}\\
       Arora--Barak--Steurer 2010\\ in tension with the UGC\\ \emph{Lecture 12}};
  % arrows
  \draw[arr] (cook)--(pcp);
  \draw[arr] (pcp)--(lc);
  \draw[arr] (lc)--(ugc);
  \draw[arr] (lc.west) -| (has.north);
  \draw[arr] (ugc)--(cons);
  \draw[arr] (has)--(two);
  \draw[darr] (ugc.east)--(abs.west);
  \draw[carr] (oai.south west)--(ugc.north east);
\end{tikzpicture}}
\caption{The route from $\NP$-completeness to Unique Games.  Solid blue arrows are established
implications.  The orange double arrow records a \emph{tension}: a subexponential algorithm for
Unique Games sits uneasily next to a conjecture asserting $\NP$-hardness.  The dashed wine-red arrow
marks a claim whose status is explained in the box below.  \emph{L$n$} abbreviates ``Lecture~$n$''.}
\label{fig:map}
\end{figure}

\begin{statusbox}{the 2026 claim}
On 23~September~2026 OpenAI released a manuscript \emph{The Unique Games Theorem} \citep{OpenAI26}.  Its abstract
reads: ``We prove the Unique Games Conjecture.  For every fixed $\eps,\delta\in(0,1/2)$, we give a
deterministic polynomial-time reduction from $3$SAT to Unique Games over a fixed finite alphabet, with
completeness at least $1-\eps$ and soundness at most $\delta$.''

\smallskip
\textbf{Status.}  At the time of writing we know of no independent expert verification of the proof, and no
refutation; we have not verified it ourselves.  Conjecture~\ref{conj:ugc} is therefore treated as
open, and every statement labelled ``assuming the UGC'' in these notes is \emph{conditional}.  Lecture~16
compares the manuscript's main theorem with the classical formulations and discusses what a careful
verification would have to check.  Until then we keep apart results established by the community,
results in a complete but unrefereed paper, and claims published by their authors only.
\end{statusbox}

\begin{table}[H]
\centering\small
\renewcommand{\arraystretch}{1.22}
\begin{tabular}{@{}r p{11.2cm}@{}}
\toprule
Year & Event\\
\midrule
1971--72 & Cook, Levin, Karp: $\NP$-completeness \citep{Cook71,Levin73,Karp72}\\
1974 & Johnson: systematic study of approximation algorithms \citep{Johnson74}\\
1988 & Ben-Or, Goldwasser, Kilian, Wigderson: multi-prover interactive proofs \citep{BGKW88}\\
1990 & Babai, Fortnow, Lund: $\mathsf{NEXP}=\mathsf{MIP}$ \citep{BFL91}\\
1991 & Feige, Goldwasser, Lov\'asz, Safra, Szegedy: PCPs imply hardness of \sat{Clique} \citep{FGLSS96}\\
1992 & Arora--Safra and Arora--Lund--Motwani--Sudan--Szegedy: the PCP theorem \citep{AroraSafra98,ALMSS98};
         Feige--Lov\'asz: two-prover one-round proof systems \citep{FeigeLovasz92}\\
1994 & Goemans--Williamson: $0.878$-approximation for \sat{Max-Cut} \citep{GW95}\\
1995 & Raz: parallel repetition \citep{Raz98}; Bellare--Goldreich--Sudan: the long code \citep{BGS98}\\
1996 & H\aa stad: \sat{Clique} is hard within $n^{1-\eps}$ \citep{Hastad99}\\
1997 & H\aa stad: optimal inapproximability for \sat{Max-E3-SAT} and \sat{Max-E3-Lin} \citep{Hastad01}\\
2002 & Khot: the Unique Games Conjecture \citep{Khot02}\\
2003 & Khot--Regev: \sat{Vertex Cover} is hard within $2-\eps$ under the UGC \citep{KhotRegev08}\\
2004--05 & Khot--Kindler--Mossel--O'Donnell: \sat{Max-Cut} \citep{KKMO07}; Khot--Vishnoi: integrality gaps \citep{KhotVishnoi15}\\
2006 & Dinur: the PCP theorem by gap amplification \citep{Dinur07}\\
2008 & Raghavendra: every CSP \citep{Raghavendra08}\\
2010 & Arora--Barak--Steurer: subexponential algorithm for Unique Games \citep{ABS15}\\
2017--18 & Khot--Minzer--Safra, Dinur--Khot--Kindler--Minzer--Safra: $2$-to-$2$ Games \citep{KMS17,DKKMS18,KMS18}\\
2026 & OpenAI manuscript claiming the UGC \citep{OpenAI26}\\
\bottomrule
\end{tabular}
\caption{A timeline of the main results (conference years; journal versions are later).  Many important
contributions are missing.}
\label{tab:timeline}
\end{table}


%% ===========================================================================
\section*{Exercises}
\addcontentsline{toc}{section}{Exercises}
%% ===========================================================================

\begin{exercise}[\stars{$\star$}]\label{ex:derand}
Derandomize the three random-assignment algorithms of Example~\ref{ex:random} with the method of
conditional expectations: fix the variables one at a time, always choosing the value that does not
decrease the conditional expected number of satisfied constraints.  Why is this efficiently computable?
\end{exercise}

\begin{exercise}[\stars{$\star$}]\label{ex:minimization}
State and prove the analogue of Lemma~\ref{lem:gap-approx} for minimization problems, where YES
instances have $\OPT\le c$ and NO instances have $\OPT>s$ with $c<s$.  What is the
inapproximability ratio?
\end{exercise}

\begin{exercise}[\stars{$\star$}]\label{ex:gauss}
Show that $\Gap{1}{s}{Max-E3-Lin}$ is in $\Ptime$ for every $s<1$.  (Hence H\aa stad's theorem
cannot have perfect completeness.)  Why does the same argument not apply to $\Gap{1}{\tfrac78+\eps}{Max-E3-SAT}$?
\end{exercise}

\begin{exercise}[\stars{$\star\star$}]\label{ex:fglss}
Let $V$ be a PCP verifier with $r$ random bits and $q$ queries, and input $x$.  Build a graph $H_x$ whose vertices are
the pairs $(\rho,a)$ with $\rho\in\{0,1\}^r$ and $a\in\{0,1\}^q$ satisfying $P_\rho(a)=1$, and in which
$(\rho,a)$ and $(\rho',a')$ are adjacent if $\rho\ne\rho'$ and the two answer vectors do not contradict each other on
positions queried by both.  Prove that the largest clique of $H_x$ has size $2^r\cdot\max_\pi\Pr_\rho[V^\pi(x)\text{ accepts}]$.
Conclude, using the PCP theorem, that it is $\NP$-hard to approximate \sat{Max-Clique} within some constant factor.
(Amplifying this to a polynomial factor needs randomness-efficient repetition; \citealp{FGLSS96}.)
\end{exercise}

\begin{exercise}[\stars{$\star$}]\label{ex:repeat}
Let $V$ be a one-oracle verifier with soundness $s$.  Prove that its $k$-fold independent repetition has soundness $s^k$.
Identify exactly where the argument uses that there is one fixed proof $\pi$, and explain why the same
argument breaks for parallel repetition of two-prover games.
\end{exercise}

\begin{exercise}[\stars{$\star\star$}]\label{ex:oddcycle}
Prove that the odd cycle game has value exactly $1-\tfrac1{2\ell}$.
\emph{Hint: for the lower bound, let both provers color the cycle with a single monochromatic edge.  For the
upper bound, write the strategies as colorings $f_A,f_B$ of the vertices and build a graph on the $2\ell$ vertices
$(v,A),(v,B)$ with ``$\ne$'' edges $(u,A)\!-\!(v,B)$ for every oriented edge $(u,v)$ of the cycle, and ``$=$'' chords
$(v,A)\!-\!(v,B)$.  The ``$\ne$'' edges form a single even cycle of length $2\ell$, and the two endpoints of any chord are
at odd distance $\ell$ along it.  So each chord together with either arc of length~$\ell$ forms a cycle ($2\ell$ of them
in total) that no coloring can satisfy completely.  Every ``$\ne$'' edge has probability $\tfrac1{4\ell}$ and lies in $\ell$
of these cycles, every chord has probability $\tfrac1{2\ell}$ and lies in $2$ of them; give each cycle weight
$\tfrac1{4\ell^2}$.}
\end{exercise}

\begin{exercise}[\stars{$\star\star$}]\label{ex:fortnow}
\emph{Repetition can beat the product.}  Consider the following game $G$.  Alice and Bob receive independent
uniform bits $x$ and $y$ and answer with bits $a$ and $b$.  They win iff $(x,y)=(0,0)$ and $a\ne b$, or
$(x,y)\in\{(0,1),(1,0)\}$ and $a=b=0$.  (For $(x,y)=(1,1)$ they always lose.)
Show that $\val(G)=\tfrac12$ but $\val(G^{\otimes2})\ge\tfrac38>\val(G)^2=\tfrac14$.
\emph{Hint: in $G^{\otimes 2}$ let both players use the same strategy, answering the pair $(a_1,a_2)$ as a function of
their input pair: $00\mapsto00$, $01\mapsto10$, $10\mapsto01$, $11\mapsto00$.}
\end{exercise}

\begin{exercise}[\stars{$\star$}]\label{ex:not-ug}
(a) Prove that \sat{Max-3-Cut} is not a Unique Game.  (b) Prove that a Unique Game with $\val=1$ and a
Label Cover instance with $\val=1$ can differ in complexity by explaining, using results of this lecture,
why the first is in $\Ptime$ and the second is $\NP$-hard.
\end{exercise}

\begin{exercise}[\stars{$\star\star$}]\label{ex:quant}
Which of the following statements implies the UGC, which is implied by it, and which is neither?
(a) For every $\eps,\delta>0$ and \emph{every} $k\ge k_0(\eps,\delta)$, $\Gap{1-\eps}{\delta}{UG}_k$ is $\NP$-hard.
(b) There exist $\eps,\delta>0$ such that for all $k$, $\Gap{1-\eps}{\delta}{UG}_k$ is $\NP$-hard.
(c) For every $\eps$ there is $\delta>0$ and $k$ such that $\Gap{1-\eps}{\delta}{UG}_k$ is $\NP$-hard.
\end{exercise}

\begin{exercise}[\stars{$\star\star$}]\label{ex:reading}
\emph{Reading.}  Read the definition of Unique Games in Section~2 of \citet{Khot10} and the
statement of Theorem~1.1 in \citet{OpenAI26}.  List all differences in the formulation (type of graph,
type of constraint, alphabet, role of determinism, order of quantifiers).  For each difference, decide
whether it is a cosmetic restatement, a strengthening, or a weakening of the classical conjecture.
\end{exercise}


%% ===========================================================================
\section*{Notes and further reading}
\addcontentsline{toc}{section}{Notes and further reading}
%% ===========================================================================

\paragraph{Where this lecture comes from.}
The equivalence of the two PCP formulations is standard; see Lecture~1 of \citet{Harsha07}, or
Lecture~1 and Lecture~3 of \citet{Trevisan06}, and Chapter~11 of \citet{AroraBarak09}.  That the formulas produced by
Cook's theorem are ``fragile''---they are almost satisfiable even for NO instances, so a
gap cannot be created by such a reduction---is the starting point of \citeauthor{Trevisan06}'s first lecture;
the gap-problem formalism of Section~\ref{sec:gap} follows \citet{Harsha07}.  The DIMACS tutorial notes \citep{DIMACS10} are an excellent
companion for Lectures 1--9.  For the Unique Games Conjecture, the survey of \citet{Khot10} and the
Bulletin article of \citet{Trevisan12} are the primary references.

\paragraph{On attributions.}
The $n^{1-\eps}$ hardness of \sat{Max-Clique} is due to \citet{Hastad99}, and was proved under the
assumption $\NP\not\subseteq\mathsf{ZPP}$; the assumption was removed by \citet{Zuckerman07}.  It is a different
paper from the one containing the $\tfrac78$ and $\tfrac12$ results \citep{Hastad01}.  The first
examples of games whose parallel repetition does not behave like a product were due to Fortnow and to Feige;
the odd cycle game and its tilings of space are studied by \citet{FKO07} and \citet{Raz11}.

\paragraph{What is not covered.}
We have not defined the \emph{long code}, nor discussed the SDP relaxations of Unique Games, nor
the connection between the UGC and the \emph{Small-Set Expansion} problem.  They enter in later lectures.  Lecture~2 develops the formalism of games and Label Cover that we have only sketched.

\chapter{Games, Label Cover, and the Language of PCPs}
\label{lec:games}

{\sffamily\large\color{softgray}One object with four names, and the price of every translation\par}
\vspace{0.5em}

\noindent
Lecture~\ref{lec:why} introduced four characters: a PCP verifier, a constraint satisfaction
problem, a two-prover game, and a Label Cover instance.  Table~\ref{tab:dict} asserted that they are
one and the same object seen from four sides.  That is true, but only up to constants, and in hardness of
approximation constants are exactly what we cannot afford to lose track of: the difference between
soundness $1-\eps_0$ and soundness $1-\eps_0/3$ is invisible to the PCP theorem and decisive for a reduction
that must end with soundness $\tfrac12+\eps$.

This lecture turns every arrow of the dictionary into a theorem and records what it costs.  Three prices
will appear: a factor $2$ in value when we pass from one oracle to two provers (Section~\ref{sec:oracle}), a
factor $q$ in the gap when we pass from $q$-ary to binary constraints (Section~\ref{sec:cv}), and a blow-up
of the alphabet.  Once the prices are known, statements such as ``Label Cover with soundness $\delta$ over
an alphabet of size $\poly(1/\delta)$'' or ``$2$-to-$2$ Games over $\F_2^\ell$'' can be read at a glance.
We end with the oldest customer of this language, the reduction of
\citet{FGLSS96} from PCPs to \sat{Max-Clique}.  Its proof fits on one page and contains the idea behind
every later reduction in this course: \emph{turn the consistency of local views of a proof into the structure of a graph}.

\begin{goalbox}{After this lecture you should be able to\ldots}
\begin{enumerate}[leftmargin=*,label=(\arabic*)]
  \item write any PCP verifier, CSP, game or Label Cover instance as a constraint graph, and compute
        the value of a small one by hand (Sections~\ref{sec:lang}, \ref{sec:cv});
  \item say which constraint types---general, projection, $2$-to-$2$, unique---a given instance has, and
        what each type buys and costs (Section~\ref{sec:ladder});
  \item convert between a one-oracle and a two-prover view of the same instance, and state exactly how
        much value is lost (Proposition~\ref{prop:cover});
  \item prove that the clause--variable game of any $q$-ary CSP has soundness $1-\eps/q$ (Theorem~\ref{thm:cv});
  \item build the FGLSS graph of a verifier and prove that its clique number is $2^r$ times the best
        acceptance probability (Theorem~\ref{thm:fglss}).
\end{enumerate}
\end{goalbox}


%% ===========================================================================
\section{Constraint graphs: one language for everything}
\label{sec:lang}
%% ===========================================================================

We begin with the most general object and specialize.

\begin{definition}[Constraint satisfaction instance]\label{def:csp}
An instance $\Phi$ of a \defn{constraint satisfaction problem} consists of a finite set $X$ of
\emph{variables}, a finite alphabet $\Sigma_x$ for every $x\in X$, and a finite multiset $\mathcal C$ of
\emph{constraints}.  A constraint $c=(S_c,P_c)$ has a \emph{scope} $S_c\subseteq X$ and a
\emph{predicate} $P_c\colon\prod_{x\in S_c}\Sigma_x\to\{0,1\}$.  The \emph{arity} of $\Phi$ is $q=\max_c|S_c|$.
An \emph{assignment} $\sigma$ gives every variable $x$ a value $\sigma(x)\in\Sigma_x$, and
\[
  \val(\Phi)\;=\;\max_{\sigma}\ \Pr_{c\sim\mathcal C}\bigl[P_c(\sigma|_{S_c})=1\bigr].
\]
\end{definition}

\noindent
(Lecture~\ref{lec:why} wrote $\OPT$ for the optimum of a Max-CSP; from now on we write $\val$, which is also
the name for the value of a game.)  The problems of Lecture~\ref{lec:why} all fit.

\begin{table}[t]
\centering\small
\renewcommand{\arraystretch}{1.3}
\begin{tabular}{@{}>{\raggedright\arraybackslash}p{3.3cm}
                   >{\raggedright\arraybackslash}p{2.1cm}
                   c
                   >{\raggedright\arraybackslash}p{3.7cm}
                   >{\raggedright\arraybackslash}p{2.2cm}@{}}
\toprule
Problem & Alphabet & Arity & Predicate & Random assignment\\
\midrule
\sat{Max-E3-SAT} & $\{0,1\}$ & $3$ & OR of three literals & $\tfrac78$\\
\sat{Max-E3-Lin}$(\F_2)$ & $\{0,1\}$ & $3$ & $x+y+z=b$ & $\tfrac12$\\
\sat{Max-Cut} & $\{0,1\}$ & $2$ & $x\ne y$ & $\tfrac12$\\
\sat{Max-3-Cut} & $\{0,1,2\}$ & $2$ & $x\ne y$ & $\tfrac23$\\
Linear Unique Games & $\Z_k$ & $2$ & $x-y=c$ & $\tfrac1k$\\
Label Cover (projection) & $\Sigma_U,\Sigma_W$ & $2$ & $b=\pi(a)$ & $\tfrac1{|\Sigma_W|}$\\
$2$-to-$2$ Games & $[k]$ & $2$ & $\pi(a)=\pi'(b)$ & $\tfrac2k$\\
\bottomrule
\end{tabular}
\caption{A zoo of constraint satisfaction problems.  The last column is the fraction of constraints satisfied
by a \emph{uniformly random assignment}; it is the ``trivial'' threshold that approximation algorithms must beat.}
\label{tab:zoo}
\end{table}

The last column of Table~\ref{tab:zoo} is not a coincidence of the examples.  It is an expectation.

\begin{lemma}[Random assignment]\label{lem:random-asg}
Let the \defn{density} of a constraint $c$ be $\delta(c)=|P_c^{-1}(1)|\big/\prod_{x\in S_c}|\Sigma_x|$.  Then
$\val(\Phi)\ge\E_{c\sim\mathcal C}\,\delta(c)$.
\end{lemma}
\begin{proof}
For a uniformly random assignment $\sigma$ and a fixed constraint $c$, the restriction $\sigma|_{S_c}$ is uniform on
$\prod_{x\in S_c}\Sigma_x$, so $c$ is satisfied with probability $\delta(c)$.  By linearity of expectation
the expected fraction of satisfied constraints is $\E_c\,\delta(c)$, and some assignment does at least as well.
\end{proof}

\noindent
This is the baseline of the scoreboard in Table~\ref{tab:scoreboard}: the first two rows say that for
\sat{Max-E3-SAT} and \sat{Max-E3-Lin} nothing efficient beats the density of the predicate, unless $\Ptime=\NP$.

\subsection{Binary constraints: constraint graphs}

From here on the arity is almost always $2$.

\begin{definition}[Constraint graph]\label{def:cg}
A \defn{constraint graph} is a CSP instance of arity $2$.  Concretely it is a finite multigraph $G=(V,E)$,
alphabets $\Sigma_v$ for $v\in V$, and, for every edge, oriented arbitrarily as $e=(u,v)$, a
\defn{relation} $C_e\subseteq\Sigma_u\times\Sigma_v$.  A \emph{labeling} $\ell$ gives each vertex $v$ a label
$\ell(v)\in\Sigma_v$; it \emph{satisfies} $e=(u,v)$ if $(\ell(u),\ell(v))\in C_e$, and
\[
  \val(G)\;=\;\max_{\ell}\ \Pr_{e=(u,v)\sim E}\bigl[(\ell(u),\ell(v))\in C_e\bigr].
\]
If $G$ is bipartite, $G=(U,W,E)$ with every edge oriented from $U$ to $W$, we call it a \defn{game}; this is
Definition~\ref{def:game}, and $\val(G)$ is the value of the game.
\end{definition}

\noindent
When an edge $e=(u,v)$ carries a projection we write $\pi_{uv}\colon\Sigma_u\to\Sigma_v$ so that
$C_e=\{(a,\pi_{uv}(a))\}$, and the value reads
$\val(G)=\max_\ell\Pr_{(u,v)\sim E}[\pi_{uv}(\ell(u))=\ell(v)]$, the formula we shall use throughout.
Unique Games (Definition~\ref{def:ug}) are the case where every $\pi_{uv}$ is a bijection of $[k]$.

The word ``Label Cover'' is attached to this optimization problem---find a labeling of maximum value---and
usually, though not always, refers to projection constraints.  We shall say \emph{Label Cover} for the
general problem and specify the constraint type when it matters.

\begin{definition}[Types of constraints]\label{def:types}
Let $C\subseteq\Sigma\times\Sigma'$ be the relation on an edge $(u,v)$, with $\Sigma=\Sigma_u$ and $\Sigma'=\Sigma_v$.
\begin{itemize}
  \item $C$ is a \defn{projection} if $C=\{(a,\pi(a))\}$ for a function $\pi\colon\Sigma\to\Sigma'$: every left label has
        exactly one partner.  It is \defn{$d$-to-$1$} if moreover every fibre $\pi^{-1}(b)$ has exactly $d$ elements.
  \item $C$ is \defn{unique} if it is the graph of a bijection, \ie a perfect matching between $\Sigma$ and $\Sigma'$.
  \item $C$ is \defn{$2$-to-$2$} if there are a set $\Gamma$ and two $2$-to-$1$ maps $\pi\colon\Sigma\to\Gamma$,
        $\pi'\colon\Sigma'\to\Gamma$ with $C=\{(a,b):\pi(a)=\pi'(b)\}$.  Every label then has exactly two partners,
        and the bipartite graph of $C$ is a disjoint union of copies of $K_{2,2}$.
\end{itemize}
An instance \emph{is} of a given type if all its constraints are.
\end{definition}

\begin{figure}[t]
\centering
\begin{tikzpicture}[x=1cm,y=1cm,every node/.style={font=\small},
   la/.style={circle,fill=accent,inner sep=1.6pt}, lb/.style={circle,fill=warn,inner sep=1.6pt},
   ed/.style={thick,black!70}]
  % (a) arbitrary relation
  \begin{scope}
    \node[font=\sffamily\footnotesize\bfseries] at (0.95,0.75) {arbitrary};
    \foreach \i in {1,...,4} {\node[la] (pa\i) at (0,-\i*0.6+0.6) {}; \node[lb] (pb\i) at (1.9,-\i*0.6+0.6) {};}
    \draw[ed] (pa1)--(pb1) (pa1)--(pb2) (pa1)--(pb3) (pa2)--(pb2) (pa3)--(pb1) (pa3)--(pb4);
    \node[font=\footnotesize] at (0.95,-2.3) {general relation};
  \end{scope}
  % (b) projection 2-to-1
  \begin{scope}[xshift=4.0cm]
    \node[font=\sffamily\footnotesize\bfseries] at (0.95,0.75) {projection};
    \foreach \i in {1,...,4} \node[la] (qa\i) at (0,-\i*0.6+0.6) {};
    \node[lb] (qb1) at (1.9,-0.3) {}; \node[lb] (qb2) at (1.9,-1.5) {};
    \draw[ed] (qa1)--(qb1) (qa2)--(qb1) (qa3)--(qb2) (qa4)--(qb2);
    \node[font=\footnotesize] at (0.95,-2.3) {here $2$-to-$1$};
  \end{scope}
  % (c) 2-to-2: two matchings
  \begin{scope}[xshift=8.0cm]
    \node[font=\sffamily\footnotesize\bfseries] at (0.95,0.75) {2-to-2};
    \foreach \i in {1,...,4} {\node[la] (ra\i) at (0,-\i*0.6+0.6) {}; \node[lb] (rb\i) at (1.9,-\i*0.6+0.6) {};}
    \draw[accent,thick] (ra1)--(rb1) (ra2)--(rb2) (ra3)--(rb3) (ra4)--(rb4);
    \draw[warn,thick,dashed] (ra1)--(rb2) (ra2)--(rb1) (ra3)--(rb4) (ra4)--(rb3);
    \node[font=\footnotesize] at (0.95,-2.3) {two matchings in one};
  \end{scope}
  % (d) unique
  \begin{scope}[xshift=12.0cm]
    \node[font=\sffamily\footnotesize\bfseries] at (0.95,0.75) {unique};
    \foreach \i in {1,...,4} {\node[la] (sa\i) at (0,-\i*0.6+0.6) {}; \node[lb] (sb\i) at (1.9,-\i*0.6+0.6) {};}
    \draw[ed] (sa1)--(sb3) (sa2)--(sb1) (sa3)--(sb4) (sa4)--(sb2);
    \node[font=\footnotesize] at (0.95,-2.3) {perfect matching};
  \end{scope}
\end{tikzpicture}
\caption{The four constraint types of Definition~\ref{def:types} as bipartite graphs between $\Sigma_u$
(left, blue) and $\Sigma_v$ (right, orange).  In a unique constraint each label determines its partner
\emph{and} vice versa; in a projection, only one direction is determined; in a $2$-to-$2$ constraint each label
leaves two candidates.  The $2$-to-$2$ relation of panel~(c) is the union of the solid and the dashed
perfect matching (Lemma~\ref{lem:split}).}
\label{fig:types}
\end{figure}

Lemma~\ref{lem:random-asg} immediately quantifies how many labels a type forces on us.

\begin{corollary}[Alphabets cannot be small]\label{cor:alph}
For a projection game, $\val(G)\ge1/|\Sigma_W|$; for a Unique Game with $k$ labels, $\val(G)\ge1/k$;
for a $2$-to-$2$ game with $k$ labels on both sides, $\val(G)\ge 2/k$.  Hence a hardness result with soundness $\delta$ needs
alphabets of size at least $1/\delta$, $1/\delta$ and $2/\delta$, respectively.
\end{corollary}
\begin{proof}
A projection has $|C|=|\Sigma|$ pairs, hence density $|\Sigma|/(|\Sigma||\Sigma'|)=1/|\Sigma'|$; a $2$-to-$2$ constraint has
$|C|=2|\Sigma|$ pairs and, when $|\Sigma|=|\Sigma'|=k$, density $2/k$.  Apply Lemma~\ref{lem:random-asg}.
\end{proof}

\begin{example}[Unique Games on a cycle]\label{ex:cycle}
Let $G$ be a cycle on vertices $0,1,\dots,\ell-1$ with the linear constraints $x_{i+1}-x_i=c_i\pmod k$ on the edges
$(i,i+1)$ (indices modulo $\ell$).  If $\sum_ic_i\equiv0$, set $x_0=0$ and propagate: all $\ell$ equations hold.
Otherwise summing all the equations would give $0\equiv\sum_ic_i\not\equiv0$, so at least one equation fails; but
deleting one edge leaves a path, which can always be satisfied.  Hence $\val(G)=1-\tfrac1\ell$.
\emph{A cycle is satisfiable iff the constraints compose to the identity around it}; this is the whole
obstruction to perfect completeness in a Unique Game, and checking it along a spanning tree is what the algorithm
of Proposition~\ref{prop:ug-easy} does.
\end{example}


%% ===========================================================================
\section{One oracle or two provers?}
\label{sec:oracle}
%% ===========================================================================

A PCP verifier that makes two queries reads two positions of \emph{one} proof.  In a game the referee sends two
questions to \emph{two} provers who answer independently.  A constraint graph on a general graph is the first,
a game the second.  The two differ in one respect: a single proof is automatically consistent with itself, whereas
two provers may give contradictory answers to the same question.

\begin{definition}[Double cover]\label{def:cover}
Let $G=(V,E)$ be a constraint graph with relations $C_e$.  Its \defn{double cover} $G^{\times2}$ is the game with left
vertices $v_A$ (Alice's questions) and right vertices $v_B$ (Bob's questions), one pair $(v_A,v_B)$ for
every $v\in V$, both with alphabet $\Sigma_v$, and for every edge $e=(u,v)$ of $G$ the two edges
$(u_A,v_B)$ with relation $C_e$ and $(v_A,u_B)$ with the transposed relation $C_e^{\top}=\{(b,a):(a,b)\in C_e\}$.
In words: the referee picks an edge of $G$ and a random orientation of it, and asks the two endpoints to the two provers.
\end{definition}

\begin{proposition}[Value of the double cover]\label{prop:cover}
For every constraint graph $G$,
\[
  \val(G)\;\le\;\val(G^{\times2})\;\le\;2\,\val(G),
\]
and both inequalities can be tight.
\end{proposition}
\begin{proof}
For the first inequality let $\ell$ be an optimal labeling of $G$ and let Alice and Bob both use it:
$\ell_A(v_A)=\ell_B(v_B)=\ell(v)$.  Whatever orientation the referee picks, the two answers are the labels of the two endpoints, so
the acceptance probability is $\val(G)$.

For the second, let $(\ell_A,\ell_B)$ be an optimal strategy in $G^{\times2}$ with value $p$.  Define a labeling of
$G$ at random: for every vertex $v$, independently and uniformly pick a side $s(v)\in\{A,B\}$ and set
$\ell(v)=\ell_{s(v)}(v)$.  Fix an edge $e=(u,v)$.  With probability $\tfrac14$ we have $s(u)=A$ and $s(v)=B$;
then $\ell$ satisfies $e$ iff the edge $(u_A,v_B)$ of $G^{\times2}$ is satisfied by $(\ell_A,\ell_B)$.  With probability
$\tfrac14$ we have $s(v)=A$ and $s(u)=B$, and $\ell$ satisfies $e$ iff the edge $(v_A,u_B)$ is satisfied.  These two events are
disjoint, so
\[
  \Pr[\ell\text{ satisfies }e]\ \ge\ \tfrac14\bigl(\mathbf 1[(u_A,v_B)\text{ sat.}]+\mathbf 1[(v_A,u_B)\text{ sat.}]\bigr).
\]
Averaging over a uniform edge $e$ of $G$, the right-hand side has expectation $\tfrac14\cdot2p=p/2$, since the $2|E|$ edges of
$G^{\times2}$ are the two orientations of the edges of $G$.  Hence some labeling has value at least $p/2$, \ie
$\val(G)\ge\val(G^{\times2})/2$.

For tightness of the first inequality take any satisfiable $G$.  For the second take \sat{Max-Cut} on the complete graph $K_n$:
the best cut contains a $\tfrac{\lceil n/2\rceil\lfloor n/2\rfloor}{\binom n2}\to\tfrac12$ fraction of the edges, but in the double
cover Alice answers $0$ everywhere and Bob $1$ everywhere, and all constraints $a\ne b$ hold (Figure~\ref{fig:cover}).
\end{proof}

\begin{figure}[t]
\centering
\begin{tikzpicture}[x=1cm,y=1cm,every node/.style={font=\small},
   al/.style={circle,draw,fill=accent!22,minimum size=0.75cm,inner sep=0pt},
   bo/.style={circle,draw,fill=warn!22,minimum size=0.75cm,inner sep=0pt},
   pl/.style={circle,draw,fill=black!6,minimum size=0.7cm,inner sep=0pt}]
  % triangle
  \node[pl] (u) at (0,1.5) {$u$};
  \node[pl] (v) at (-1.3,-0.5) {$v$};
  \node[pl] (w) at (1.3,-0.5) {$w$};
  \draw[thick] (u)--(v)--(w)--(u);
  \node[align=center,font=\footnotesize] at (0,-1.7) {\sat{Max-Cut} on a triangle\\ $\val=\tfrac23$};
  \draw[-{Stealth[length=2.2mm]},thick] (2.5,0.5)--(4.1,0.5) node[midway,above,font=\footnotesize] {double cover};
  % hexagon
  \begin{scope}[xshift=7.4cm,yshift=0.35cm]
    \node[al] (uA) at (90:1.55) {$u_A$};
    \node[bo] (vB) at (30:1.55) {$v_B$};
    \node[al] (wA) at (-30:1.55) {$w_A$};
    \node[bo] (uB) at (-90:1.55) {$u_B$};
    \node[al] (vA) at (-150:1.55) {$v_A$};
    \node[bo] (wB) at (150:1.55) {$w_B$};
    \draw[thick] (uA)--(vB)--(wA)--(uB)--(vA)--(wB)--(uA);
    \node[align=center,font=\footnotesize] at (0,-2.05) {the induced game\\ $\val=1$: Alice says $0$, Bob says $1$};
  \end{scope}
\end{tikzpicture}
\caption{The double cover of a triangle is a hexagon, which is bipartite.  Self-consistency of a single labeling is what
the triangle tests; the two provers of the hexagon need not be consistent with each other, and they win with probability~$1$.}
\label{fig:cover}
\end{figure}

\begin{remark}
The factor~$2$ is harmless when the soundness we care about is small: a hardness result for constraint graphs with completeness $c$ and soundness $s$
gives one for games with completeness $c$ and soundness $2s$.  It is fatal when $s\ge\tfrac12$.  This is why the PCP theorem is converted into a game by the
clause--variable construction of Section~\ref{sec:cv} (soundness $1-\eps/q$) and \emph{not} by the double cover
(soundness $2(1-\eps)$, which exceeds $1$ and says nothing as soon as $\eps<\tfrac12$).  For fixed orientations the loss is a factor $4$ instead of~$2$, which is the
bound quoted in \citet{Khot10}.
\end{remark}

A game is conversely a perfectly good PCP.

\begin{lemma}[Answer tables]\label{lem:tables}
Let $G=(U,W,E)$ be a game with alphabets $\Sigma_U,\Sigma_W$ of size at most $\sigma$.  Consider the verifier $V_G$ whose proof is a
pair of tables $A\colon U\to\Sigma_U$ and $B\colon W\to\Sigma_W$, that picks a uniformly random edge $(u,w)$, reads $A(u)$ and $B(w)$, and accepts
iff $(A(u),B(w))\in C_{(u,w)}$.  Then $\max_{A,B}\Pr[V_G\text{ accepts}]=\val(G)$.  If proofs must be binary strings, each label
is written with $\lceil\log_2\sigma\rceil$ bits, and the verifier makes $2\lceil\log_2\sigma\rceil$ binary queries, uses $\log_2|E|$ random bits
and has the same completeness and soundness.
\end{lemma}
\begin{proof}
A proof $(A,B)$ is exactly a labeling of $G$, and the acceptance probability is the fraction of satisfied edges.
\end{proof}

\noindent
So $\Gap{1}{\delta}{LabelCover}_\sigma$ is a PCP with $O(\log\sigma)$ queries and soundness~$\delta$.  Is that better than what we
could already do?  By Lemma~\ref{lem:tables} and Theorem~\ref{thm:lc} the parameters are $O(\log(1/\delta))$ queries and
soundness~$\delta$ with $O(\log n)$ random bits---and independent repetition of the plain verifier of Proposition~\ref{prop:pcp-equiv}
achieves the very same trade-off.  The point of games is not their parameters but their \emph{structure}: the $2$ queries
come in two groups, and the answer in one group is a function of the answer in the other.  This is the structure that the long-code
tests of Lectures~7--11 exploit.

\begin{proposition}[Two binary queries are not enough]\label{prop:2sat}
If a language $L$ has a PCP verifier with $r=O(\log n)$ random bits, \emph{two} binary queries, perfect completeness and soundness $s<1$,
then $L\in\Ptime$.  Hence, unless $\Ptime=\NP$, a two-query PCP for $\NP$ with perfect completeness needs alphabets of growing size.
\end{proposition}
\begin{proof}
On input $x$ enumerate the $2^r=\poly(n)$ random strings $\rho$.  Each defines a predicate $P_\rho$ on two Boolean proof bits,
and every Boolean predicate on two variables is a $2$-CNF (forbid each rejected pair of values by one clause of width $2$).
So the verifier defines a $2$-CNF $\psi_x$ of polynomial size.  If $x\in L$ the honest proof satisfies $\psi_x$; if $x\notin L$ every
proof is rejected for some $\rho$, so $\psi_x$ is unsatisfiable.  $2$SAT is in $\Ptime$.
\end{proof}


%% ===========================================================================
\section{From verifiers and CSPs to games}
\label{sec:cv}
%% ===========================================================================

Lecture~\ref{lec:why} turned the clauses of a $3$-CNF into a game by asking one prover about a clause and the other
about a variable of it.  The construction works for every CSP, and its analysis is the same one-paragraph argument.

\begin{definition}[Constraint--variable game]\label{def:cvgame}
Let $\Phi$ be a CSP instance of arity $q$.  The \defn{constraint--variable game} $G_\Phi$ is the following game.  The referee draws a
uniformly random constraint $c$ and a uniformly random variable $x\in S_c$, and sends $c$ to Alice and $x$ to Bob.  Alice
answers an assignment $a\in\prod_{y\in S_c}\Sigma_y$ to the scope of $c$, Bob answers a symbol $b\in\Sigma_x$, and they win iff
$P_c(a)=1$ and $a_x=b$.
\end{definition}

\noindent
(The edge $(c,x)$ has weight $\tfrac1{|\mathcal C||S_c|}$; weights are harmless since we may repeat edges.  If every predicate is satisfiable we restrict
Alice's alphabet to the satisfying assignments of $c$, which changes nothing since a non-satisfying answer loses outright.  Then $G_\Phi$ is a \emph{projection game}:
$\pi_{(c,x)}(a)=a_x$, with $|\Sigma_U|\le\max_x|\Sigma_x|^q$ and $\Sigma_W=\Sigma_x$.)

\begin{theorem}[Soundness of the constraint--variable game]\label{thm:cv}
For every CSP instance $\Phi$ of arity $q$,
\[
  \val(\Phi)\ \le\ \val(G_\Phi)\ \le\ 1-\frac{1-\val(\Phi)}{q}.
\]
In particular $\val(\Phi)=1$ iff $\val(G_\Phi)=1$, and if every assignment violates an $\eps$-fraction of the constraints, then
$\val(G_\Phi)\le1-\eps/q$.
\end{theorem}
\begin{proof}
\emph{Lower bound.}  Let $\sigma$ be an optimal assignment of $\Phi$.  Bob answers $\sigma(x)$; Alice, for a constraint $c$ that $\sigma$
satisfies, answers $\sigma|_{S_c}$, and for a constraint it violates, anything.  On every constraint satisfied by $\sigma$ the provers win
for every $x\in S_c$, hence the winning probability is at least $\val(\Phi)$.

\emph{Upper bound.}  Fix deterministic strategies.  Bob's answers define an assignment $b$ of $\Phi$ with $b(x)=\ell_W(x)$.
Let $c$ be a constraint violated by $b$, and let $a=\ell_U(c)$ be Alice's answer.  If $P_c(a)=0$ the provers lose on every
edge at~$c$.  Otherwise $a\ne b|_{S_c}$ (because $P_c(b|_{S_c})=0$), so there is some $x\in S_c$ with $a_x\ne b(x)$, and the referee
asks that $x$ with probability at least $1/|S_c|\ge1/q$.  So the provers lose at $c$ with probability at least $1/q$.  Since $b$ violates at
least a $1-\val(\Phi)$ fraction of the constraints, the total losing probability is at least $(1-\val(\Phi))/q$.
\end{proof}

\noindent
This is Proposition~\ref{prop:cv} for $q=3$ and a CNF.  The bound is attained, \eg for the formula below.

\begin{example}[A formula and its game]\label{ex:phi0}
Let $\varphi_0$ have the four clauses $C_1=(x_1\vee x_2)$, $C_2=(\neg x_1\vee x_2)$, $C_3=(\neg x_2\vee x_3)$, $C_4=(\neg x_2\vee\neg x_3)$.
If $x_2=1$ then $C_1,C_2$ hold and exactly one of $C_3,C_4$ does; if $x_2=0$ then $C_3,C_4$ hold and exactly one of $C_1,C_2$ does.
So \emph{every} assignment satisfies exactly three of the four clauses: $\val(\varphi_0)=\tfrac34$.  Theorem~\ref{thm:cv} with $q=2$ gives
$\val(G_{\varphi_0})\le1-\tfrac{1/4}{2}=\tfrac78$.  This is tight: let Bob answer any assignment $b$, say $b=(0,0,0)$, which violates $C_1$ only; Alice answers $b|_{S_c}$ on $C_2,C_3,C_4$ and
$(x_1,x_2)=(1,0)$ on $C_1$, a satisfying assignment that still agrees with $b$ on $x_2$.  The provers lose only if the referee picks $C_1$ and then the variable $x_1$, an event of
probability $\tfrac14\cdot\tfrac12=\tfrac18$.
\end{example}

\begin{corollary}[Every verifier is a game]\label{cor:verifier-game}
Let $V$ be a PCP verifier with $r$ random bits and $q$ queries, perfect completeness and soundness~$s$ on a language~$L$.  Let $\Phi_V(x)$ be the CSP with one
constraint per random string~$\rho$ (scope: the positions queried at~$\rho$, predicate: $P_{x,\rho}$).  Then the game
$G_{\Phi_V(x)}$, in which Alice is asked a random string $\rho$ and Bob a position queried at $\rho$, is a projection game with $|\Sigma_U|\le2^q$, $|\Sigma_W|=2$, and
\[
  x\in L\ \Rightarrow\ \val=1,\qquad x\notin L\ \Rightarrow\ \val\le1-\frac{1-s}{q}.
\]
\end{corollary}
\begin{proof}
$\val(\Phi_V(x))$ is the best acceptance probability over proofs, which is $1$ or at most $s$; apply Theorem~\ref{thm:cv}.
\end{proof}

\noindent
Combined with $\NP=\PCP[O(\log n),O(1)]$, Corollary~\ref{cor:verifier-game} says that it is $\NP$-hard to distinguish projection games of value $1$ from games of value at most $1-\eps_1$, for a constant $\eps_1>0$:
the PCP theorem in the language of games, stated without mentioning $3$SAT.

Table~\ref{tab:price} collects the translations between the four columns of the dictionary and what each costs.

\begin{table}[t]
\centering\small
\renewcommand{\arraystretch}{1.3}
\begin{tabular}{@{}>{\raggedright\arraybackslash}p{3.8cm}
                   >{\raggedright\arraybackslash}p{4.1cm}
                   >{\raggedright\arraybackslash}p{4.6cm}@{}}
\toprule
Translation & Effect on the value & Price\\
\midrule
verifier $\to$ CSP & $\val(\Phi_V)=\max_\pi\Pr[V^\pi\text{ acc.}]$ exactly & none; $2^r$ constraints\\
CSP $\to$ verifier & the same, exactly & none (pick a random constraint)\\
constraint graph $\to$ game (double cover, Prop.~\ref{prop:cover}) & $\val(G)\le\val(G^{\times2})\le2\val(G)$ & soundness at most doubles\\
game $\to$ verifier (answer tables, Lemma~\ref{lem:tables}) & exactly the value of the game & $2\lceil\log\sigma\rceil$ binary queries\\
$q$-ary CSP $\to$ projection game (Thm.~\ref{thm:cv}) & $1-\frac{1-\val(\Phi)}q\ge\val(G_\Phi)\ge\val(\Phi)$ & gap shrinks by a factor $q$; alphabet $|\Sigma|^q$\\
verifier $\to$ clique (Thm.~\ref{thm:fglss}) & $\omega(H_x)=2^r\max_\pi\Pr[V^\pi\text{ acc.}]$ & $2^{r+q}$ vertices\\
\bottomrule
\end{tabular}
\caption{The price of every translation in the dictionary of Table~\ref{tab:dict}.}
\label{tab:price}
\end{table}


%% ===========================================================================
\section{The ladder of constraint types}
\label{sec:ladder}
%% ===========================================================================

Why are we so careful about the type of a constraint?  Because it determines both what an algorithm can do and what a hardness reduction needs.
We saw in Proposition~\ref{prop:ug-easy} that the \emph{unique} type makes perfect completeness easy: a label on one vertex forces the label on its neighbours in both
directions, so a satisfying labeling is found by propagation along a spanning tree.  For a projection the forcing goes in one direction only, and for a general relation in neither.
A $2$-to-$2$ constraint sits exactly between: after one step of propagation \emph{two} candidates remain, after $t$ steps up to $2^t$.

This is the right way to read the hierarchy of conjectures and theorems that organizes the rest of the course.  Every row of Table~\ref{tab:ladder} is a
statement about the complexity of a gap problem for constraint graphs of this type.

\begin{table}[t]
\centering\small
\renewcommand{\arraystretch}{1.3}
\begin{tabular}{@{}>{\raggedright\arraybackslash}p{2.6cm}
                   >{\raggedright\arraybackslash}p{3.6cm}
                   >{\raggedright\arraybackslash}p{3.3cm}
                   >{\raggedright\arraybackslash}p{2.9cm}@{}}
\toprule
Constraints & Gap problem & Status & Main use\\
\midrule
projection (large alphabet) & $\Gap{1}{\delta}{LabelCover}$, every $\delta>0$ & theorem (Thm.~\ref{thm:lc}) & H\aa stad's tight $3$-bit results (Lecture~\ref{lec:hastad})\\
$2$-to-$2$ & $\Gap{1-\eps}{\eps}{2-to-2}_k$ for some $k=k(\eps)$ & theorem \citep{KMS17,DKKMS18,KMS18} & \sat{Vertex Cover} within $\sqrt2-\eps$ (Lectures~13--15)\\
unique & $\Gap{1-\eps}{\delta}{UG}_k$ for some $k=k(\eps,\delta)$ & conjecture \citep{Khot02}; claimed in \citep{OpenAI26}, not verified & tight bounds for \sat{Max-Cut}, \sat{Vertex Cover}, every CSP (Lectures~8--11)\\
unique, completeness $1$ & $\Gap{1}{\delta}{UG}$ & in $\Ptime$ (Prop.~\ref{prop:ug-easy}) & ---\\
\bottomrule
\end{tabular}
\caption{The ladder of constraint types.  Moving down the table the constraints become more rigid and the hardness
results become stronger \emph{per label}, but harder to prove.}
\label{tab:ladder}
\end{table}

The $2$-to-$2$ theorem is called ``half-way'' to the UGC for a reason that can be proved in a few lines, and the proof is a good test of the language we have built.

\begin{lemma}[A $2$-to-$2$ relation is two bijections]\label{lem:split}
Let $C\subseteq\Sigma\times\Sigma'$ be a relation in which every label of $\Sigma$ and every label of $\Sigma'$ has exactly two partners; in particular
$C$ may be a $2$-to-$2$ constraint in the sense of Definition~\ref{def:types}.  Then $C$ is the disjoint union $M\sqcup M'$ of two perfect matchings, \ie of the graphs of two bijections.
\end{lemma}
\begin{proof}
The bipartite graph $(\Sigma\sqcup\Sigma',C)$ is simple and $2$-regular, hence a disjoint union of cycles, each of even length at least~$4$.  Colour the edges of each cycle alternately red and blue.
Every vertex lies on exactly one red and one blue edge, so the red edges form a perfect matching $M$ and the blue edges another one, $M'$.
\end{proof}

\begin{corollary}[Splitting a $2$-to-$2$ game]\label{cor:half}
Let $G$ be a constraint graph whose relations are all as in Lemma~\ref{lem:split}.  Let $\mathcal U_G$ be the Unique Games instance on the same vertices
in which every edge $e$ of $G$ is replaced by two parallel edges, one carrying each of the two bijections of $C_e$.  Then for every labeling $\ell$, the fraction of edges of $\mathcal U_G$
satisfied by $\ell$ is exactly \emph{half} the fraction of edges of $G$ satisfied by $\ell$.  Consequently $\val(\mathcal U_G)=\tfrac12\val(G)$.
\end{corollary}
\begin{proof}
A labeling $\ell$ satisfies $e=(u,v)$ in $G$ iff $(\ell(u),\ell(v))\in C_e=M_e\sqcup M'_e$, \ie iff it satisfies exactly one of the two parallel edges of $\mathcal U_G$ (the matchings
are disjoint, so never both); otherwise it satisfies neither.  So $\ell$ satisfies as many edges of $\mathcal U_G$ as of $G$, while $\mathcal U_G$ has twice as many edges.
\end{proof}

\noindent
Parallel edges are harmless: they only change the weights.  Hence the $2$-to-$2$ theorem gives the following unconditional
statement, which differs from the Unique Games Conjecture only in the \emph{completeness}.

\begin{corollary}[A half-way Unique Games Conjecture, from Table~\ref{tab:ladder}]\label{cor:halfway}
For every $\eps>0$ there is a $k=k(\eps)$ such that it is $\NP$-hard to distinguish Unique Games instances with $k$ labels and $\val\ge\tfrac{1-\eps}2$ from those with $\val\le\tfrac\eps2$.
\end{corollary}

\noindent
This is the sense in which the $2$-to-$2$ theorem proves the UGC ``half-way'' \citep{Khot19}: completeness $\approx\tfrac12$ instead of $1-\eps$.  Khot remarks that the arguments usually cited as evidence against the UGC apply to this
version too.  The claim of \citet{OpenAI26} is a claim to supply the missing half; we treat it as unverified (Lecture~16).

\begin{remark}[Max-3-Cut is $2$-regular]
The relation $x\ne y$ on three labels has exactly two partners for every label, so Lemma~\ref{lem:split} applies: it is the union of the two translations $y=x+1$ and $y=x+2$ modulo~$3$.
(It is not a union of $K_{2,2}$'s, since the bipartite graph is a $6$-cycle, so it is not $2$-to-$2$ in the sense of Definition~\ref{def:types}; see Exercise~\ref{ex:3cut}.)
\end{remark}


%% ===========================================================================
\section{From PCPs to cliques: the FGLSS graph}
\label{sec:fglss}
%% ===========================================================================

All of the above treated the proof as an oracle to be queried.  The reduction of \citet{FGLSS96} does something different: it turns the \emph{set of everything the verifier might see} into a
graph, and lets consistency do the work.  We fix a verifier $V$ as in Definition~\ref{def:pcp}, and an input $x$.  For $\rho\in\{0,1\}^r$ write $I_\rho=(i_1(\rho),\dots,i_q(\rho))$ for the positions it queries (possibly with repetitions).

\begin{definition}[FGLSS graph]\label{def:fglss}
A \defn{local view} is a pair $(\rho,a)$ with $\rho\in\{0,1\}^r$ and $a\in\{0,1\}^q$ such that $P_{x,\rho}(a)=1$ and $a_j=a_{j'}$ whenever $i_j(\rho)=i_{j'}(\rho)$.
Two local views $(\rho,a)$ and $(\rho',a')$ are \defn{compatible} if $\rho\neq\rho'$ and $a_j=a'_{j'}$ whenever $i_j(\rho)=i_{j'}(\rho')$.  The \defn{FGLSS graph} $H_x$ has the local views as vertices and an edge between
any two compatible views.  It has at most $2^{r+q}$ vertices.
\end{definition}

In words: a vertex says ``on random string $\rho$, the proof contained the accepting answers $a$''; two vertices are adjacent when they could both occur in one and the same proof.

\begin{theorem}[FGLSS]\label{thm:fglss}
$\displaystyle\omega(H_x)\;=\;2^r\cdot\max_{\pi\in\{0,1\}^N}\Pr_\rho\bigl[V^\pi(x)\text{ accepts}\bigr].$
\end{theorem}
\begin{proof}
($\ge$) Fix a proof $\pi$ and let $R=\{\rho:V^\pi(x)\text{ accepts on }\rho\}$.  For $\rho\in R$ the pair $(\rho,\pi|_{I_\rho})$ is a local view (accepting, and
consistent since it is read from the single string $\pi$).  Any two of these views agree with $\pi$ and have different $\rho$, so they are compatible: a clique of size $|R|=2^r\Pr_\rho[V^\pi\text{ accepts}]$.

($\le$) Let $K$ be a clique.  Adjacent vertices have different $\rho$, so $K$ contains at most one view per~$\rho$.  Define a proof $\pi$ by setting $\pi_i=a_j$ whenever $(\rho,a)\in K$ and $i_j(\rho)=i$, and $\pi_i=0$ for unconstrained positions.
This is well defined: two occurrences of position $i$ within one view agree by the consistency condition of a local view, and two occurrences in different views agree by compatibility.
Now for every $(\rho,a)\in K$ we have $\pi|_{I_\rho}=a$, so $P_{x,\rho}(\pi|_{I_\rho})=P_{x,\rho}(a)=1$: $V^\pi$ accepts on~$\rho$.  Hence $V^\pi$ accepts on at least $|K|$ random strings.
\end{proof}

\begin{figure}[t]
\centering
\begin{tikzpicture}[x=1cm,y=1cm,every node/.style={font=\small},
   vtx/.style={circle,draw=black!70,minimum size=0.78cm,inner sep=0pt,font=\footnotesize\ttfamily}]
  \node[vtx,fill=accent!22] (n001) at (111.00:3.00) {01};
  \node[vtx,fill=accent!22] (n010) at (135.00:3.00) {10};
  \node[vtx,fill=accent!22] (n011) at (159.00:3.00) {11};
  \node[vtx,fill=warn!22] (n100) at (21.00:3.00) {00};
  \node[vtx,fill=warn!22] (n101) at (45.00:3.00) {01};
  \node[vtx,fill=warn!22] (n111) at (69.00:3.00) {11};
  \node[vtx,fill=claim!18] (n200) at (-69.00:3.00) {00};
  \node[vtx,fill=claim!18] (n201) at (-45.00:3.00) {01};
  \node[vtx,fill=claim!18] (n211) at (-21.00:3.00) {11};
  \node[vtx,fill=black!14] (n300) at (-159.00:3.00) {00};
  \node[vtx,fill=black!14] (n301) at (-135.00:3.00) {01};
  \node[vtx,fill=black!14] (n310) at (-111.00:3.00) {10};
  \begin{scope}[on background layer]
    \draw[black!35] (n001)--(n101);
    \draw[black!35] (n001)--(n211);
    \draw[black!35] (n001)--(n310);
    \draw[black!35] (n010)--(n200);
    \draw[black!35] (n010)--(n201);
    \draw[black!35] (n010)--(n300);
    \draw[black!35] (n010)--(n301);
    \draw[black!35] (n011)--(n211);
    \draw[black!35] (n100)--(n200);
    \draw[black!35] (n100)--(n201);
    \draw[black!35] (n100)--(n300);
    \draw[black!35] (n100)--(n301);
    \draw[black!35] (n101)--(n211);
    \draw[black!35] (n101)--(n310);
    \draw[black!35] (n111)--(n211);
    \draw[black!35] (n200)--(n300);
    \draw[black!35] (n201)--(n301);
  \end{scope}
  \draw[accent!85!black,line width=1.8pt] (n011)--(n111);
  \draw[accent!85!black,line width=1.8pt] (n011)--(n310);
  \draw[accent!85!black,line width=1.8pt] (n111)--(n310);
  \node[vtx,fill=accent!22,draw=accent!85!black,line width=1.4pt] at (n011) {11};
  \node[vtx,fill=warn!22,draw=accent!85!black,line width=1.4pt] at (n111) {11};
  \node[vtx,fill=black!14,draw=accent!85!black,line width=1.4pt] at (n310) {10};
  \node[align=center,font=\footnotesize] at (135:4.55) {$\rho=1$: $x_1\vee x_2$\\ answers $x_1x_2$};
  \node[align=center,font=\footnotesize] at (45:4.55)  {$\rho=2$: $\neg x_1\vee x_2$\\ answers $x_1x_2$};
  \node[align=center,font=\footnotesize] at (-45:4.55) {$\rho=3$: $\neg x_2\vee x_3$\\ answers $x_2x_3$};
  \node[align=center,font=\footnotesize] at (-135:4.55){$\rho=4$: $\neg x_2\vee\neg x_3$\\ answers $x_2x_3$};
\end{tikzpicture}
\caption{The FGLSS graph of the verifier for $\varphi_0$ (Example~\ref{ex:phi0}) that picks one of the four clauses and reads its two variables.
Each group of three vertices contains the three satisfying answers of its clause; edges join compatible views.  The heavy triangle is a maximum clique: it consists of the views
of $\rho=1,2,4$ that agree with the assignment $x=(1,1,0)$, which satisfies $C_1,C_2,C_4$ and violates $C_3$.
No view of $\rho=3$ is compatible with all three.}
\label{fig:fglss}
\end{figure}

\begin{example}[The FGLSS graph of $\varphi_0$]\label{ex:fglss}
Take $V$ to be the verifier that expects an assignment $\pi\in\{0,1\}^3$ of $\varphi_0$ (Example~\ref{ex:phi0}), picks one of the four clauses ($r=2$) and reads its two variables ($q=2$).
Each clause has three satisfying answers, so $H_x$ has $4\cdot3=12$ vertices and (by direct count) $20$ edges; Figure~\ref{fig:fglss}.  Every assignment satisfies exactly three of
the four clauses, so $\max_\pi\Pr[\text{accept}]=\tfrac34$ and Theorem~\ref{thm:fglss} predicts $\omega(H_x)=4\cdot\tfrac34=3$, as one finds by inspection (there are $8$ maximum cliques, one for each
assignment).
\end{example}

\begin{pitfall}[unbreakable]{a clique is a consistent set, not just an accepting one}
Two vertices of $H_x$ may both be accepting views and still be non-adjacent: they disagree about one position of the proof.  A clique does not say that many random strings accept
\emph{somewhere}; it says that they accept \emph{the same proof}.  This is the same idea as the difference between the two provers of a game, who need not be consistent, and a single proof, which is:
Theorem~\ref{thm:fglss} is the one-oracle statement, and it is exact; no factor $2$ appears.
\end{pitfall}

\begin{corollary}[Constant-factor hardness for Clique]\label{cor:clique-const}
Unless $\Ptime=\NP$, no polynomial-time algorithm outputs, for every graph, a clique of more than half the maximum size.
\end{corollary}
\begin{proof}
By Theorem~\ref{thm:pcp1} there is a verifier for $\sat{3SAT}$ with $r=O(\log n)$, $q=O(1)$, perfect completeness and soundness $\tfrac12$.  Build $H_x$ in time polynomial in $n$ ($2^{r+q}=\poly(n)$ vertices).
If $x\in\sat{3SAT}$ then $\omega(H_x)=2^r$; otherwise $\omega(H_x)\le2^{r-1}$.  By Lemma~\ref{lem:gap-approx} no polynomial-time algorithm achieves ratio $\alpha>\tfrac12$ unless $\Ptime=\NP$.
\end{proof}

To pass from the constant $2$ to larger factors we amplify the gap.  Here graph products do exactly what independent repetition does for proofs.  The \defn{strong product} $G\boxtimes H$ has vertex set $V(G)\times V(H)$, and
two distinct vertices $(g,h)$, $(g',h')$ are adjacent iff ($g=g'$ or $g\sim g'$) and ($h=h'$ or $h\sim h'$).

\begin{proposition}[Clique number is multiplicative]\label{prop:product}
$\omega(G\boxtimes H)=\omega(G)\,\omega(H)$.
\end{proposition}
\begin{proof}
If $A\subseteq V(G)$ and $B\subseteq V(H)$ are cliques then $A\times B$ is a clique, giving $\ge$.  Conversely let $K$ be a clique of $G\boxtimes H$ and let $A$ and $B$ be its
projections to $V(G)$ and $V(H)$.  Two distinct elements of $A$ come from two vertices of $K$ that are adjacent, and since their first coordinates differ they must be adjacent in $G$; so $A$ is a clique, and likewise $B$.
Since $K\subseteq A\times B$ we get $|K|\le|A||B|\le\omega(G)\omega(H)$.
\end{proof}

\begin{corollary}[Hardness of Clique at larger factors]\label{cor:clique-quasi}
(a) For every constant $c$, it is $\NP$-hard to find a clique of size at least $1/c$ times the maximum.
(b) For every $\eps>0$, if on $N$-vertex graphs one can find in polynomial time a clique of size at least $2^{-(\log N)^{1-\eps}}$ times the maximum, then $\NP\subseteq\mathsf{DTIME}(2^{\polylog n})$.
\end{corollary}
\begin{proof}
Let $H=H_x$ be the graph of Corollary~\ref{cor:clique-const}, with $M=\poly(n)$ vertices, say $M\le n^{a}$.  The $k$-fold power $H^{\boxtimes k}$ has $M^k$ vertices and clique number $2^{rk}$ (YES) or at most $2^{(r-1)k}$ (NO), by Proposition~\ref{prop:product}: a gap of $2^k$.
For (a) take $k=\lceil\log_2c\rceil+1$, a constant.  For (b) take $k=\lceil(\log n)^{1/\eps}\rceil$: the graph has $N=n^{O(k)}$ vertices, is built in time $\poly(N)=2^{O((\log n)^{1+1/\eps})}$, and since
$\log N=\Theta(k\log n)$ the gap is $2^k=2^{\Theta((\log N)^{1/(1+\eps)})}$, which exceeds $2^{(\log N)^{1-\eps}}$ for large $n$ because $\tfrac1{1+\eps}>1-\eps$.  An algorithm as in (b) would therefore decide $\sat{3SAT}$ in quasi-polynomial time.
\end{proof}

\noindent
This is the shape of the argument of \citet{FGLSS96}, who had a PCP for $\NP$ with randomness and queries $O(\log n\log\log n)$ rather than $O(\log n)$ and $O(1)$; Theorem~\ref{thm:fglss}
applied to their proof system gives graphs with $n^{O(\log\log n)}$ vertices, and hence hardness under the assumption $\NP\not\subseteq\mathsf{DTIME}(n^{O(\log\log n)})$.  The PCP theorem removes the loss.

\begin{bigpicture}{Where the polynomial factors come from}
Corollary~\ref{cor:clique-quasi} reaches $2^{(\log N)^{1-\eps}}$ but not $N^{\Omega(1)}$, because $k$ independent repetitions cost $rk$ random bits, and $2^{rk}$ vertices is quasi-polynomial once $k$ grows with $n$.  The remedy is
\emph{randomness-efficient repetition}: replace the $k$ independent random strings by a $k$-step random walk on a constant-degree expander, which needs only $r+O(k)$ random bits and still drives the soundness down
exponentially in $k$ (Exercise~\ref{ex:expander}).  With $k=\Theta(\log n)$ one gets graphs of polynomial size and gap $N^{\Omega(1)}$, which is the conclusion of Theorem~6 of \citet{ALMSS98} (we do not claim that their proof proceeds this way).
The exponent was improved by the \emph{free bit} analysis of \citet{BGS98} and finally to $N^{1-\eps}$ by \citet{Hastad99}; the latter proof uses Label Cover, parallel repetition and the long code, the toolkit of Lectures~5 and~7.
\end{bigpicture}


%% ===========================================================================
\section{What Label Cover is for}
\label{sec:for}
%% ===========================================================================

The FGLSS reduction reads a PCP and builds a graph.  For problems with a rigid local structure---\sat{Max-E3-SAT}, \sat{Max-E3-Lin}, and later \sat{Max-Cut}---a different recipe is needed, and Label Cover is its first ingredient.

\begin{bigpicture}{The recipe of the next lectures}
\textbf{Step 1 (Lectures 1--2, 5).}  Prove that a Label Cover instance (projection constraints, soundness $\delta$, alphabet $\poly(1/\delta)$) is $\NP$-hard to approximate.
This is the \emph{outer verifier}.
\textbf{Step 2 (Lecture~\ref{lec:hastad}).}  Replace every vertex $v$ of the instance by a table $f_v$ that is supposed to be the \emph{long-code} encoding of the label of~$v$, and every edge $e=(u,v)$ by a
\emph{local test} between $f_u$ and $f_v$ that reads very few entries; this is the \emph{inner verifier}.  The tests have the shape of the target problem: three bits for \sat{Max-E3-Lin}, two bits for \sat{Max-Cut}.
\textbf{Step 3 (analysis).}  If the Label Cover instance has a labeling of value~$1$, encoding it passes all tests (completeness).  If a family of tables passes the tests with
probability above the claimed soundness, \emph{decode} each table into a short list of candidate labels; a random choice from the lists is a labeling that satisfies many edges, contradicting the soundness of
Label Cover.  The decoding is where the \emph{type} of the constraints matters: a projection lets the table at $v$ be read from the table at $u$, and the \emph{unique} type of the UGC lets decoding lose nothing
when the two tables are asked only two bits.
\end{bigpicture}

Label Cover is also the source of hardness for problems that look nothing like constraint satisfaction.  The same two-step recipe gives the logarithmic hardness of \sat{Set Cover}
\citep{LundYannakakis94,Feige98,DinurSteurer14}; and Label Cover under its original name was introduced in the context of lattice and code problems \citep{ABSS97}.


%% ===========================================================================
\section*{Exercises}
\addcontentsline{toc}{section}{Exercises}
%% ===========================================================================

\begin{exercise}[\stars{$\star$}]\label{ex:holonomy}
Let $G$ be a connected graph with linear constraints $x_v-x_u=c_e\pmod k$ on the edges $e=(u,v)$.  Prove that $\val(G)=1$ iff around every cycle the signed sum of the $c_e$ is $0$ mod $k$ (an edge traversed against
its orientation contributes $-c_e$).  \emph{Hint: spanning tree.}  Use this to give another proof of Proposition~\ref{prop:ug-easy} for linear constraints.
\end{exercise}

\begin{exercise}[\stars{$\star\star$}]\label{ex:cover-tight}
Compute $\val(G)$ and $\val(G^{\times2})$ for \sat{Max-Cut} on $K_n$ and show $\val(G^{\times2})/\val(G)=2-\tfrac2n$ for even $n$.  Show that for \sat{Max-Cut} on an odd cycle $C_\ell$ the ratio is $\tfrac\ell{\ell-1}$.  Conclude that the factor $2$ in Proposition~\ref{prop:cover} cannot be replaced by any constant smaller than~$2$.
\end{exercise}

\begin{exercise}[\stars{$\star\star$}]\label{ex:cv-tight}
(a) Let $\varphi$ be the $3$-CNF consisting of all eight clauses on three variables (every sign pattern).  Show that $\val(\varphi)=\tfrac78$ and $\val(G_\varphi)=\tfrac{23}{24}=1-\tfrac{1/8}3$, so the bound of Theorem~\ref{thm:cv} is tight.
(b) Show that the bound is also tight for \sat{Max-Cut} on a triangle, viewed as a CSP of arity $2$: $\val(\Phi)=\tfrac23$ and $\val(G_\Phi)=\tfrac56$.
\end{exercise}

\begin{exercise}[\stars{$\star\star$}]\label{ex:dreg}
Show that a relation $C\subseteq\Sigma\times\Sigma'$ in which every label of both sides has exactly $d$ partners is a disjoint union of $d$ perfect matchings.  \emph{Hint: Hall's theorem; remove one perfect matching and induct.}
Deduce the analogue of Corollary~\ref{cor:half}: from a $d$-to-$d$ game with $\val=1-\eps$ one gets a Unique Game with $\val=(1-\eps)/d$ on the same graph, with every edge repeated $d$ times.
\end{exercise}

\begin{exercise}[\stars{$\star$}]\label{ex:3cut}
(a) Verify that $x\ne y$ on $\Z_3$ is $2$-regular but not a disjoint union of copies of $K_{2,2}$.  (b) Write it as the union of the Unique constraints $y=x+1$ and $y=x+2$.  (c) For \sat{Max-3-Cut} on $K_4$ compute $\val(G)$ and the value of
the split Unique Games instance of Corollary~\ref{cor:half}.
\end{exercise}

\begin{exercise}[\stars{$\star\star\star$}]\label{ex:expander}
\emph{Randomness-efficient repetition.}  Let $A$ be the normalized adjacency matrix of a $d$-regular graph on the vertex set $\{0,1\}^r$ with second largest eigenvalue in absolute value at most $\lambda$, and let $B\subseteq\{0,1\}^r$ have density $\mu$.
(a) Prove that $\|P_BAP_B\|\le\mu+\lambda$, where $P_B$ is the coordinate projection onto $B$.  \emph{Hint: write a vector supported on $B$ as $\alpha\mathbf1/\sqrt{2^r}+y$ with $y\perp\mathbf1$ and bound $\alpha^2\le\mu\|u\|^2$ by Cauchy--Schwarz.}
(b) Deduce that a uniformly started random walk $\rho_1,\dots,\rho_k$ stays inside $B$ throughout with probability at most $\mu(\mu+\lambda)^{k-1}$.
(c) Let $V$ be a verifier with perfect completeness and soundness $\tfrac12$ and suppose an explicit family of such expanders with $\lambda\le\tfrac14$ and constant degree~$d$ exists.  Define $V_k$ that reads the walk's $k$ random strings, makes all $kq$ queries and
accepts iff all accept.  Show that $V_k$ uses $r+(k-1)\log d$ random bits and has soundness $(3/4)^{k-1}$, and conclude that with $k=\Theta(\log n)$ the FGLSS graph has $N=\poly(n)$ vertices and a gap $N^{\Omega(1)}$.
\end{exercise}

\begin{exercise}[\stars{$\star\star$}]\label{ex:reading2}
\emph{Reading.}  Section~2.3 of \citet{DIMACS10} reduces the PCP theorem to Label Cover by putting the random strings of the verifier on the left and the proof positions on the right.  Compare this reduction with Corollary~\ref{cor:verifier-game}.
Apply both to the verifier that picks a random clause of a $3$-CNF and queries its three variables, and identify the games.
\end{exercise}


%% ===========================================================================
\section*{Notes and further reading}
\addcontentsline{toc}{section}{Notes and further reading}
%% ===========================================================================

\paragraph{Where this lecture comes from.}
The dictionary between verifiers, CSPs and games is standard; we followed the treatment of the PCP theorem as a statement about Label Cover in Section~2.3 of the DIMACS notes \citep{DIMACS10}.  The name \emph{Label Cover}
comes from \citet{ABSS97}; two-prover one-round proof systems are due to \citet{FeigeLovasz92}; the clause--variable game is the canonical example.  The $d$-to-$1$ constraints are from \citet{Khot02}, and the $2$-to-$2$ constraints and the
``half-way'' remark are from \citet{Khot19}; the passage from a Unique Game to the bipartite two-prover game, with a loss of a factor of at most~$4$, is noted in \citet{Khot10}; we sharpened the factor to~$2$ for
the symmetric double cover (Proposition~\ref{prop:cover}).

\paragraph{On attributions.}
The Clique reduction and the exact formula $\omega(H_x)=2^r\cdot\max_\pi\Pr[\text{accept}]$ are due to \citet{FGLSS96}.  The hardness of approximating Clique within $N^\eps$ for some $\eps>0$, assuming only $\Ptime\ne\NP$, is Theorem~6 of
\citet{ALMSS98}, improving the $2^{\sqrt{\log N}}$ of \citet{AroraSafra98}; the factor $N^{1-\eps}$ is \citet{Hastad99}, as mentioned in Lecture~\ref{lec:why}.  The multiplicativity of the clique number under the strong
product is folklore.

\paragraph{What is not covered.}
We have not proved any hardness result for Label Cover itself beyond what Lecture~\ref{lec:why} did, and we have not discussed the \emph{size} of Label Cover instances (the Projection Games Conjecture asks for a nearly linear size).
Raz's theorem, which supplies the soundness, is quoted here; Lecture~\ref{lec:raz} proves the special case that Label Cover needs (expanding projection games) and states the general theorem.

\chapter{The Arora--Safra Route: Polynomials, Sum-Check and Composition}
\label{lec:as}

{\sffamily\large\color{softgray}How a global algebraic statement becomes a handful of local queries\par}
\vspace{0.5em}

\noindent
Lectures~\ref{lec:why} and~\ref{lec:games} took the PCP theorem as a black box and asked what it gives us.
This lecture opens the box---or rather, opens the first one, the proof of \citet{AroraSafra98}.  The second
proof, the one that reaches $O(1)$ queries \citep{ALMSS98}, is Lecture~4; the third, combinatorial one is
Lecture~6.

Be warned about what Arora and Safra actually proved.  Their theorem is \emph{not} $\NP=\PCP[O(\log n),O(1)]$.  It is
\[
  \NP=\PCP[\log n,\log n]\qquad\text{and}\qquad
  \NP=\PCP\bigl[\log n,\ \log^{1/2+\eps}n\bigr]\ \text{ for every fixed }\eps>0,
\]
a verifier that reads a \emph{sublogarithmic} number of proof bits.  The step from $\log^{1/2+\eps}n$ to $O(1)$ was taken
soon afterwards by \citet{ALMSS98} (both papers appeared at FOCS 1992), using the same composition lemma and a sharper low-degree test (Section~\ref{sec:as-limits}).
We therefore state precisely the theorem proved here, and we resist the common shorthand ``Arora--Safra proved the PCP theorem''.

What makes the proof worth a full lecture is that it is built from \emph{modules} with clean interfaces, and we can
state and prove most of them.  The cast:
\begin{itemize}
  \item a \emph{code}: low-degree polynomials, which turn ``$f$ is almost right everywhere'' into ``$f$ is exactly right at a random point'';
  \item an \emph{arithmetization}: satisfiability of a $3$-CNF becomes ``a certain polynomial sums to $0$ over a cube'';
  \item the \emph{sum-check}: a table that lets a verifier check such a sum by reading $O(m)$ entries;
  \item the \emph{low-degree test}: a table that lets a verifier check that a function is close to a polynomial by reading $O(m)$ entries;
  \item \emph{composition}: the lemma that lets a verifier with many queries to long symbols be rewritten as one that reads
        fewer, shorter symbols.
\end{itemize}
Each module has an input, an output and a parameter change.  The point of the lecture is to track them.

\begin{goalbox}{After this lecture you should be able to\ldots}
\begin{enumerate}[leftmargin=*,label=(\arabic*)]
  \item state exactly what \citet{AroraSafra98} proved and which constants separate it from the PCP theorem (Section~\ref{sec:as-land});
  \item prove the Schwartz--Zippel lemma, and use it to decode a function that is close to a low-degree polynomial
        (Section~\ref{sec:poly});
  \item arithmetize a $3$-CNF as a vanishing condition on a cube, and compute the single nonzero point of $g_A$ on a
        small example (Section~\ref{sec:arith});
  \item prove the soundness $1-dl/q$ of the sum-check by induction and compute the cheating probability of a given table
        (Section~\ref{sec:sumcheck});
  \item prove that a function all of whose axis-parallel lines are low-degree is a polynomial, and explain why the real test
        needs a \emph{robust} version of this fact (Section~\ref{sec:ldt});
  \item prove the composition lemma, and carry out the parameter chase that ends at $\log^{1/2+\eps}n$ queried bits
        (Sections~\ref{sec:compose}--\ref{sec:chase}).
\end{enumerate}
\end{goalbox}


%% ===========================================================================
\section{What was proved, and when}
\label{sec:as-land}
%% ===========================================================================

Recall the notation of Definition~\ref{def:pcp}: $\PCP[r,q]$ is the class of languages with a non-adaptive verifier using
$O(r)$ random bits and $O(q)$ queries, perfect completeness and soundness~$\tfrac12$.  Trivially
$\PCP[r,q]\subseteq\mathsf{NTIME}\bigl(2^{O(r)}q+\poly(n)\bigr)$, so $\PCP[O(\log n),q]\subseteq\NP$ for every $q=\poly(n)$ (we used this in
Section~\ref{sec:pcp}).  The whole question is how \emph{small} $q$ can be.

\begin{table}[t]
\centering\small
\renewcommand{\arraystretch}{1.3}
\begin{tabular}{@{}>{\raggedright\arraybackslash}p{3.5cm}
                   >{\raggedright\arraybackslash}p{4.4cm}
                   >{\raggedright\arraybackslash}p{4.4cm}@{}}
\toprule
Result & Randomness and queries & Consequence for \sat{Max-Clique}\\
\midrule
$\NP=\PCP[0,\poly(n)]$ (definition) & none, all of the witness & none\\
\citet{BFLS91} & $\NP\subseteq\PCP[\polylog n,\polylog n]$ (implicit) & ---\\
\citet{FGLSS96} & $\PCP[\log n\log\log n,$\newline $\log n\log\log n]$ & hard within any constant unless $\NP\subseteq\mathsf{DTIME}(n^{O(\log\log n)})$\\
\citet{AroraSafra98} & $\PCP[\log n,\log n]$ and\newline $\PCP[\log n,\log^{1/2+\eps}n]$ & within $2^{c\log^{1/2-\eps}N}$ is $\NP$-hard\\
\citet{ALMSS98} & $\PCP[\log n,1]$ & within $N^{\eps_1}$ is $\NP$-hard\\
\bottomrule
\end{tabular}
\caption{The race for the number of queries (all verifiers non-adaptive; $N$ is the number of vertices of the clique instance).
The row of Arora and Safra is the subject of this lecture.  Each row implies the ones above it, since $\PCP[r,q]\subseteq\PCP[r',q']$ whenever $r\le r'$ and $q\le q'$.}
\label{tab:race}
\end{table}

\noindent
Two further statements from \citet{AroraSafra98} will be used.

\begin{theorem}[Arora--Safra]\label{thm:as}
For every fixed $\eps>0$, $\NP=\PCP\bigl[\log n,\ \log^{1/2+\eps}n\bigr]$.
\end{theorem}

\begin{corollary}[Arora--Safra]\label{cor:as-clique}
For all constants $c,\eps>0$, if \sat{Max-Clique} on $N$-vertex graphs can be approximated within $2^{c\log^{1/2-\eps}N}$
in polynomial time, then $\Ptime=\NP$.
\end{corollary}

\noindent
Corollary~\ref{cor:as-clique} is the \emph{reason} the paper was written: it removes the ``almost'' from the title of \citet{FGLSS96}, who had
shown that approximating Clique within any constant factor implies $\NP\subseteq\mathsf{DTIME}(n^{O(\log\log n)})$.  We derive it in Section~\ref{sec:clique}.

\begin{pitfall}{three different statements share the name ``PCP theorem''}
(1) $\NP=\PCP[\log n,\log n]$ and $\NP=\PCP[\log n,\log^{1/2+\eps}n]$: Arora--Safra, FOCS 1992.
(2) $\NP=\PCP[\log n,1]$: Arora--Lund--Motwani--Sudan--Szegedy, FOCS 1992.  This is what Lecture~\ref{lec:why} called the PCP theorem
(Theorem~\ref{thm:pcp1}); the \emph{approximation} consequences for \sat{Max-3SAT} need~(2), not~(1): through the reduction of Proposition~\ref{prop:pcp-equiv} the constant $\eps_0$ is $2^{-O(q)}$, so with
$\log^{1/2+\eps}n$ queried bits one only gets a gap that shrinks with~$n$, not a constant (Arora and Safra, Section~6).
(3) The combinatorial proof of the same statement as~(2) by \citet{Dinur07}.
Arora and Safra's own Section~6 is explicit about this division of labour.
\end{pitfall}


%% ===========================================================================
\section{Polynomials as codes}
\label{sec:poly}
%% ===========================================================================

Fix a finite field $\F_q$.  For integers $h,m\ge1$ let
\[
  \mathcal P^m_h\;=\;\{\,p\in\F_q[x_1,\dots,x_m]\ :\ \deg_{x_i}p\le h\text{ for every }i\,\}
\]
be the polynomials of degree at most $h$ \emph{in each variable}, viewed as functions $\F_q^m\to\F_q$.  (Each $p\in\mathcal P^m_h$ has total degree at most $hm$.)
For $f,g\colon\F_q^m\to\F_q$ write $\Delta(f,g)=\Pr_{x\in\F_q^m}[f(x)\ne g(x)]$ for their (normalized) distance, and say that $f$ is
\defn{$\delta$-close} to $\mathcal P^m_h$ if $\Delta(f,p)\le\delta$ for some $p\in\mathcal P^m_h$.

\begin{lemma}[Schwartz--Zippel]\label{lem:sz}
A nonzero $p\in\mathcal P^m_d$ vanishes on at most a $dm/q$ fraction of the points of $\F_q^m$.
\end{lemma}
\begin{proof}
Induction on $m$.  For $m=1$ a nonzero univariate polynomial of degree $\le d$ has at most $d$ roots.
For $m>1$ write $p=\sum_{i=0}^{k}x_1^i\,p_i(x_2,\dots,x_m)$ with $k\le d$ and $p_k\ne0$.  By induction $p_k$ vanishes on at most a
$d(m-1)/q$ fraction of $\F_q^{m-1}$.  For every point $(x_2,\dots,x_m)$ where $p_k\ne0$, the polynomial $p(\cdot,x_2,\dots,x_m)$ is a
nonzero univariate polynomial of degree $k\le d$, hence vanishes for at most $k$ values of $x_1$.  So the fraction of zeros is at most
$\frac{d(m-1)}q+\frac kq\le\frac{dm}q$.
\end{proof}

\begin{corollary}[Codes of large distance]\label{cor:dist}
Distinct $p,p'\in\mathcal P^m_h$ satisfy $\Delta(p,p')\ge1-hm/q$.  If $hm<q/2$ and $f$ is $\delta$-close to $\mathcal P^m_h$ for some $\delta<\tfrac14$,
then $f$ is $\delta$-close to a \emph{unique} $\tilde f\in\mathcal P^m_h$, and no other polynomial of $\mathcal P_h^m$ is within $\tfrac12-\delta$ of~$f$.
\end{corollary}
\begin{proof}
The first claim is Lemma~\ref{lem:sz} applied to $p-p'$.  If $p,p'$ were both within $\delta$ of $f$, the triangle inequality would give
$\Delta(p,p')\le2\delta<\tfrac12<1-hm/q$, a contradiction.  The last claim is the same computation with $\tfrac12-\delta$.
\end{proof}

\noindent
Corollary~\ref{cor:dist} is the reason polynomials are used: a function that is $1\%$ wrong still \emph{determines} the polynomial it is trying to be.
Next, how to turn a string into a polynomial.

\begin{lemma}[Polynomial extension]\label{lem:ext}
Let $h<q$, identify $\{0,1,\dots,h\}$ with the first $h+1$ elements of $\F_q$, and let $s\colon\{0,\dots,h\}^m\to\F_q$ be any function.  Then there is a \emph{unique}
$\hat s\in\mathcal P^m_h$ with $\hat s=s$ on $\{0,\dots,h\}^m$, namely
\[
  \hat s(x)\;=\;\sum_{u\in\{0,\dots,h\}^m}s(u)\,\Lambda_u(x),\qquad
  \Lambda_u(x)=\prod_{i=1}^m\lambda_{u_i}(x_i),\qquad
  \lambda_a(x)=\prod_{b\in\{0..h\},\,b\ne a}\frac{x-b}{a-b}.
\]
\end{lemma}
\begin{proof}
Each $\lambda_a$ has degree $h$, equals $1$ at $a$ and $0$ at the other points of $\{0..h\}$; hence $\Lambda_u$ is $1$ at $u$ and $0$
elsewhere on the cube, and $\hat s\in\mathcal P^m_h$ agrees with $s$ on the cube.  For uniqueness it suffices to show that $p\in\mathcal P^m_h$ vanishing on the cube is
zero.  Write $p=\sum_{j\le h}x_m^jp_j(x_1,\dots,x_{m-1})$.  For every fixed $x'\in\{0..h\}^{m-1}$ the univariate polynomial $p(x',\cdot)$ has degree $\le h$ and $h+1$ roots,
so it is zero, \ie $p_j(x')=0$ for all $j$.  Hence each $p_j$ vanishes on $\{0..h\}^{m-1}$ and by induction $p_j=0$.
\end{proof}

\begin{example}[The multilinear extension of XOR]\label{ex:xor}
For $m=2$, $h=1$ and $s(0,0)=s(1,1)=0$, $s(1,0)=s(0,1)=1$, Lemma~\ref{lem:ext} gives
\[
  \hat s(x_1,x_2)=x_1(1-x_2)+(1-x_1)x_2=x_1+x_2-2x_1x_2 .
\]
Over $\F_7$ it takes the value $3+5-30\equiv6$ at $(3,5)$.  The function $s$ is determined by $4$ values; its extension is a table with $q^2=49$ entries, and two distinct such extensions differ on at least $1-2/7=5/7$ of the $49$ points
(Corollary~\ref{cor:dist}), so even a corrupted table pins down $s$.
\end{example}

\noindent
\textbf{Encoding strings.}  Let $n=(h+1)^m$.  Identify $\{1,\dots,n\}$ with the cube $\{0,\dots,h\}^m$ (base-$(h+1)$ digits) and encode $a\in\{0,1\}^n\subset\F_q^n$ by
$\hat a\in\mathcal P^m_h$.  The encoding is injective, and the codeword is a table of $q^m$ field elements; if $q=\poly(h)$ then
$q^m=h^{O(m)}=\poly(n)$ since $(h+1)^m=n$.  This is the sense in which the encoding is ``efficient''.

\begin{pitfall}{a table that is close to a polynomial is not a polynomial}
The verifier will not hold a polynomial; it holds a \emph{table} $f\colon\F_q^m\to\F_q$ supplied by an adversary.  The low-degree test (Section~\ref{sec:ldt}) can only certify that
$f$ is $\delta$-close to some $\tilde f\in\mathcal P^m_h$ for a small constant~$\delta$.  Any later step that needs values of $\tilde f$ at a point $x$ reads $f(x)$ and is correct unless $x$
lies in the $\delta$-fraction where $f\ne\tilde f$.  The design principle of the whole verifier is therefore: \emph{never evaluate $\tilde f$ at a point the adversary can predict.}
Every point at which $f$ is read must be uniformly distributed, so that the error probability is $\le\delta$ per read.
\end{pitfall}


%% ===========================================================================
\section{Arithmetization: satisfiability as a vanishing condition}
\label{sec:arith}
%% ===========================================================================

Let $\varphi$ be a $3$-CNF with $n=(h+1)^m$ variables and $n$ clauses (pad with dummy variables and clauses).  Let $\mathbb K=\{0,1,\dots,h\}\subset\F_q$.  Identify variables and clauses with
points of the cube $\mathbb K^m$.  Every clause $c\in\mathbb K^m$ has an ordered triple of variables $v_1(c),v_2(c),v_3(c)\in\mathbb K^m$ and signs
$\varsigma_j(c)\in\{0,1\}$ ($\varsigma_j(c)=1$ if the $j$th literal is unnegated).  Define
\[
  \chi_j(c,v)=\begin{cases}1&\text{if }v=v_j(c)\\0&\text{otherwise}\end{cases}\quad\bigl(\mathbb K^m\times\mathbb K^m\to\{0,1\}\bigr),\qquad
  \varsigma_j\colon\mathbb K^m\to\{0,1\},
\]
and let $\hat\chi_j\in\mathcal P^{2m}_h$, $\hat\varsigma_j\in\mathcal P^m_h$ be their polynomial extensions (Lemma~\ref{lem:ext}); they depend only on
$\varphi$ and are computable in time $\poly(n,q^m)$.

For $A\in\mathcal P^m_h$, which we think of as the extension of a Boolean assignment, define $g_A\colon\F_q^{4m}\to\F_q$ by
\[
  g_A(z,w_1,w_2,w_3)\;=\;\prod_{j=1}^3\hat\chi_j(z,w_j)\cdot\bigl(\hat\varsigma_j(z)-A(w_j)\bigr),\qquad z,w_1,w_2,w_3\in\F_q^m .
\]

\begin{lemma}[Algebraic view of $3$SAT]\label{lem:arith}
Let $A\in\mathcal P^m_h$.
\begin{enumerate}[label=(\alph*)]
  \item If $a\in\{0,1\}^n$ satisfies $\varphi$ and $A=\hat a$, then $g_A(y)=0$ for every $y\in\mathbb K^{4m}=\{0,\dots,h\}^{4m}$.  If $a\in\{0,1\}^n$ does not satisfy $\varphi$, then $g_{\hat a}$
        is nonzero at some point of $\mathbb K^{4m}$.
  \item Conversely, for \emph{every} $A$: if $g_A$ vanishes on $\mathbb K^{4m}$ then $\varphi$ is satisfiable.  (The assignment $a'_v=\mathbf 1[A(v)=1]$ satisfies it.)
  \item $g_A\in\mathcal P^{4m}_{6h}$, and evaluating $g_A$ at a point $(z,w_1,w_2,w_3)$ requires only the values of $A$ at the three points
        $w_1,w_2,w_3\in\F_q^m$ (plus $\poly(n)$ work independent of $A$).  If the point is uniform in $\F_q^{4m}$ then each $w_j$ is uniform in $\F_q^m$.
\end{enumerate}
\end{lemma}
\begin{proof}
Fix $(c,v_1,v_2,v_3)\in\mathbb K^{4m}$.  On the cube the extensions $\hat\chi_j,\hat\varsigma_j$ equal $\chi_j,\varsigma_j$.  If $(v_1,v_2,v_3)\ne(v_1(c),v_2(c),v_3(c))$ then some $\chi_j(c,v_j)=0$ and $g_A=0$
at this point.  Otherwise
$g_A(c,v_1,v_2,v_3)=\prod_{j=1}^3\bigl(\varsigma_j(c)-A(v_j)\bigr)$, which vanishes iff $A(v_j)=\varsigma_j(c)$ for some $j$.

(a) For Boolean $a$ and $A=\hat a$ we have $A(v_j)=a_{v_j}$, and $a_{v_j(c)}=\varsigma_j(c)$ says exactly that the $j$th literal of $c$ is true.  So $g_A$ vanishes at $(c,v_1(c),v_2(c),v_3(c))$ iff $a$ satisfies $c$.

(b) If $A(v_j(c))=\varsigma_j(c)\in\{0,1\}$ then $a'_{v_j(c)}=\varsigma_j(c)$ (for $\varsigma_j=1$ because $A=1$; for $\varsigma_j=0$ because $A=0\ne1$).  So every clause has a literal made true by $a'$.

(c) Each factor is a product of two polynomials of degree $\le h$ in every variable; the three factors share~$z$, so the degree is $\le6h$ in each $z_i$, and $\le2h$ in each $w_{j,i}$.  The rest is read off
the formula; the $w_j$ are disjoint blocks of coordinates of a uniform point, hence uniform.
\end{proof}

\begin{pitfall}{the arithmetization certifies the \emph{rounded} assignment}
Part~(b) says that a table $A$ with $g_A\equiv0$ on the cube is satisfiable-by-rounding, not that its values on the cube are bits.  Indeed take any satisfying $a$ and change $a_v$ to $5$ for a variable $v$ that
occurs in clauses whose other literals are true: $g_A$ still vanishes on the cube.  This does no harm for the \emph{soundness} of a verifier for $\sat{3SAT}$ (an unsatisfiable $\varphi$ is rejected by every $A$, by (b)),
but it matters in Section~\ref{sec:normal}, where the verifier will be asked about the \emph{decoded string} of~$A$.  There one needs the further vanishing condition $A(A-1)\equiv0$ on the cube, and one more sum-check.
\citet{AroraSafra98} leave this implicit; the cost is a constant factor.
\end{pitfall}

\begin{example}[A formula, an assignment, one nonzero point]\label{ex:arith}
Take $h=1$, $m=2$ (so $n=4$ variables $x_0,\dots,x_3$, named by their binary index, and $4$ clauses) and
\[
  C_0=x_3\vee\neg x_2\vee\neg x_1,\quad C_1=\neg x_3\vee x_2\vee x_0,\quad C_2=x_1\vee x_0\vee\neg x_2,\quad C_3=\neg x_0\vee\neg x_2\vee\neg x_3 .
\]
The assignment $a=(x_0,x_1,x_2,x_3)=(0,0,0,1)$ falsifies exactly $C_1$ (here $x_3=1$, $x_2=0$, $x_0=0$).  Computing $g_{\hat a}$ at all $2^8=256$ points of $\{0,1\}^{8}$
over $\F_{31}$ gives a single nonzero value:
\[
  g_{\hat a}\bigl(\underbrace{(0,1)}_{C_1},\underbrace{(1,1)}_{x_3},\underbrace{(1,0)}_{x_2},\underbrace{(0,0)}_{x_0}\bigr)
  \;=\;(0-1)(1-0)(1-0)\;=\;-1,
\]
the product of $\varsigma_j-a_{v_j}$ for the three false literals.  So a single violated clause becomes a single nonzero point of a $256$-point cube:
a needle in a haystack.  We cannot hope to find it by sampling, and this is exactly why we multiply by the random \emph{zero-testers} of the next lemma, which turn
``somewhere nonzero'' into ``the weighted sum is nonzero''.
\end{example}

\begin{lemma}[Zero-testers]\label{lem:zt}
For $t=(t_1,\dots,t_{k})\in\F_q^{k}$ let $R_t\in\mathcal P^{k}_h$ be the polynomial $R_t(y)=\prod_{i=1}^{k}I_{t_i}(y_i)$,
where $I_s\in\mathcal P^1_h$ is the polynomial with $I_s(a)=s^a$ for $a=0,\dots,h$ (Lemma~\ref{lem:ext}).  Then for every $f\colon\{0..h\}^{k}\to\F_q$ that is not identically zero,
\[
  \Pr_{t\in\F_q^{k}}\Bigl[\sum_{y\in\{0..h\}^{k}}R_t(y)f(y)=0\Bigr]\ \le\ \frac{hk}{q}.
\]
In particular this is less than $\tfrac18$ when $q>8hk$ ($q>32hm$ for $k=4m$).
\end{lemma}
\begin{proof}
For $y\in\{0..h\}^{k}$ we have $R_t(y)=\prod_it_i^{y_i}$.  Hence $F(t):=\sum_yR_t(y)f(y)=\sum_yf(y)\prod_it_i^{y_i}$ is a polynomial in $t$ of degree $\le h$ in every variable, and it is
nonzero because distinct $y$ give distinct monomials and some $f(y)\ne0$.  By Lemma~\ref{lem:sz}, $\Pr_t[F(t)=0]\le hk/q$.
\end{proof}

\begin{example}[Zero-testers on Example~\ref{ex:arith}]\label{ex:zt}
Here $k=4m=8$ and $g_{\hat a}|_{\mathbb K^{4m}}$ has a single nonzero point, $y=(0,1\,|\,1,1\,|\,1,0\,|\,0,0)$ with value $-1$.  So $F(t)=-t_2t_3t_4t_5$ (the coordinates where $y$ has a~$1$), which vanishes iff one of these four
coordinates of $t$ is $0$:
\[
  \Pr_t[F(t)=0]\;=\;1-\bigl(\tfrac{30}{31}\bigr)^4\;\approx\;0.123\ \ \le\ \frac{hk}q=\frac8{31}\approx0.258 .
\]
(A simulation with $2\cdot10^5$ samples gives $0.123$.)  The field is too small for the hypothesis $q>8hk=64$ of the lemma, but the bound $hk/q$ itself holds.  Note how sparse $F$ is: a single violated clause
is a one-monomial polynomial, and the zero-tester detects it unless $t$ is unlucky in one of four coordinates.
\end{example}

\begin{corollary}[The summed polynomials $P_t$]\label{cor:P}
For $t\in\F_q^{4m}$ and $A\in\mathcal P^m_h$ put $P^A_t=R_t\cdot g_A\in\mathcal P^{4m}_{7h}$ ($R_t$ the zero-testers with $k=4m$, $q>32hm$).  Then
\begin{itemize}
  \item if $A=\hat a$ for a satisfying assignment $a$, then $\sum_{y\in\mathbb K^{4m}}P^A_t(y)=0$ for \emph{every} $t$;
  \item if $\varphi$ is \emph{unsatisfiable}, then for \emph{every} $A\in\mathcal P^m_h$ we have $\sum_{y\in\mathbb K^{4m}}P^A_t(y)\ne0$ for all but at most an $\tfrac18$ fraction of the $t$.
\end{itemize}
Moreover $P^A_t(y)$ at a uniformly random $y\in\F_q^{4m}$ is computed from three uniformly distributed values of $A$.
\end{corollary}
\begin{proof}
First bullet: Lemma~\ref{lem:arith}(a).  Second bullet: by Lemma~\ref{lem:arith}(b), $g_A\not\equiv0$ on $\mathbb K^{4m}$ when $\varphi$ is unsatisfiable; apply Lemma~\ref{lem:zt} to $f=g_A|_{\mathbb K^{4m}}$.
\end{proof}

\begin{bigpicture}{What arithmetization buys}
A $3$-CNF is a list of $n$ constraints, and a single violated clause is a needle in a haystack that a verifier reading $O(1)$ clauses cannot find.  After arithmetization the same fact is one equation,
``$\sum_{y\in\mathbb K^{4m}}P^A_t(y)=0$'', about a single low-degree polynomial for a random~$t$.  The sum has $(h+1)^{4m}=n^4$ terms, but the polynomial has structure---low degree in every variable---and
the \emph{sum-check} of the next section exploits that structure to verify the sum by evaluating the polynomial at \emph{one} random point.
\end{bigpicture}


%% ===========================================================================
\section{The sum-check protocol, as a table}
\label{sec:sumcheck}
%% ===========================================================================

The sum-check is due to \citet{LFKN92}; we give the non-adaptive, table-based form of \citet{AroraSafra98}.

Let $B\in\mathcal P^{\ell}_d$ and $H\subseteq\F_q$.  For $1\le i\le\ell$ and $a_1,\dots,a_{i-1}\in\F_q$ put
\[
  B_{a_1,\dots,a_{i-1}}(x_i)\;=\;\sum_{x_{i+1},\dots,x_\ell\in H}B(a_1,\dots,a_{i-1},x_i,x_{i+1},\dots,x_\ell),
\]
a univariate polynomial of degree $\le d$ in $x_i$ (for $i=1$ we write $B_\varnothing(x_1)$).  A \defn{table of partial sums} is any assignment $g$ of a univariate polynomial
$g_{a_1,\dots,a_{i-1}}$ of degree $\le d$ to every prefix $(a_1,\dots,a_{i-1})\in\F_q^{i-1}$, $1\le i\le\ell$; it has $\sum_i q^{i-1}\le q^\ell$ entries, each of $O(d\log q)$ bits.
The \emph{honest} table is $g_{a_1..a_{i-1}}=B_{a_1..a_{i-1}}$.

\begin{definition}[Sum-check test]\label{def:sumcheck}
Input: $B$ (available through oracle access), $H$, a claimed value $c\in\F_q$ and a table $g$.
Choose $a_1,\dots,a_\ell\in\F_q$ independently and uniformly; then accept iff
\begin{align*}
  &\sum_{x\in H}g_\varnothing(x)=c,\\
  &\sum_{x\in H}g_{a_1,\dots,a_{i-1}}(x)=g_{a_1,\dots,a_{i-2}}(a_{i-1})\quad(2\le i\le\ell),\\
  &g_{a_1,\dots,a_{\ell-1}}(a_\ell)=B(a_1,\dots,a_\ell).
\end{align*}
\end{definition}

\noindent
The test reads $\ell$ table entries and one value of $B$ at the uniform point $a=(a_1,\dots,a_\ell)$.  The entries to be read depend only on $a$, not on
earlier answers: the test is non-adaptive.  It uses $\ell\log q=\log|\F_q^\ell|$ random bits.

\begin{lemma}[Sum-check]\label{lem:sumcheck}
Let $B\in\mathcal P^\ell_d$, $c\in\F_q$.
\begin{enumerate}[label=(\alph*)]
  \item If $\sum_{x\in H^\ell}B(x)=c$, the honest table is accepted with probability~$1$.
  \item If $\sum_{x\in H^\ell}B(x)\ne c$, then \emph{every} table is accepted with probability at most $d\ell/q$.
\end{enumerate}
\end{lemma}
\begin{proof}
(a) Each check is an identity between partial sums.

(b) Induction on $\ell$.  For $\ell=1$ the test accepts iff $\sum_Hg_\varnothing=c$ and $g_\varnothing(a_1)=B(a_1)$.  Since $\sum_HB\ne c$ we have $g_\varnothing\ne B$, and two distinct polynomials of
degree $\le d$ agree in at most $d$ points, so the acceptance probability is at most $d/q$.

For $\ell>1$ let $B_\varnothing(x_1)=\sum_{x_2..x_\ell\in H}B(x_1,x_2,\dots,x_\ell)$.  Then $\sum_{x\in H}B_\varnothing(x)=\sum_{H^\ell}B\ne c$.
\emph{Case 1:} $g_\varnothing=B_\varnothing$.  Then $\sum_Hg_\varnothing\ne c$ and the first check fails with probability~$1$.
\emph{Case 2:} $g_\varnothing\ne B_\varnothing$.  They agree in at most $d$ points $a_1$.  For every other $a_1$, the claim ``$\sum_{x_2..x_\ell\in H}B(a_1,x_2,\dots,x_\ell)=g_\varnothing(a_1)$'' is
false.  Fixing $a_1$, the remaining checks form exactly the sum-check test for the polynomial $B(a_1,\cdot)\in\mathcal P^{\ell-1}_d$, the claim $c'=g_\varnothing(a_1)$ and the sub-table $g_{a_1,\ast}$;
by induction it accepts with probability $\le d(\ell-1)/q$.  Hence the overall acceptance probability is at most
\[
  \frac dq\cdot1+\Bigl(1-\frac dq\Bigr)\frac{d(\ell-1)}q\ \le\ \frac{d\ell}q.\qedhere
\]
\end{proof}

\noindent
(The recursion $p_\ell\le d/q+(1-d/q)p_{\ell-1}$, $p_1\le d/q$, solves to the slightly sharper $p_\ell\le1-(1-d/q)^\ell\le d\ell/q$.)

\begin{example}[A cheating table, exactly]\label{ex:sc}
Let $q=7$, $d=2$, $\ell=2$, $H=\{0,1\}$ and $B(x_1,x_2)=x_1x_2+x_2^2+1$.  Then
$B_\varnothing(x_1)=B(x_1,0)+B(x_1,1)=1+(x_1+2)=x_1+3$, and the true sum is $B_\varnothing(0)+B_\varnothing(1)=3+4=7\equiv0$.
Suppose the prover wants to convince us of the false claim $c=1$.  An exhaustive computation over all tables (every polynomial of degree $\le2$ at every node of the tree) gives
\[
  \max_g\Pr[\text{accept}]\;=\;\frac{24}{49}\;=\;1-\Bigl(\frac57\Bigr)^2\;\approx\;0.490\ \ <\ \frac{d\ell}q=\frac47\approx0.571 .
\]
The maximum is the same for every false claim $c\ne0$, and equals exactly the sharper bound $1-(1-d/q)^\ell$ of the remark after the proof: at the first level the cheating polynomial can have
$d=2$ points $a_1$ at which it agrees with $B_\varnothing$, and at each of the other points the second level leaves a further $d/q$.  The bound $d\ell/q$ of the lemma is its linearization.
\end{example}

\begin{pitfall}{a table of partial sums is a prover strategy written out}
In the interactive sum-check of \citet{LFKN92} the prover answers round by round, and its $i$th message may depend on the challenges $a_1,\dots,a_{i-1}$.  A deterministic strategy of such a prover is
\emph{exactly} a table of partial sums: the entry $g_{a_1,\dots,a_{i-1}}$ is its $i$th message.  So the lemma says the same thing in both models.  What the table form adds is that the verifier
may choose all of $a$ first and then read the $\ell$ entries on the path to $a$ in one batch: the test is non-adaptive.  This is needed for the composition lemma (Section~\ref{sec:compose}).
\end{pitfall}

\subsection{Putting arithmetization and sum-check together}

Combine Corollary~\ref{cor:P} and Lemma~\ref{lem:sumcheck} with $B=P^A_t$, $d=7h$, $\ell=4m$, $H=\{0..h\}$, $c=0$.  A verifier that holds $f\colon\F_q^m\to\F_q$ (in the role of $A$)
and a collection of sum-check tables $\{g^{(t)}\}_{t\in\F_q^{4m}}$ does the following:
\begin{enumerate}
  \item pick $t\in\F_q^{4m}$ uniformly ($4m\log q$ random bits) and compute from $\varphi$ and $t$ everything about $P^f_t$ except the values of $f$;
  \item run the sum-check on $P^f_t$ with claim $c=0$ and table $g^{(t)}$.
\end{enumerate}
If $f$ extends a satisfying assignment, there are tables making this accept with probability~$1$.  If $\varphi$ is unsatisfiable and $f$ is \emph{any} polynomial in $\mathcal P^m_h$, then by
Corollary~\ref{cor:P} and Lemma~\ref{lem:sumcheck} the test accepts with probability at most $\tfrac18+28hm/q$ (here $d\ell/q=7h\cdot4m/q$).  In the next section we remove the assumption that $f$ is a polynomial.


%% ===========================================================================
\section{The low-degree test}
\label{sec:ldt}
%% ===========================================================================

The sum-check of the previous section was proved for polynomials.  The verifier holds a table $f$.  We need a \emph{local} test that accepts every
$f\in\mathcal P^m_h$ and rejects every $f$ far from $\mathcal P^m_h$.

\subsection{The exact statement, and why it is not enough}

An \defn{axis-parallel line} of direction $i$ in $\F_q^m$ is $\ell=\{(a_1,\dots,a_{i-1},x,a_{i+1},\dots,a_m):x\in\F_q\}$.  There are $m\,q^{m-1}$ of them.
If $p\in\mathcal P^m_h$ then its restriction to every axis-parallel line is a univariate polynomial of degree $\le h$.  The converse is true, and its proof is a good warm-up.

\begin{lemma}[Axis-parallel lines suffice, exact version]\label{lem:axis}
Let $f\colon\F_q^m\to\F_q$ be such that the restriction of $f$ to every axis-parallel line is a polynomial of degree $\le h$ (with $h<q$).  Then $f\in\mathcal P^m_h$.
\end{lemma}
\begin{proof}
Induction on $m$; $m=1$ is the hypothesis.  Fix distinct $c_0,\dots,c_h\in\F_q$.  For each $j$ the function $f(\cdot,c_j)$ on $\F_q^{m-1}$ satisfies the hypothesis in $m-1$ variables, so it is a polynomial
$p_j\in\mathcal P^{m-1}_h$.  Let $\lambda_j$ be the univariate Lagrange polynomials of degree $h$ for the nodes $c_0,\dots,c_h$ and define
$G(x',x_m)=\sum_{j=0}^h\lambda_j(x_m)\,p_j(x')\in\mathcal P^m_h$.  For each fixed $x'$, the functions $f(x',\cdot)$ and $G(x',\cdot)$ are both univariate polynomials of degree $\le h$ and they agree at the
$h+1$ points $c_0,\dots,c_h$; hence they are equal.  So $f=G$ everywhere.
\end{proof}

\noindent
As a sanity check, take $q=5$, $h=1$, $m=2$.  A pruned exhaustive search through the $5^{25}$ functions $\F_5^2\to\F_5$ finds exactly $625=5^{(h+1)^2}$ with all rows and all columns of degree $\le1$:
precisely the $5^4$ polynomials of $\mathcal P^2_1$, as the lemma predicts.

\begin{bigpicture}{Exact versus robust}
Lemma~\ref{lem:axis} is a statement about \emph{perfect} functions.  A verifier can only look at a few lines, and can at best find that \emph{most} of the lines it looked at are low-degree and \emph{mostly} agree where they meet.
The theorem the test needs is the \emph{robust} form: if the restrictions to most axis-parallel lines are consistent with each other and with~$f$, then $f$ is \emph{close} to a polynomial.
All the difficulty of the Arora--Safra analysis sits here, and it is the reason the field size is cubic in the degree.
\end{bigpicture}

\subsection{The test of Arora and Safra}

A \defn{decomposition} is a table $g$ that assigns to every axis-parallel line $\ell$ a univariate polynomial $g[\ell]$ of degree $\le h$ (an entry of $O(h\log q)$ bits).  For $1\le i\le m$ and
$a\in\F_q^m$ let $g_i(a)=g[\ell_i(a)](a_i)$, where $\ell_i(a)$ is the line of direction $i$ through $a$.  A \defn{checkpoint} is a pair $(i,a)$ with $1\le i<m$; the decomposition is \defn{consistent} at it if
$g_i(a)=g_{i+1}(a)$.  Note that a genuine decomposition of a polynomial $p$ ($g[\ell]=p|_\ell$) is consistent everywhere and $g_i=p$ for all $i$.

\begin{definition}[Low-degree test]\label{def:ldt}
Given $f\colon\F_q^m\to\F_q$, the closeness parameter $\delta<0.1$ and a decomposition $g$:
\begin{enumerate}[label=(\roman*)]
  \item choose $\lceil4/\delta\rceil$ points $a\in\F_q^m$ uniformly and independently, and check $f(a)=g_1(a)$ at each;
  \item choose $\lceil8m/\delta\rceil$ checkpoints $(i,a)$ uniformly and independently, and check $g_i(a)=g_{i+1}(a)$ at each.
\end{enumerate}
Accept iff all checks pass.
\end{definition}

\noindent
The test reads $O(1/\delta)$ values of $f$ and $O(m/\delta)$ entries of $g$: $O(m)$ queries for constant $\delta$, non-adaptively.  (The randomness is $O(m\log q)$ bits per sample, which is reduced to $O(\log q^m)$ in total
by pairwise-independent sampling; we omit it.)  It accepts every $f\in\mathcal P^m_h$ together with its true decomposition.

\begin{lemma}[Robust consistency; Arora--Safra, Lemma~4.6.4]\label{lem:robust}
Let $m\ge2$ and $q\ge100\,h^3m^2$.  If $g_1$ is not $\eps$-close to $\mathcal P^m_h$, then the fraction of checkpoints at which $g$ is inconsistent is at least
$\gamma^{2(m-1)}\min\{0.1,\eps\}/(m-1)$, where $\gamma=1-\sqrt{h/q}$.  Since $\sqrt{h/q}\le1/(10m)$ we have $\gamma^{2(m-1)}\ge1-2(m-1)/(10m)\ge0.8$, so this fraction is at least $\min\{0.1,\eps\}/(2m)$.
\end{lemma}

\noindent
We do not prove Lemma~\ref{lem:robust}.  The proof is an induction on $m$.  Its base case $m=2$ is the \emph{bivariate} lemma that \citet{ALMSS98} quote as Theorem~69 (due to Arora--Safra): if the row polynomial and the column polynomial of a
function $F\colon\F_q^2\to\F_q$ each agree with $F$ on all but an $\eta$ fraction of points ($\eta\le\eta_0$) and $q\ge c\,h^3$, then $F$ is $4\eta$-close to a polynomial of degree $\le h$ in each variable.  This is the robust analogue
of the exact Lemma~\ref{lem:axis} with $m=2$, and it is where the \emph{cubic} dependence on~$h$ comes from.  \citet{PolishchukSpielman94} later gave a proof that needs only $q=\Omega(h)$.

\begin{corollary}[Soundness of the test]\label{cor:ldt}
Let $q\ge100\,h^3m^2$ and $m\ge2$.  If $f$ is not $\delta$-close to $\mathcal P^m_h$ ($\delta<0.1$), then for \emph{every} decomposition $g$ the test of Definition~\ref{def:ldt} rejects with probability at least $3/4$.
\end{corollary}
\begin{proof}
If $\Delta(f,g_1)\ge\delta/2$, part~(i) fails to reject with probability at most $(1-\delta/2)^{4/\delta}\le e^{-2}\approx0.135<\tfrac14$.
Otherwise $\Delta(f,g_1)<\delta/2$, so $g_1$ is not $\delta/2$-close to $\mathcal P^m_h$ (else $f$ would be $\delta$-close by the triangle inequality).  By Lemma~\ref{lem:robust} with $\eps=\delta/2\le0.05$ a fraction at least $\delta/(4m)$ of the
checkpoints is inconsistent, and part~(ii) misses all of them with probability at most $(1-\delta/4m)^{8m/\delta}\le e^{-2}<\tfrac14$.
\end{proof}

\begin{remark}
With $\delta=0.01$ the two failure probabilities are $\le0.1347$ and $\le0.1353$ respectively (computed from the formulas above).  The older analysis of this test (for $h=1$ in \citealp{FGLSS96}) required $O(mh)$ queries; Arora and Safra's algebraic argument gets $O(m)$, at the price of $q=\Omega(h^3m^2)$.  The number of queries to the low-degree table is exactly what is later
reduced to a constant by \citet{ALMSS98}; see Section~\ref{sec:as-limits}.
\end{remark}


%% ===========================================================================
\section{The basic verifier}
\label{sec:vpoly}
%% ===========================================================================

We now assemble the three modules.  Let $n=(h+1)^m$ with $h\ge\log n$ and $q=\Theta(h^3m^2)$ a prime power (note that $m<h$ forces $q=O(h^5)$).  The verifier $V_{\rm poly}$ for a $3$-CNF $\varphi$ with $n$ variables expects the proof to consist of

\begin{center}\small
\renewcommand{\arraystretch}{1.25}
\begin{tabular}{@{}>{\raggedright\arraybackslash}p{2.6cm}
                   >{\raggedright\arraybackslash}p{6.0cm}
                   >{\raggedright\arraybackslash}p{4.0cm}@{}}
\toprule
component & content & size\\
\midrule
$f$ & a function $\F_q^m\to\F_q$ (the claimed extension of a satisfying assignment) & $q^m=n^{O(1)}$ field elements\\
$g$ & a decomposition for the low-degree test of $f$ & $mq^{m-1}$ polynomials of degree $h$\\
$\{\tau^{(t)}\}_{t\in\F_q^{4m}}$ & a table of partial sums for the sum-check of $P_t^{f}$, one per $t$ & $\le2q^{8m}=n^{O(1)}$ polynomials of degree $7h$\\
\bottomrule
\end{tabular}
\end{center}

\noindent
and works as follows.

\begin{definition}[The verifier $V_{\rm poly}$]
Given the proof components above, $V_{\rm poly}$ does the following.
\begin{enumerate}
  \item Run the low-degree test (Definition~\ref{def:ldt}) on $f$ and $g$ with $\delta=0.01$; reject if it rejects.
  \item Choose $t\in\F_q^{4m}$ uniformly and run the sum-check (Definition~\ref{def:sumcheck}) for $B=P^{f}_t$ (that is, with the values of $f$ in place of $A$), $H=\{0..h\}$, $\ell=4m$, $d=7h$, $c=0$, and table $\tau^{(t)}$;
        reject if it rejects.  Whenever the sum-check needs $f$ at a point $w$, it reads $f(w)$.
  \item Otherwise accept.
\end{enumerate}
\end{definition}

\begin{theorem}[The basic verifier; Arora--Safra, Theorem~3.5(a)]\label{thm:vpoly}
$V_{\rm poly}$ is a non-adaptive verifier for $\sat{3SAT}$ with perfect completeness, soundness $<\tfrac14$, using $O(\log n)$ random bits and $O(m)$ queries,
each to an entry of $O(h\log h)$ bits, and with decision time $\poly(h)$.
\end{theorem}
\begin{proof}
\emph{Completeness.}  If $\varphi$ is satisfiable, let $f$ be the extension of a satisfying assignment, $g$ its true decomposition, and $\tau^{(t)}$ the honest table for $P_t^f$.  The low-degree test accepts by completeness, and
the sum-check accepts by Lemma~\ref{lem:sumcheck}(a), because $\sum_{y\in\mathbb K^{4m}}P^f_t(y)=0$ by Corollary~\ref{cor:P}.

\emph{Soundness.}  Let $\varphi$ be unsatisfiable and fix any proof.  If $f$ is not $0.01$-close to $\mathcal P^m_h$, step~1 rejects with probability $\ge3/4$ (Corollary~\ref{cor:ldt}).
So assume $f$ is $0.01$-close, and let $\tilde f\in\mathcal P^m_h$ be its unique nearest polynomial (Corollary~\ref{cor:dist}).  Step~2 accepts only if one of three events happens:
\begin{enumerate}[label=(\roman*)]
  \item $\sum_{\mathbb K^{4m}}P^{\tilde f}_t=0$: by Corollary~\ref{cor:P} this has probability $\le\tfrac18$ over $t$ (this is where unsatisfiability is used, for the arbitrary polynomial $\tilde f$);
  \item the sum is nonzero but the sum-check does not detect it: probability $\le d\ell/q=28hm/q$ by Lemma~\ref{lem:sumcheck}(b), which is applied to the genuine polynomial $B=P_t^{\tilde f}$;
  \item the sum-check reads $f(w_j)$ at one of three points $w_1,w_2,w_3$ with $f(w_j)\ne\tilde f(w_j)$: each $w_j$ is uniform (Lemma~\ref{lem:arith}(c)), so by the union bound this has probability $\le3\cdot0.01$.
\end{enumerate}
Altogether the acceptance probability is at most $\tfrac18+0.03+28hm/q<\tfrac14$ (for $q\ge h^3m^2$ and $h\ge\log n$ the last term is $o(1)$).

\emph{Parameters.}  The random bits are those of $t$, the point $a\in\F_q^{4m}$ of the sum-check, and the samples of the low-degree test: $O(m\log q)=O(m\log h)=O(\log n)$ in total, since $(h+1)^{m}=n$.
The low-degree test reads $O(m)$ entries and the sum-check $4m$ entries of $\tau^{(t)}$, plus three values of $f$.  Entries have $O(h\log q)=O(h\log h)$ bits.  Before reading the proof, the verifier computes from $\varphi$ and $(t,a)$ everything about
$P^f_t(a)$ that does not involve $f$: $\hat\chi_j,\hat\varsigma_j,R_t$ at the relevant points, which takes $\poly(n)$ time; after reading, it does $\poly(m,h,\log q)=\poly(h)$ field operations.  All positions queried are determined by the random string alone: the verifier draws all its randomness first, so it can read the entries of the
low-degree test and of the sum-check in one batch.  This is the point of using sum-check \emph{tables} rather than the interactive protocol.
\end{proof}

\begin{remark}
Note how the verifier avoids a circularity.  The sum-check needs $\tilde f$ at three points, but only $f$ is available.  The three points are \emph{uniform}, hence outside the $1\%$ of disagreement with probability $\ge0.97$.
This is the ``never evaluate at a predictable point'' principle of Section~\ref{sec:poly}, and it is the reason the arithmetization of Lemma~\ref{lem:arith}(c) was designed so that each $w_j$ is uniform.
\end{remark}

Theorem~\ref{thm:vpoly} is already a PCP: choosing $h+1=n^{1/m}$ for a constant $m$ it has $O(1)$ queries.  The catch is the \emph{size} of the queried symbols: $O(h\log h)=n^{\Omega(1/m)}$ bits each.
To get a verifier that reads \emph{bits} one must re-encode every queried symbol and check it without reading it.  That is the job of composition.


%% ===========================================================================
\section{Normal form verifiers and split-encoded assignments}
\label{sec:normal}
%% ===========================================================================

Suppose an \emph{outer} verifier $V_1$ has a decision circuit $C_w$ (on the symbols it queried under random string $w$), of size $T$ much smaller than the proof.
Instead of reading the queried symbols $P[i_1(w)],\dots,P[i_Q(w)]$ in full, we would like to check ``$C_w$ accepts them'' by running a second, \emph{inner} verifier $V_2$ on the formula of $C_w$.
Two requirements:
\begin{itemize}
  \item the inner verifier must work when the \emph{input assignment is spread over several separate pieces}, each of which is held in some encoded form in the proof (the $Q$ queried symbols and a scratch string);
  \item the inner verifier must read only a few \emph{bits} of each piece, since reading the pieces in full is what we want to avoid.
\end{itemize}
This is formalized in the following definition (a simplified form of Definition~3.2 of \citet{AroraSafra98}).

\begin{definition}[Normal form verifier]\label{def:normal}
An \defn{$(r,s,q,t)$-constrained normal form verifier} $V_2$ has an associated encoding scheme $\sigma$ (bit strings to codewords of minimum distance $\delta_{\min}$ over an alphabet of $2^{O(s)}$ symbols).
Given a $3$-CNF $\psi$ with $N$ variables and an integer $p$ dividing $N$, it is given access to a proof $z_1\sharp\cdots\sharp z_p\sharp\pi$, where each $z_i$ is a word over the alphabet and $\pi$ is auxiliary, and must satisfy:
\begin{itemize}
  \item \emph{(completeness)} if each $z_i=\sigma(a_i)$ for $a_i\in\{0,1\}^{N/p}$ and $a_1\circ\dots\circ a_p$ satisfies $\psi$, there is a $\pi$ with acceptance probability~$1$;
  \item \emph{(soundness)} if some $z_i$ is not $\delta_{\min}/3$-close to a codeword, or if each $z_i$ is close to a codeword $\sigma(a_i)$ but $a_1\circ\dots\circ a_p$ does not satisfy $\psi$, then every $\pi$ is accepted with probability $<\tfrac12$;
  \item \emph{(resources)} it uses $O(r(N))$ random bits and alphabet symbols of $O(s(N))$ bits, makes $O(p+q(N))$ non-adaptive queries, and has decision time $O(p\,t(N))$.
\end{itemize}
\end{definition}

\noindent
The unusual feature is the query count $O(p+q(N))$ rather than $O(p\cdot q(N))$: the inner verifier reads $O(1)$ symbols of each piece plus $q(N)$ symbols of the \emph{shared} auxiliary table.
Proposition~3.3 of \citet{AroraSafra98} records the trivial case $p=1$: a normal form verifier gives a verifier for $\sat{3SAT}$ with the same resources.

This notion is the ancestor of the \emph{PCP of proximity} \citep{BGHSV06} and of the \emph{assignment tester} \citep{DR04}, which isolate exactly the property used here---``a tester that looks at a few places of an encoded input
and a proof''---and which are what \citet{Dinur07} composes with in Lecture~6.

\subsection{How $V_{\rm poly}$ becomes a normal form verifier}

By Theorem~\ref{thm:vpoly}, $V_{\rm poly}$ is a verifier with parameters $(\log n,\ h\log h,\ m,\ \poly(h))$ for $r,s,q,t$.  Two additions turn it into a normal form verifier, and
both are instances of the exact-versus-robust pattern.

\paragraph{The encoding.}  Take $\sigma$ to be the polynomial extension: $\sigma\colon\{0,1\}^N\to\mathcal P^{m}_h$ with $N=(h+1)^m$, a code of minimum distance $\ge1-hm/q>\tfrac12$.

\paragraph{Split encodings.}  Suppose $N=(h+1)^m$ and $p=(h+1)^{m-\ell}$, so that each of the $p$ pieces has $N/p=(h+1)^\ell$ bits.  The proof contains $f\colon\F_q^m\to\F_q$ (the encoding of the whole assignment, for the sum-check as before)
and $f_1,\dots,f_p\colon\F_q^\ell\to\F_q$ (the encodings of the pieces), and must check the \defn{concatenation property}: for $i\in\mathbb K^{m-\ell}$ and $u\in\mathbb K^\ell$,
\begin{equation}\label{eq:concat}
  f(i,u)=f_i(u).
\end{equation}

\begin{lemma}[Concatenation]\label{lem:concat}
Let $f\in\mathcal P^m_h$ and $f_1,\dots,f_p\in\mathcal P^\ell_h$ (indexed by $i\in\mathbb K^{m-\ell}$).  Then~\eqref{eq:concat} holds if and only if, as polynomials,
\[
  f(x_1,\dots,x_m)\;=\;\sum_{i\in\mathbb K^{m-\ell}}\Lambda_i(x_1,\dots,x_{m-\ell})\,f_i(x_{m-\ell+1},\dots,x_m),
\]
where $\Lambda_i$ is the Lagrange basis polynomial of Lemma~\ref{lem:ext} for the cube $\mathbb K^{m-\ell}$.
\end{lemma}
\begin{proof}
Let $G$ denote the right-hand side; $G\in\mathcal P^m_h$.  On a cube point $(i',u)$ we have $\Lambda_i(i')=\mathbf 1[i=i']$, so $G(i',u)=f_{i'}(u)$.  If~\eqref{eq:concat} holds, $f$ and $G$ agree on $\mathbb K^m$ and are equal by the
uniqueness in Lemma~\ref{lem:ext}.  Conversely, if $f=G$ then $f(i',u)=G(i',u)=f_{i'}(u)$.
\end{proof}

\noindent
So the verifier tests the identity $f(x)=\sum_i\Lambda_i(x')f_i(x'')$ at $O(1)$ random points $x=(x',x'')\in\F_q^m$.  Each test reads $f(x)$ and one value $f_i(x'')$ for each of the $p$ pieces: $O(1)$ values per piece, $O(p)$ in all,
which is the source of the $O(p+q)$ in Definition~\ref{def:normal}.  If the functions are polynomials but the identity fails, then $f-G$ is a nonzero polynomial and the test catches it with probability $\ge1-hm/q$ per sample (Lemma~\ref{lem:sz}).
Arora and Safra prove the \emph{robust} version (their Theorem~4.4.2): if $f$ is $0.01$-close and the test accepts with probability $>\tfrac14$, then each $f_i$ is $0.03$-close and the concatenation property holds for the nearest polynomials.
The proof rests on a second robust lemma (their Lemma~5.1.1, proved with the same field-size hypothesis as Lemma~\ref{lem:robust}): a function $G$ that is $0.02$-close to a polynomial and all of whose
restrictions $x'\mapsto G(x',x'')$, for fixed $x''$, are polynomials of degree $\le h$, has the property that every restriction $x''\mapsto G(x',x'')$ is $0.03$-close to a polynomial---namely to the corresponding
restriction of the nearest polynomial $\tilde G$.  Applied to $G=\sum_i\Lambda_if_i$ it yields that each $f_i$ is $0.03$-close, with $\tilde f_i=\tilde f(i,\cdot)$.

\begin{pitfall}{decoding must produce bits}
In the definition of a normal form verifier the decoded pieces $a_i=\sigma^{-1}(\text{nearest codeword to }z_i)$ must be \emph{bit strings}, and $\psi$ is evaluated on their concatenation.
By the pitfall after Lemma~\ref{lem:arith}, the arithmetization alone only certifies satisfiability of the rounded assignment; to meet the definition the verifier adds a Booleanity check
($A(A-1)$ vanishing on the cube) to the sum-check, which changes only the constants.
\end{pitfall}


%% ===========================================================================
\section{The composition lemma}
\label{sec:compose}
%% ===========================================================================

Let $V_1$ be a verifier with random bits $R$, alphabet of $2^{S}$ symbols, $Q$ queries, decision time $T$, perfect completeness and soundness $\tfrac14$.  For a random string $w\in\{0,1\}^R$ it queries positions
$i_1(w),\dots,i_Q(w)$ of its proof $P$, and accepts iff a circuit $C_w$ of size $O(T^2)$ accepts $P[i_1(w)]\circ\dots\circ P[i_Q(w)]$.  By the Cook--Levin theorem for circuits this is equivalent to the satisfiability of a $3$-CNF
$c_w$ of size $\tau=O(T^2)$ over the bits of the $Q$ symbols and of a witness string $y_w$.  We append the strings $y_w$ ($w\in\{0,1\}^R$) to $P$ as extra symbols and let $V_1$ read $y_w$ as a $(Q{+}1)$st query.
After this harmless change, writing $i_{Q+1}(w)$ for the position of $y_w$,
\begin{equation}\label{eq:decide}
  V_1\text{ accepts }P\text{ on }w\iff P[i_1(w)]\circ\dots\circ P[i_{Q+1}(w)]\text{ satisfies }c_w .
\end{equation}

\begin{lemma}[Composition; Arora--Safra, Lemma~3.4]\label{lem:compose}
Let $V_1$ be as above, with parameters $(R,S,Q,T)$, and let $V_2$ be an $(r,s,q,t)$-constrained normal form verifier (Definition~\ref{def:normal}).
Then the same language has a verifier $V_{\rm new}$ with perfect completeness and soundness $<\tfrac12$, and with
\[
  \begin{aligned}
  \text{random bits } &R+O(r(\tau)),\qquad &\text{alphabet bits } &O(s(\tau)),\\
  \text{queries } &O\bigl(Q+q(\tau)\bigr),\qquad &\text{decision time } &O\bigl(Q\,t(\tau)\bigr),
  \end{aligned}\qquad\tau=O(T^2).
\]
\end{lemma}

\begin{figure}[t]
\centering
\begin{tikzpicture}[x=1cm,y=1cm,every node/.style={font=\small},
   wd/.style={draw,minimum width=0.62cm,minimum height=0.62cm,inner sep=0pt,fill=white},
   hl/.style={wd,fill=accent!30},
   hy/.style={wd,fill=warn!35},
   hp/.style={wd,fill=claim!22}]
  \foreach \i in {0,...,13} \node[wd] at (\i*0.66,0) {};
  \node[hl] at (1*0.66,0) {};  \node[hl] at (5*0.66,0) {};  \node[hl] at (9*0.66,0) {};
  \node[hy] at (12*0.66,0) {};
  \foreach \i in {15,...,19} \node[wd] at (\i*0.66,0) {};
  \node[hp] at (17*0.66,0) {};
  \node[anchor=east] at (-0.5,0) {$P_{\rm new}$};
  \draw[decorate,decoration={brace,amplitude=3pt,mirror}] (-0.28,-0.5) -- (13*0.66+0.28,-0.5)
        node[midway,below=4pt,font=\footnotesize] {$Y$ words: $\sigma(P[j])$ and $\sigma(y_w)$};
  \draw[decorate,decoration={brace,amplitude=3pt,mirror}] (15*0.66-0.28,-0.5) -- (19*0.66+0.28,-0.5)
        node[midway,below=4pt,font=\footnotesize] {$2^R$ words $\pi_w$};
  % verifiers
  \node[draw,rounded corners=3pt,fill=black!5,minimum width=3.2cm,minimum height=0.8cm] (V1) at (3.0,2.4) {$V_1$: picks $w$, forms $c_w$};
  \node[draw,rounded corners=3pt,fill=black!5,minimum width=4.2cm,minimum height=0.8cm] (V2) at (9.6,2.4) {$V_2$ on $c_w$ (normal form)};
  \draw[-{Stealth[length=2mm]}] (V1)--(V2) node[midway,above,font=\footnotesize] {$c_w$};
  \foreach \x in {1,5,9,12} \draw[accent,-{Stealth[length=1.6mm]}] (V2.south) to[out=-110,in=90] (\x*0.66,0.35);
  \draw[claim,-{Stealth[length=1.6mm]}] (V2.south) to[out=-70,in=90] (17*0.66,0.35);
\end{tikzpicture}
\caption{The composed verifier $V_{\rm new}$ (Lemma~\ref{lem:compose}).  The proof consists of encodings $\sigma(P[j])$ of the symbols of $V_1$'s proof (including the witness strings $y_w$) and one auxiliary word
$\pi_w$ per random string $w$ of $V_1$.  On random string $w$ the verifier $V_1$ would read the $Q$ blue words and the orange word in full; $V_{\rm new}$ instead hands them to $V_2$ as a split-encoded assignment to $c_w$
and reads only a few symbols of each.}
\label{fig:compose}
\end{figure}

\begin{proof}
Replacing $V_1$ and $V_2$ by two independent repetitions of themselves we may assume that both have soundness $\tfrac14$ (this affects only constants).
Let $\sigma$ be the encoding of $V_2$ and write $Y$ for the number of symbols of $V_1$'s (extended) proof.

\emph{The verifier.}  $V_{\rm new}$ expects a proof $P_{\rm new}$ of $Y+2^R$ words; the word $P_{\rm new}[j]$ for $j\le Y$ is supposed to be $\sigma(P[j])$, and $P_{\rm new}[Y+w]$ is an auxiliary proof $\pi_w$.
On input $x$ it (1)~chooses $w\in\{0,1\}^R$ and computes $c_w$ and $i_1(w),\dots,i_{Q+1}(w)$ without reading the proof; (2)~chooses a random string $w'$ for $V_2$ and runs $V_2$ on $c_w$ with the
$(Q{+}1)$-part split-encoded assignment $z_j=P_{\rm new}[i_j(w)]$ ($1\le j\le Q+1$) and auxiliary proof $\pi_w=P_{\rm new}[Y+w]$.  It reads only those symbols of the $z_j$ and of $\pi_w$ that $V_2$ queries under $w'$, and accepts iff $V_2$ does.

\emph{Resources.}  The random bits are those of $w$ and $w'$.  Everything else is that of $V_2$ run with $p=Q+1$ parts on a formula of size~$\tau$: alphabet $O(s(\tau))$, queries $O(p+q(\tau))$, decision time $O(p\,t(\tau))$.
The queries depend only on $(w,w')$, so $V_{\rm new}$ is non-adaptive; this is why non-adaptivity was built into every definition.

\emph{Completeness.}  If $x\in L$ take a proof $P$ with $V_1$ accepting on every $w$, and put $P_{\rm new}[j]=\sigma(P[j])$.  By~\eqref{eq:decide} the concatenation of the $P[i_j(w)]$ satisfies $c_w$, so by completeness of $V_2$ there is $\pi_w$ accepted with
probability~$1$; place it at $Y+w$.

\emph{Soundness.}  Let $x\notin L$ and suppose $P_{\rm new}$ is accepted with probability $\ge\tfrac12$.  Define $P[j]=\sigma^{-1}(N(P_{\rm new}[j]))$, where $N$ is the nearest codeword (an arbitrary string if it is undefined).
Let $p_1=\Pr_w[V_1\text{ accepts }P\text{ on }w]<\tfrac14$.  If $V_1$ rejects $P$ on $w$ then by~\eqref{eq:decide} the decoded concatenation $P[i_1(w)]\circ\dots\circ P[i_{Q+1}(w)]$ does not satisfy $c_w$, so
either some $z_j$ is not close to a codeword or the decoded assignment is unsatisfying.  By the soundness of $V_2$, $V_{\rm new}$ accepts on such $w$ with probability $<\tfrac14$, \emph{whatever $\pi_w$ is}.  Hence
\[
  \Pr[V_{\rm new}\text{ accepts}]\ <\ p_1+(1-p_1)\cdot\tfrac14\ <\ \tfrac14+\tfrac34\cdot\tfrac14\ =\ \tfrac7{16}\ <\ \tfrac12,
\]
a contradiction.
\end{proof}

\begin{pitfall}{what composition does \emph{not} do}
Composition does not make a \emph{hard} instance easier; it makes the \emph{queries} cheaper by moving the cost into the proof length and the verifier's computation.
Also, the lemma needs $V_2$ to \emph{check an encoded input}.  A verifier with the same parameters that must read its input in the clear would give nothing: the whole point is that $V_{\rm new}$ never reads a $z_j$ in full.
This is the reason normal form (Section~\ref{sec:normal}), and not just a verifier with small parameters, is the right interface.
\end{pitfall}


%% ===========================================================================
\section{The parameter chase}
\label{sec:chase}
%% ===========================================================================

We now put the pieces together.  Throughout, $n$ is the formula size, $L=\log n$, and all logarithms are base~$2$.

\begin{enumerate}[label=\textbf{Step \arabic*.},leftmargin=*]
  \item \emph{Outer verifier.}  Apply Theorem~\ref{thm:vpoly} with $h+1=2^{\sqrt L}$ and $m=\sqrt L$ ($h\ge\log n$ holds for large $n$).  Call it $V^{(0)}$:
        \[
          (R,S,Q,T)\ =\ \bigl(O(L),\ 2^{O(\sqrt L)},\ \sqrt L,\ 2^{O(\sqrt L)}\bigr).
        \]
  \item \emph{Inner verifier.}  By Section~\ref{sec:normal} (the concatenation test and the Booleanity check added), the same verifier with the same parameter choice is, at input size $N$, an
        $(r,s,q,t)$-constrained normal form verifier with $(r,s,q,t)(N)=(\log N,\ 2^{O(\sqrt{\log N})},\ \sqrt{\log N},\ 2^{O(\sqrt{\log N})})$.
  \item \emph{Compose with $V^{(0)}$ as inner verifier.}  Apply Lemma~\ref{lem:compose}.  Since $T=2^{O(\sqrt L)}$, $\tau=O(T^2)$ has $\log\tau=O(\sqrt L)$, and
        \[
          \begin{aligned}
          \text{random bits}&\;:\; L+O(\log\tau)=O(L),\\
          \text{alphabet bits}&\;:\; 2^{O(\sqrt{\log\tau})}=2^{O(L^{1/4})},\\
          \text{queries}&\;:\; O\bigl(\sqrt L+\sqrt{\log\tau}\bigr)=O(\sqrt L),\\
          \text{decision time}&\;:\; \sqrt L\cdot2^{O(\sqrt{\log\tau})}=2^{O(L^{1/4})}.
          \end{aligned}
        \]
  \item \emph{Iterate.}  Each further composition with $V^{(0)}$ replaces the exponent $e$ in $2^{O(L^{e})}$ by $e/2$ and leaves queries and randomness unchanged up to constant factors.  After $k$ compositions the verifier has
        alphabet bits and decision time $2^{O(L^{2^{-(k+1)}})}$, $O(\sqrt L)$ queries, $O(L)$ random bits.
  \item \emph{Final composition with $V_{\rm savior}$.}  Apply Theorem~\ref{thm:vpoly} with $h=\Theta(\log N)$ and $m=\log N/\log\log N$ at input size $N$: it is a normal form verifier with
        $(r,s,q,t)(N)=(\log N,\ \log N\log\log N,\ \log N/\log\log N,\ \poly\log N)$.  Composing with it, where $\log\tau=O(L^{2^{-(k+1)}})$, gives
        \begin{align*}
          \text{alphabet bits}&=\log\tau\log\log\tau=L^{2^{-(k+1)}+o(1)},\\
          \text{queries}&=O(\sqrt L),\\
          \text{decision time}&=\sqrt L\cdot\poly(\log\tau)=\poly(L).
        \end{align*}
\end{enumerate}

\begin{table}[t]
\centering\small
\renewcommand{\arraystretch}{1.25}
\begin{tabular}{@{}lllll@{}}
\toprule
verifier & random bits & queries (symbols) & bits per symbol & decision time\\
\midrule
$V^{(0)}$ (Step 1) & $O(L)$ & $\sqrt L$ & $2^{O(L^{1/2})}$ & $2^{O(L^{1/2})}$\\
$V^{(1)}=V^{(0)}\circ V^{(0)}$ & $O(L)$ & $O(\sqrt L)$ & $2^{O(L^{1/4})}$ & $2^{O(L^{1/4})}$\\
$V^{(2)}=V^{(1)}\circ V^{(0)}$ & $O(L)$ & $O(\sqrt L)$ & $2^{O(L^{1/8})}$ & $2^{O(L^{1/8})}$\\
$V^{(k)}$ & $O(L)$ & $O(\sqrt L)$ & $2^{O(L^{2^{-(k+1)}})}$ & $2^{O(L^{2^{-(k+1)}})}$\\
$V^{(k)}\circ V_{\rm savior}$ & $O(L)$ & $O(\sqrt L)$ & $L^{2^{-(k+1)}+o(1)}$ & $\poly(L)$\\
\bottomrule
\end{tabular}
\caption{The parameter chase of Section~\ref{sec:chase} ($L=\log n$).  Each composition square-roots the exponent of the alphabet size; the number of queries is stuck at $\sqrt L$ because it is inherited from
the outer verifier.}
\label{tab:chase}
\end{table}

\begin{proof}[Proof of Theorem~\ref{thm:as}]
Given $\eps>0$ choose $k$ with $2^{-(k+1)}\le\eps/2$, a constant.  The verifier of Table~\ref{tab:chase} after $k$ compositions and the final one with $V_{\rm savior}$ uses $O(L)$ random bits,
makes $O(\sqrt L)$ queries to symbols of $L^{\eps/2+o(1)}$ bits each, so reads $L^{1/2+\eps/2+o(1)}\le L^{1/2+\eps}$ bits of the proof, has perfect completeness and soundness $\tfrac12$.
Hence $\sat{3SAT}\in\PCP[\log n,\log^{1/2+\eps}n]$ and so $\NP\subseteq\PCP[\log n,\log^{1/2+\eps}n]$.  The reverse inclusion was noted at the beginning of the lecture.
\end{proof}

\noindent
Each step of the chase hides a constant in $\tau=O(T^2)$ and in $O(Q+q(\tau))$; the next section is about why that matters.

\begin{bigpicture}{What is proved here and what is quoted}
\textbf{Proved in this lecture, from scratch:} Schwartz--Zippel, polynomial extension, the arithmetization (Lemma~\ref{lem:arith}) and zero-testers (Lemma~\ref{lem:zt}), the sum-check (Lemma~\ref{lem:sumcheck}), the exact low-degree
statement (Lemma~\ref{lem:axis}), the concatenation identity (Lemma~\ref{lem:concat}), the composition lemma (Lemma~\ref{lem:compose}), and the whole parameter chase.
\textbf{Quoted:} the two \emph{robust} algebraic lemmas, Lemma~\ref{lem:robust} (consistency of a decomposition) and the robust concatenation theorem (Theorem~4.4.2 of \citealp{AroraSafra98}), both of which rely on
the bivariate lemma of Arora and Safra and on the cubic field size.  Everything else in the text is bookkeeping around them.
\end{bigpicture}


%% ===========================================================================
\section{\texorpdfstring{Consequence: Clique within $2^{\log^{1/2-\eps}N}$ is $\NP$-hard}{Consequence: hardness of approximating Clique}}
\label{sec:clique}
%% ===========================================================================

\begin{proof}[Proof of Corollary~\ref{cor:as-clique}]
Let $0<\eps_1<\eps$.  By Theorem~\ref{thm:as} there is a verifier $V$ for $\sat{3SAT}$ with $O(\log n)$ random bits and $q=\log^{1/2+\eps_1}n$ queries, perfect completeness and soundness $\tfrac12$.  Amplify it by
the random walk on an explicit constant-degree expander on the set of random strings (Exercise~\ref{ex:expander}; \citealp{AKS87}): $k=\lceil\log n/q\rceil$ steps use $O(\log n)+O(k)=O(\log n)$ random bits, make $kq=O(\log n)$ queries, and
give soundness $2^{-\Omega(k)}=2^{-\Omega(\log^{1/2-\eps_1}n)}$.  The FGLSS graph $H_x$ of the amplified verifier (Definition~\ref{def:fglss}) has at most $2^{O(\log n)}=n^{O(1)}=:N$ vertices, and by Theorem~\ref{thm:fglss}
\[
  x\in\sat{3SAT}\Rightarrow\omega(H_x)=2^{r},\qquad x\notin\sat{3SAT}\Rightarrow\omega(H_x)\le2^{r}\cdot2^{-\Omega(\log^{1/2-\eps_1}n)} .
\]
The gap is $2^{\Omega(\log^{1/2-\eps_1}N)}$, which exceeds $2^{c\log^{1/2-\eps}N}$ for every constant $c$ and large $N$ because $\eps_1<\eps$.  An approximation algorithm within $2^{c\log^{1/2-\eps}N}$ would therefore
decide $\sat{3SAT}$ in polynomial time.
\end{proof}

\noindent
Compare with Corollary~\ref{cor:clique-quasi}, which reached $2^{\Omega((\log N)^{1/(1+\eps)})}$ from a constant-query PCP and quasi-polynomial time.  Here the PCP has $\log^{1/2+\eps_1}n$ queries, but the amplification is
randomness-efficient, so the reduction takes \emph{polynomial} time: \emph{this} is the improvement that earned the paper its title.


%% ===========================================================================
\section{\texorpdfstring{Why it stops at $\log^{1/2+\eps}n$, and how to reach $O(1)$}{Why it stops at sublogarithmic query bits, and how to reach O(1)}}
\label{sec:as-limits}
%% ===========================================================================

The chase of Table~\ref{tab:chase} has two bottlenecks.

\paragraph{Queries are inherited, never reduced.}  By Lemma~\ref{lem:compose} the new number of queries is $O(Q+q(\tau))\ge Q$.  Composition shrinks the \emph{symbols} but never the number of symbols queried,
and $V_{\rm poly}$ has $Q=m=\log n/\log(h+1)$.  To have few queries one needs $h$ large, hence symbols of $\Omega(h)$ bits, hence \emph{many} compositions to make them short.

\paragraph{Constants compound.}  With $h+1=n^{1/m}$ and $m$ constant, the alphabet has $n^{\Omega(1/m)}$ bits and shrinking it to $\polylog n$ takes $\Theta(\log\log n)$ compositions, each
multiplying the hidden constants in $O(Q+q(\tau))$ by a factor $c>1$.  After $\Theta(\log\log n)$ rounds the product is $c^{\Theta(\log\log n)}=\polylog n$ queries.  This is exactly the caution in
the paper: ``We do not wish to invoke the Composition Lemma more than $O(1)$ times because it hid constant factors in its $O(\cdot)$ notation and those constant factors could grow very fast.''  The $\log^{1/2+\eps}n$ of Theorem~\ref{thm:as}
is what survives when one limits oneself to $O(1)$ compositions.

\paragraph{The remedy of ALMSS.}  According to Section~6 of \citet{AroraSafra98}, \citet{ALMSS98} proceed in three steps.
\begin{enumerate}
  \item They construct a new outer verifier that reads only $O(1)$ entries of the polynomial tables (not $O(m)$) and has polylogarithmic decision time.  The two ingredients are a more efficient low-degree test, based on the
        main lemma of Arora--Safra (the bivariate Lemma~\ref{lem:robust}'s base case) and on the work of \citet{RubinfeldSudan92}, and a \emph{parallelization} step in the spirit of \citet{LapidotShamir91} and \citet{FeigeLovasz92}
        (a close relative of the constraint--variable game of Lecture~\ref{lec:games}, in which the answers to all queries of a random string are bundled into one symbol), which collapses several queries into one.
  \item They compose this verifier with itself.  The result reads $O(1)$ symbols of size $O(\log\log n)$.
  \item They compose that with a verifier, built from ideas of \citet{BLR90}, which is \emph{wasteful in randomness} (polynomial in the proof size) but makes only $O(1)$ \emph{bit} queries.  Since by now the input
        to the last verifier has size $\polylog n$, its wastefulness is harmless.  The result is a $(\log n,1)$-restricted verifier.
\end{enumerate}
This is Lecture~4.  The last of the three verifiers is based on the \emph{Hadamard code} and linearity testing \citep{BLR90}; the \emph{long code} of \citet{BGS98}, the workhorse of Lecture~7, is a more redundant relative of the Hadamard code.

\begin{bigpicture}{The shape of the argument, in one paragraph}
A satisfying assignment is encoded by a low-degree polynomial, a code in which a $1\%$-corrupted word still determines the codeword.  Satisfiability becomes one vanishing sum, which the sum-check verifies at a single random
point of a polynomial that is itself computed from three random points of the table.  A low-degree test certifies that the table is close to a polynomial.  Because the verifier only ever asks for \emph{uniformly distributed}
points, corruption of a small fraction of the table is invisible to it.  This yields a verifier with $O(\log n)$ random bits and few, but long, queries.  Composition then re-checks each long symbol \emph{in encoded form} with a second such verifier, trading
symbol length for proof length.  Iterating, one reaches sublogarithmic query bits; getting to $O(1)$ needs a better outer verifier and a Hadamard-code inner verifier (Lecture~4), or Dinur's entirely different amplification (Lecture~6).
\end{bigpicture}


%% ===========================================================================
\section*{Exercises}
\addcontentsline{toc}{section}{Exercises}
%% ===========================================================================

\begin{exercise}[\stars{$\star\star$}]\label{ex:sz-tight}
Show that Lemma~\ref{lem:sz} is nearly tight: for $d<q$ and every $m$ there is a nonzero $p\in\mathcal P^m_d$ vanishing on exactly a $1-(1-d/q)^m$ fraction of $\F_q^m$.  Compare $1-(1-d/q)^m$ with $dm/q$.
\end{exercise}

\begin{exercise}[\stars{$\star$}]\label{ex:maj}
Compute the multilinear extension $\hat s$ of the majority function $s\colon\{0,1\}^3\to\{0,1\}$ and evaluate it at $(2,2,2)$ over $\F_7$.  Verify that $\hat s$ has degree $1$ in each variable but total degree $3$.
\end{exercise}

\begin{exercise}[\stars{$\star$}]\label{ex:arith2}
In Example~\ref{ex:arith} take instead $a=(1,1,1,1)$.  Which clause is violated?  Find all nonzero points of $g_{\hat a}$ on $\mathbb K^{8}$ and their values, and compute $\Pr_t[F(t)=0]$ exactly as in Example~\ref{ex:zt}.
\end{exercise}

\begin{exercise}[\stars{$\star\star$}]\label{ex:sc-tight}
Let $B\equiv0$ be the zero polynomial in $\ell$ variables, $H\subseteq\F_q$, $c\ne0$.  Show that the best table convinces the sum-check with probability \emph{exactly} $1-(1-d/q)^\ell$ provided that there is a polynomial
of degree $\le d$ with $d$ distinct roots in $\F_q$ whose sum over $H$ is $c$.  Exhibit one for $q=7$, $d=2$, $H=\{0,1\}$, $c=1$.  (\emph{Hint: scaling a polynomial by a nonzero constant keeps its roots and scales its sum over $H$.})
\end{exercise}

\begin{exercise}[\stars{$\star$}]\label{ex:onedir}
Show that Lemma~\ref{lem:axis} fails if lines of one direction are omitted: for $m=2$ let $f(x_1,x_2)=\varphi(x_1)$ where $\varphi\colon\F_q\to\F_q$ is a function of degree greater than $h$.  Every row $f(x_1,\cdot)$ is constant,
but $f\notin\mathcal P^2_h$.  Prove that $\Delta(f,\mathcal P^2_h)=\Delta(\varphi,\mathcal P^1_h)$, so $f$ can be as far from $\mathcal P^2_h$ as a univariate function can be from $\mathcal P^1_h$.
\emph{Hint: for a fixed polynomial $p$ average over $x_2$.}
\end{exercise}

\begin{exercise}[\stars{$\star\star$}]\label{ex:chase}
Suppose an inner verifier has $s(N)=t(N)=N^{1/3}$, $q(N)=O(1)$, $r(N)=\log N$ and the outer verifier has $(R,S,Q,T)=(O(\log n),n^{1/3},c_0,n^{1/3})$.  Compute the parameters after $k$ compositions
and find the least $k$ for which the alphabet is $\polylog n$.  If each composition multiplies the number of queries by a constant $c>1$, what is the final number of queries?  Compare with Step~4 of Section~\ref{sec:chase}.
\end{exercise}

\begin{exercise}[\stars{$\star\star$}]\label{ex:bool}
Design the Booleanity check mentioned in the pitfalls of Sections~\ref{sec:arith} and~\ref{sec:normal}: give a polynomial in $\mathcal P^{m}_{2h}$ that vanishes on $\mathbb K^m$ iff $A$ is $\{0,1\}$-valued
there, and show how a second sum-check (with its own zero-tester, $k=m$ variables) detects a violation with probability at least $1-3hm/q$.  How many extra queries does the verifier make?
\end{exercise}

\begin{exercise}[\stars{$\star\star\star$}]\label{ex:reading3}
\emph{Reading.}  Compare the low-degree test of Section~4.6 of \citet{AroraSafra98} with the one analyzed in Section~7 of \citet{ALMSS98}: what is tested (axis-parallel lines or random lines), how many queries
are made, how large the field must be, and which of the two lemmas is the bivariate base case.  Which of these differences is responsible for the drop from $O(m)$ to $O(1)$ queries?
\end{exercise}


%% ===========================================================================
\section*{Notes and further reading}
\addcontentsline{toc}{section}{Notes and further reading}
%% ===========================================================================

\paragraph{Where this lecture comes from.}
Everything is from \citet{AroraSafra98}: Sections~3 and~4 give the verifier and the composition lemma, Section~5 the algebraic lemmas that we quoted, and Section~6 the account of what followed.  The algebraic characterization of satisfying polynomials (Lemma~\ref{lem:arith}) is credited there to \citet{BFLS91}, in a form similar in spirit to \citet{BFL91}, and the zero-testers (Lemma~\ref{lem:zt}) to
\citet{BFLS91} and \citet{FGLSS96}; the sum-check is due to \citet{LFKN92}.  The names ``zero-tester'' and ``table of partial sums'', the form of the composition lemma
and our parameter table follow \citet{AroraSafra98}; the numerical examples (Examples~\ref{ex:arith}, \ref{ex:zt}, \ref{ex:sc}) were computed for these notes.

\paragraph{On attributions.}
We distinguish three statements that are often confused: Arora--Safra proved sub-logarithmic queries, Arora--Lund--Motwani--Sudan--Szegedy constant queries.  Section~6 of \citet{AroraSafra98} is generous in crediting the second
paper and the earlier observation of Szegedy and Sudan that a verifier with small decision time yields a gap for \sat{Max-3SAT}.  The modular formulation of composition (PCPs of proximity, assignment testers) is
due to \citet{BGHSV06} and \citet{DR04}; Arora and Safra's normal form verifier is the special case used here.  The analysis of the low-degree test was subsequently improved by \citet{RubinfeldSudan92},
\citet{PolishchukSpielman94} and \citet{AroraSudan97}; none of these is needed for what we do.

\paragraph{A caveat on the text.}
Two small points in the presentation of \citet{AroraSafra98} that we were careful with: the arithmetization certifies rounded satisfiability, so Booleanity must be checked separately when decoding
(pitfall after Lemma~\ref{lem:arith}); and in the remark after Lemma~4.6.4 the paper writes $\gamma^{2(l-1)}\ge\tfrac15$, which is true but too weak for the conclusion drawn from it for $l>5$; the correct bound is $\ge\tfrac45$ (Bernoulli's inequality), and
it is this that makes the choice of $8m/\delta$ checkpoints in the test sufficient.

\paragraph{What is not covered.}
We have not proved the robust lemmas, nor the analysis of $V_{\rm savior}$ at small input size, and we have not treated the $\PCP[\log n,(\log\log n)^2]$ refinement mentioned in the paper.
The proof of the constant-query theorem is Lecture~4, and the combinatorial alternative is Lecture~6.

\chapter{Constant Queries: Hadamard Codes, Parallelization and the ALMSS Proof}
\label{lec:alm}

{\sffamily\large\color{softgray}Reorder the compositions, and a wasteful verifier finishes the job\par}
\vspace{0.5em}

\noindent
Lecture~\ref{lec:as} ended with a verifier that reads $\log^{1/2+\eps}n$ bits.  Its bottleneck was not the
algebra but the bookkeeping of the composition lemma: composition shrinks \emph{symbols} but never the \emph{number}
of symbols read, and the starting verifier reads $\Omega(m)=\Omega(\log n/\log h)$ of them.
\citet{ALMSS98}, writing at the same time as Arora and Safra (both papers appeared at FOCS 1992), found a way
around this by changing the \emph{order} of the operations.  Their theorem is the one that is usually meant by the PCP theorem.

\begin{theorem}[Arora--Lund--Motwani--Sudan--Szegedy]\label{thm:alm}
There is a constant $q$ such that $\NP=\bigcup_{c>0}\PCP\bigl[c\log n,\ q\bigr]$.
\end{theorem}

\noindent
The proof has three ideas, and each gets a section of this lecture.
\begin{enumerate}
  \item \emph{A wasteful verifier with $O(1)$ bit queries.}  Encode the assignment with the \emph{Hadamard code} (all
        parities) and its tensor square.  The proof has length $2^{\Theta(n^2)}$, but a verifier can check it with
        $O(1)$ queries using the linearity test of Blum, Luby and Rubinfeld, self-correction, and a Freivalds-type test
        (Sections~\ref{sec:had-code}--\ref{sec:had-ver}).  Used directly this is useless: the proof is exponential.
        Used as the \emph{innermost} verifier of a composition, on a circuit of size $\polylog\log n$, the exponential is harmless.
  \item \emph{Parallelization.}  Starting from the Arora--Safra verifier, which makes $\polylog n$ queries, collapse
        all its queries into $O(1)$ queries by asking for the restriction of the encoded proof to a \emph{curve}
        through the $\polylog n$ points, at the price of larger answers (Sections~\ref{sec:lines}--\ref{sec:par}).
  \item \emph{Reorder the composition.}  Compose \emph{the parallelized verifier with itself} once (answer size
        $\polylog n\to\polylog\log n$), then with the Hadamard verifier (answer size $\to1$).
        Only two compositions are needed, so the constants do not compound (Section~\ref{sec:assemble}).
\end{enumerate}

\begin{goalbox}{After this lecture you should be able to\ldots}
\begin{enumerate}[leftmargin=*,label=(\arabic*)]
  \item state the outer/inner verifier formalism of ALMSS and prove its composition theorem, and say how it
        relates to the normal form verifiers of Lecture~\ref{lec:as} (Section~\ref{sec:outin});
  \item prove the linearity test of Blum--Luby--Rubinfeld by Fourier analysis, with the sharp bound
        ``rejection probability $\ge$ distance to the nearest linear function'' (Section~\ref{sec:had-code});
  \item prove the soundness of the Freivalds-type tensor test, and assemble the Hadamard verifier, computing its
        number of queries and acceptance probabilities on a $5$-gate circuit (Section~\ref{sec:had-ver});
  \item prove the properties of self-correction on lines and of the curve-check, and explain how a curve turns
        $\polylog n$ queries into one (Sections~\ref{sec:lines}--\ref{sec:curves});
  \item carry out the three-step parameter chase of Theorem~\ref{thm:alm} and say which ingredients we quote rather than prove (Section~\ref{sec:assemble}).
\end{enumerate}
\end{goalbox}


%% ===========================================================================
\section{Outer verifiers, inner verifiers, and composition again}
\label{sec:outin}
%% ===========================================================================

\subsection{The formalism}

Lecture~\ref{lec:as} proved the composition lemma for \emph{normal form verifiers}.  ALMSS use a formalism that
separates the two roles more cleanly and states soundness in terms of decoding.  Both are the same idea, and we shall see
precisely how they correspond.

\begin{definition}[Outer verifier]\label{def:outer}
An \defn{$(r,p,c,a)$ outer verifier} $V$ for inputs of length $n$ expects as proof an oracle $\pi$ whose entries are strings
of $a(n)$ bits.  On input $x$ it (i)~draws $R\in\{0,1\}^{r(n)}$; (ii)~computes in time $\poly(n)$ a circuit $C_R$ of size $c(n)$ and
$p(n)$ positions $q_1,\dots,q_{p(n)}$; (iii)~reads $a_i=\pi(q_i)$; (iv)~accepts iff $a_1\circ\cdots\circ a_p$ satisfies $C_R$.
A language $L$ is in $\mathrm{RPCP}(r,p,c,a,e)$ if such a verifier accepts some oracle with probability~$1$ when $x\in L$ and every oracle with probability at
most~$e$ when $x\notin L$.
\end{definition}

\noindent
Since the verifier reads $p(n)\,a(n)$ bits, $\mathrm{RPCP}(r,p,c,a,\tfrac12)\subseteq\PCP[r,\,pa]$.  What distinguishes an outer verifier from an arbitrary verifier is
\emph{succinctness}: its decision is a small circuit evaluated on the concatenated answers.  This is the interface through which an inner verifier is plugged in.

An inner verifier must check ``the concatenation of the queried answers satisfies $C$'' \emph{without reading the answers}, so these must be given in an encoded form.

\begin{definition}[Encoding scheme]\label{def:enc}
An \defn{$(\ell,a)$ encoding scheme} is a family $E=\{E_m\}$ of maps $E_m\colon\{0,1\}^m\to\{f\colon[\ell(m)]\to\{0,1\}^{a(m)}\}$ together with a
\defn{decoding} $E^{-1}_m$ in the other direction with $E^{-1}_m(E_m(x))=x$ for all $x$.  (Decoding is arbitrary on words that are not codewords.)
\end{definition}

\begin{definition}[Inner verifier system]\label{def:inner}
A \defn{$(k;r,p,c,a;e)$ inner verifier system} is a triple $(V,E,E^{-1})$.  The input is a circuit $C$ of size $n$ with $km$ input nodes; the verifier expects a proof
consisting of $k+1$ oracles $X_1,\dots,X_k$ and $\pi$, each with answer size $a(n)$.  It draws $r(n)$ random bits, computes a circuit $C'$ of size $c(n)$ and
$p(n)$ positions in the oracles, reads them, and accepts iff the concatenated answers satisfy $C'$.
\begin{itemize}
  \item \emph{Completeness.}  If $x_1\circ\cdots\circ x_k$ ($x_i\in\{0,1\}^m$) satisfies $C$, then there is $\pi$ such that $V$ accepts $(E(x_1),\dots,E(x_k),\pi)$ with probability~$1$.
  \item \emph{Soundness.}  If for some oracles $X_1,\dots,X_k,\pi$ the verifier accepts with probability $>e$, then $E^{-1}(X_1)\circ\cdots\circ E^{-1}(X_k)$ satisfies $C$.
\end{itemize}
\end{definition}

\noindent
Compare with Definition~\ref{def:normal}: the $X_i$ are the pieces $z_i$ of the split-encoded assignment, $\pi$ is the auxiliary proof, and soundness here is the
conjunction of the two soundness clauses there (``some piece is far from a codeword'' and ``the decoded assignment does not satisfy'') phrased through the \emph{decoding} map, whose behaviour on non-codewords is
part of the scheme.  The number of queries is $p$, in place of $O(p+q)$.

\begin{proposition}[An inner verifier is an outer verifier]\label{prop:in-out}
If a $(k;r,p,c,a;e)$ inner verifier system exists then $\sat{CKT-SAT}\in\mathrm{RPCP}(r,p,c,a,e)$.
\end{proposition}
\begin{proof}
Concatenate $X_1,\dots,X_k,\pi$ into one oracle.  If $C$ is satisfiable, completeness gives a proof accepted with probability~$1$.  If a proof is accepted with probability $>e$, soundness shows that
$C$ is satisfiable.  Hence for unsatisfiable $C$ every proof is accepted with probability $\le e$.
\end{proof}

\begin{proposition}[Repetition]\label{prop:rep}
From a $(k;r,p,c,a;e)$ inner verifier system one obtains, for every $l\ge1$, a $(k;lr,lp,lc,a;e^l)$ inner verifier system.
\end{proposition}
\begin{proof}
Run $V$ $l$ times with independent random strings and accept iff all runs accept.  For fixed oracles the runs are independent, so the acceptance probability is the $l$-th power of that of a single run:
if it is $\ge e^l$ then a single run accepts with probability $\ge e$, and soundness of $V$ applies.  Completeness is unchanged.  (ALMSS write $la$ for the answer size; $a$ suffices, since the same oracles are queried.)
\end{proof}

\subsection{Composition}

\begin{theorem}[Composition; ALMSS Theorem~13, after Arora--Safra]\label{thm:compose}
Suppose a $(p;r_1,p_1,c_1,a_1;e_1)$ inner verifier system exists.  Then for all functions $r,c,a$ and $0<e<1$,
\[
  \mathrm{RPCP}\bigl(r,\,p,\,c,\,a,\,e\bigr)\ \subseteq\ \mathrm{RPCP}\bigl(r+r_1(c),\ p_1(c),\ c_1(c),\ a_1(c),\ e+e_1-ee_1\bigr).
\]
\end{theorem}

\begin{proof}
Let $V_1$ be an $(r,p,c,a)$ outer verifier for $L$ and $(V_2,E,E^{-1})$ the inner system.  Fix $x$ of length $n$ and let $Y$ be the number of entries of $V_1$'s oracle.
The new verifier $V$ expects an oracle $\pi$ that consists of $Y+2^{r(n)}$ \emph{sub-oracles}: $\pi[1],\dots,\pi[Y]$ (supposedly $E(\pi'(1)),\dots,E(\pi'(Y))$ for a proof $\pi'$ of $V_1$) and $\pi[Y+R]$ for every $R\in\{0,1\}^{r(n)}$ (an auxiliary proof for the run $R$).

\emph{Program.}  Draw $R\in\{0,1\}^{r(n)}$ and simulate steps (i)--(ii) of $V_1$, obtaining a circuit $C_R$ of size $c(n)$ and positions $Q_1,\dots,Q_p$.  Draw $R_2\in\{0,1\}^{r_1(c(n))}$ and run $V_2$ on input $C_R$ with random string $R_2$ and oracles
$X_i=\pi[Q_i]$ ($i\le p$) and $\pi_{\rm aux}=\pi[Y+R]$.  Accept iff $V_2$ accepts.  The number of random bits is $r(n)+r_1(c(n))$, and all other parameters are those of $V_2$ on an input of size $c(n)$; in particular the verdict is a circuit of size $c_1(c(n))$ of the $p_1(c(n))$ answers.

\emph{Completeness.}  Let $\pi'$ be accepted by $V_1$ with probability~$1$ and set $\pi[j]=E(\pi'(j))$.  For each $R$ the answers $\pi'(Q_1)\circ\cdots\circ\pi'(Q_p)$ satisfy $C_R$, so by completeness of $V_2$ there is an auxiliary proof for $R$ which we place in $\pi[Y+R]$; then $V$ accepts with probability~$1$.

\emph{Soundness.}  Suppose $x\notin L$ and $V$ accepts some $\pi$ with probability $>e+e_1-ee_1$.  Define $\pi'(j)=E^{-1}(\pi[j])$ for $j\le Y$.  Let
$\mathcal R=\{R:\Pr_{R_2}[V\text{ accepts on }(R,R_2)]>e_1\}$ and $\rho=|\mathcal R|/2^{r(n)}$.  Then $\Pr[V\text{ accepts}]\le\rho+(1-\rho)e_1$, which is increasing in $\rho$; if $\rho\le e$ this is at most $e+(1-e)e_1=e+e_1-ee_1$, a contradiction.
So $\rho>e$.  For $R\in\mathcal R$, the inner verifier $V_2$ accepts the oracles $\pi[Q_1],\dots,\pi[Q_p],\pi[Y+R]$ with probability $>e_1$, so by its soundness
$E^{-1}(\pi[Q_1])\circ\cdots\circ E^{-1}(\pi[Q_p])=\pi'(Q_1)\circ\cdots\circ\pi'(Q_p)$ satisfies $C_R$, \ie $V_1$ accepts $\pi'$ on $R$.  Hence $V_1$ accepts $\pi'$ with probability $>e$, contradicting soundness of $V_1$.
\end{proof}

\noindent
Two observations.  First, the soundness arithmetic is slightly different from Lemma~\ref{lem:compose} ($e+e_1-ee_1$ versus $p_1+(1-p_1)\tfrac14$); both say ``the cheater either fools the outer verifier or fools the inner verifier on the runs where the outer verifier rejects''.  Second, the roles of the parameters:
the number of queries of the composed verifier is $p_1(c)$, that of the \emph{inner} verifier on the \emph{small} circuit.  This is what changes the game: if the \emph{inner} verifier has $O(1)$ queries,
so does the composition, whatever the outer verifier's query count $p$ was.  In Lecture~\ref{lec:as} the inner verifiers were algebraic verifiers whose query count $q(\tau)$ grew with $\tau$ (as $\sqrt{\log\tau}$, and as $\log\tau/\log\log\tau$ for $V_{\rm savior}$).

\begin{figure}[t]
\centering
\begin{tikzpicture}[x=1cm,y=1cm,every node/.style={font=\small},
   st/.style={draw,rounded corners=3pt,minimum width=3.15cm,minimum height=1.35cm,align=center,fill=accent!7,font=\footnotesize},
   ar/.style={-{Stealth[length=2.2mm]},thick,accent!80!black}]
  \node[st] (a) at (0,0)    {\textbf{Arora--Safra}\\ $\log^{c_2}n$ queries\\ $\log^{c_3}n$-bit answers};
  \node[st] (b) at (3.8,0)  {\textbf{parallelized}\\ $O(1)$ queries\\ $\polylog n$-bit answers};
  \node[st] (c) at (7.6,0)  {\textbf{composed with}\\ \textbf{itself}\\ $O(1)$ queries\\ $\polylog\log n$-bit answers};
  \node[st,fill=warn!12] (d) at (11.4,0) {\textbf{composed with}\\ \textbf{Hadamard}\\ $O(1)$ queries\\ $1$-bit answers};
  \draw[ar] (a)--(b);
  \draw[ar] (b)--(c);
  \draw[ar] (c)--(d);
  \node[font=\footnotesize,accent!80!black] at (1.9,-1.05) {\S\ref{sec:par}};
  \node[font=\footnotesize,accent!80!black] at (5.7,-1.05) {Thm~\ref{thm:compose}};
  \node[font=\footnotesize,accent!80!black] at (9.5,-1.05) {Thm~\ref{thm:compose}};
\end{tikzpicture}
\caption{The architecture of the ALMSS proof.  Random bits stay $O(\log n)$ throughout; only the \emph{answer size} and the
\emph{decision circuit} are shrunk by the compositions, and the number of queries is made constant once and for all in the first arrow.}
\label{fig:alm-arch}
\end{figure}


%% ===========================================================================
\section{The Hadamard code and its local tests}
\label{sec:had-code}
%% ===========================================================================

Throughout this section all arithmetic is over $\F_2$ and $z\cdot y=\sum_iz_iy_i$.

\begin{definition}[Hadamard / parity encoding]\label{def:hadamard}
For $z\in\F_2^n$ let $L_z\colon\F_2^n\to\F_2$, $L_z(y)=z\cdot y$.  The \defn{Hadamard encoding} of $z$ is the table of all $2^n$ values of $L_z$.
Decoding maps a function $f\colon\F_2^n\to\F_2$ to the string $z$ minimizing $\Delta(f,L_z)$.
\end{definition}

\noindent
Every linear function $\F_2^n\to\F_2$ is some $L_z$.  For $z\ne z'$ the function $L_z+L_{z'}=L_{z+z'}$ is a nonzero linear function, hence equals~$1$ on exactly half of $\F_2^n$:
the code has relative distance $\tfrac12$, so a function within distance $<\tfrac14$ of a codeword decodes to it.
This is the code with the \emph{worst rate} one could imagine ($2^n$ bits for $n$) and the \emph{best local structure}.

\subsection{Fourier analysis over $\F_2^n$ in one page}

Write a function $F\colon\F_2^n\to\F_2$ as $f=(-1)^F\colon\F_2^n\to\{\pm1\}$, and for $S\in\F_2^n$ let $\chi_S(x)=(-1)^{S\cdot x}$, so that $\chi_S=(-1)^{L_S}$.  Define
\[
  \widehat f(S)\;=\;\E_{x}\bigl[f(x)\chi_S(x)\bigr]\;=\;1-2\Delta(F,L_S).
\]
(The second equality: $f\chi_S=+1$ where $F=L_S$ and $-1$ elsewhere.)

\begin{lemma}[Orthonormality and Parseval]\label{lem:fourier}
(a)~$\E_x[\chi_S\chi_T]=\mathbf 1[S=T]$.  (b)~$f=\sum_S\widehat f(S)\chi_S$.  (c)~$\sum_S\widehat f(S)^2=\E[f^2]=1$.
\end{lemma}
\begin{proof}
(a) $\chi_S\chi_T=\chi_{S+T}$, and for $U\ne0$ the function $x\mapsto U\cdot x$ takes each value on half of $\F_2^n$, so $\E[\chi_U]=0$.
(b) The $2^n$ functions $\chi_S$ are orthonormal for the inner product $\langle g,h\rangle=\E[gh]$ on the $2^n$-dimensional space of real functions, hence a basis, and the coefficient of $\chi_S$ in $f$ is $\langle f,\chi_S\rangle$.
(c) Expand $\E[f^2]$ using (a) and~(b).
\end{proof}

\subsection{Linearity testing}

\begin{definition}[Linearity test; Blum--Luby--Rubinfeld]\label{def:blr}
Given $F\colon\F_2^n\to\F_2$: choose $x,y\in\F_2^n$ uniformly and independently; accept iff $F(x)+F(y)=F(x+y)$.  (Three queries, $2n$ random bits.)
\end{definition}

\begin{theorem}[Linearity test, sharp form]\label{thm:blr}
If $F$ is linear the test accepts with probability~$1$.  In general
\begin{gather*}
  \Pr[\text{accept}]\;=\;\frac12+\frac12\sum_S\widehat f(S)^3,\\
  \min_{S}\Delta(F,L_S)\ \le\ \Pr[\text{reject}]\ \le\ 3\min_S\Delta(F,L_S).
\end{gather*}
\end{theorem}
\begin{proof}
Since $f(x)f(y)f(x+y)=1$ iff the test accepts, $\Pr[\text{accept}]=\tfrac12+\tfrac12\E[f(x)f(y)f(x+y)]$.  Expanding each factor by Lemma~\ref{lem:fourier}(b),
\[
  \E_{x,y}\bigl[f(x)f(y)f(x+y)\bigr]=\sum_{S,T,U}\widehat f(S)\widehat f(T)\widehat f(U)\,\E_x[\chi_S\chi_U(x)]\,\E_y[\chi_T\chi_U(y)]=\sum_S\widehat f(S)^3,
\]
using $\chi_U(x+y)=\chi_U(x)\chi_U(y)$ and Lemma~\ref{lem:fourier}(a).  For the lower bound on rejection: $\widehat f(S)^3=\widehat f(S)\cdot\widehat f(S)^2\le\max_T\widehat f(T)\cdot\widehat f(S)^2$, so by Parseval $\sum_S\widehat f(S)^3\le\max_T\widehat f(T)$ and
\[
  \Pr[\text{reject}]=\frac12-\frac12\sum_S\widehat f(S)^3\ \ge\ \frac{1-\max_T\widehat f(T)}2=\min_T\Delta(F,L_T).
\]
For the upper bound, let $L$ be a nearest linear function; if $F(x)+F(y)\ne F(x+y)$ then one of $F(x),F(y),F(x+y)$ differs from $L$, and each of these three points is uniform, so $\Pr[\text{reject}]\le3\Delta(F,L)$.
\end{proof}

\begin{remark}
\citet{ALMSS98} quote the theorem of \citet{BLR90} in the weaker form ``rejection probability $\delta<2/9$ implies $2\delta$-close to linear'' and remark that the exact constant may differ from the original;
the Fourier-analytic argument above, from the later literature on linearity testing \citep{BCHKS96}, gives $\delta$-closeness for \emph{every} $\delta$, which improves the soundness constants below.
Exhaustively, over all $2^8$ functions on $\F_2^3$ and all $2^{16}$ on $\F_2^4$ the identity above holds exactly and the largest ratio $\min_S\Delta/\Pr[\text{reject}]$ is $2/3$ and $8/9$ respectively.
\end{remark}

\subsection{Self-correction and concatenation}

A function that is $\eps$-close to linear is correct on most points but the verifier needs its value at \emph{specific} points (for instance a coefficient vector of a circuit).
Linearity lets us turn an arbitrary point into two uniform ones.

\begin{proposition}[Self-correction; BLR]\label{prop:selfcorr}
If $F$ is $\eps$-close to the linear function $L$ then for \emph{every} $x\in\F_2^n$, $\Pr_{y}\bigl[F(x+y)+F(y)\ne L(x)\bigr]\le2\eps$.
\end{proposition}
\begin{proof}
Both $x+y$ and $y$ are uniformly distributed, so $F(x+y)=L(x+y)$ and $F(y)=L(y)$ each fail with probability $\le\eps$, and when both hold $F(x+y)+F(y)=L(x)+L(y)+L(y)=L(x)$.
\end{proof}

\noindent
Self-correction is why the verifier may be given a \emph{wrong} table: it never queries the table at a point of its own choosing, only at points that are individually uniform.  This is the same principle as in Section~\ref{sec:poly}.

A linear function on $\F_2^{n_1+\cdots+n_k}$ is the same thing as a $k$-tuple of linear functions: $L(x_1,\dots,x_k)=\sum_iL_i(x_i)$ with $L_i(u)=L(0,\dots,u,\dots,0)$ ($u$ in the $i$-th block).
This is the \emph{concatenation} we need to compare an oracle $A$ for the whole assignment with oracles $X_i$ for its pieces.

\begin{proposition}[Concatenation test]\label{prop:linconcat}
Let $f\colon\F_2^{n_1+\dots+n_k}\to\F_2$ and $f'\colon\F_2^{n_i}\to\F_2$ be $\eps$-close to linear $g,g'$ respectively.  Test: choose $u\in\F_2^{n_i}$ and $y\in\F_2^{n_1+\dots+n_k}$ uniformly,
let $\iota_i(u)$ be $u$ placed in block $i$, and accept iff $f(\iota_i(u)+y)+f(y)=f'(u)$.  Then (a)~if $g_i=g'$ the test accepts with probability~$1$ when $f=g$, $f'=g'$; (b)~if $g_i\ne g'$ it accepts with probability at most $\tfrac12+3\eps$.
\end{proposition}
\begin{proof}
(b) Since $g_i$ and $g'$ are distinct linear functions, $g(\iota_i(u))=g_i(u)\ne g'(u)$ for half of all $u$.  Accepting requires that this event fails, or that self-correction of $f$ at $\iota_i(u)$ returns a wrong value (probability $\le2\eps$ by Proposition~\ref{prop:selfcorr}), or that $f'(u)\ne g'(u)$ (probability $\le\eps$).
\end{proof}


%% ===========================================================================
\section{Verifying a circuit with Hadamard codes}
\label{sec:had-ver}
%% ===========================================================================

\subsection{A circuit as a system of quadratic equations}

Let $C$ be an algebraic circuit (gates $+$ and $\cdot$ over $\F_2$, fan-in~$2$) with $n$ gates numbered $1,\dots,n$ and $km$ input nodes $1,\dots,km$; gate $n$ is the output.  For an input $x\in\F_2^{km}$ let
$z=C\text{-aug}(x)\in\F_2^n$ be the vector of all gate values (so the first $km$ coordinates are $x$).  Associate to every non-input gate $j$ with input gates $j_1,j_2$ the equation
\[
  P_j(z)=z_{j_1}+z_{j_2}+z_j=0\ (\text{addition}),\qquad P_j(z)=z_{j_1}z_{j_2}+z_j=0\ (\text{multiplication}),
\]
and add $P_{\rm out}(z)=z_n+1=0$.  These are $n-km+1$ polynomials of degree $\le2$.

\begin{lemma}\label{lem:aug}
$C(x)=1$ iff there is $z\in\F_2^n$ with $z_{1..km}=x$ and $P_j(z)=0$ for all $j$.  Moreover then $z=C\text{-aug}(x)$.
\end{lemma}
\noindent(Immediate: the gate equations force $z$ to be the gate-value vector, and $P_{\rm out}$ says the output is~$1$.)

Because $z_a^2=z_a$ over $\F_2$, every polynomial of degree $\le2$ is $P(z)=\sum_{a,b}c_{ab}z_az_b+c_0$, \ie
\begin{equation}\label{eq:quad}
  P(z)=\bigl\langle z\otimes z,\ c\bigr\rangle+c_0,\qquad(z\otimes z)_{ab}=z_az_b,\quad c\in\F_2^{n^2}.
\end{equation}
This is the key observation: \emph{evaluating every quadratic function of $z$ is a single linear query to the Hadamard encoding $L_{z\otimes z}$ of $z\otimes z\in\F_2^{n^2}$.}
The proof therefore consists of $L_z$ (to compare with the input pieces) and $L_{z\otimes z}$ (to evaluate the equations); what must be tested is that the second really is the tensor square of the first.

\begin{lemma}[Tensor test; Freivalds]\label{lem:freivalds}
Let $M\in\F_2^{n\times n}$ be nonzero.  For uniform $u,v\in\F_2^n$, $\Pr[u^{\top}Mv=1]\ge\tfrac14$.  (In fact it equals $\tfrac12-2^{-\mathrm{rank}(M)-1}$; Exercise~\ref{ex:frank}.)
\end{lemma}
\begin{proof}
The map $v\mapsto Mv$ is linear and nonzero, so $Mv\ne0$ with probability $\ge\tfrac12$.  For a fixed nonzero vector $w$, $u^\top w$ is a uniform bit.  Multiply.
\end{proof}

\subsection{The verifier $V^{\oplus}$}

\begin{definition}[Hadamard verifier]\label{def:vplus}
Input: a circuit $C$ as above.  Proof: functions $X_1,\dots,X_k\colon\F_2^m\to\F_2$, $A\colon\F_2^n\to\F_2$ and $B\colon\F_2^{n^2}\to\F_2$.  (Expected: $X_i=L_{x_i}$, $A=L_z$, $B=L_{z\otimes z}$ for $z=C\text{-aug}(x_1\circ\cdots\circ x_k)$.)
The verifier performs \emph{all} of the following tests with independent randomness, and accepts iff all pass.
\begin{enumerate}
  \item \emph{Linearity:} the test of Definition~\ref{def:blr} on $A$, on $B$ and on each $X_i$.
  \item \emph{Concatenation:} for $i=1,\dots,k$ the test of Proposition~\ref{prop:linconcat} on $(A,X_i)$, with block $i$ the $i$-th chunk of $m$ input coordinates.
  \item \emph{Tensor:} draw $z_1,z_2\in\F_2^n$ and $w\in\F_2^{n^2}$ uniformly; check $A(z_1)A(z_2)=B(z_1\otimes z_2+w)+B(w)$.
  \item \emph{Circuit:} draw $r\in\F_2^{n-km+1}$ and $w\in\F_2^{n^2}$ uniformly; let $c=\sum_jr_jc^{(j)}$ and $c_0=\sum_jr_jc^{(j)}_0$ be the coefficients, in the sense of~\eqref{eq:quad}, of $\sum_jr_jP_j$; check $B(c+w)+B(w)=c_0$.
\end{enumerate}
\end{definition}

\begin{theorem}[The Hadamard inner verifier; ALMSS Theorem~15]\label{thm:vplus}
$(V^\oplus,E^\oplus,(E^\oplus)^{-1})$ is a $(k;\,r,\,6k+12,\,O(k),\,1;\,\tfrac{19}{20})$ inner verifier system with $r=4n^2+(k+5)n+2km+1$ random bits.
\end{theorem}
\begin{proof}
\emph{Resources.}  Linearity: $3(k+2)$ queries; concatenation: $3k$; tensor: $4$; circuit: $2$; total $6k+12$.  Random bits: $(2n+2n^2+2km)+k(m+n)+(2n+n^2)+(n-km+1+n^2)=4n^2+(k+5)n+2km+1$.
The verdict is the AND of $O(k)$ parity checks of the answer bits, a circuit of size $O(k)$.

\emph{Completeness.}  If $x_1\circ\cdots\circ x_k$ satisfies $C$, put $z=C\text{-aug}(\cdot)$, $X_i=L_{x_i}$, $A=L_z$, $B=L_{z\otimes z}$.  Linearity and (since $\iota_i(u)\cdot z=u\cdot x_i$) concatenation tests hold; the tensor test holds as $L_z(z_1)L_z(z_2)=(z\cdot z_1)(z\cdot z_2)=\langle z\otimes z,z_1\otimes z_2\rangle$;
the circuit test holds as $\langle z\otimes z,c\rangle=\sum_jr_jP_j(z)+c_0=c_0$ by Lemma~\ref{lem:aug}.

\emph{Soundness.}  Suppose the proof is accepted with probability $>\tfrac{19}{20}$.  The tests use independent randomness, so each individual test accepts with probability $>\tfrac{19}{20}$.
\begin{enumerate}[label=(\roman*)]
  \item By Theorem~\ref{thm:blr} each of $X_1,\dots,X_k,A,B$ is $\eps$-close, for some $\eps<\tfrac1{20}$, to a linear function, say $L_{x_i}$, $L_z$, $L_b$, with $b\in\F_2^{n^2}$ regarded as an $n\times n$ matrix $M$.
  \item \emph{Tensor.}  Suppose $M\ne zz^\top$.  By Lemma~\ref{lem:freivalds} applied to $M+zz^\top$, with probability $\ge\tfrac14$ over $z_1,z_2$ we have $L_z(z_1)L_z(z_2)\ne L_b(z_1\otimes z_2)$.  The test accepts only if this event fails, or $A(z_1)$ or $A(z_2)$ is wrong (probability $\le\eps$ each, both points uniform), or self-correction of $B$ at $z_1\otimes z_2$ is wrong ($\le2\eps$).  So it accepts with probability $\le\tfrac34+4\eps<\tfrac{19}{20}$, a contradiction.  Hence $L_b=L_{z\otimes z}$.
  \item \emph{Circuit.}  Let $z_1,\dots,z_k$ be the consecutive $m$-bit blocks of the first $km$ coordinates of $z$.  If some $P_j(z)\ne0$ then for uniform $r$ we have $\sum_jr_jP_j(z)=1$ with probability exactly $\tfrac12$, \ie $\langle z\otimes z,c\rangle\ne c_0$ with probability $\tfrac12$; the test accepts only if this event fails or the self-correction of $B$ at $c$ errs ($\le2\eps$): at most $\tfrac12+2\eps<\tfrac{19}{20}$, a contradiction.  So $P_j(z)=0$ for all $j$, and by Lemma~\ref{lem:aug} $C$ accepts $z_1\circ\cdots\circ z_k$ and $z=C\text{-aug}(z_1\circ\cdots\circ z_k)$.
  \item \emph{Concatenation.}  By Proposition~\ref{prop:linconcat}(b) and $\tfrac12+3\eps<\tfrac{19}{20}$, the $i$-th block function of $L_z$ equals $L_{x_i}$, \ie $z_i=x_i$.  Since $\eps<\tfrac14$, $(E^\oplus)^{-1}(X_i)=x_i$.
\end{enumerate}
Therefore $(E^\oplus)^{-1}(X_1)\circ\cdots\circ(E^\oplus)^{-1}(X_k)=z_1\circ\cdots\circ z_k$ satisfies $C$.
\end{proof}

\begin{remark}
\citet{ALMSS98} obtain soundness $\tfrac{35}{36}$ with the weaker linearity theorem (closeness $2\delta$ for rejection $\delta<\tfrac29$); with Theorem~\ref{thm:blr} the same argument gives $\tfrac{19}{20}$, as above.  Both are constants less than~$1$, which is all that is used.
\end{remark}

\begin{example}[Two circuits on five gates]\label{ex:had}
Take $k=1$, $m=2$ (inputs $z_1,z_2$), $n=5$.  The proof has $|X_1|=4$, $|A|=2^5=32$ and $|B|=2^{25}=33{,}554{,}432$ bits, the verifier reads $18$ of them and uses $4\cdot25+6\cdot5+2\cdot2+1=135$ random bits.

\smallskip\noindent\emph{(a) A satisfiable circuit.}  Gates $3=1\cdot2$, $4=1+2$, $5=3+4$, output $5$: $C(x)=x_1\vee x_2$.  Its equations are $z_3+z_1z_2$, $z_4+z_1+z_2$, $z_5+z_3+z_4$, $z_5+1$.  For $x=(1,0)$ the augmented vector is $z=(1,0,0,1,1)$, and the honest proof is accepted with probability~$1$.

\smallskip\noindent\emph{(b) An unsatisfiable circuit.}  Gates $3=1+2$, $4=1\cdot2$, $5=3\cdot4$, output $5$.  Over $\F_2$, $(x_1+x_2)x_1x_2=x_1x_2+x_1x_2=0$, so $C\equiv0$; the equations are $z_3+z_1+z_2$, $z_4+z_1z_2$, $z_5+z_3z_4$, $z_5+1$.
No $z$ satisfies all four, and some $z$ violates exactly one.  We now find the best cheating proof \emph{among linear oracles} ($X_1=L_x$, $A=L_z$, $B=L_b$), for which every test can be computed exactly.
Take $X_1$ the prefix of $z$ (so the concatenation test passes) and let $M$ be the matrix of $b$.  With $N=M+zz^\top$, the tensor test accepts with probability $\tfrac12+2^{-\mathrm{rank}N-1}$ (Exercise~\ref{ex:frank}), and the circuit test accepts with probability $1$ if
$\langle M,c^{(j)}\rangle=c^{(j)}_0$ for all four equations, and with probability exactly $\tfrac12$ otherwise.
\begin{itemize}
  \item $\mathrm{rank}\,N=0$ ($B$ is honest): tensor $1$; but for every $z$ some equation fails, so circuit $\tfrac12$.  Total $\tfrac12$.
  \item $\mathrm{rank}\,N=1$: tensor $\tfrac34$.  An exhaustive search over all $32\cdot31\cdot31$ choices of $(z,u,v)$ with $N=uv^\top$ finds the circuit test passing with probability $1$ for, \eg, $z=0$, $u=e_3+e_5$, $v=e_4+e_5$: then $u^\top C^{(j)}v=(0,0,0,1)=c^{(j)}_0$.  Total $\tfrac34$.
  \item $\mathrm{rank}\,N\ge2$: tensor $\le\tfrac58$.
\end{itemize}
So the best linear cheating proof is accepted with probability exactly $\tfrac34$, well below the guarantee $\tfrac{19}{20}$ of Theorem~\ref{thm:vplus}.  Note \emph{which} test it cheats on: $B$ is not the tensor square of $A$ at all (here $A\equiv0$), and passes the circuit test by being a rank-one tensor that happens to solve the four linear conditions.  Without the tensor test the verifier would accept the unsatisfiable circuit with probability~$1$.
\end{example}

\begin{bigpicture}{Why this verifier is a good \emph{inner} verifier and a useless outer one}
The proof has length $2^{n^2}$ and the verifier uses $\Theta(n^2)$ random bits, where $n$ is the size of the circuit.  For $n=\poly(N)$ this is exponential, so $V^\oplus$ can not be used as a PCP for $\sat{3SAT}$ of size $N$.
It reads \emph{six bits per piece} regardless of $n$, and is exact about what it needs: the input is a \emph{split-encoded} assignment ($X_1,\dots,X_k$), exactly the interface demanded by Definition~\ref{def:inner}.
So $V^\oplus$ is the ideal last verifier of a recursion, applied to circuits that are tiny: if the previous verifier has a decision circuit of size $\polylog\log N$, then $2^{n^2}$ is $2^{(\log\log N)^{O(1)}}=N^{o(1)}$.
\end{bigpicture}


%% ===========================================================================
\section{\texorpdfstring{Lines: low-degree testing and correction with $O(1)$ queries}{Lines: low-degree testing and correction with O(1) queries}}
\label{sec:lines}
%% ===========================================================================

For the rest of this lecture $\F$ is a finite field, and $\mathcal Q^m_d$ denotes the polynomials $\F^m\to\F$ of \emph{total} degree at most $d$.
(In Lecture~\ref{lec:as} we used the per-variable degree $h$, for which total degree is $\le hm$.  A line sees the total degree, so that is the natural notion here.)

\begin{lemma}[Schwartz--Zippel, total degree]\label{lem:sz-tot}
Two distinct $p,p'\in\mathcal Q^m_d$ agree on at most a $d/|\F|$ fraction of $\F^m$.
\end{lemma}
\begin{proof}
It suffices to show that a nonzero $p\in\mathcal Q^m_d$ has at most a $d/|\F|$ fraction of zeros.  For $m=1$ there are at most $d$ roots.  For $m>1$ write $p=\sum_{i\le k}x_m^ip_i(x_1,\dots,x_{m-1})$ with $p_k\ne0$, so $\deg p_k\le d-k$.
By induction $p_k$ vanishes on at most a $(d-k)/|\F|$ fraction; for every other $(x_1..x_{m-1})$ the polynomial in $x_m$ has degree $k$ and at most $k$ roots.  Total: $\le(d-k)/|\F|+k/|\F|=d/|\F|$.
\end{proof}

\noindent
In particular, if $d/|\F|\le\tfrac12$ and $\delta<\tfrac14$ then $f$ is $\delta$-close to at most one $p\in\mathcal Q^m_d$ (Remark~22 of ALMSS).

\paragraph{The encoding.}  Let $H\subseteq\F$ with $|H|=h+1$ and let $n\le|H|^m$.  The proof of Lemma~\ref{lem:ext} works verbatim for any $H\subseteq\F$ with $|H|=h+1$: every function $s\colon H^m\to\F$, in particular every bit string $s\in\{0,1\}^n$ padded to length $|H|^m$,
has an extension $\hat s$ of total degree $\le mh$.  The polynomial encoding $P_{n,m,H,\F}(s)=\hat s$ is the table of $\hat s$ on $\F^m$; the decoder
returns, for $f\colon\F^m\to\F$, the values on $H^m$ of the nearest polynomial of $\mathcal Q^m_{d}$ ($d$ a degree bound specified by the verifier; $d\ge mh$), truncated to $n$ bits.  We need $|\F|\ge2d$ so that nearest polynomials are unique.

\subsection{The low-degree test with a table of lines}

The Arora--Safra test of Lecture~\ref{lec:as} reads $O(m)$ entries.  The test of \citet{RubinfeldSudan96}, as strengthened by Arora--Safra's bivariate lemma, reads $O(1)$.
Its trick is to use \emph{all} lines, not only axis-parallel ones.  Let $\ell_{x,h}(t)=x+th$ be the line with base $x$ and direction $h\in\F^m$.  For $p\in\mathcal Q^m_d$ the restriction $t\mapsto p(x+th)$ is a univariate polynomial of degree $\le d$.

\begin{definition}[Line table, Poly-test]\label{def:polytest}
A \defn{line table} is a function $B\colon\F^{2m}\to\mathcal Q^1_d$ assigning to every pair $(x,h)$ a univariate polynomial $B(x,h)$ of degree $\le d$.  The honest line table of $p$ is $\mathrm{lines}_p(x,h)=p(x+t\,h)$ (as a polynomial in $t$).
\emph{Poly-test}$(f,B)$: choose $x,h\in\F^m$ and $t\in\F$ uniformly; accept iff $f(x+th)=B(x,h)(t)$.  It makes two queries.
\end{definition}

\begin{theorem}[Low-degree test; ALMSS Theorem~45]\label{thm:polytest}
There are constants $\delta_0>0$ and $\alpha<\infty$ such that for every $d$ and field with $|\F|\ge\alpha(d+1)^3$ and every $\delta\le\delta_0$: (1)~$\mathrm{Poly\text{-}test}(p,\mathrm{lines}_p)$ accepts with probability~$1$ for $p\in\mathcal Q^m_d$;
(2)~if $\mathrm{Poly\text{-}test}(f,B)$ accepts with probability $\ge1-\delta$ for some $B$, then $f$ is $2\delta$-close to $\mathcal Q^m_d$.
\end{theorem}

\noindent
We quote it.  Part~(1) is trivial; part~(2) is proved in Section~A of ALMSS by combining the analysis of \citet{RubinfeldSudan96} (whose number of queries is independent of $m$ but whose field size depends on more than $d$) with the bivariate lemma of Arora--Safra (the base case of Lemma~\ref{lem:robust}).
The point to retain is the quantifiers: \emph{two} queries, independent of $m$, $d$ and $\delta$, in exchange for a field of size $\Omega(d^3)$.

\subsection{Correcting a polynomial at a point}

\begin{definition}[Poly-self-corr]\label{def:polycorr}
\emph{Poly-self-corr}$(A,B,x)$ with $A\colon\F^m\to\F$ and a line table $B$: choose $h\in\F^m$ and $t\in\F\setminus\{0\}$ uniformly; let $q=B(x,h)$; if $q(t)\ne A(x+th)$ reject, else output $q(0)$.  Two queries.
\end{definition}

\begin{proposition}[ALMSS Proposition~46]\label{prop:polycorr}
Let $p\in\mathcal Q^m_d$.  (1)~$\mathrm{Poly\text{-}self\text{-}corr}(p,\mathrm{lines}_p,x)$ outputs $p(x)$ with probability~$1$.  (2)~If $\Delta(A,p)=\eps$ then for every $x\in\F^m$ and \emph{every} line table $B$,
\[
  \Pr_{h,t}\bigl[\text{the procedure does not reject and its output is}\ne p(x)\bigr]\ \le\ 2\sqrt\eps+\frac d{|\F|-1}.
\]
\end{proposition}
\begin{proof}
(1) is clear.  (2) For fixed $x$ and $t\ne0$, the point $x+th$ is uniform when $h$ is uniform, so $\Pr_{h,t}[A(x+th)\ne p(x+th)]=\eps$.  Call $h$ \emph{bad} if $\Pr_t[A(x+th)\ne p(x+th)]>\sqrt\eps$; by Markov $\Pr_h[\text{bad}]<\sqrt\eps$.
Fix a good $h$ and write $q=B(x,h)$.  If $q=p(x+\cdot\,h)$ then the output is $p(x)$ or the procedure rejects.  If $q\ne p(x+\cdot\,h)$ then the two degree-$\le d$ polynomials agree in at most $d$ values of $t$ (so for $d/(|\F|-1)$ of $t\ne0$), and a wrong output
without rejection needs $q(t)=A(x+th)$, which happens when $A(x+th)=p(x+th)=q(t)$ (probability $\le d/(|\F|-1)$) or when $A(x+th)\ne p(x+th)$ (probability $\le\sqrt\eps$ since $h$ is good).  Summing, $\sqrt\eps+\sqrt\eps+d/(|\F|-1)$.
\end{proof}

\subsection{Concatenating polynomials}

The analogue of Proposition~\ref{prop:linconcat} for polynomial encodings: from $p_1,\dots,p_k\in\mathcal Q^m_d$ build one polynomial in $m+1$ variables whose slices are the $p_i$.  Fix an enumeration $\zeta_1,\zeta_2,\dots$ of the elements of $\F$; the $\zeta_i$ serve as indices for the $k$ pieces.

\begin{definition}[Concatenated polynomial]\label{def:polyconcat}
For $k\le|\F|$ let $\mathrm{concat}_{p_1,\dots,p_k}(x,y)=\sum_{i=1}^k\lambda_i(y)p_i(x)\in\mathcal Q^{m+1}_{d+k-1}$, where $\lambda_i$ is the Lagrange polynomial of degree $k-1$ with $\lambda_i(\zeta_j)=\mathbf 1[i=j]$.  It satisfies $\mathrm{concat}(x,\zeta_i)=p_i(x)$.
\end{definition}

\begin{lemma}[Concatenation test; ALMSS Lemma~49]\label{lem:polyconcat}
Let $A\colon\F^{m+1}\to\F$ be $\eps$-close to $p\in\mathcal Q^{m+1}_d$, let $p_i=p(\cdot,\zeta_i)$, and let $C\colon\F^m\to\F$.  \emph{Poly-corr-concat-test}$(\zeta_i,A,B,C)$: choose $x\in\F^m$; accept iff
$\mathrm{Poly\text{-}self\text{-}corr}(A,B,(x,\zeta_i))$ does not reject and outputs $C(x)$.  (Three queries.)  (1)~If $A=p$, $B=\mathrm{lines}_p$, $C=p_i$ it accepts with probability~$1$.  (2)~If for some $B$ it accepts with probability $\rho$, then $\Delta(C,p_i)<1-\rho+2\sqrt\eps+d/(|\F|-1)$.
\end{lemma}
\begin{proof}
(2) Let $\Delta=\Delta(C,p_i)$.  The test accepts only if $C(x)=p(x,\zeta_i)$ (probability $1-\Delta$ over $x$) or if the self-correction outputs something other than $p(x,\zeta_i)$ without rejecting; the latter has probability $\le2\sqrt\eps+d/(|\F|-1)$ for \emph{every} point, in particular for $(x,\zeta_i)$ (Proposition~\ref{prop:polycorr}).
\end{proof}

\noindent
Note that the correction is performed at the \emph{non-uniform} points $(x,\zeta_i)$; this is why the line $\ell_{(x,\zeta_i),h}$ with random $h$ is used: the point actually queried in $A$ is uniform.


%% ===========================================================================
\section{\texorpdfstring{Curves: collapsing $l$ queries into one}{Curves: collapsing l queries into one}}
\label{sec:curves}
%% ===========================================================================

The Arora--Safra verifier asks $l=\polylog n$ queries.  The $l$ query points $z_1,\dots,z_l\in\F^{w+1}$ are determined by the random string $R$, so we may pre-compute them, and ask for the values of the polynomial $p$ at all of them \emph{at once}.
A line through two points is easy; through $l$ points in general position we need a curve of degree $l-1$.

\begin{definition}[Curves]\label{def:curve}
A \defn{curve of degree $s$} in $\F^m$ is a map $\mathcal C\colon\F\to\F^m$ whose coordinates are univariate polynomials of degree $\le s$.  Given points $z_1,\dots,z_l\in\F^m$ and the first $l$ elements $\zeta_1,\dots,\zeta_l$ of $\F$, let $\mathrm{curve}_{z_1..z_l}$ be the curve of degree $l-1$ obtained by Lagrange interpolation in every coordinate, so that $\mathrm{curve}(\zeta_j)=z_j$.
\end{definition}

\begin{proposition}\label{prop:curve-deg}
If $p\in\mathcal Q^m_d$ and $\mathcal C$ is a curve of degree $s$ then $t\mapsto p(\mathcal C(t))$ is a univariate polynomial of degree $\le ds$.  In particular for $\mathcal C=\mathrm{curve}_{z_1..z_l}$ the degree is $\le d(l-1)$.
\end{proposition}
\begin{proof}
Each monomial of $p$ has total degree $\le d$, and substituting polynomials of degree $\le s$ gives a polynomial of degree $\le ds$.
\end{proof}

\begin{figure}[t]
\centering
\begin{tikzpicture}[x=1cm,y=1cm,every node/.style={font=\small},
   pt/.style={circle,fill=accent,inner sep=1.6pt}]
  \begin{scope}
    \draw[rounded corners=6pt,fill=black!3] (-0.3,-1.9) rectangle (5.7,1.9);
    \node[anchor=north west,font=\footnotesize] at (-0.2,1.85) {$\F^{w+1}$};
    \node[pt,label=above:{$z_1$}] (a) at (0.6,0.6) {};
    \node[pt,label=above:{$z_2$}] (b) at (1.9,-0.9) {};
    \node[pt,label=above:{$z_3$}] (c) at (3.2,0.9) {};
    \node[pt,label=above:{$z_l$}] (d) at (4.9,-0.4) {};
    \node at (4.0,0.1) {$\cdots$};
    \draw[warn,very thick] plot[smooth,tension=0.8] coordinates {(0.6,0.6) (1.2,-0.1) (1.9,-0.9) (2.6,-0.2) (3.2,0.9) (4.1,0.5) (4.9,-0.4)};
    \node[warn,font=\footnotesize] at (2.9,-1.55) {$\mathrm{curve}_{z_1..z_l}$ (degree $l-1$)};
  \end{scope}
  \node[align=left,anchor=west,font=\footnotesize] at (6.2,0.8) {\textbf{Without curves:} $l$ queries $A(z_1),\dots,A(z_l)$,\\ each of which must itself be self-corrected.};
  \node[align=left,anchor=west,font=\footnotesize] at (6.2,-0.55) {\textbf{With curves:} one query $q=O(z_1,\dots,z_l)$,\\ a polynomial of degree $d(l-1)$, then one random\\ $t$ and one self-corrected value at $\mathrm{curve}(t)$.\\ Output: $q(\zeta_1),\dots,q(\zeta_l)$.};
\end{tikzpicture}
\caption{Aggregating queries.  The oracle $O$ lists, for every tuple of points, the restriction of the (encoded) proof to the curve through them.  The verifier \emph{checks} this restriction at one random point of the curve, and \emph{reads off} the $l$ values it wanted from the coefficients.}
\label{fig:curve}
\end{figure}

\begin{definition}[Curve-corr-check]\label{def:curvecheck}
Given $A\colon\F^m\to\F$, a line table $B$, an oracle $O\colon\F^{ml}\to\mathcal Q^1_{d(l-1)}$ and points $z_1,\dots,z_l\in\F^m$: let $q=O(z_1,\dots,z_l)$; choose $t\in\F$ uniformly; if $q(t)\ne\mathrm{Poly\text{-}self\text{-}corr}(A,B,\mathrm{curve}_{z_1..z_l}(t))$ (including the case that the latter rejects) then reject; else output $(q(\zeta_1),\dots,q(\zeta_l))$.
It makes three queries (one to $O$, two inside Poly-self-corr) independently of $l$.
\end{definition}

\begin{lemma}[ALMSS Lemma~57]\label{lem:curvecheck}
Let $p\in\mathcal Q^m_d$ and $A$ be $\eps$-close to $p$.  (1)~If $B=\mathrm{lines}_p$ and $O(z_1..z_l)=p|_{\mathrm{curve}_{z_1..z_l}}$ then the output is $(p(z_1),\dots,p(z_l))$ with probability~$1$.
(2)~For all $B$, $O$ and $z_1,\dots,z_l$,
\[
  \Pr\bigl[\text{no rejection and output}\ne(p(z_1),\dots,p(z_l))\bigr]\ \le\ 2\sqrt\eps+\frac d{|\F|-1}+\frac{d(l-1)}{|\F|}.
\]
\end{lemma}
\begin{proof}
(1) is clear since $q(\zeta_j)=p(\mathrm{curve}(\zeta_j))=p(z_j)$.  (2) Let $\mathcal C=\mathrm{curve}_{z_1..z_l}$ and $q=O(z_1..z_l)$.  If $q=p\circ\mathcal C$ then the output, if any, is $(p(z_j))_j$.  Otherwise $q$ and $p\circ\mathcal C$ are distinct polynomials of degree $\le d(l-1)$ and agree for at most $d(l-1)/|\F|$ of the $t$.
Not rejecting requires $q(t)$ to equal the self-corrected value at $\mathcal C(t)$.  That value is $p(\mathcal C(t))$ except with probability $\le2\sqrt\eps+d/(|\F|-1)$ (Proposition~\ref{prop:polycorr}, which is valid at every point, in particular at $\mathcal C(t)$ for each fixed $t$).
So non-rejection requires $q(t)=p(\mathcal C(t))$ (probability $\le d(l-1)/|\F|$) or a wrong self-correction.
\end{proof}

\begin{pitfall}{a curve is not a random point}
A point $\mathrm{curve}(t)$ on a curve through \emph{given} points is not uniformly distributed in $\F^m$; it is determined by $z_1,\dots,z_l$.  An adversary who knows the verifier's program may corrupt $A$ exactly on the points that the verifier
will want to read, and querying $A$ directly at $\mathrm{curve}(t)$ would give no protection from the $\eps$-fraction where $A$ is wrong.  This is why the check goes through Poly-self-corr, which replaces the point by $\mathrm{curve}(t)+t'h$ for a random direction $h$: a uniformly distributed point.
\end{pitfall}


%% ===========================================================================
\section{The parallelized verifier}
\label{sec:par}
%% ===========================================================================

The starting point is the Arora--Safra verifier in the inner-verifier formalism.

\begin{theorem}[Arora--Safra, in the form of ALMSS Theorem~39; quoted]\label{thm:seq}
For every constant $k$ there are constants $c_1,a,b$ and $e_{\rm seq}<1$ and an inner verifier system $V^{\rm seq}$ with parameters
\[
  (k;\ r,p,c,a;\ e)\;=\;(k;\ c_1\log n,\ \log^an,\ \log^bn,\ 1;\ e_{\rm seq}),
\]
\ie $c_1\log n$ random bits, $\log^an$ queries, a decision circuit of size $\log^bn$ and bit answers.
\end{theorem}

\noindent
This is the verifier of Section~\ref{sec:normal} with $h=\Theta(\log n)$ and $m=\log n/\log\log n$ (the verifier $V_{\rm savior}$ of Section~\ref{sec:chase}), read as an inner verifier system: its encoding is the polynomial extension, it makes $O(k+m)=\polylog n$ queries, and its decision is a circuit of size $\poly(h)=\polylog n$.  (ALMSS write the exponents of the number of queries and of the circuit size as $c_2$ and $c_3$ in the statement but use $c_3$ for the number of queries in Section~7.5; we call them $a$ and $b$.)  Its oracles are bit strings of length $N\le\log^an\cdot2^{c_1\log n}=\poly(n)$.

\paragraph{Parameters of $V^{\rm par}$.}  Let $w=\lceil\log N/\log\log N\rceil$, $H\subseteq\F$ with $|H|=\log N$ (so $|H|^w\ge N$ and the polynomial encodings have total degree at most $w(|H|-1)$), $d=w|H|+k$, and $|\F|=\log^sN$ with $s=\max\{7,a+3\}$.  Then $|\F|\ge\alpha(d+1)^3$ for large $n$ (as $d=O(\log^2N)$),
and $d\,l/|\F|=o(1)$ where $l=\log^an$ is the number of queries of $V^{\rm seq}$.  The encoding of $V^{\rm par}$ is $E^{\rm par}=P_{N,w,H,\F}\circ E^{\rm seq}$ (first encode with $V^{\rm seq}$'s scheme, then as a polynomial).  Its decoder first finds the nearest polynomial of degree $\le d$
(the degree bound against which the test certifies closeness; for codewords of degree $\le w|H|$ this changes nothing) and then applies $(E^{\rm seq})^{-1}$.

\begin{definition}[The verifier $V^{\rm par}$]\label{def:vpar}
The input is a circuit $C$ with $km$ input nodes.  The proof consists of $X_1,\dots,X_k\colon\F^w\to\F$ and of $X_{k+1}$ with four suboracles:
\begin{center}\small
\renewcommand{\arraystretch}{1.2}
\begin{tabular}{@{}lll@{}}
\toprule
oracle & type & expected content\\
\midrule
$X_i$ ($i\le k$) & $\F^w\to\F$ & $P(E^{\rm seq}(x_i))$: encoding of the $i$-th piece\\
$Z$ & $\F^w\to\F$ & $P(\pi)$: encoding of the auxiliary proof $\pi$ of $V^{\rm seq}$\\
$A$ & $\F^{w+1}\to\F$ & $\mathrm{concat}_{X_1,\dots,X_k,Z}$\\
$B$ & $\F^{2(w+1)}\to\mathcal Q^1_d$ & $\mathrm{lines}_A$ (to test and correct $A$)\\
$O$ & $\F^{(w+1)l}\to\mathcal Q^1_{d(l-1)}$ & $A|_{\mathrm{curve}(\cdot)}$ (to aggregate queries)\\
\bottomrule
\end{tabular}
\end{center}
\begin{enumerate}
  \item \emph{Low degree:} run Poly-test$(A,B)$.
  \item \emph{Concatenation:} for $i=1,\dots,k$ run the test of Lemma~\ref{lem:polyconcat} for $(\zeta_i,A,B,X_i)$, and for $(\zeta_{k+1},A,B,Z)$.
  \item \emph{Simulation of $V^{\rm seq}$:} draw $R\in\{0,1\}^{c_1\log n}$; let $(i_1,q_1),\dots,(i_l,q_l)$ be the queries of $V^{\rm seq}$ on $(C,R)$ (oracle $i_j$, position $q_j\in[N]$); let $\#\colon[N]\to H^w$ be the canonical map and $z_j=(\#(q_j),\zeta_{i_j})\in\F^{w+1}$.
        Run Curve-corr-check$(A,B,O,z_1,\dots,z_l)$ to obtain $(a_1,\dots,a_l)$; reject if some $a_j\notin\{0,1\}$; reject if $V^{\rm seq}$'s decision circuit on $(C,R)$ rejects the answers $(a_1,\dots,a_l)$.
\end{enumerate}
Accept iff nothing rejects.
\end{definition}

\begin{theorem}[Parallelization; ALMSS Theorem~14]\label{thm:par}
$V^{\rm par}$ is a $(k;\ O(\log n),\ 3k+8,\ \polylog n,\ \polylog n;\ e)$ inner verifier system for some constant $e<1$ depending on $k$.
\end{theorem}
\begin{proof}
\emph{Resources.}  Queries: $2$ (Poly-test) $+\,3(k+1)$ (concatenation tests) $+\,3$ (curve check) $=3k+8$.  (ALMSS count $3k+5$, omitting the test for $Z$; this is immaterial.)
Random bits: each test draws $O(w)$ field elements, \ie $O(w\log|\F|)=O(\log N)$ bits since $w\log|\F|=s\log N$, so the total is $O(k\log N)+r=O(\log n)$ for fixed $k$.  Answer sizes are those of the entries of $B$ and $O$: $O(dl\log|\F|)=\polylog n$ bits.
Poly-test, the correction and the curve-check involve $\poly(d,l,m)$ field operations; the decision of $V^{\rm seq}$ is a circuit of size $\log^bn$; so the decision is a circuit of size $\polylog n$ of the answers.

\emph{Completeness.}  Given $x_1\circ\cdots\circ x_k$ satisfying $C$, let $\pi$ be the proof of $V^{\rm seq}$ accepted with probability~$1$ and set the oracles to the expected contents of the table.  Every test is then an identity between polynomials, and in Step~3 the answers $a_j$ are exactly the bits $V^{\rm seq}$ would read, which it accepts.

\emph{Soundness.}  Let $\eta=2\sqrt{0.002}+\frac d{|\F|-1}+\frac{dl}{|\F|}<0.1$ (for large $n$) and $e=\max\{0.999,\ 1-\delta_0,\ 1-(1-\eta)(1-e_{\rm seq})\}<1$.  Suppose $V^{\rm par}$ accepts $(X_1..X_k,Z,A,B,O)$ with probability $>e$.
\begin{enumerate}[label=(\roman*)]
  \item Step~1 accepts with probability $>\max\{0.999,1-\delta_0\}$, so by Theorem~\ref{thm:polytest} there is $p\in\mathcal Q^{w+1}_d$ with $\eps:=\Delta(A,p)<2\min\{0.001,\delta_0\}\le0.002$.
  \item Each concatenation test accepts with probability $>0.999$, so by Lemma~\ref{lem:polyconcat} $\Delta(X_i,p_i)<0.001+2\sqrt{0.002}+o(1)<\tfrac14$ for $p_i=p(\cdot,\zeta_i)$ ($i\le k$) and likewise for $Z$ and $p_{k+1}$.  Since $d/|\F|\le\tfrac12$, $p_i$ is the unique polynomial of $\mathcal Q^w_d$ within $\tfrac14$ of $X_i$, so the decoder returns, from $X_i$, the string $Y_i$ of values of $p_i$ on $H^w$ (after $(E^{\rm seq})^{-1}$).  Let $Y_{k+1}$ be the values of $p_{k+1}$ on $H^w$.
  \item Let $S$ be the set of $R$ for which the tuple $\bigl(p(z_j)\bigr)_{j\le l}$ is \emph{not} a tuple of bits accepted by the decision circuit of $V^{\rm seq}$ on $(C,R)$.  Since $p(\#(q_j),\zeta_{i_j})=p_{i_j}(\#(q_j))$ is the $q_j$-th value of $Y_{i_j}$, this tuple is what $V^{\rm seq}$ would read from $Y_1,\dots,Y_{k+1}$ on $R$ (replace non-bit entries of the $Y_i$ by~$0$ to obtain bit strings; this does not change the behaviour on $R\notin S$).  Fix $R\in S$.
        By Lemma~\ref{lem:curvecheck}, Step~3 returns $(p(z_j))_j$ or rejects, except with probability $\le\eta$; and on $(p(z_j))_j$ Step~3(f) rejects since $R\in S$.  So Step~3 accepts with probability $\le\eta$ given $R\in S$.
  \item Hence $\Pr[V^{\rm par}\text{ accepts}]\le\Pr[S]\,\eta+(1-\Pr[S])=1-(1-\eta)\Pr[S]$.  Since this exceeds $e\ge1-(1-\eta)(1-e_{\rm seq})$, we get $\Pr[S]<1-e_{\rm seq}$: $V^{\rm seq}$ accepts $Y_1,\dots,Y_{k+1}$ (as bit strings) with probability $>e_{\rm seq}$.
        By its soundness $(E^{\rm seq})^{-1}(Y_1)\circ\cdots\circ(E^{\rm seq})^{-1}(Y_k)$ satisfies $C$, and this string is the decoding of $X_1,\dots,X_k$ under $E^{\rm par}$.\qedhere
\end{enumerate}
\end{proof}

\begin{remark}[Where each technique is used]
\emph{Low-degree test and concatenation tests} certify that the (arbitrary) tables $X_i,Z$ are close to polynomials that are the slices of one polynomial $p$: they play the role that linearity and concatenation play in Section~\ref{sec:had-ver}.
\emph{Self-correction} is used so that every read of $A$ is at a uniform point.  \emph{Curves} are the new ingredient, absent from Lecture~\ref{lec:as}: they are the reason the number of queries is independent of $l$.
\end{remark}


%% ===========================================================================
\section{Assembling the proof: three steps, constant parameters}
\label{sec:assemble}
%% ===========================================================================

\begin{theorem}[ALMSS Theorem~17]\label{thm:alm17}
There is a constant $C$ such that for every $L\in\NP$ there is a constant $c_L$ with $L\in\mathrm{RPCP}(c_L\log n,\ C,\ 2C,\ 1,\ \tfrac12)$.
\end{theorem}

\noindent
Since an $(r,p,c,a)$ outer verifier reads $pa$ bits, Theorem~\ref{thm:alm17} gives $L\in\PCP[c_L\log n,C]$, which is Theorem~\ref{thm:alm}.

\begin{proof}
It suffices to treat $\sat{CKT-SAT}$, which is $\NP$-complete.  All inner verifiers below are brought to error $\tfrac1{16}$ by Proposition~\ref{prop:rep}; this multiplies each parameter by a constant.

\medskip
\noindent\textbf{Step 1: constant queries, large answers.}  Apply Theorem~\ref{thm:par} with $k=1$ (its proof starts from the verifier of Theorem~\ref{thm:seq}).  By Proposition~\ref{prop:in-out},
\[
  \sat{CKT-SAT}\in\mathrm{RPCP}\bigl(c_1\log n,\ p_1,\ \polylog n,\ \polylog n,\ \tfrac1{16}\bigr)\qquad(p_1\text{ a constant}).\tag{1}
\]

\medskip
\noindent\textbf{Step 2: shrink the answers by composing with the same kind of verifier.}  Apply Theorem~\ref{thm:par} again, now with $k:=p_1$ \emph{the constant of Step~1}.  This gives a $(p_1;\,r_2,p_2,c_2,a_2;\tfrac1{16})$ inner system with $r_2=O(\log n)$, $p_2=O(1)$ and
$c_2,a_2=\polylog n$.  Compose with~(1) (Theorem~\ref{thm:compose}); the circuit size to feed it is $c(n)=\polylog n$, so $r_2(c(n))=O(\log\log n)$ and $c_2(c(n)),a_2(c(n))=\polylog(\polylog n)=(\log\log n)^{O(1)}$:
\[
  \sat{CKT-SAT}\in\mathrm{RPCP}\bigl(c_1\log n+O(\log\log n),\ p_2,\ (\log\log n)^{O(1)},\ (\log\log n)^{O(1)},\ \tfrac18\bigr).\tag{2}
\]

\medskip
\noindent\textbf{Step 3: shrink the answers to one bit with the Hadamard verifier.}  Apply Theorem~\ref{thm:vplus} with $k:=p_2$.  This gives a $(p_2;\,gn^2,h,2h,1;\tfrac1{16})$ inner system, with $g,h$ constants.  Compose with~(2); its circuit size is $c(n)=(\log\log n)^{O(1)}$, so the random bits become
$c_1\log n+O(\log\log n)+g\,c(n)^2=c_1\log n+(\log\log n)^{O(1)}=O(\log n)$:
\[
  \sat{CKT-SAT}\in\mathrm{RPCP}\bigl(O(\log n),\ h,\ 2h,\ 1,\ \tfrac{3}{16}\bigr).
\]
The errors compose by $e\mapsto e+\tfrac1{16}-\tfrac e{16}$: from $\tfrac1{16}$ to $\tfrac{31}{256}<\tfrac18$ in step~2 and to $\tfrac{721}{4096}<\tfrac3{16}$ in step~3 (Exercise~\ref{ex:alm-reps}).  This is below $\tfrac12$; take $C=h$.
\end{proof}

\begin{table}[t]
\centering\small
\renewcommand{\arraystretch}{1.25}
\begin{tabular}{@{}l>{\raggedright\arraybackslash}p{3.2cm}lll@{}}
\toprule
stage & random bits & queries & answer bits & decision circuit\\
\midrule
Thm~\ref{thm:seq} ($k=1$) & $c_1\log n$ & $\log^an$ & $1$ & $\log^bn$\\
(1) parallelized & $O(\log n)$ & $p_1=O(1)$ & $\polylog n$ & $\polylog n$\\
(2) composed with $V^{\rm par}_{k=p_1}$ & $O(\log n)+O(\log\log n)$ & $p_2=O(1)$ & $(\log\log n)^{O(1)}$ & $(\log\log n)^{O(1)}$\\
(3) composed with $V^\oplus_{k=p_2}$ & $O(\log n)+(\log\log n)^{O(1)}$ & $h=O(1)$ & $1$ & $2h$\\
\bottomrule
\end{tabular}
\caption{The parameter chase of Theorem~\ref{thm:alm17}.  Compare with Table~\ref{tab:chase}: there, \emph{every} composition kept $\sqrt{\log n}$ queries and shrank the exponent of the alphabet by half; here the queries are made constant before any composition, and two compositions suffice.}
\label{tab:alm-chase}
\end{table}

\begin{bigpicture}{Why the order matters}
In Lecture~\ref{lec:as} the number of queries after a composition is $O(Q+q(\tau))$: it never goes below the outer verifier's $Q$.  So the \emph{first} verifier must already have $O(1)$ queries,
and parallelization achieves that at the price of larger answers.  The second observation is that the \emph{size} of the answer is irrelevant to the number of queries of the composed verifier, so large answers can be cut down later by inner verifiers of \emph{any} speed.
Finally, because each inner verifier is fed an outer verifier with a constant number $p_i$ of queries, its parameter $k=p_i$ is a constant, so the constants $p_2$ and $h$ are \emph{independent of $n$}: nothing compounds.
The Hadamard verifier closes the recursion because its exponential size is paid on a circuit of size $(\log\log n)^{O(1)}$.
\end{bigpicture}

\begin{remark}[Constructivity]
Following the proofs one sees that, given a satisfying assignment, the accepted proof is computable in polynomial time, and conversely a satisfying assignment can be recovered in polynomial time from any proof accepted with probability $\ge\tfrac12$ (Remark~18 of ALMSS: decoding polynomial and Hadamard codes is efficient).
This constructive form is what lets a PCP reduction be \emph{a Karp reduction with a witness map}, used in the Label Cover reductions of Lectures~5--7.
\end{remark}

\begin{remark}[How large is $q$?]
The constant is not made explicit in the paper: it depends on $\delta_0$ and $\alpha$ of Theorem~\ref{thm:polytest}, which are not computed.  Section~8 of ALMSS reports that the smallest $q$ that can be obtained went, in subsequent work, from $29$ (Bellare--Goldwasser--Lund--Russell) to $22$ (Feige--Kilian) to $16$ (Bellare--Goldreich--Sudan) to at most $9$ (the best that can be inferred directly from H\aa stad's work); all of these are for soundness~$\tfrac12$.
H\aa stad's optimal results (Lecture~\ref{lec:hastad}) concern \emph{three}-bit verifiers with soundness close to the trivial value: completeness $1-\eps$ and soundness $\tfrac12+\eps$ (accept iff a parity of the three bits equals a prescribed value), and perfect completeness with soundness $\tfrac78+\eps$ (accept iff a $3$-clause is satisfied).  They are tight because a random proof achieves $\tfrac12$, respectively $\tfrac78$.
The Hadamard verifier of this lecture on its own reads $6k+12$ bits per run, and $55$ repetitions bring $\tfrac{19}{20}$ down to $\tfrac1{16}$ (Exercise~\ref{ex:alm-reps}); this is a first indication of where the constant comes from.
\end{remark}


%% ===========================================================================
\section{Consequences, and the link to games}
\label{sec:alm-cons}
%% ===========================================================================

\paragraph{Max-3SAT.}  By Proposition~\ref{prop:pcp-equiv} (hardness direction) the theorem gives a constant $\eps_0>0$ such that $\Gap{1}{1-\eps_0}{3SAT}$ is $\NP$-hard; with $q$ queries and soundness $\tfrac12$ one may take $\eps_0=1/(2(q-2)2^q)$, and ALMSS's own computation gives $1/(q\,2^{q+1})$ in their Theorem~3 (where they credit the idea to Arora--Motwani--Safra--Sudan--Szegedy).
By the L-reductions of Papadimitriou and Yannakakis this transfers to every \sat{Max-SNP}-hard problem: none has a PTAS unless $\Ptime=\NP$.

\paragraph{Max-Clique.}  Let $V$ be the verifier of Theorem~\ref{thm:alm} ($r=c\log n$ random bits, $q$ queries, soundness $\tfrac12$).  Replace it by the verifier $V_k$ that uses a $k$-step random walk on a constant-degree expander on $\{0,1\}^r$
(Exercise~\ref{ex:expander}): $r+O(k)$ random bits, $kq$ queries, soundness $\le(\tfrac34)^{k-1}$.  Take $k=\lceil\log n\rceil$.  Then the FGLSS graph $H_x$ has $N=2^{r+O(k)+kq}=n^{O(1)}$ vertices and, by Theorem~\ref{thm:fglss}, clique number $2^{r'}$ if $x\in\sat{3SAT}$ and $\le2^{r'}(3/4)^{k-1}$ otherwise:
a gap of $(4/3)^{k-1}=n^{\Omega(1)}=N^{\eps_1}$.  So approximating the clique number within $N^{\eps_1}$ is $\NP$-hard \citep[Theorem~6]{ALMSS98}.  This is exactly the route indicated by their Theorem~2 (FGLSS plus ``recycled'' randomness),
and it supersedes Corollary~\ref{cor:as-clique}: the exponent is now polynomial.  H\aa stad's $N^{1-\eps}$ \citep{Hastad99} is discussed, though not proved, in Lecture~\ref{lec:hastad}.

\paragraph{Label Cover.}  Apply Corollary~\ref{cor:verifier-game} to the verifier of Theorem~\ref{thm:alm}: the constraint--variable game is a projection game with $|\Sigma_U|\le2^q$, $|\Sigma_W|=2$ and
\[
  x\in\sat{3SAT}\Rightarrow\val=1,\qquad x\notin\sat{3SAT}\Rightarrow\val\le1-\tfrac1{2q}.
\]
This is the starting point of Theorem~\ref{thm:lc}: ALMSS point out (Section~1.2) that their Theorem~4 ``provides the first known construction of a $2$-prover $1$-round proof system that uses logarithmic random bits and constant number of communication bits''.
Lecture~5 will show how Raz's theorem drives $1-\frac1{2q}$ to $\delta$ at the cost of an alphabet of size $\poly(1/\delta)$.

\paragraph{Transparent proofs of mathematics.}  The set of pairs $(\varphi,1^n)$ such that $\varphi$ has a proof of length $n$ in ZFC is in $\NP$.  By Theorem~\ref{thm:alm} every theorem therefore has a ``proof'' that a verifier using $O(\log n)$ random bits checks by reading $O(1)$ bits, and by the constructive form of the remark above the transformation from an ordinary proof to such a proof is polynomial.
This is the sense of the pitch in Lecture~\ref{lec:why} (``proofs that can be checked by reading only a few bits'').


%% ===========================================================================
\section{What the proof does not do, and what replaces it}
\label{sec:alm-limits}
%% ===========================================================================

Three features of the proof are worth noticing before moving on.

\begin{enumerate}
  \item \emph{There is no mechanism that makes a gap grow, other than repetition.}  Each module---linearity, tensor, low-degree, concatenation, curve-check---detects a corrupted proof with some \emph{constant} probability, and soundness $\tfrac12$ is reached by independent repetition (Proposition~\ref{prop:rep}).  The final constants are those of the weakest module (here $\delta_0$ of Theorem~\ref{thm:polytest}).  A \emph{gap-amplifying} step, the heart of Lecture~6, is a different idea.
  \item \emph{The only part with an explicit constant is the Hadamard verifier.}  The rest passes through the low-degree test with the field-size condition $|\F|\ge\alpha(d+1)^3$ and the non-explicit $\delta_0$.
  \item \emph{The algebra is used twice, for two different reasons.}  Polynomials give \emph{local} testability with few queries (low-degree test, self-correction, curves).  Hadamard codes give \emph{constant-size} testability at the price of exponential length.
\end{enumerate}

\citet{Dinur07} (Lecture~6; see Figure~\ref{fig:map}) replaces all of the polynomial machinery by a combinatorial operation on constraint graphs that amplifies the gap by a constant factor at each step, using expander graphs and random walks, so that after $O(\log n)$ steps the gap is constant.
The \emph{composition} step survives, but only with a verifier of \emph{constant} input size.  Dinur remarks that this one verifier can be found by exhaustive search, provided its existence has been established independently,
and that one may use ``Long-code based or Hadamard-code based'' constructions for it; the Hadamard-type verifier of Section~\ref{sec:had-ver} is the prototype.  (Dinur composes with an \emph{assignment tester}, \ie a PCP of proximity, which is a slightly stronger notion than our inner verifier; Section~6 of her paper gives the long-code instance.)
In short: Lecture~4 builds the theorem from algebra, Lecture~6 from combinatorics plus one small, constant-size verifier of the kind constructed in Section~\ref{sec:had-ver}.


%% ===========================================================================
\section*{Exercises}
\addcontentsline{toc}{section}{Exercises}
%% ===========================================================================

\begin{exercise}[\stars{$\star$}]\label{ex:maj3}
Let $F\colon\F_2^3\to\F_2$ be the majority function.  Compute all Fourier coefficients of $f=(-1)^F$, the acceptance probability of the linearity test via Theorem~\ref{thm:blr}, and the distance to the nearest linear function.  Check against a direct count of the $64$ pairs $(x,y)$.
Compare $\Pr[\text{reject}]$ with the two bounds $\Delta$ and $3\Delta$ of the theorem.
\end{exercise}

\begin{exercise}[\stars{$\star$}]\label{ex:frank}
Prove that for $M\in\F_2^{n\times n}$ of rank $r\ge1$, $\Pr_{u,v}[u^\top Mv=1]=\tfrac12-2^{-r-1}$.  \emph{Hint: change bases so that $M=\mathrm{diag}(1,\dots,1,0,\dots,0)$; then $u^\top Mv=\sum_{i\le r}u_iv_i$, and the bias of an inner product of $r$-bit strings is $2^{-r}$.}
Conclude that Lemma~\ref{lem:freivalds} is tight for $r=1$.
\end{exercise}

\begin{exercise}[\stars{$\star$}]\label{ex:quadeq}
Show that every $3$-CNF with $n$ variables and $m$ clauses can be written as a system of $2m$ quadratic equations over $\F_2$ in $n+m$ variables, each equation of the form $P(z)=0$ with $\deg P\le2$, which is satisfiable iff the formula is.
\emph{Hint: the clause $\ell_1\vee\ell_2\vee\ell_3$ is violated iff $(1+\ell_1)(1+\ell_2)(1+\ell_3)=1$; introduce $y=(1+\ell_1)(1+\ell_2)$.}  Why does the verifier of Section~\ref{sec:had-ver} work for this input directly, without a circuit?
\end{exercise}

\begin{exercise}[\stars{$\star\star$}]\label{ex:had-cheat}
Return to Example~\ref{ex:had}(b).  (a)~List all $z\in\F_2^5$ that violate exactly one of the four equations.  (b)~Verify that $z=0$, $u=e_3+e_5$, $v=e_4+e_5$ and $B=L_{uv^\top}$ pass the circuit test with probability~$1$.
(c)~Prove that no linear cheating proof (all of $X_1,A,B$ linear) is accepted with probability more than $\tfrac34$.  (d)~Find a \emph{non-linear} $A$ that is accepted with higher probability than your best linear cheating strategy, or explain why Theorem~\ref{thm:vplus} still guarantees at most $\tfrac{19}{20}$.
\end{exercise}

\begin{exercise}[\stars{$\star\star$}]\label{ex:alm-reps}
(a)~How many independent repetitions bring the error $\tfrac{19}{20}$ of $V^\oplus$ to at most $\tfrac1{16}$?  What is the resulting number of queries as a function of $k$?  (b)~By Theorem~\ref{thm:compose} the error after composing three verifiers each with error $\tfrac1{16}$ is $e+e_1-ee_1$ applied twice; compute it exactly.
\end{exercise}

\begin{exercise}[\stars{$\star\star$}]\label{ex:oracle-size}
The oracle $O$ of $V^{\rm par}$ is defined on $\F^{(w+1)l}$, a superpolynomial set.  Explain why the \emph{proof} is nevertheless of polynomial size.  How many entries of $O$ can be read, as a function of $n$?  Why is it legitimate for the soundness argument that $O$ is arbitrary on all other entries?
\end{exercise}

\begin{exercise}[\stars{$\star\star$}]\label{ex:curve-line}
Suppose we try to aggregate $l$ queries with a \emph{line} instead of a curve: ask for $p$ restricted to a line through $z_1$ and $z_2$ and so on.  (a)~Show that for $l$ points in general position no line contains them.  (b)~What do curves of degree $s$ cost in the degree of the restriction and in the soundness error $d s/|\F|$?  (c)~Why is it essential that $|\F|\gg dl$?
\end{exercise}

\begin{exercise}[\stars{$\star\star\star$}]\label{ex:alm-reading}
\emph{Reading.}  Compare Lemma~\ref{lem:compose} (Lecture~\ref{lec:as}) with Theorem~\ref{thm:compose}.  Which soundness clause of Definition~\ref{def:normal} corresponds to ``$E^{-1}$ of the oracles satisfies $C$'' in Definition~\ref{def:inner}?  Where does the inner system of this lecture use \emph{decoding} while the normal-form verifier of Lecture~\ref{lec:as} used \emph{closeness}?
Then identify, in the proof of Theorem~\ref{thm:par}, the single place where the unique-decoding radius is used, and what would break if $\Delta(X_i,p_i)$ could only be bounded by $\tfrac12+\eps$.
\end{exercise}


%% ===========================================================================
\section*{Notes and further reading}
\addcontentsline{toc}{section}{Notes and further reading}
%% ===========================================================================

\paragraph{Where this lecture comes from.}
Everything is from \citet{ALMSS98}: the outer/inner formalism and composition are Sections~2.1; the Hadamard verifier is Sections~5--6; parallelization is Section~7; the \sat{Max-3SAT} consequence is Section~3.
The linearity test and self-correction are due to \citet{BLR90}; the tensor test is a Freivalds-type matrix-product check \citep{Freivalds79}; the Fourier-analytic form of the linearity analysis follows the standard modern treatment, for which see the detailed study of the test in characteristic two by \citet{BCHKS96}.
The low-degree test with a table of lines and its analysis combine \citet{RubinfeldSudan96} and the bivariate lemma of \citet{AroraSafra98}; the correction procedure goes back to Beaver--Feigenbaum and Lipton and, in the form with constant queries, to Gemmell--Lipton--Rubinfeld--Sudan--Wigderson, as cited in ALMSS.
The numerical Example~\ref{ex:had} was computed for these notes.

\paragraph{On attributions.}
The conference versions of \citet{AroraSafra98} and \citet{ALMSS98} both appeared at FOCS~1992 (pp.~2--12 and 13--22).  The constant-query theorem is rightly credited to ALMSS, but its proof rests on the composition technique and the polylogarithmic verifier of Arora and Safra;
ALMSS say so in their Section~1.2.  The \emph{linearity-testing / Hadamard} half of the proof was in turn inspired, as they note, by work on program checking.

\paragraph{Small points in the text of \citet{ALMSS98} that we handled differently.}
Five points, as we read the paper; none affects the theorem, and we list them so that a reader comparing with the original is not surprised.
(i)~Step~4(c) of $V^\oplus$ compares the self-corrected value of $B$ at $\mathrm{coeff}_P$ with $0$, while $\mathrm{coeff}_p$ is defined only for $p(0)=0$ and the output polynomial (zero iff the output gate accepts, \ie $z_n+1$ over $\F_2$) has a constant term; we compare with the constant term of $\sum_jr_jP_j$ (Definition~\ref{def:vplus}), which repairs this.
(ii)~The number of queries of $V^{\rm par}$ is stated as $3k+5$; counting the concatenation test for $Z$ gives $3k+8$.
(iii)~The decoder of $E^{\rm par}$ is defined with the nearest polynomial of degree $mh$, whereas the soundness proof produces the nearest polynomial of the degree $d$ that Poly-test certifies; we define the decoder with degree $d$, which agrees with the original on all honest codewords.
(iv)~The linearity theorem is quoted with closeness $2\delta$; the Fourier argument gives $\delta$, which we used to get soundness $\tfrac{19}{20}$ rather than $\tfrac{35}{36}$.
(v)~In Section~7.5 the exponent $c_3$ is used both for the circuit size (as in the statement of Theorem~39) and for the number of queries $l$; the field size $\log^{\max\{2+c_3,6\}}N$ must dominate $dl$, so it is the exponent of the number of queries that matters.

\paragraph{What is not covered.}
We have not proved Theorem~\ref{thm:polytest} (the low-degree test with lines) or Theorem~\ref{thm:seq} beyond what Lecture~\ref{lec:as} gave, and we have not discussed the $\PCP[\log n,(\log\log n)^2]$ intermediate result.
Lecture~5 turns to the \emph{soundness} side: Raz's parallel repetition theorem, which the constant-query PCP of this lecture feeds with a projection game of value $1-\frac1{2q}$.

\chapter{Parallel Repetition: Driving the Soundness Down}
\label{lec:raz}

{\sffamily\large\color{softgray}Playing a game $k$ times at once is not the same as playing it $k$ times\par}
\vspace{0.5em}

\providecommand{\lamp}{\lambda_{+}}

\noindent
Lecture~\ref{lec:why} ended with a promise.  The PCP theorem gives a game of value $1$ or $1-\eta_0$; repeating it $k=O(\log\frac1\delta)$ times in parallel gives a game of value $1$ or at most $\delta$,
at the price of an alphabet of size $\poly(1/\delta)$.  The proof there was two lines long: ``apply Theorem~\ref{thm:raz}''.  This lecture opens that black box.

The difficulty is easy to state.  A verifier that repeats itself \emph{sequentially} multiplies its acceptance probabilities (Proposition~\ref{prop:pcp-equiv}).
Two provers cannot be questioned sequentially $k$ times without $k$ rounds of interaction, so the referee sends all $k$ questions at once, and each prover sees all $k$ of its questions before answering any of them.
It may therefore answer coordinate~$1$ using what it learned about coordinate~$2$.  Does that help the provers?  It does (Section~\ref{sec:notpower}), and for years it was open whether the value even tends to $0$; a non-constructive proof that it does came first, as Raz notes in his introduction.
\citet{Raz98} proved that it tends to $0$ exponentially fast in $k$, with a proof of about forty pages.

We shall do three things.  First, we meet the examples that explain why the answer is not obvious: a two-answer game whose two-fold repetition beats the square of its value, and a game of value $\frac12$ whose two-fold repetition still has value~$\frac12$.
Second, we prove a repetition theorem \emph{completely}, for projection games whose ``collision graph'' is an expander.  The proof, due to \citet{DinurSteurer14}, is linear algebra: it takes a few pages, and it is enough to establish the Label Cover theorem
(Theorem~\ref{thm:lc}) with explicit constants, using a one-line trick that makes any game expanding.  Third, we state what is known without the expansion hypothesis, sketch how Raz and Holenstein prove it, and explain why repetition can never,
by itself, prove the Unique Games Conjecture.

\begin{goalbox}{After this lecture you should be able to\ldots}
\begin{enumerate}[leftmargin=*,label=(\arabic*)]
  \item compute the $k$-fold value of a small game by exhaustive search, and exhibit a strategy that beats $\val(G)^k$ (Section~\ref{sec:notpower});
  \item define the collision value $\|G\|$ and prove $\val(G)^2\le\|G\|^2\le\val(G)$ (Section~\ref{sec:operator});
  \item prove that $\|G\otimes H\|\le\|G\|$ and, more sharply, $\|G\otimes H\|\le\lamp(G)\,\|H\|$ (Sections~\ref{sec:tensor}--\ref{sec:product});
  \item prove that $\lamp(G)\le1-\gamma\eta/54$ for a $\gamma$-expanding projection game of value $\le1-\eta$, and deduce the Label Cover theorem with an explicit number of repetitions (Sections~\ref{sec:expander}--\ref{sec:lc});
  \item state the repetition theorems of Raz, Holenstein, Rao and Dinur--Steurer, say what fails (strong repetition) and why that obstructs a repetition-based proof of the Unique Games Conjecture (Sections~\ref{sec:general}--\ref{sec:ugc-rep}).
\end{enumerate}
\end{goalbox}


%% ===========================================================================
\section{Repeating a game is not raising its value to a power}
\label{sec:notpower}
%% ===========================================================================

We use the games of Lecture~\ref{lec:why} (Definition~\ref{def:game}) with one harmless generalization: the referee draws the edge from an arbitrary probability distribution~$\mu$ rather than the uniform one.
All games in Lectures~\ref{lec:why}--\ref{lec:games} have uniform~$\mu$, and a rational~$\mu$ can always be made uniform by repeating edges (see Remark~\ref{rem:multiplicity} below).

\begin{definition}[Weighted projection game]\label{def:wgame}
A \defn{projection game} $G$ has question sets $U$ (Alice) and $W$ (Bob), label sets $\Sigma_U,\Sigma_W$, a finite set $E$ of edges $e=(u_e,w_e)\in U\times W$ (parallel edges allowed) with a probability distribution $\mu$ on $E$, and for every edge a map $\pi_e\colon\Sigma_U\to\Sigma_W$.
The referee draws $e\sim\mu$, sends $u_e$ to Alice and $w_e$ to Bob, and accepts answers $a\in\Sigma_U$, $b\in\Sigma_W$ iff $\pi_e(a)=b$.
We write $\mu_U,\mu_W$ for the marginals of $\mu$ on questions (assumed positive), $\mu(\cdot\mid w)$ for $\mu$ conditioned on $w_e=w$, and $\val(G)$ for the maximum acceptance probability over deterministic strategies.
The \defn{product} $G\otimes H$ has questions $U\times U'$ and $W\times W'$, labels $\Sigma_U\times\Sigma_{U'}$ and $\Sigma_W\times\Sigma_{W'}$, edges $E\times E'$ with the product distribution, and constraint maps $\pi_e\times\pi_{e'}$; $G^{\otimes k}$ is the $k$-fold product.
\end{definition}

\noindent
This is the $k$-fold parallel repetition of Lecture~\ref{lec:why}: the referee draws $k$ independent edges, sends all $k$ Alice-questions to Alice and all $k$ Bob-questions to Bob, and accepts iff every coordinate is accepted.
Note that \emph{totality} of $\pi_e$ is part of the definition: every label of Alice has exactly one accepted partner.  Unique Games are the case of bijective maps.

\begin{remark}\label{rem:multiplicity}
If $\mu$ takes rational values, write $\mu(e)=n_e/N$; replacing $e$ by $n_e$ parallel copies gives an equivalent game with the uniform distribution on $N$ edges.
\end{remark}

\subsection{The easy inequalities}

\begin{lemma}[Elementary facts]\label{lem:rep-basic}
For every game $G$ and $k\ge2$: $\ \val(G)^k\ \le\ \val(G^{\otimes k})\ \le\ \val(G^{\otimes(k-1)})$.  In particular $\val(G)=1$ implies $\val(G^{\otimes k})=1$.
\end{lemma}
\begin{proof}
Lower bound: both provers play an optimal strategy of $G$ independently in each coordinate; the $k$ coordinates are independent.
Upper bound: take a deterministic strategy $(f,g)$ for $G^{\otimes k}$ and write $e=(e_1,\dots,e_k)$.  Winning all $k$ coordinates implies winning the first $k-1$, so
$\val(G^{\otimes k})\le\E_{e_k}\Pr_{e_1..e_{k-1}}[\text{first }k-1\text{ coordinates won}\mid e_k]$, hence some fixed edge $e_k=(u_k,w_k)$ attains at least this average.  With $u_k$ and $w_k$ hard-wired, $f(\cdot,u_k)$ and $g(\cdot,w_k)$ are strategies for $G^{\otimes(k-1)}$ that win with probability at least $\val(G^{\otimes k})$.
\end{proof}

\noindent
Sequential repetition would give $\val(G)^k$ exactly.  Is the first inequality an equality?  Here is the smallest example showing that it is not.

\begin{example}[CHSH]\label{ex:chsh}
Let $U=W=\{0,1\}$ and $\Sigma_U=\Sigma_W=\{0,1\}$, with the uniform distribution on the four pairs $(x,y)$ and the constraint $a\oplus b=x\wedge y$, \ie $\pi_{(x,y)}(a)=a\oplus xy$.  Every $\pi_{(x,y)}$ is a bijection, so this is a Unique Game with two labels.

\emph{Its value is $\frac34$.}  The four constraints $f(x)\oplus g(y)=xy$ cannot all hold: adding them, the left sides cancel in pairs and give $0$, while the right sides add up to $0+0+0+1=1$.  The strategy $f=g=0$ satisfies three.

\emph{Its two-fold repetition has value $\frac58>\big(\frac34\big)^2=\frac9{16}$.}  Questions are now $x=(x_1,x_2)$ and $y=(y_1,y_2)$, answers are pairs, and both coordinates must satisfy $a_i\oplus b_i=x_iy_i$.  Consider the strategy
\[
  a(x)=\begin{cases}(1,0)&x=(1,1)\\(0,0)&\text{otherwise}\end{cases}\qquad
  b(y)=\begin{cases}(0,1)&y=(1,1)\\(0,0)&\text{otherwise.}\end{cases}
\]
Table~\ref{tab:chsh-win} shows that it wins on $10$ of the $16$ question pairs: $10/16=\frac58$.
Exhaustive search over all $4^4\cdot4^4$ strategies confirms that nothing does better; for $k=3$ an exhaustive search over Alice's $8^8$ strategies, with Bob responding optimally, gives $\val(G^{\otimes3})=\frac{31}{64}$, which exceeds \mbox{$\big(\frac34\big)^3=\frac{27}{64}$}.
\end{example}

\begin{table}[t]
\centering\small
\renewcommand{\arraystretch}{1.15}
\begin{tabular}{@{}c|cccc@{}}
\toprule
$x\backslash y$ & $00$ & $01$ & $10$ & $11$\\
\midrule
$00$ & \checkmark & \checkmark & \checkmark & $\cdot$\\
$01$ & \checkmark & $\cdot$ & \checkmark & \checkmark\\
$10$ & \checkmark & \checkmark & $\cdot$ & $\cdot$\\
$11$ & $\cdot$ & $\cdot$ & \checkmark & \checkmark\\
\bottomrule
\end{tabular}
\caption{The winning question pairs of the two-fold CHSH strategy of Example~\ref{ex:chsh}: $10$ of $16$.  A single coordinate of the same strategy would win only $\frac34$ of the time on its own; the strategy gains by letting each prover's answer in coordinate~$i$ depend on the \emph{other} coordinate's question.}
\label{tab:chsh-win}
\end{table}

\noindent
The gain is real but small: $\frac58$ against $\frac9{16}$.  The next example shows that without structure, repetition can fail to help at all.

\begin{example}[Feige's non-interactive agreement game]\label{ex:feige}
Let $x,y\in\{0,1\}$ be independent uniform questions to Alice and Bob, and let the label set of both be $\{A_0,A_1,B_0,B_1\}$.  The provers win iff they give the \emph{same} label and it is either $A_x$ (``we agree on Alice's question'') or $B_y$ (``we agree on Bob's question'').

\emph{Its value is $\frac12$.}  Fix deterministic strategies $f,g$.  For a given Bob-question $y$, the label $g(y)$ can win against at most one $x$ if $g(y)=A_s$ (namely $x=s$), against at most two if $g(y)=B_y$ (namely those $x$ with $f(x)=B_y$), and against none otherwise.
If some $y$ is won twice then $g(y)=B_y$ and $f\equiv B_y$, and then the other $y'$ is never won; otherwise each $y$ is won at most once.  In either case at most $2$ of the $4$ question pairs are won.  Always answering $A_0$ achieves $\frac12$.

\emph{Its two-fold repetition has value $\frac12$ as well.}  Let Alice, on $(x_1,x_2)$, answer $(A_{x_1},B_{x_1})$, and Bob, on $(y_1,y_2)$, answer $(A_{y_2},B_{y_2})$.  In coordinate~$1$ the pair $(A_{x_1},A_{y_2})$ is accepted iff $y_2=x_1$ (the label $A_{x_1}$ is Alice's own question); in coordinate~$2$ the pair $(B_{x_1},B_{y_2})$ is accepted iff $y_2=x_1$ (the label $B_{y_2}$ is Bob's own question).  Both conditions are the same event $\{x_1=y_2\}$, which has probability $\frac12$.
Together with Lemma~\ref{lem:rep-basic}, $\val(G^{\otimes2})=\val(G)=\frac12$.  An exhaustive search over all strategies confirms both numbers.
\end{example}

\begin{pitfall}{repetition does not have to help, and does not help as the product says}
Example~\ref{ex:feige} is a game for which \emph{no} bound of the form $\val(G^{\otimes2})\le\varphi(\val(G))$ with $\varphi(x)<x$ can hold; a repetition theorem must therefore speak about $k$ large, and about the size of the answer sets or the structure of the constraints.
Example~\ref{ex:chsh} shows that even for Unique Games the true value of the repeated game lies strictly between $\val(G)^k$ and $\val(G)$.  Neither game is exotic: the first is a free game on two questions, the second is the XOR game behind Bell's inequality.
The Feige game is \emph{not} a projection game in our sense (a label $A_{1-x}$ of Alice has no accepted partner at all), and Lemma~\ref{lem:val-coll} below uses totality; this is no accident.
\end{pitfall}

We will return to these two examples in Section~\ref{sec:general}.  For now the lesson is: to prove that $\val(G^{\otimes k})\to0$ we cannot reason one coordinate at a time.  We need a quantity of $G$ that \emph{is} multiplicative-like under products and still controls the value.


%% ===========================================================================
\section{A projection game as a linear operator}
\label{sec:operator}
%% ===========================================================================

\subsection{Strategies, the operator $P_G$, and the success probability}

Let $G$ be a projection game.  An \defn{assignment} for Alice is a function $f\colon U\times\Sigma_U\to[0,1]$ with $\sum_af(u,a)=1$ for every $u$; $f(u,a)$ is the probability that Alice answers $a$ to $u$.
A deterministic \defn{labeling} $\ell\colon U\to\Sigma_U$ is the assignment $f_\ell(u,a)=\mathbf 1[a=\ell(u)]$.  Assignments for Bob are defined in the same way on $W\times\Sigma_W$.
On real functions $f,f'$ on $U\times\Sigma_U$ and $g,g'$ on $W\times\Sigma_W$ we use the inner products
\[
  \langle f,f'\rangle=\E_{u\sim\mu_U}\sum_af(u,a)f'(u,a),\qquad
  \langle g,g'\rangle=\E_{w\sim\mu_W}\sum_bg(w,b)g'(w,b),
\]
and the norms $\|f\|=\langle f,f\rangle^{1/2}$.
The \defn{operator of the game} sends a function on $U\times\Sigma_U$ to a function on $W\times\Sigma_W$:
\[
  (P_Gf)(w,b)\;=\;\E_{e\sim\mu(\cdot\mid w)}\ \sum_{a\colon\pi_e(a)=b}f(u_e,a).
\]
In words: to evaluate $P_Gf$ at $w$, pick a random neighbour of $w$, let Alice answer according to $f$, and project her answer to a label of Bob.  Because every $\pi_e$ is a \emph{total} map, $\sum_b(P_Gf)(w,b)=\E_{e\mid w}\sum_af(u_e,a)=1$ when $f$ is an assignment:
\emph{$P_Gf$ is again an assignment}, this time for Bob.

\begin{lemma}[The success probability is a bilinear form]\label{lem:bilinear}
If Alice plays the assignment $f$ and Bob plays $g$, the referee accepts with probability $\langle g,P_Gf\rangle$.  Consequently $\val(G)=\max\langle g,P_Gf\rangle$ over assignments, and the maximum is attained by deterministic ones.
\end{lemma}
\begin{proof}
The acceptance probability is $\E_{e\sim\mu}\sum_{a,b}f(u_e,a)g(w_e,b)\mathbf 1[\pi_e(a)=b]$.  Condition on $w_e=w$ to get $\E_w\sum_bg(w,b)\,\E_{e\mid w}\sum_{a\colon\pi_e(a)=b}f(u_e,a)=\langle g,P_Gf\rangle$.
This expression is bilinear in $(f,g)$, and the set of assignments is a product of simplices whose extreme points are the deterministic ones.
\end{proof}

\subsection{The collision value}

The optimal $g$ for a given $f$ is obvious (answer the most likely label of $P_Gf$), so the interesting object is $P_Gf$ itself.  Its length is the quantity that will behave well under products.

\begin{definition}[Collision value]\label{def:coll}
$\|G\|\;=\;\max\{\|P_Gf\|:\ f\text{ an assignment for Alice}\}$.
\end{definition}

\noindent
Since $f\mapsto\|P_Gf\|$ is convex, the maximum is attained at a deterministic labeling.  The name comes from the following experiment, in which Bob has been eliminated.

\begin{lemma}[The collision game]\label{lem:collgame}
$\|P_Gf\|^2$ is the probability of acceptance in this experiment: draw $w\sim\mu_W$ and two edges $e,e'\sim\mu(\cdot\mid w)$ independently; let two independent copies of Alice answer $a\sim f(u_e,\cdot)$ and $a'\sim f(u_{e'},\cdot)$; accept iff $\pi_e(a)=\pi_{e'}(a')$.
\end{lemma}
\begin{proof}
Expand the square and exchange the sum over $b$ with the sums over $a,a'$:
\begin{align*}
  \|P_Gf\|^2&=\E_w\sum_b\Big(\E_{e\mid w}\sum_{a\colon\pi_e(a)=b}f(u_e,a)\Big)^2\\
  &=\E_w\ \E_{e,e'\mid w}\ \sum_{a,a'}f(u_e,a)\,f(u_{e'},a')\,\mathbf 1\big[\pi_e(a)=\pi_{e'}(a')\big].\qedhere
\end{align*}
\end{proof}

\noindent
For the clause--variable game of a $3$-CNF this is the following experiment (Figure~\ref{fig:collision}): choose a random variable $x$ and two random clauses containing it; the two clause-provers must give satisfying assignments of their clauses that \emph{agree on $x$}.
The same symmetrization appears, implicitly, throughout the PCP literature: for example when a ``line versus point'' low-degree test is turned into a ``line versus line'' test \citep[footnote to Section~2.3]{DinurSteurer14}.

\begin{figure}[t]
\centering
\begin{tikzpicture}[x=1cm,y=1cm,every node/.style={font=\small},
   cl/.style={draw,rounded corners=2pt,minimum width=5.3cm,minimum height=0.95cm,fill=accent!8,align=center},
   vr/.style={draw,circle,minimum size=0.9cm,fill=warn!10},
   ok/.style={font=\bfseries\small,color=green!50!black},
   no/.style={font=\bfseries\small,color=claim}]
  \node[font=\sffamily\small\bfseries,color=accent,anchor=west] at (-2.7,2.55) {(a)\ \ a collision};
  \node[cl] (c1) at (0,1.5) {$C_1=(x_1\vee x_2\vee\neg x_3)$\\ label $(x_1,x_2,x_3)=(1,0,0)$};
  \node[cl] (c2) at (0,0) {$C_2=(\neg x_1\vee x_3\vee x_4)$\\ label $(x_1,x_3,x_4)=(0,0,0)$};
  \node[vr] (x3) at (5.2,0.75) {$x_3$};
  \draw[-{Stealth[length=2.2mm]},thick] (c1.east) -- node[above,sloped,font=\footnotesize] {$x_3=0$} (x3);
  \draw[-{Stealth[length=2.2mm]},thick] (c2.east) -- node[below,sloped,font=\footnotesize] {$x_3=0$} (x3);
  \node[ok,anchor=west] at (5.9,0.75) {agree: accept};
  \begin{scope}[yshift=-3.9cm]
  \node[font=\sffamily\small\bfseries,color=accent,anchor=west] at (-2.7,2.55) {(b)\ \ no collision};
  \node[cl] (d1) at (0,1.5) {$C_1=(x_1\vee x_2\vee\neg x_3)$\\ label $(x_1,x_2,x_3)=(1,0,0)$};
  \node[cl] (d2) at (0,0) {$C_2=(\neg x_1\vee x_3\vee x_4)$\\ label $(x_1,x_3,x_4)=(1,1,0)$};
  \node[vr] (y3) at (5.2,0.75) {$x_3$};
  \draw[-{Stealth[length=2.2mm]},thick] (d1.east) -- node[above,sloped,font=\footnotesize] {$x_3=0$} (y3);
  \draw[-{Stealth[length=2.2mm]},thick] (d2.east) -- node[below,sloped,font=\footnotesize] {$x_3=1$} (y3);
  \node[no,anchor=west] at (5.9,0.75) {disagree: reject};
  \end{scope}
\end{tikzpicture}
\caption{The collision game of the clause--variable game of Figure~\ref{fig:cv}.  The referee picks the variable $x_3$ and two clauses containing it; the two copies of Alice answer with satisfying assignments of their clauses (both are satisfying in (a) and (b)) and are accepted iff their projections to $x_3$ coincide.  Bob does not appear.}
\label{fig:collision}
\end{figure}

\begin{lemma}[Collision value versus value]\label{lem:val-coll}
$\ \val(G)^2\ \le\ \|G\|^2\ \le\ \val(G).$
\end{lemma}
\begin{proof}
Let $f,g$ be deterministic and optimal, so $\val(G)=\langle g,P_Gf\rangle\le\|g\|\,\|P_Gf\|\le\|G\|$, because $\|g\|^2=\E_w\sum_bg(w,b)^2=1$.
For the other inequality let $f$ attain $\|G\|$.  Since $P_Gf$ is an assignment for Bob, Lemma~\ref{lem:bilinear} gives $\|G\|^2=\langle P_Gf,P_Gf\rangle\le\val(G)$.
\end{proof}

\noindent
In particular $\val(G)=1$ iff $\|G\|=1$.  The squaring is harmless for us: ``exponentially small'' survives a square root, at the cost of a factor $2$ in the number of repetitions.
(The proof of the second inequality is the place where totality of the maps $\pi_e$ is used; for a partial projection one must increase some entries of $P_Gf$ to turn it into an assignment, and this is how \citet{DinurSteurer14} proceed.)

\begin{example}[CHSH, by hand]\label{ex:chsh-coll}
In the collision game of Example~\ref{ex:chsh} the referee picks a uniform $y$ and two independent uniform $x,x'$; a deterministic Alice answers $f(x),f(x')$; the pair is accepted iff $f(x)\oplus xy=f(x')\oplus x'y$.
If $x=x'$ (probability $\frac12$) this always holds.  If $x\ne x'$ it holds iff $f(0)\oplus f(1)=y$, which happens for exactly one of the two values of $y$.  So $\|P_Gf\|^2=\frac12+\frac12\cdot\frac12=\frac34$ for every deterministic $f$, \ie
$\|G\|^2=\frac34=\val(G)$.  In Lemma~\ref{lem:val-coll} the left inequality is strict here ($\frac9{16}<\frac34$) and the right one is an equality.
\end{example}


%% ===========================================================================
\section{The collision value under products}
\label{sec:tensor}
%% ===========================================================================

For two functions on question--label pairs the product game acts coordinatewise.  Write $L(X)$ for the real functions on $X$ with the inner product above.

\begin{lemma}[Products are tensor products]\label{lem:tensor}
$P_{G\otimes H}=P_G\otimes P_H$: for $F\colon(U\times U')\times(\Sigma_U\times\Sigma_{U'})\to\R$,
\[
  (P_{G\otimes H}F)\big((w,w'),(b,b')\big)=\E_{e\sim\mu(\cdot\mid w)}\ \E_{e'\sim\mu'(\cdot\mid w')}\ \sum_{a\colon\pi_e(a)=b}\ \sum_{a'\colon\pi_{e'}(a')=b'}F\big((u_e,u_{e'}),(a,a')\big).
\]
\end{lemma}
\begin{proof}
Given $(w,w')$, the edge of $G\otimes H$ is a product of an edge of $G$ and an edge of $H$, drawn independently from $\mu(\cdot\mid w)$ and $\mu'(\cdot\mid w')$.
\end{proof}

\noindent
For the value this tensor structure is of little use, since the value of a product is not the product of the values (Section~\ref{sec:notpower}).  For the \emph{collision} value we can at least prove the monotonicity that is trivial for the value.

\begin{lemma}[Playing in parallel cannot raise the collision value]\label{lem:coll-mono}
For all projection games $G,H$: $\ \|G\otimes H\|\le\|G\|$ and $\|G\otimes H\|\le\|H\|$.
\end{lemma}
\begin{proof}
By symmetry it suffices to prove the first.  Let $F$ be an assignment for $G\otimes H$.  For $u'\in U'$ and $w'\in W'$ define functions on $U\times\Sigma_U$:
\[
  F_{u'}(u,a)=\sum_{a'}F\big((u,u'),(a,a')\big),\qquad F_{w'}=\E_{e'\sim\mu'(\cdot\mid w')}F_{u_{e'}}.
\]
$F_{u'}$ is an assignment for $G$ (non-negative, and $\sum_aF_{u'}(u,a)=\sum_{a,a'}F=1$), hence so is the average $F_{w'}$.
Fix $(w,w')$ and $b$, and put $x_{b'}=(P_{G\otimes H}F)((w,w'),(b,b'))\ge0$ as in Lemma~\ref{lem:tensor}.  Summing over $b'$ removes the constraint on $a'$:
\[
  \sum_{b'}x_{b'}=\E_{e\mid w}\ \E_{e'\mid w'}\sum_{a\colon\pi_e(a)=b}F_{u_{e'}}(u_e,a)=\E_{e\mid w}\sum_{a\colon\pi_e(a)=b}F_{w'}(u_e,a)=(P_GF_{w'})(w,b).
\]
For non-negative numbers $\sum_{b'}x_{b'}^2\le\big(\sum_{b'}x_{b'}\big)^2$.  Therefore
\[
  \|P_{G\otimes H}F\|^2=\E_{w,w'}\sum_{b,b'}x_{b'}^2\ \le\ \E_{w'}\ \E_w\sum_b\big((P_GF_{w'})(w,b)\big)^2=\E_{w'}\|P_GF_{w'}\|^2\ \le\ \|G\|^2.\qedhere
\]
\end{proof}

\noindent
Lemma~\ref{lem:coll-mono} says that $\|G^{\otimes k}\|$ is non-increasing in $k$ (compare Lemma~\ref{lem:rep-basic}).  But it does not say that it \emph{decreases}.  For that we need to compare $G\otimes H$ with $H$ alone, and this is where the argument becomes interesting.


%% ===========================================================================
\section{The product theorem}
\label{sec:product}
%% ===========================================================================

Fix $G$.  We look for a number $\rho_G<1$, depending only on $G$, such that $\|G\otimes H\|\le\rho_G\|H\|$ \emph{for every} projection game~$H$.
Given such a $\rho_G$, $\|G^{\otimes k}\|\le\rho_G^{k-1}\|G\|$ follows at once, and the value of $G^{\otimes k}$ is bounded by Lemma~\ref{lem:val-coll}.

Suppose $F$ is a best assignment for $G\otimes H$ and consider first the ``null'' version of $G$ that never rejects.  This is the trivial game below; comparing $G\otimes H$ with it, the part coming from $H$ will cancel.

\begin{definition}[The trivial game and the parameter $\lamp(G)$]\label{def:lamp}
The \defn{trivial game} $T_G$ has the same Alice side as $G$ (questions $U$, labels $\Sigma_U$, marginal $\mu_U$), Bob's questions are again $U$, Bob has the single label $\ast$, and the edges are $(u,u)$ with probability $\mu_U(u)$.  Every strategy wins.  Its operator computes the mass $(P_{T_G}h)(u,\ast)=\sum_ah(u,a)$ of a function $h\ge0$, so $\|P_{T_G}h\|^2=\E_u\big(\sum_ah(u,a)\big)^2$.
Define
\[
  \lamp(G)\;=\;\sup\Big\{\frac{\|P_Gh\|}{\|P_{T_G}h\|}\ :\ h\colon U\times\Sigma_U\to\R_{\ge0},\ h\ne0\Big\}.
\]
\end{definition}

\noindent
The supremum is over all non-negative functions, not only assignments, so $\|G\|\le\lamp(G)$ (take $h$ an assignment: $\|P_{T_G}h\|=1$).  Also $\lamp(G)\le1$, because for $m(u)=\sum_ah(u,a)$,
$\|P_Gh\|^2=\E_w\sum_b\big(\E_{e\mid w}\sum_{a\colon\pi_e(a)=b}h(u_e,a)\big)^2\le\E_w\big(\E_{e\mid w}m(u_e)\big)^2\le\E_um(u)^2$ by $\sum x_b^2\le(\sum x_b)^2$ and Jensen.

\begin{theorem}[Product theorem; Dinur--Steurer]\label{thm:product}
For all projection games $G,H$: $\ \|G\otimes H\|\le\lamp(G)\cdot\|H\|$.
\end{theorem}
\begin{proof}
Let $F$ be an assignment for $G\otimes H$ with $\|P_{G\otimes H}F\|=\|G\otimes H\|$.  Define a non-negative function $h$ of $x=(u,a)$ and $j=(w',b')$ by applying only the $H$-half of the operator:
\[
  h(x;j)\;=\;\E_{e'\sim\mu'(\cdot\mid w')}\ \sum_{a'\colon\pi_{e'}(a')=b'}F\big((u,u_{e'}),(a,a')\big).
\]
By Lemma~\ref{lem:tensor}, $P_{G\otimes H}F=(P_G\otimes\mathrm{Id})h$ and $P_{T_G\otimes H}F=(P_{T_G}\otimes\mathrm{Id})h$.
The assignment $F$ has the same Alice side in $G\otimes H$ and in $T_G\otimes H$, so it is an assignment for $T_G\otimes H$ as well, and by Lemma~\ref{lem:coll-mono}
\[
  \|(P_{T_G}\otimes\mathrm{Id})h\|=\|P_{T_G\otimes H}F\|\ \le\ \|T_G\otimes H\|\ \le\ \|H\|.
\]
(The left side is positive: for an assignment $F$ the entries of $(P_{T_G}\otimes\mathrm{Id})h$ at $(u,w')$ sum over $b'$ to $1$.)  Dividing,
\[
  \frac{\|G\otimes H\|}{\|H\|}\ \le\ \frac{\|(P_G\otimes\mathrm{Id})h\|}{\|(P_{T_G}\otimes\mathrm{Id})h\|}.
\]
Both numerator and denominator are sums over the columns $j$ of $h$: with $h_j=h(\cdot\,;j)\ge0$ and weights $m_j=\mu_{W'}(w')$,
$\|(P_G\otimes\mathrm{Id})h\|^2=\sum_jm_j\|P_Gh_j\|^2$ and $\|(P_{T_G}\otimes\mathrm{Id})h\|^2=\sum_jm_j\|P_{T_G}h_j\|^2$.
A ratio of sums with positive denominators is at most the largest ratio of corresponding terms, and for each $j$ the ratio is at most $\lamp(G)^2$ by definition (columns with $h_j=0$ contribute $0$ to both sums).
\end{proof}

\noindent
The mechanism is worth a sentence.  The function $h$ summarizes everything that $H$ did to the strategy; the factorization $P_{G\otimes H}=(P_G\otimes\mathrm{Id})(\mathrm{Id}\otimes P_H)$ lets us discard $H$ and keep only the best $h$.
The denominator, supplied by the trivial game, is what turns an inequality about two interacting games into one about~$G$ alone.

\begin{corollary}[Repetition bound in terms of $\lamp$]\label{cor:lamp}
For every projection game $G$ and $k\ge1$: $\ \val(G^{\otimes k})\ \le\ \|G^{\otimes k}\|\ \le\ \lamp(G)^k$.
\end{corollary}
\begin{proof}
Apply Theorem~\ref{thm:product} with $H=G^{\otimes(k-1)}$ repeatedly, then use $\|G\|\le\lamp(G)$ and Lemma~\ref{lem:val-coll}.
\end{proof}

\begin{example}[CHSH, continued]\label{ex:chsh-lamp}
Let $h=m\cdot f_\ell$ for a labeling $\ell$ of Alice's two questions and weights $m=(m_0,m_1)\ge0$.  The computation of Example~\ref{ex:chsh-coll}, now with weights, gives $\|P_Gh\|^2=\frac14(m_0^2+m_1^2)+\frac14m_0m_1$ for \emph{every} $\ell$ (the two questions collide for exactly one of the two values of $y$), while $\|P_{T_G}h\|^2=\frac12(m_0^2+m_1^2)$.
The largest value of the ratio of these two quadratic forms is $\frac34$, attained at $m=(1,1)$.  Deterministic $h$ suffice for the supremum (Step~1 of the proof of Theorem~\ref{thm:expander} below), so $\lamp(G)^2=\frac34$.  Corollary~\ref{cor:lamp} reads $\val(G^{\otimes k})\le(\frac34)^{k/2}$.

This bound is tight for the \emph{collision value}, for every $k$: the product assignment $f^{\otimes k}$ has $P_{G^{\otimes k}}f^{\otimes k}=(P_Gf)^{\otimes k}$, whose norm is $\|P_Gf\|^k=(\sqrt3/2)^k$, so $\|G^{\otimes k}\|=(\sqrt3/2)^k=\lamp(G)^k$.
It is not tight for the value: Table~\ref{tab:chsh-sandwich} puts the chain $\val(G)^k\le\val(G^{\otimes k})\le\|G^{\otimes k}\|=\lamp(G)^k$ next to the true values.
\end{example}

\begin{table}[t]
\centering\small
\renewcommand{\arraystretch}{1.2}
\begin{tabular}{@{}rcccc@{}}
\toprule
$k$ & $\val(G)^k$ & $\val(G^{\otimes k})$ & $\|G^{\otimes k}\|$ & $\lamp(G)^k$\\
\midrule
$1$ & $0.7500$ & $0.7500$ & $0.8660$ & $0.8660$\\
$2$ & $0.5625$ & $0.6250$ & $0.7500$ & $0.7500$\\
$3$ & $0.4219$ & $0.4844$ & $0.6495$ & $0.6495$\\
\bottomrule
\end{tabular}
\caption{CHSH: the true values of Example~\ref{ex:chsh} (second column; exhaustive search) are sandwiched between the na\"ive power (first column) and the collision value (third column), which equals $(\sqrt3/2)^k=\lamp(G)^k$ by Example~\ref{ex:chsh-lamp}.  The first inequality of Lemma~\ref{lem:rep-basic} is strict for $k\ge2$.}
\label{tab:chsh-sandwich}
\end{table}

\begin{pitfall}{$\lamp(G)$ can equal $1$ for a game that is far from satisfiable}
Take $G$ to be the disjoint union of a perfectly satisfiable game and a game of value near $0$, each drawn with probability $\frac12$.  A function $h$ supported on the satisfiable half has $\|P_Gh\|=\|P_{T_G}h\|$, so $\lamp(G)=1$ although $\val(G)\approx\frac12$.
Theorem~\ref{thm:product} is then vacuous: the parameter $\lamp(G)$ cannot see that half of the game is hopeless, because one global denominator ignores how the game is connected.  The remedy of \citet{DinurSteurer14} for general games is to replace $T_G$ by one trivial game per question (their relaxation $\val_+$);
for the games that matter for Label Cover there is a simpler remedy, the subject of the next section: make the game expand.
\end{pitfall}


%% ===========================================================================
\section{Expanding games: the proof}
\label{sec:expander}
%% ===========================================================================

\subsection{The Markov chain of a game}

In the collision game of Lemma~\ref{lem:collgame} the referee draws two edges $e,e'$ at a common Bob-question.  Let $J_G(u,u')=\Pr[u_e=u,\ u_{e'}=u']$ in that experiment.  It is a symmetric matrix on $U$ whose row sums are $\mu_U(u)$, because each of $u_e,u_{e'}$ alone has law $\mu_U$.
Let $A_G(u,u')=J_G(u,u')/\mu_U(u)$ be the corresponding random walk: from $u$, move to a random Bob-question $w\sim\mu(\cdot\mid u)$ and come back to a random neighbour $u'\sim\mu(\cdot\mid w)$.  Write $\langle x,y\rangle_U=\E_{u\sim\mu_U}x(u)y(u)$ for functions on $U$ and $\mathrm{Var}(x)=\langle x,x\rangle_U-\langle x,\mathbf 1\rangle_U^2$.
Then $\langle x,A_Gy\rangle_U=\sum_{u,u'}J_G(u,u')x(u)y(u')$, so $A_G$ is self-adjoint with respect to $\mu_U$, and
\[
  \langle x,A_Gx\rangle_U=\E_w\Big(\E_{e\sim\mu(\cdot\mid w)}x(u_e)\Big)^2\ \ge\ 0 .
\]
Hence $A_G$ has eigenvalues $1=\lambda_1\ge\lambda_2\ge\dots\ge0$, and the quantity
\[
  \gamma(G)\;=\;1-\lambda_2(A_G)\ \in[0,1]
\]
is the \defn{spectral gap} of the game; $G$ is \defn{$\gamma$-expanding} if $\gamma(G)\ge\gamma$.  Since $\langle x,(I-A_G)x\rangle_U=\E_w\mathrm{Var}_{e\mid w}\,x(u_e)$ (the first term is $\E_ux(u)^2$ and the second is the square of the mean in each neighbourhood), the spectral gap is the best constant in a Poincar\'e inequality:
\begin{equation}\label{eq:poincare}
  \gamma(G)\cdot\mathrm{Var}_{u\sim\mu_U}(x)\ \le\ \E_w\,\mathrm{Var}_{e\sim\mu(\cdot\mid w)}\,x(u_e)\qquad\text{for all }x\colon U\to\R.
\end{equation}
In words: a function that varies a lot over $U$ must vary inside many neighbourhoods; a disconnected game has $\gamma=0$.

\begin{example}\label{ex:gaps}
(a) The CHSH game is \emph{free}: $\mu=\mu_U\otimes\mu_W$.  Then $J_G(u,u')=\mu_U(u)\mu_U(u')$, $A_G$ maps every function to its mean, and $\gamma=1$.  The same holds for every free game (Exercise~\ref{ex:free}).
(b) In the clause--variable game of Figure~\ref{fig:cv} (three clauses, four variables, nine equally likely edges) one finds $J(C_i,C_i)=\frac4{27}$ and $J(C_i,C_j)=\frac5{54}$ for $i\ne j$, so $A_G=\frac16I+\frac5{18}\mathbf 1\mathbf 1^\top$ has second eigenvalue $\frac16$ and $\gamma=\frac56$ (Exercise~\ref{ex:fig-gap}).
(c) For the odd cycle game on $\ell$ vertices, $\gamma=1-\cos^4(\pi/\ell)\approx2\pi^2/\ell^2$ (Exercise~\ref{ex:cycle-gap}); the triangle has $\gamma=\frac{15}{16}$.
\end{example}

\subsection{Expansion makes the product theorem effective}

\begin{theorem}[$\lamp$ is close to the collision value on expanding games; Dinur--Steurer]\label{thm:expander}
Let $G$ be a $\gamma$-expanding projection game and let $\eps>0$ satisfy $\lamp(G)>1-\eps$.  Then $\ \|G\|^2\ge1-\dfrac{54\,\eps}{\gamma}.$
\end{theorem}
\begin{proof}
Pick $h\ge0$ with $\|P_Gh\|>(1-\eps)\|P_{T_G}h\|$, and scale it so that $\|P_{T_G}h\|=1$.  Write $m(u)=\sum_ah(u,a)$; then $\E_um(u)^2=1$ and $\|P_Gh\|>1-\eps$.

\emph{Step 1: we may assume $h$ is concentrated on one label per question.}  Choose independently for every $u$ with $m(u)>0$ a label $a_u$ with probability $h(u,a)/m(u)$, and set $h'(u,a_u)=m(u)$, $h'(u,a)=0$ otherwise.  Then $\E h'=h$ and $m$ is unchanged, so $\|P_{T_G}h'\|=1$; by convexity of the norm, $\E\|P_Gh'\|\ge\|P_G\E h'\|=\|P_Gh\|$, so some outcome $h'$ has $\|P_Gh'\|\ge\|P_Gh\|$.  (This is where non-negativity of $h$ is used.)
Henceforth $h=m\cdot f_\ell$ for a labeling $\ell\colon U\to\Sigma_U$.

\emph{Step 2: the quadratic form.}  For $u,u'$ with $J(u,u')>0$ let $Q(u,u')\in[0,1]$ be the probability, in the experiment of Lemma~\ref{lem:collgame}, that $\pi_e(\ell(u))=\pi_{e'}(\ell(u'))$ given $u_e=u$, $u_{e'}=u'$.  The computation of Lemma~\ref{lem:collgame}, now with the weights $m$, gives
\begin{equation}\label{eq:coll-quad}
  \|P_Gh\|^2=\sum_{u,u'}J(u,u')\,m(u)m(u')\,Q(u,u'),\qquad
  \|G\|^2\ \ge\ \|P_Gf_\ell\|^2=\sum_{u,u'}J(u,u')\,Q(u,u').
\end{equation}
Write $\E_J$ for the average over $(u,u')\sim J$; so $\|G\|^2\ge\E_JQ$ and $\E_J[mm'\,Q]>(1-\eps)^2$, where $mm'$ abbreviates $m(u)m(u')$.

\emph{Step 3: $m$ is nearly constant.}  Since $Q\le1$ and $\E_J[mm']=\langle m,A_Gm\rangle_U$,
\[
  \langle m,(I-A_G)m\rangle_U=1-\E_J[mm']\le1-\E_J[mm'Q]<1-(1-\eps)^2\le2\eps .
\]
By~\eqref{eq:poincare}, and since $\mathrm{Var}_{e\mid w}x(u_e)$ averaged over $w$ is exactly $\langle x,(I-A_G)x\rangle_U$,
$\ \mathrm{Var}(m)\le2\eps/\gamma$.  Put $\bar m=\E_um(u)$; then $\bar m^2=1-\mathrm{Var}(m)$.  If $\eps/\gamma>\frac16$ the claim is trivial (its right side is negative), so assume $\eps/\gamma\le\frac16$; then $\mathrm{Var}(m)\le\frac13$ and $\bar m\ge\frac23$.

\emph{Step 4: the conclusion.}  For $s,t\ge0$ and $c\ge\frac23$,
\begin{equation}\label{eq:pointwise}
  1\ \le\ 9st+9(s-c)^2+9(t-c)^2 .
\end{equation}
Indeed, if $s<\frac13$ then $(s-c)^2\ge(\frac23-\frac13)^2=\frac19$; likewise for $t$; and if both are at least $\frac13$ then $9st\ge1$.
Apply~\eqref{eq:pointwise} with $c=\bar m$, $s=m(u)$, $t=m(u')$, multiply by $1-Q(u,u')\in[0,1]$ and average over $J$:
\[
  1-\E_JQ\ \le\ 9\,\E_J\big[(1-Q)mm'\big]+9\,\E_J\big[(1-Q)(m(u)-\bar m)^2\big]+9\,\E_J\big[(1-Q)(m(u')-\bar m)^2\big].
\]
Both marginals of $J$ equal $\mu_U$ and $0\le1-Q\le1$, so each of the last two terms is at most $9\,\mathrm{Var}(m)\le18\eps/\gamma$.  The first is $9(\E_J[mm']-\E_J[mm'Q])\le9(1-(1-\eps)^2)\le18\eps$, using $\E_J[mm']=\langle m,A_Gm\rangle_U\le\langle m,m\rangle_U=1$.
Hence $1-\|G\|^2\le1-\E_JQ\le18\eps+36\eps/\gamma\le54\eps/\gamma$, as $\gamma\le1$.
\end{proof}

\begin{corollary}[Parallel repetition for expanding projection games]\label{cor:expander}
Let $G$ be a $\gamma$-expanding projection game with $\val(G)\le1-\eta$.  Then $\lamp(G)\le1-\gamma\eta/54$, and for all $k\ge1$
\[
  \val(G^{\otimes k})\ \le\ \|G^{\otimes k}\|\ \le\ \Big(1-\frac{\gamma\eta}{54}\Big)^k .
\]
\end{corollary}
\begin{proof}
Let $\eps>0$ with $\lamp(G)>1-\eps$.  By Theorem~\ref{thm:expander} and Lemma~\ref{lem:val-coll}, $1-54\eps/\gamma\le\|G\|^2\le\val(G)\le1-\eta$, so $\eps\ge\gamma\eta/54$.  As this holds for every $\eps>1-\lamp(G)$, we get $1-\lamp(G)\ge\gamma\eta/54$.  The rest is Corollary~\ref{cor:lamp}.
\end{proof}

\noindent
The rate is linear in $\eta$, not quadratic as in the general bounds of Section~\ref{sec:general}; the price is the factor $\gamma$.  Numerical experiments for these notes (876 random projection games with $\gamma>0.001$ and $\lamp<1$, half of them planted to be nearly satisfiable) never produced a value of $(1-\|G\|^2)\gamma/\eps$ above $3.7$, where $\eps=1-\lamp(G)$; the proved bound is $54$.  The constant is certainly not tight, and we have not tried to optimize it.

\begin{pitfall}{what the expansion hypothesis is doing}
Steps~3--4 are the only places where $\gamma$ enters.  A function $h$ with $\|P_Gh\|\approx\|P_{T_G}h\|$ could, a priori, be a large multiple of a perfect labeling on a tiny part of $U$ and nearly vanish elsewhere; then the weights $m$ would be far from constant and $\|G\|$ could be small even though $\lamp(G)\approx1$.
That is exactly the disconnected game of the pitfall above.  Expansion forbids it: a function with small Dirichlet form $\langle m,(I-A_G)m\rangle_U$ is close to a constant.
\end{pitfall}


%% ===========================================================================
\section{The Label Cover theorem, with explicit parameters}
\label{sec:lc}
%% ===========================================================================

The hypothesis of Corollary~\ref{cor:expander} is easy to arrange: we may mix in a harmless piece of game that connects everything.

\begin{lemma}[Making a game expand]\label{lem:mix}
Let $G$ be a projection game, $0\in\Sigma_W$ a fixed label.  Let $G'$ be the game with one new Bob-question $w_0$ whose edge distribution is
\[
  \mu'=\tfrac12\mu\ +\ \tfrac12\sum_{u\in U}\mu_U(u)\,\delta_{(u,w_0)},
\]
where the edge $(u,w_0)$ carries the constant map $\pi(a)=0$.  Then (i)~$\val(G')=\frac12+\frac12\val(G)$, and (ii)~$\gamma(G')\ge\frac12$.  The question sets $U$, the label sets, and the maps of the old edges are unchanged.
\end{lemma}
\begin{proof}
(i) Any strategy wins the new edges with probability at most $\frac12$, their total weight, and the old edges with probability at most $\frac12\val(G)$.  Both bounds are attained together: Bob answers $0$ at $w_0$ (which wins every new edge, whatever Alice does) and plays an optimal strategy of $G$ elsewhere.
(ii) The marginal of $\mu'$ on $U$ is $\mu_U$ again.  In the collision experiment of $G'$, with probability $\frac12$ the Bob-question is $w_0$ and $u,u'$ are independent copies of $\mu_U$; otherwise it is as in $G$.  Thus $A_{G'}=\frac12A_G+\frac12\mathbf 1\mu_U^\top$.  On functions orthogonal to the constants $A_{G'}$ acts as $\frac12A_G$, so $\lambda_2(A_{G'})=\frac12\lambda_2(A_G)\le\frac12$.
\end{proof}

\begin{remark}[Uniformity]
If $G$ is uniform on $|E|$ edges, then $\mu_U(u)=\deg(u)/|E|$ and $G'$ is the \emph{uniform} game on $2|E|$ edges obtained by adding $\deg(u)$ copies of $(u,w_0)$ for every $u$.  So no weights are needed, and Remark~\ref{rem:multiplicity} applies to $G'^{\otimes k}$.
\end{remark}

\begin{theorem}[Label Cover theorem, with constants]\label{thm:lc-explicit}
Suppose it is $\NP$-hard to distinguish projection games (with label sets of sizes $7$ and $2$) of value $1$ from those of value at most $1-\eta_0$; by Theorem~\ref{thm:pcp2} and Proposition~\ref{prop:cv} this holds with $\eta_0=\eps_0/3$.  Then for every $\delta\in(0,1)$ and
\[
  k\ =\ \Big\lceil\frac{216}{\eta_0}\ln\frac1\delta\Big\rceil,
\]
the map $G\mapsto(G')^{\otimes k}$ is a polynomial-time reduction (for $\delta$ fixed) to distinguishing projection games with label sets of sizes $7^k$ and $2^k$ of value $1$ from value at most~$\delta$.  In particular $7^k,2^k=\poly(1/\delta)$.
\end{theorem}
\begin{proof}
If $\val(G)=1$ then $\val(G')=1$ and $\val((G')^{\otimes k})=1$ by Lemma~\ref{lem:rep-basic}.  If $\val(G)\le1-\eta_0$ then $\val(G')\le1-\eta_0/2$ by Lemma~\ref{lem:mix}(i), and $G'$ is $\frac12$-expanding by (ii).  Corollary~\ref{cor:expander} with $\gamma=\frac12$ and $\eta=\eta_0/2$ gives
$\val((G')^{\otimes k})\le(1-\eta_0/216)^k\le e^{-k\eta_0/216}\le\delta$.  The product has $(2|E|)^k$ edges, polynomial for fixed $\delta$.
\end{proof}

\noindent
This is Theorem~\ref{thm:lc} again, now with a proof that is complete except for the PCP theorem, with $k=O(\log(1/\delta))$ and exponent $\poly(1/\delta)$ of order $1/\eta_0$.
The proof in Lecture~\ref{lec:why} cited Rao's bound, which needs no expansion but has rate $\eta^2$ instead of $\eta$.  Here the number of repetitions is linear in $1/\eta_0$ rather than quadratic, at the price of the expansion step of Lemma~\ref{lem:mix} and of a worse absolute constant; and nothing beyond linear algebra is used.

\begin{pitfall}{the Label Cover theorem does not give Unique Games}
Taking a product of projection maps preserves their fibre sizes: if each $\pi_e$ is $d$-to-$1$ then each $\pi_{e_1}\times\dots\times\pi_{e_k}$ is $d^k$-to-$1$, and a product of bijections is a bijection.
The games produced by Theorem~\ref{thm:lc-explicit} are therefore $d$-to-$1$ with $d$ growing with $k$ (for the clause--variable game, the fibres have $3$ or $4$ elements; the mixed-in edges are constant maps).
Not one of them is a Unique Game, and repetition cannot make them one.  We return to this in Section~\ref{sec:ugc-rep}.
\end{pitfall}


%% ===========================================================================
\section{What is true in general, and what is not}
\label{sec:general}
%% ===========================================================================

Expansion was a convenience.  The theorems below need no such hypothesis; they are all stated for $\val(G)\le1-\eta$ unless noted.

\begin{table}[t]
\centering\small
\renewcommand{\arraystretch}{1.3}
\begin{tabular}{@{}l>{\raggedright\arraybackslash}p{3.6cm}>{\raggedright\arraybackslash}p{4.9cm}@{}}
\toprule
result & games & bound on $\val(G^{\otimes k})$\\
\midrule
\citet{Raz98}; \citet{Holenstein09} & all games, answer sets of size $\le\sigma$ & $(1-\eta^{c})^{\Omega(k/\log\sigma)}$, with $c=32$, resp.\ $c=3$ (Theorem~\ref{thm:raz})\\
\citet{Rao11} & projection games & $(1-\Omega(\eta^2))^k$\\
\citet{DinurSteurer14} & projection games with $\val(G)\le\rho$, any $\rho\in(0,1)$ & $\big(\frac{2\sqrt\rho}{1+\rho}\big)^{k/2}$; this is $\le(1-\eta^2/16)^k$ for $\rho=1-\eta$ and $\le(4\rho)^{k/4}$ for small $\rho$\\
Corollary~\ref{cor:expander} & $\gamma$-expanding projection games & $(1-\gamma\eta/54)^k$\\
\bottomrule
\end{tabular}
\caption{Parallel repetition theorems.  The first three are quoted; the last was proved in this lecture.}
\label{tab:rep-thms}
\end{table}

\noindent
The new feature of the Dinur--Steurer bound is that $\rho$ may be as small as one likes; before it, it was not known that repetition decreases a value that is already sub-constant \citep[Section~1]{DinurSteurer14}.  Their proof is the general version of the argument of this lecture, with the relaxation $\val_+$ in place of $\lamp$.

\paragraph{Strong repetition.}  The bounds of Table~\ref{tab:rep-thms} for projection games decay like $(1-\Omega(\eta^2))^k$.  \defn{Strong parallel repetition} is the hope that $(1-\Omega(\eta))^k$ holds, as it would if the na\"ive guess were roughly right.
It is \emph{false} in general, even for Unique Games on two labels: the odd cycle game of Lecture~\ref{lec:why} has $\eta=\frac1{2\ell}$, yet $G^{\otimes k}$ keeps value $1-O(\sqrt k/\ell)$ up to $k\approx\ell^2$ \citep{Raz11}, and this is essentially tight \citep{FKO07}.
It is true for restricted classes: free projection games \citep{BRRRS09} and projection games on expanders \citep{RR12}.  Corollary~\ref{cor:expander} has the same flavour, with a rate linear in $\eta$ but proportional to the gap $\gamma$; for free games $\gamma=1$ (Exercise~\ref{ex:free}).
The two facts are consistent: the odd cycle has $\gamma\approx2\pi^2/\ell^2$ (Exercise~\ref{ex:cycle-gap}), so Corollary~\ref{cor:expander} only promises decay after $k\gtrsim\ell^3$ repetitions.
As \citet{DinurSteurer14} note, a strong repetition bound for every Unique Game ``would imply a reduction from max cut to unique games''; Raz's counterexample rules out this route.

\begin{bigpicture}{The shape of the information-theoretic proof (a sketch; none of this is proved here)}
The theorems of Raz and Holenstein are proved by contradiction, one coordinate at a time.  Fix a strategy for $G^{\otimes k}$, and for $C\subseteq[k]$ let $W_C$ be the event that all coordinates in $C$ are won.
The central claim is: \emph{if $\Pr[W_C]$ is not too small (say $\ge2^{-\beta k}$), there is a coordinate $i\notin C$ with $\Pr[W_i\mid W_C]\le\val(G)+o(1)$}.  Iterating, $\Pr[W_{C\cup\{i\}}]\le(\val(G)+o(1))\Pr[W_C]$ until the probability falls below $2^{-\beta k}$, which gives exponential decay.
Why the claim holds: conditioning on an event of probability $2^{-\beta k}$ can distort the $k$ pairs of questions by at most $\beta k$ bits of information in total, so for most coordinates $i$ the pair $(x_i,y_i)$ is still almost distributed as in $G$.
The two provers can then \emph{simulate} a strategy for the single game $G$ at coordinate $i$: they embed the real questions at coordinate $i$, and use shared randomness (correlated sampling) to fill in the other coordinates, together with the answers, from the conditional distribution given $W_C$, each looking only at its own question.
The simulated strategy wins $G$ with probability about $\Pr[W_i\mid W_C]$, hence this is at most $\val(G)+o(1)$.  The dependence on the size $\sigma$ of the answer sets enters through the same information accounting.
\end{bigpicture}

\noindent
Compare the two proofs.  Raz and Holenstein follow a \emph{strategy}: they transform a strategy for the product into one for a single copy.  The proof of this lecture never mentions a strategy of the product; it bounds the \emph{norm} of an operator and uses a spectral property of a graph.
The first is more general; the second shorter, and it yields the explicit rate $\gamma\eta$ that makes the role of the graph visible.


%% ===========================================================================
\section{What repetition cannot do for Unique Games}
\label{sec:ugc-rep}
%% ===========================================================================

Repetition turned the PCP theorem into Label Cover.  The Unique Games Conjecture asks for a Label Cover \emph{with bijections}, completeness $1-\zeta$ and soundness $\delta$, for every $\zeta,\delta>0$.  It is natural to try the same recipe, and it is instructive to see exactly where it breaks.

\begin{bigpicture}{Three reasons why repetition does not prove the Unique Games Conjecture}
\textbf{1.\ The type of the constraints is preserved.}  A product of $d$-to-$1$ maps is $d^k$-to-$1$ and a product of bijections is a bijection (Pitfall in Section~\ref{sec:lc}).  The PCP theorem gives projections that are not bijections (fibres of size $3$ and $4$ for $3$-CNF), and repetition cannot repair that.

\smallskip
\textbf{2.\ Repetition never improves completeness.}  By Lemma~\ref{lem:rep-basic}, $\val(G^{\otimes k})\le\val(G)$ for every instance, so whatever completeness the base instances have, repetition cannot raise it.  Perfect completeness is not available as a base, since Unique Games of value $1$ can be recognized in polynomial time (Proposition~\ref{prop:ug-easy}).
The UGC asks for completeness $1-\zeta$ for \emph{every} $\zeta>0$; a repetition-based proof would have to start from a hardness result for Unique Games with the same completeness and a weaker, constant soundness, that is, from a result of the same kind as the conjecture.

\smallskip
\textbf{3.\ Soundness does not decay at the na\"ive rate.}  Even from such a base problem, $(1-\Omega(\eta))^k$ is false for Unique Games \citep{Raz11,FKO07}: the odd cycle games keep value $1-O(\sqrt k\,\eta)$ for $k$ up to about $1/\eta^2$, so in the worst case about $1/\eta^2$ repetitions are needed before the value drops appreciably.
Starting from completeness $1-\zeta$ and value at most $1-\eta$ on the no-instances, $k$ repetitions give completeness at most $1-\zeta$ and, by Rao's bound, soundness at most $e^{-\Omega(k\eta^2)}$.

\smallskip
These are not proofs that the conjecture is hard to prove; they say that the proof, if it exists, must use a different idea.  Lectures~8--12 use the conjecture as a hypothesis, and Lectures~13--15 prove its $2$-to-$2$ relative, which is \emph{not} a statement about Unique Games but comes close.
\end{bigpicture}


%% ===========================================================================
\section*{Exercises}
\addcontentsline{toc}{section}{Exercises}
%% ===========================================================================

\begin{exercise}[\stars{$\star\star$}]\label{ex:rep-prod}
\emph{Products and an experiment.}  (a)~Prove $\val(G)\val(H)\le\val(G\otimes H)\le\min\{\val(G),\val(H)\}$.  (b)~Prove $\|G\otimes H\|\ge\|G\|\,\|H\|$ by exhibiting a product assignment.
(c)~Write a program that computes $\val(G)$, $\|G\|$ and $\lamp(G)$ for small random projection games, and test Theorem~\ref{thm:product}, and compare $\|G\otimes H\|$ with $\|G\|\,\|H\|$.
\emph{Hints for (c):} for a labeling $\ell$ put $c_\ell(w,b,u)=\sum_{e=(u,w)\colon\pi_e(\ell(u))=b}\mu(e)$; then $\|P_Gf_\ell\|^2=\sum_{w,b}\big(\sum_uc_\ell(w,b,u)\big)^2/\mu_W(w)$, and $\lamp(G)^2$ is the maximum over $\ell$ of the top eigenvalue of $D^{-1/2}M_\ell D^{-1/2}$, with $D=\mathrm{diag}(\mu_U)$ and $M_\ell(u,u')=\sum_{w,b}c_\ell(w,b,u)c_\ell(w,b,u')/\mu_W(w)$.
\emph{Warning:} most small random games are satisfiable, hence have $\|G\|=1$, and all the inequalities are trivial for them; filter for $\|G\|<1$ (dense graphs with random projection maps work).  Does $\|G\otimes H\|<\|G\|\,\|H\|$ ever occur in your experiments?  These notes make no claim either way.
\end{exercise}

\begin{exercise}[\stars{$\star$}]\label{ex:fig-gap}
For the clause--variable game of the formula of Figure~\ref{fig:cv} compute the matrices $J$ and $A_G$ and verify $\gamma(G)=\frac56$.  Apply Lemma~\ref{lem:mix} and verify that $\gamma(G')=\frac{11}{12}$ (so the bound $\frac12$ of the lemma is far from tight here).
\end{exercise}

\begin{exercise}[\stars{$\star\star$}]\label{ex:cycle-gap}
Let $G$ be the odd cycle game on $\ell$ vertices.  (a)~Show that $A_G=K^2$ for $K=\frac12I+\frac14(S+S^{-1})$, where $S$ is the cyclic shift, and that $A_G$ has eigenvalues $\cos^4(\pi j/\ell)$, $0\le j<\ell$.  Conclude $\gamma(G)=1-\cos^4(\pi/\ell)$, which is $\frac{15}{16}$ for the triangle.
(b)~With $\eta=\frac1{2\ell}$, Corollary~\ref{cor:expander} gives $\val(G^{\otimes k})\le(1-\gamma\eta/54)^k$.  Show that this is at least $\frac12$ for $k\le\ell^3$ and $\ell$ large, so that it does not contradict Raz's example (which keeps the value above $\frac12$ for $k\le c\ell^2$).  This is a consistency check, not a claim that the corollary is tight.
\end{exercise}

\begin{exercise}[\stars{$\star\star$}]\label{ex:disjoint}
Let $G_1$ be a satisfiable projection game, $G_2$ any projection game with $\val(G_2)<1$ on question sets disjoint from those of $G_1$, and let $G$ draw an edge of $G_1$ or of $G_2$ with probability $\frac12$ each.  Show that $\lamp(G)=1$ and $\gamma(G)=0$, while $\|G\|^2\le\val(G)\le\frac12+\frac12\val(G_2)<1$.
Conclude that no statement of the form of Theorem~\ref{thm:expander} can hold without a hypothesis on $\gamma$.  Which step of its proof fails?
\end{exercise}

\begin{exercise}[\stars{$\star$}]\label{ex:ug-mix}
Let $G$ be a Unique Game.  (a)~Show that the game $G'$ of Lemma~\ref{lem:mix} is not a Unique Game, and explain why the proof of Theorem~\ref{thm:lc-explicit} therefore gives no information about Unique Games.
(b)~Suppose the new edges $(u,w_0)$ carried bijections $\sigma_u$ instead of the constant map.  Show that Lemma~\ref{lem:mix}(i) then fails: find a satisfiable $G$ for which $\val(G')<1$.  \emph{Hint: a path $u_1$---$w$---$u_2$ with identity maps, and $\sigma_{u_1}=\mathrm{id}$, $\sigma_{u_2}=$ the transposition of the two labels.}  Why does this destroy the argument?
\end{exercise}

\begin{exercise}[\stars{$\star\star$}]\label{ex:free}
A game is \emph{free} if $\mu=\mu_U\otimes\mu_W$.  Show that $J_G=\mu_U\otimes\mu_U$ and $\gamma(G)=1$ for every free game.  Deduce $\val(G^{\otimes k})\le(1-\eta/54)^k$ for every free projection game of value at most $1-\eta$.  Check on CHSH that $\val(G^{\otimes k})>\val(G)^k$ for $k=2,3$, and that it nevertheless respects the bound.
\end{exercise}

\begin{exercise}[\stars{$\star\star\star$}]\label{ex:reading-ds}
\emph{Reading.}  Read Section~4 of \citet{DinurSteurer14}.  Identify the place where our proof of Theorem~\ref{thm:expander} uses expansion, and explain how their ``partial assignments'' together with correlated sampling replace it for games with $\gamma=0$.
Then compare the dependence on $\eta$ in Corollary~\ref{cor:expander} (linear) with that of Corollary~1.2 of their paper (quadratic) and say what accounts for the difference.
\end{exercise}


%% ===========================================================================
\section*{Notes and further reading}
\addcontentsline{toc}{section}{Notes and further reading}
%% ===========================================================================

\paragraph{Where this lecture comes from.}
Sections~\ref{sec:operator}--\ref{sec:lc} follow Sections~2 and~3 and Claim~A.1 of \citet{DinurSteurer14} closely: the collision value, Claims~2.3 and~2.4 (our Lemmas~\ref{lem:val-coll} and~\ref{lem:coll-mono}), the parameter $\lamp$ and Theorems~3.2 and~3.4 (our Theorems~\ref{thm:product} and~\ref{thm:expander}), and the expansion trick.
The intuition that one should compare $G\otimes H$ with $T\otimes H$ to ``cancel $H$'' is theirs.  Section~\ref{sec:general} and Table~\ref{tab:rep-thms} summarize the literature as described in the introduction of that paper.

\paragraph{What we changed.}
(i)~\emph{Orientation.}  Dinur and Steurer let the player with the large alphabet be Bob; we let it be Alice, as in Lectures~\ref{lec:why} and~\ref{lec:games}, so the roles of ``$u$'' and ``$v$'' are swapped throughout and the new vertex of the expansion trick sits on Bob's side.
(ii)~\emph{Totality.}  Their projections may be partial; ours are total.  This removes the step ``increase some entries of $P_Gf$ to turn it into an assignment'' in the proof of $\|G\|^2\le\val(G)$, and it is why Feige's game (Example~\ref{ex:feige}) is outside our class.
(iii)~\emph{Constants.}  They write $\|G\|>1-O(\eps/\gamma)$; the proof ends in a bound on $\|G\|^2$, which is the form we state and the one we need, with the explicit constant $54$.
(iv)~\emph{Presentation.}  We avoid $\val_+$ altogether, since the expanding case does not need it.

\paragraph{On attributions.}
Raz's theorem appeared in preliminary form at STOC~1995 and in full as \citet{Raz98}; Holenstein's simplification and Rao's projection-game bound are cited as in Lecture~\ref{lec:why}.
Feige's non-interactive agreement game is from \citet{Feige91}, as presented in Appendix~B of \citet{DinurSteurer14}, where $\val(G)=\val(G^{\otimes2})=\frac12$ is stated; we wrote out the short proofs and confirmed both values by exhaustive search.
The CHSH game is the XOR game of Clauser--Horne--Shimony--Holt; its classical values $\frac34,\frac58,\frac{31}{64}$ for $k=1,2,3$ were computed for these notes by exhaustive search, with $k=3$ checked by two separate programs.
The odd cycle game and its repetition are due to \citet{FKO07} and \citet{Raz11}.

\paragraph{What is not covered.}
The proof of the repetition theorem for non-expanding games (the relaxation $\val_+$, the Cheeger-type inequality, correlated sampling), the information-theoretic proofs of Raz and Holenstein, the tight examples for the dependence on the alphabet size, and the applications to Set Cover and to sub-constant soundness (see \citealp{DinurSteurer14}).
Lecture~6 turns to the other classical way of amplifying a gap, due to \citet{Dinur07}; unlike repetition it keeps the alphabet small.

\chapter{Gap Amplification: Dinur's Combinatorial Proof}
\label{lec:dinur}

{\sffamily\large\color{softgray}Amplify a gap by making the graph talk to itself\par}
\vspace{0.5em}

\noindent
Lectures~\ref{lec:as} and~\ref{lec:alm} proved the PCP theorem with algebra: polynomials, the low-degree test, curves, Hadamard codes.
The proof is short for what it achieves, but every step moves through an encoding, and the soundness constant at the end is the weakest constant of five modules.
This lecture gives the other proof, due to \citet{Dinur07}, in which there is no algebra at all: a constraint graph, a graph power, and a spectral gap.
The theorem it proves is the same, with one visible parameter --- the amplification factor per round --- replacing the five hidden ones.

The idea can be told in one paragraph.  A constraint graph with $\textsc{unsat}$-value $\eps$ is a system of local tests of which an $\eps$~fraction must fail.
Raise the graph to a power $t$: the new ``edges'' are $t$-step paths of the old graph, and each vertex now carries a table of opinions about its neighbourhood.
A path that passes through a failing edge is itself likely to fail, and by expansion a random path passes through a failing edge with probability $\Omega(\sqrt t)\eps$.
So the $\textsc{unsat}$-value grows by $\sqrt t$ at the price of a larger alphabet; and the alphabet can be bought back by composition, which is allowed to be exponentially wasteful there because it is only ever applied to \emph{constant-size} constraints.
Iterating $O(\log n)$ times drives a gap of $1/n$ up to a constant.

\begin{goalbox}{After this lecture you should be able to\ldots}
\begin{enumerate}[leftmargin=*,label=(\arabic*)]
  \item define a constraint graph and $\textsc{unsat}$, and compute both on small examples (Section~\ref{sec:cg});
  \item explain why powering amplifies, and prove the two elementary estimates that make the amplification lemma work on a \emph{regular} graph (Sections~\ref{sec:power}--\ref{sec:amp});
  \item state the amplification lemma and the preprocessing and composition lemmas, and carry out the arithmetic of the main theorem (Sections~\ref{sec:prep}, \ref{sec:amp}, \ref{sec:main});
  \item prove the PCP theorem from the main theorem by iterating it $O(\log n)$ times, and say what the size does at each step (Section~\ref{sec:pcp6});
  \item compare the two proofs of the PCP theorem and say what each one is good for (Section~\ref{sec:compare}).
\end{enumerate}
\end{goalbox}


%% ===========================================================================
\section{Constraint graphs}
\label{sec:cg}
%% ===========================================================================

\begin{definition}[Constraint graph; \citet{Dinur07}]\label{def:cg6}
A \defn{constraint graph} $G=\langle(V,E),\Sigma,\mathcal C\rangle$ consists of a finite undirected graph $(V,E)$ (parallel edges and self-loops allowed), a finite alphabet $\Sigma$, and for every edge $e\in E$ a \defn{constraint} $c(e)\subseteq\Sigma\times\Sigma$.
Every edge is oriented for the purpose of writing the constraint down: $c(e)$ is a set of ordered pairs, and the constraint of the reversed edge is the reversed relation.
An \defn{assignment} is a map $\sigma\colon V\to\Sigma$.  For an edge $e=(u,v)$ written in its oriented form, $\sigma$ \defn{satisfies} $e$ if $(\sigma(u),\sigma(v))\in c(e)$.  We put
\[
  \textsc{unsat}_\sigma(G)=\Pr_{e\in E}\big[\sigma\text{ does not satisfy }e\big],\qquad
  \textsc{unsat}(G)=\min_\sigma\textsc{unsat}_\sigma(G).
\]
The \defn{size} of $G$ is $|V|+|E|\cdot|\Sigma|^2$ (the length of its description), and $\val(G)=1-\textsc{unsat}(G)$ is the value in the sense of Lectures~\ref{lec:why} and~\ref{lec:games}.
\end{definition}

\noindent
A constraint graph of arity $2$ and alphabet $\Sigma$ is a two-prover game with questions on both sides equal to $V$ (Exercise~\ref{ex:cg-game}); what is new here is not the object but the \emph{operations} we shall perform on it.
Note that the edges are counted as a multiset, so a $d$-regular graph has a uniform edge distribution and $\textsc{unsat}_\sigma(G)$ is the fraction of edges violated.

\begin{proposition}\label{prop:np-cg}
It is $\NP$-hard to decide whether an instance of constraint graph satisfiability with $|\Sigma|=3$ has $\textsc{unsat}(G)=0$.
\end{proposition}
\begin{proof}
$3$-colourability: the vertices are the vertices of the given graph, $\Sigma=\{1,2,3\}$, and every edge carries the inequality constraint $c(e)=\{(a,b):a\ne b\}$.
A graph is $3$-colourable iff this system has $\textsc{unsat}=0$.
\end{proof}

\noindent
So $\NP$-hardness of the gap problem $\Gap{0}{\alpha}{\textsc{CSP}}$ for some constant $\alpha>0$ \emph{is} the PCP theorem (Proposition~\ref{prop:pcp-equiv}).  The rest of the lecture constructs it.

\begin{example}[Two small constraint graphs]\label{ex:two-small}
(a)~The triangle with $\Sigma=\{0,1\}$ and inequality on all three edges has $\textsc{unsat}=\frac13$: a two-colouring of a triangle violates exactly one edge, and every assignment violates at least one.

(b)~The same triangle with a self-loop added at each vertex, the self-loops carrying the always-true constraint, has six edges of which only the three original ones can fail, so $\textsc{unsat}=\frac16$: the loops change the normalisation.
This second form is the one the amplification lemma applies to, and it is worth noting that the change from $\frac13$ to $\frac16$ is purely a matter of what we divide by.
\end{example}

\begin{remark}[Why self-loops]\label{rem:loops}
The amplification lemma needs every vertex to have a self-loop carrying a satisfiable constraint.  The reason is visible in the proof: the lemma compares the distribution of the endpoint of a random walk of length $\ell$ with that of a walk of length $t/2$, and with a self-loop present these two distributions can be coupled tightly (Lemma~\ref{lem:binom}).  Without self-loops the lemma is false already at $t=2$: for the triangle of Example~\ref{ex:two-small}(a), give each vertex the label that assigns to every vertex of its ball the parity of that vertex's index.  Every walk of even length alternates, and each edge of the walk joins vertices of different parity, so all edges of the walk are satisfied and the two tables agree at the middle: this labelling of $G^2$ violates nothing, although $\textsc{unsat}(G)=\frac13$ (Exercise~\ref{ex:parity}).
\end{remark}


%% ===========================================================================
\section{Preprocessing: regularity and expansion}
\label{sec:prep}
%% ===========================================================================

The amplification lemma assumes the graph is regular with a good spectral gap.  Both can be assumed without loss, at the price of a constant factor in the $\textsc{unsat}$-value.

\begin{lemma}[Preprocessing; \citet{Dinur07}, Lemma~3.1]\label{lem:prep}
There are constants $0<\lambda<d$ and $\beta_1>0$ such that every constraint graph $G$ can be transformed in polynomial time into a constraint graph $G'=\mathrm{prep}(G)$ with
\[
  \text{(i) } G'\text{ is }d\text{-regular, has a self-loop at every vertex, and }\lambda(G')\le\lambda,
\]
\[
  \text{(ii) } \beta_1\,\textsc{unsat}(G)\le\textsc{unsat}(G')\le\textsc{unsat}(G),
\]
where $\lambda(G')$ is the second largest eigenvalue of the adjacency matrix of~$G'$ (with self-loops counted once).  Also $\mathrm{size}(G')=O(\mathrm{size}(G))$ and $G'$ has the same alphabet as $G$.
\end{lemma}

\begin{proof}[Sketch]
Two steps, both standard.

\emph{Step 1: constant degree (\citet{PapadimitriouYannakakis91}).}  Replace every vertex $v$ of degree $d_v$ by $d_v$ new vertices $[v]$, joined by the edges of a fixed $d_0$-regular expander on $d_v$ vertices with the equality constraint, and put the $d_v$ original edges of $v$ on the $d_v$ new vertices, one each.  The result is $(d_0+1)$-regular; the equality constraints force the copies of $v$ to agree, by a plurality argument, whenever the assignment does not want to pay a constant fraction of edges (Exercise~\ref{ex:prep-degree}).

\emph{Step 2: expansion.}  Given a $d$-regular $G$, add a $d_0$-regular expander on the same vertex set with the always-true constraint, and one self-loop per vertex, again with the always-true constraint.  The result is $(d+d_0+1)$-regular and always satisfied on the new edges, so the $\textsc{unsat}$-value only changes by the change of normalisation: $\textsc{unsat}(G')=O(1)\cdot\textsc{unsat}(G)$ from below.
Its expansion is that of the expander we added, so $\lambda(G')\le d+d_0+1-h_0^2/(d+d_0+1)<\deg(G')$ (Exercise~\ref{ex:prep-expand}).
\end{proof}

\begin{remark}\label{rem:prep-losses}
Each step changes the $\textsc{unsat}$-value by a constant factor:
the first can lose a factor like $h_0/(2d)$ and the second the ratio $\frac{d}{d+d_0+1}$ between the old and the new number of edges.  Both are absorbed into $\beta_1$.
Completeness is untouched: $\textsc{unsat}(G)=0$ implies $\textsc{unsat}(G')=0$ in both steps, since the extension of a satisfying assignment (constants on the copies, and the old values on the old vertices) uses only always-satisfied or equality edges.
\end{remark}


%% ===========================================================================
\section{Powering}
\label{sec:power}
%% ===========================================================================

Here is the operation that creates the gap.  It is easiest to understand it in the special case of a colouring problem first.

\begin{definition}[Power of a constraint graph; \citet{Dinur07}, Section~3.2]\label{def:power}
Let $G=\langle(V,E),\Sigma,\mathcal C\rangle$ be a constraint graph and let $t\ge2$ be even, $h=t/2$.
The \defn{$t$-th power} $G^t=\langle(V,E^t),\Sigma^t,\mathcal C^t\rangle$ has the same vertex set; its alphabet is
\[
  \Sigma^t\;=\;\Big\{\text{functions }f\colon B(v,h)\to\Sigma\Big\}\quad\text{with }B(v,h)=\{x\in V:\ \mathrm{dist}(v,x)\le h\},
\]
the set of \emph{opinions of $v$ about its $h$-neighbourhood} (we write $\Sigma^t$ for the alphabet size $|\Sigma|^{|B(v,h)|}$, which is bounded by $|\Sigma|^{d^{h}}$ in a $d$-regular graph).
For every $t$-walk $p=(x_0,x_1,\dots,x_{2h})$ in $G$ --- \emph{one edge of $G^t$ per walk}, joining $x_0$ and $x_{2h}$ --- the constraint is satisfied by a pair of labels $(f,g)$ iff
\begin{enumerate}[label=(\roman*),leftmargin=*]
  \item $f$ is defined on $x_0,\dots,x_h$ and $g$ on $x_h,\dots,x_{2h}$, and they \defn{agree at the middle}, $f(x_h)=g(x_h)$; and
  \item the merged colouring $a(x_i)=f(x_i)$ for $i\le h$, $a(x_i)=g(x_i)$ for $i\ge h$, satisfies every edge $(x_i,x_{i+1})$ of the walk.
\end{enumerate}
\end{definition}

\noindent
Both conditions are needed.  Condition~(ii) alone would let the two tables disagree at the middle vertex; condition~(i) alone would let them agree there while contradicting the edges of the walk.

\begin{figure}[t]
\centering
\begin{tikzpicture}[x=1cm,y=1cm,every node/.style={font=\small},
   vx/.style={circle,draw,minimum size=0.62cm,fill=accent!8}]
  \node[font=\sffamily\small\bfseries,color=accent,anchor=west] at (-1.15,1.05) {(a)\ \ the walk $x_0x_1x_2$};
  \node[vx] (x0) at (0,0) {$x_0$};
  \node[vx] (x1) at (1.9,0) {$x_1$};
  \node[vx] (x2) at (3.8,0) {$x_2$};
  \draw[thick] (x0)--(x1)--(x2);
  \node[font=\footnotesize,above=2pt of x0] {$f$};
  \node[font=\footnotesize,above=2pt of x1] {agree};
  \node[font=\footnotesize,above=2pt of x2] {$g$};
  \node[font=\footnotesize,align=center,anchor=west,text width=5.6cm] at (5.0,0)
     {$f$ lives on $B(x_0,h)=\{x_0,x_1\}$\\
      $g$ lives on $B(x_2,h)=\{x_1,x_2\}$\\
      both must name $x_h=x_1$ the same way};
  \begin{scope}[yshift=-3.0cm]
  \node[font=\sffamily\small\bfseries,color=accent,anchor=west] at (-1.15,1.05) {(b)\ \ $h=2$: each ball has four vertices};
  \node[draw,rounded corners=2pt,fill=accent!7,align=left,font=\footnotesize] (t1) at (1.5,0)
      {$f\colon B(x_0,2)\to\Sigma$,\quad $f=(0,\cdot,\cdot,\cdot)$\\
       $B(x_0,2)=\{x_0,x_1,x_2,x_3\}$};
  \node[draw,rounded corners=2pt,fill=warn!8,align=left,font=\footnotesize] (t2) at (8.0,0)
      {$g\colon B(x_3,2)\to\Sigma$,\quad $g=(\cdot,\cdot,\cdot,0)$\\
       $B(x_3,2)=\{x_0,x_1,x_2,x_3\}$};
  \node[font=\footnotesize,align=center] at (4.75,-1.45) {the edge of $G^t$ joins $x_0$ and $x_3$;\\ they must agree on $x_2$};
  \end{scope}
\end{tikzpicture}
\caption{Powering.  A label of a vertex is a table of opinions about its whole $h$-neighbourhood, and every $2h$-walk of the old graph becomes a single edge of the new graph, from its first to its last vertex.  The two tables are compared at the middle vertex $x_h$ and each covers one half of the walk; in (a) the edge $(x_0,x_2)$ of $G^t$ is satisfied only if $f$ and $g$ name $x_1$ the same way and the two edges of the walk are satisfied.}
\label{fig:power}
\end{figure}

\begin{remark}[The size of $G^t$]\label{rem:size-power}
The number of edges of $G^t$ is the number of $t$-walks in $G$, which for a $d$-regular $G$ is $|E|\,d^{t-1}$.  With $t$ constant this is a constant factor: \emph{linear} blow-up.
The alphabet, on the other hand, is $|\Sigma|^{d^{h}}$, which is constant but very large (for $d=h=2$ and $|\Sigma|=2$ it is $2^{4}=16$ per vertex; the exponent is exponential in $t$).  Two-thirds of the work below goes into removing it.
\end{remark}

\begin{remark}[Powering versus parallel repetition]\label{rem:power-vs-rep}
Two differences from the repetition of Lecture~\ref{lec:raz} are worth stating at once.  First, the number of \emph{questions} is unchanged: a vertex of $G^t$ is a vertex of $G$, and the number of vertices does not grow, whereas $G^{\otimes k}$ has $|V|^k$ vertices.
Second, the two provers of $G^t$ are questioned about two \emph{vertices}, as before, and the constraints still have arity $2$.  What grows is the size of each answer, and that growth is the amplifier: a label of $v$ must commit to a colour for every vertex of $B(v,h)$, so the two provers are forced to take a position on the whole of a length-$t$ walk.
\end{remark}


%% ===========================================================================
\section{The amplification lemma}
\label{sec:amp}
%% ===========================================================================

\begin{lemma}[Amplification; \citet{Dinur07}, Lemma~3.5 and Section~5]\label{lem:amp}
Let $0<\lambda<d$ and $|\Sigma|$ be constants.  There is a constant $\beta_2=\beta_2(\lambda,d,|\Sigma|)>0$ such that for every even $t\in\N$ and every $d$-regular constraint graph $G$ with a self-loop at every vertex and $\lambda(G)\le\lambda$,
\[
  \textsc{unsat}(G^t)\ \ge\ \beta_2\sqrt t\cdot\min\Big(\textsc{unsat}(G),\ \frac1t\Big).
\]
\end{lemma}

\noindent
So as long as $\textsc{unsat}(G)\le1/t$ --- which is the situation at the beginning, when the gap is $1/n$ --- the $\textsc{unsat}$-value of the power is $\Omega(\sqrt t)$ times larger, while the number of \emph{variables} is unchanged and the size only grows by the constant factor $d^{t-1}$.
The `$\min$' is there because a game cannot have $\textsc{unsat}$-value bigger than $1$, and with $t$ repetitions a single bad edge can be hit at most $1/t$ of the time in the worst case; it is harmless because we apply the lemma only while $\textsc{unsat}\le1/t$.

\subsection{Why amplifying works: an example, and the extraction of a labelling}

Before the proof, the example that shows the operation is not vacuous.

\begin{example}[The triangle]\label{ex:tri-power}
Let $G$ be the triangle with two colours and inequality on the three edges, plus a self-loop with the always-true constraint at each vertex (Example~\ref{ex:two-small}(b)), so that $\textsc{unsat}(G)=\frac16$.
The label set of $G^2$ assigns to each vertex a value at itself and at each of its two neighbours, and the edges of $G^2$ correspond to walks $x_0x_1x_2$, \ie to triples in which $x_0,x_1,x_2$ are the three distinct vertices of the triangle (a closed walk such as $x_0x_1x_0$ also occurs).
An exhaustive search over all $8^3=512$ labellings of $G^2$ gives
\[
  \frac{\textsc{unsat}(G^2)}{1/6}=\frac{8}{27}\approx1.78,\qquad
  \frac{\textsc{unsat}(G^4)}{1/6}\approx3.36,\qquad
  \frac{\textsc{unsat}(G^6)}{1/6}\approx4.28,
\]
against $\sqrt t=1.41,2,2.45$.  The amplification is real and grows like $\sqrt t$; the constant $\beta_2$ hides a numerical factor, as it must, and it is much smaller than the observed one.
Exhaustive search is possible here because the alphabet of $G^t$ has only $2^3=8$ elements per vertex for the triangle: the ball $B(v,1)$ has $3$ vertices and $B(v,3)=V$ has $3$ vertices, and there are $3^3$ labellings.
Without the self-loops the $t=2$ power is \emph{satisfiable} (Remark~\ref{rem:loops}), which is the reason for the hypothesis.
\end{example}

\noindent
We now prove the lemma.  The proof has a clean shape: given a labelling of $G^t$, extract from it a labelling of $G$; compare the two.

\begin{definition}[Popular opinion]\label{def:popular}
Let $\vec\sigma\colon V\to\Sigma^t$ be an assignment for $G^t$, so $\vec\sigma(w)$ is a table of opinions about $B(w,h)$.  Write $\vec\sigma(w)_v$ for its value at $v$.
For $v\in V$ let $X_v$ be the random value obtained by taking a random walk of length $h$ from $v$ (with self-loops allowed, as the graph has them) and reporting the opinion $\vec\sigma(w)_v$ of its endpoint $w$ about $v$.
Define the \defn{popular-opinion assignment} $\sigma\colon V\to\Sigma$ by
\[
  \sigma(v)\;=\;\arg\max_{a\in\Sigma}\ \Pr[X_v=a].
\]
\end{definition}

\noindent
We abbreviate $p_v=\Pr[X_v=\sigma(v)]\ \ge\ 1/|\Sigma|$.

\begin{proof}[Proof of Lemma~\ref{lem:amp}]
Let $\vec\sigma$ be a best assignment for $G^t$, so that $\textsc{unsat}(G^t)=\textsc{unsat}_{\vec\sigma}(G^t)$, and let $\sigma$ be the popular-opinion assignment.

Let $F\subseteq E$ be the set of edges of $G$ violated by $\sigma$, restricted if necessary so that $|F|/|E|=\min(\textsc{unsat}_\sigma(G),1/t)$:  since $\textsc{unsat}_\sigma(G)\ge\textsc{unsat}(G)$ and the min is at most $1/t$, we have
\begin{equation}\label{eq:F}
  \frac{|F|}{|E|}\ \ge\ \min\Big(\textsc{unsat}(G),\frac1t\Big).
\end{equation}

A walk $p=(x_0,\dots,x_{2h})$ is \defn{hit by its $i$-th edge} if $(x_{i-1},x_i)\in F$ \emph{and} the opinions of the two ends about that edge are correct, \ie $\vec\sigma(x_0)_{x_{i-1}}=\sigma(x_{i-1})$ and $\vec\sigma(x_{2h})_{x_i}=\sigma(x_i)$.
Let $N(p)$ be the number of indices $i$ in the middle block
\[
  I=\{i:\ h-\sqrt t<i\le h+\sqrt t\}
\]
by which $p$ is hit.  A hit is a genuine disagreement: the merged labelling of $p$ induced by $\vec\sigma(x_0)$ and $\vec\sigma(x_{2h})$ then violates the edge $(x_{i-1},x_i)$, so the edge $p$ of $G^t$ is \emph{not} satisfied.
Hence
\begin{equation}\label{eq:N-positive}
  \Pr_p[N(p)>0]\ \le\ \textsc{unsat}(G^t).
\end{equation}

The rest of the proof is the estimate
\begin{equation}\label{eq:first-moment}
  \E_p[N(p)]\ \ge\ \Omega(1)\cdot|I|\cdot\frac{|F|}{|E|}\ =\ \Omega(\sqrt t)\cdot\frac{|F|}{|E|},
  \qquad\text{and}\qquad
  \E_p[N(p)^2]\ \le\ O(\sqrt t)\cdot\frac{|F|}{|E|}.
\end{equation}
Given~\eqref{eq:first-moment}, the Cauchy--Schwarz type inequality of \citet{Dinur07} (Exercise~\ref{ex:fact25}) gives
\[
  \Pr[N>0]\ \ge\ \frac{\E[N]^2}{\E[N^2]}\ \ge\ \Omega(\sqrt t)\cdot\frac{|F|}{|E|}.
\]
Combining with~\eqref{eq:N-positive} and~\eqref{eq:F} proves the lemma.

\medskip
\noindent\emph{First moment.}  By linearity, $\E[N]=\sum_{i\in I}\Pr[N_i>0]$ where $N_i$ is the indicator of being hit by the $i$-th edge.  Draw $p$ by choosing first a uniform edge $(u,v)\in E$, then a random walk of length $i-1$ backwards from $u$ to $x_0$, then a random walk of length $2h-i$ forwards from $v$ to $x_{2h}$; since $G$ is regular this is a uniformly random walk.
By the definition of a hit and the independence of the two ends given the edge,
\[
  \Pr[N_i>0]=\frac{|F|}{|E|}\cdot q_u\cdot q_v,\qquad q_u:=\Pr_{x_0}\big[\vec\sigma(x_0)_u=\sigma(u)\big],\ q_v:=\Pr_{x_{2h}}\big[\vec\sigma(x_{2h})_v=\sigma(v)\big].
\]
Now $q_u$ is the probability that a random walk of length $i-1$ from $u$ ends at a vertex whose opinion about $u$ is the popular one, \ie $q_u=\Pr[X_{u,i-1}=\sigma(u)]$ in the notation of Definition~\ref{def:popular}.  The key estimate is
\begin{equation}\label{eq:key}
  \Pr[X_{u,\ell}=\sigma(u)]\ \ge\ \tau\cdot\Pr[X_u=\sigma(u)]=\tau\,p_u\qquad\text{for every }\ell\text{ with }|\ell-h|\le\sqrt t,
\end{equation}
for a constant $\tau>0$ depending only on $d$ and $|\Sigma|$: the two distributions of the endpoint of a walk of length $\ell$ and of length $h$ are close, because one can insert or delete $h-\ell$ steps that use the marked self-loop (this is proved at the end of the section, Lemma~\ref{lem:binom}).
Since $p_u\ge1/|\Sigma|$, we get $q_u\ge\tau/|\Sigma|$ and likewise $q_v$, hence $\Pr[N_i>0]\ge\frac{|F|}{|E|}\cdot\frac{\tau^2}{|\Sigma|^2}$ for every $i\in I$, and
$\E[N]\ge|I|\cdot\frac{|F|}{|E|}\cdot\frac{\tau^2}{|\Sigma|^2}=2\sqrt t\cdot\frac{\tau^2}{|\Sigma|^2}\cdot\frac{|F|}{|E|}$.

\medskip
\noindent\emph{Second moment.}  With $Z(p)=|\{i\in I:\ (x_{i-1},x_i)\in F\}|$ we have $N\le Z$, so $\E[N^2]\le\E[Z^2]$.  Expanding,
\[
  \E[Z^2]=\sum_{i\in I}\E[Z_i]+2\sum_{\substack{i<j\\ i,j\in I}}\E[Z_iZ_j] \le |I|\cdot\frac{|F|}{|E|}+|I|^2\Big(\frac{|F|}{|E|}\Big)^2+2\sum_{\substack{i<j\\ i,j\in I}}\Pr[Z_j>0\mid Z_i>0]\cdot\frac{|F|}{|E|}.
\]
For the conditional probabilities, the graph is an expander: the walk forgets its past, and the probability that it re-enters $F$ after $j-i$ more steps is at most $\frac{|F|}{|E|}+\lambda^{j-i}$ (this is Proposition~\ref{prop:walk-f}, the standard mixing estimate, applied to the normalized adjacency matrix, whose second eigenvalue in absolute value is $\lambda/d$).
Therefore
\[
  \sum_{\substack{i<j\\ i,j\in I}}\E[Z_iZ_j]\le|I|^2\Big(\frac{|F|}{|E|}\Big)^2+|I|\,\frac{|F|}{|E|}\sum_{i\ge1}\Big(\frac{\lambda}{d}\Big)^{i}=O\Big(\sqrt t\cdot\frac{|F|}{|E|}\Big)
\]
because $|I|=2\sqrt t$, $\frac{|F|}{|E|}\le\frac1t$, and $\sum_i(\lambda/d)^i=O(1)$.
\end{proof}

\begin{proposition}[Random walks on an expander re-enter $F$ at the trivial rate]\label{prop:walk-f}
Let $G$ be $d$-regular with second largest eigenvalue $\lambda$, and let $F\subseteq E$ with $|F|/|E|=\mu$.  The probability that a random walk which starts at a uniformly random edge of $F$ is again at an edge of $F$ after $i+1$ further steps is at most $\mu+(\lambda/d)^i$.
\end{proposition}
\begin{proof}
Let $K$ be the distribution of the vertex reached by the first edge of $F$, let $A$ be the normalized adjacency matrix, and let $y_v$ be the probability that a random step from $v$ is inside $F$.  The probability in question is $\langle y,A^ix\rangle$, where $x$ is the density of $K$.  Splitting $x=x^\parallel+x^\perp$ into its mean $x^\parallel=\frac1n\mathbf 1$ and its zero-mean part,
\[
  \langle y,A^ix\rangle=\langle y,A^ix^\parallel\rangle+\langle y,A^ix^\perp\rangle\ \le\ \mu+\|y\|\,\|A^ix^\perp\|\ \le\ \mu+(\lambda/d)^i,
\]
using $\E_vy_v=\mu$, $\|y\|\le1$ and $\|A^ix^\perp\|\le(\lambda/d)^i\|x^\perp\|\le(\lambda/d)^i$.
\end{proof}

\begin{lemma}[Inserting self-loops does not move the endpoint much]\label{lem:binom}
Let $G$ be $d$-regular with a self-loop at every vertex, let $h\ge1$, and let $\vec\sigma$ be an assignment for $G^t$ with $t=2h$.  For every $u$ and every $\ell$ with $|\ell-h|\le\sqrt t$,
\[
  \Pr[X_{u,\ell}=\sigma(u)]\ \ge\ \tau\cdot\Pr[X_{u,h}=\sigma(u)]
\]
for a constant $\tau=\tau(d,|\Sigma|)>0$, where $X_{u,\ell}$ is the opinion about $u$ of the endpoint of a random $\ell$-step walk from $u$.
\end{lemma}
\begin{proof}
Mark one self-loop at every vertex.  A random $\ell$-step walk is described by choosing which $k$ of its steps use the marked loop, and running a walk on the remaining $d-1$ edges otherwise.  Correspondingly,
\[
  \Pr[X_{u,\ell}=a]=\sum_{k=0}^{\ell}\Pr[B_{\ell,p}=k]\cdot\Pr[X'_{u,k}=a],\qquad p=1-\tfrac1d,
\]
where $X'_{u,k}$ is the opinion obtained from a $k$-step walk conditioned to avoid the marked loops.  The binomial distributions $B_{h,p}$ and $B_{\ell,p}$ are comparable on the range $|k-ph|\le c\sqrt h$ for a constant $\tau$ depending only on $d$: that is, $\Pr[B_{h,p}=k]\ge\tau\Pr[B_{\ell,p}=k]$ there, because $|h-\ell|\le\sqrt t$ and both are binomials with the same mean up to $O(\sqrt t)$ and the same variance $\Theta(t)$.
Choose $c$ so that $B_{h,p}$ is in that range with probability $\ge1-\frac1{2|\Sigma|}$; then
\begin{align*}
  \Pr[X_{u,\ell}=a]\ &\ge\ \tau\!\!\sum_{|k-ph|\le c\sqrt h}\!\!\Pr[B_{h,p}=k]\Pr[X'_{u,k}=a]\\
  &\ge\ \tau\Big(\Pr[X_{u,h}=a]-\frac1{2|\Sigma|}\Big)\ \ge\ \frac\tau2\Pr[X_{u,h}=a],
\end{align*}
the last step because $\Pr[X_{u,h}=a]\ge\frac1{|\Sigma|}$ when $a=\sigma(u)$ is the maximizer.  Taking $\tau:=\tau/2$ proves the claim.
\end{proof}

\begin{remark}[Where each hypothesis is used]\label{rem:hypotheses}
Regularity is used twice: to identify the uniform distribution on walks with the sequential description in the first moment, and in Proposition~\ref{prop:walk-f}, where the stationary distribution must be uniform.
A self-loop at every vertex is used only in Lemma~\ref{lem:binom}.  The bound $\lambda(G)<d$ is what makes $\sum_i(\lambda/d)^i$ converge.  The alphabet enters only through $p_v\ge1/|\Sigma|$ and through the constant $\tau$ of Lemma~\ref{lem:binom}.
\end{remark}


%% ===========================================================================
\section{Composition: buying back the alphabet}
\label{sec:comp}
%% ===========================================================================

After powering, the alphabet is $\Sigma^{d^{h}}$: constant, but enormous.  Composition replaces every constraint of $G^t$ by a small constraint graph over a fixed alphabet $\Sigma_0$, without losing more than a constant factor of the $\textsc{unsat}$-value.

The crucial point is the \emph{size} of the objects we must handle.  A constraint of $G^t$ is a relation on $\Sigma^{t}\times\Sigma^{t}$ where the alphabet is constant, so the constraint has constant size, and there are only $2^{2^{O(1)}}$ distinct constraints.  A transformation that is allowed to be exponentially slow on constant-size input is therefore allowed to be \emph{anything}: it may enumerate all assignments, or search all circuits of bounded size, as long as we know one exists.
This is the sense in which the argument is not circular, and it is the reason the algebraic PCP of Lectures~\ref{lec:as}--\ref{lec:alm} reappears: we use it once, on constant-size inputs, to obtain the reduction $P$ whose existence we need.

\begin{definition}[Assignment tester; \citet{Dinur07}, Definition~3.1, after \citet{DinurReingold04} and \citet{BGHSV06}]\label{def:at}
An \defn{assignment tester} with alphabet $\Sigma_0$ and rejection probability $\epsilon>0$ is a map $P$ that takes a Boolean circuit $\Phi$ on variable set $X$ to a constraint graph $G=\langle(V,E),\Sigma_0,\mathcal C\rangle$ with $X\subseteq V$, such that for every $a\colon X\to\{0,1\}$:
\begin{itemize}
  \item (completeness) if $a$ satisfies $\Phi$, there is $b\colon V\setminus X\to\Sigma_0$ with $\textsc{unsat}_{a\cup b}(G)=0$;
  \item (soundness) if $a$ does not satisfy $\Phi$, then for every $b\colon V\setminus X\to\Sigma_0$, $\ \textsc{unsat}_{a\cup b}(G)\ge\epsilon\cdot\mathrm{dist}(a,\mathrm{SAT}(\Phi))$,
        where $\mathrm{dist}(a,\mathrm{SAT}(\Phi))$ is the Hamming distance from $a$ to the nearest satisfying assignment.
\end{itemize}
No bound is imposed on the running time of $P$ or on $\mathrm{size}(G)$.
\end{definition}

\noindent
The difference from a PCP verifier is only the ``proximity'' in the soundness clause: the rejection probability is proportional to the distance of the input from satisfiability, not merely positive.  This is what makes composition work at many overlapping constraints simultaneously; \citet{Dinur07} credits the modular treatment to \citet{DinurReingold04} and \citet{BGHSV06}.

\begin{lemma}[Existence, and how to get one]\label{lem:at-exists}
For every constant $\Sigma_0$ there is an assignment tester with alphabet $\Sigma_0$ and rejection probability $\epsilon>0$ depending only on $\Sigma_0$.  It can be taken to be computable in time exponential in the size of the input circuit.
\end{lemma}
\begin{proof}[Sketch]
Two routes, and it is worth seeing both.
\emph{(a) By the PCP theorem of Lectures~\ref{lec:as}--\ref{lec:alm}.}  Theorem~\ref{thm:alm17} gives a verifier for $\textsc{CKT-SAT}$ that reads $q=O(1)$ bits of a proof; from it, the FGLSS or the inner-verifier construction of Lecture~\ref{lec:alm} gives a reduction from a circuit to constraints over a constant alphabet, and the soundness clause of an inner verifier is exactly the proximity statement of Definition~\ref{def:at} (Exercise~\ref{ex:at-from-pcp}).
\emph{(b) Directly, from the Long Code.}  \citet{Dinur07} gives a self-contained construction in her Section~6: encode the assignment $a\in\{0,1\}^s$ by the table $A_a(f)=f(a)$ over all functions $f\colon\{0,1\}^s\to\{0,1\}$, and run dictatorial tests on it.  With $s$ constant this table has constant size, and the tests can be enumerated.
The construction is exponential in $s$, which is harmless.
\end{proof}

\begin{remark}[Why we may use the algebraic PCP here]\label{rem:no-circularity}
The proof of the PCP theorem in this lecture uses as a subroutine a \emph{fixed} reduction $P$ for inputs of \emph{constant} size.  This is not circular: the theorem we are proving is about inputs of size $n$, and $P$ is never applied to such an input.
Lectures~\ref{lec:as}--\ref{lec:alm} establish the existence of $P$ by constructing it explicitly.  There is no logical dependence on the larger theorem, only on its constant-size instance, which one could also verify by exhaustive search once existence is known.
\end{remark}

\begin{lemma}[Composition; \citet{Dinur07}, Lemma~3.6]\label{lem:comp6}
Let $P$ be an assignment tester with constant rejection probability $\epsilon>0$ and constant alphabet $\Sigma_0$.  There is a constant $\beta_3>0$, depending only on $P$, such that every constraint graph $G=\langle(V,E),\Sigma,\mathcal C\rangle$ with $|\Sigma|=O(1)$ can be transformed in polynomial time into a constraint graph $G\circ P$ with alphabet $\Sigma_0$,
\[
  \mathrm{size}(G\circ P)=O\big(\mathrm{size}(G)\big),\qquad
  \beta_3\cdot\textsc{unsat}(G)\le\textsc{unsat}(G\circ P)\le\textsc{unsat}(G).
\]
\end{lemma}
\begin{proof}[Sketch]
Two steps.

\emph{Robustization.}  A constraint $c(e)\subseteq\Sigma\times\Sigma$ must be presented to $P$ as a circuit on Boolean variables.  Encode each label in $\Sigma$ by a codeword of an error-correcting code $e\colon\Sigma\to\{0,1\}^\ell$ of constant rate and relative distance $\rho>0$; replace each vertex $v$ by $\ell$ Boolean variables $[v]$, and each constraint by the Boolean constraint $\tilde c$ on $[u]\cup[v]$ that is satisfied iff the two blocks are legal codewords of labels that satisfied $c$ (Exercise~\ref{ex:robustization}).
The code is needed so that a far-from-satisfying input to $P$ corresponds to an assignment that is far from satisfying, with a quantitative bound; this is the technical point that \citet{Dinur07} refers to as discussed at length in \citet{DinurReingold04}.

\emph{Composition.}  Run $P$ on each of the (constantly many distinct) constraints $\tilde c$; the output $G_{\tilde c}$ has $[u]\cup[v]$ among its vertices, and we take the union over all $e\in E$.  Because the same variables $[v]$ appear in the outputs for different edges, the union is not disjoint and the invocations are consistent; this is precisely what the assignment-tester property buys.

The soundness inequality follows from the soundness of $P$.  If $\sigma'$ is an assignment for $G\circ P$, extract $\sigma$ by decoding each block $[v]$ to the nearest codeword.  A fraction $\textsc{unsat}_\sigma(G)$ of constraints rejects $\sigma$; for each such edge $e$, the restriction of $\sigma'$ to $[u]\cup[v]$ is at Hamming distance $\ge\rho/2$ from a satisfying input\footnote{Deleting the $\rho$-distance between codewords: if $\sigma$ and $\sigma'$ disagree on the decoded labels, then the two blocks are $\rho$-far from each other's codewords, so at least one of them is $\rho/2$-far from its own.} of $\tilde c$, so by the soundness of $P$ a fraction at least $\epsilon\rho/2$ of the constraints of $G_{\tilde c}$ rejects.
Summing over the rejecting edges gives $\textsc{unsat}(G\circ P)\ge(\epsilon\rho/2)\,\textsc{unsat}(G)$, \ie $\beta_3=\epsilon\rho/2$.
\end{proof}

\noindent
The upper bound $\textsc{unsat}(G\circ P)\le\textsc{unsat}(G)$ needs no work: a satisfying assignment for $G$ extends, edge by edge, to a satisfying assignment of every $G_{\tilde c}$ by the completeness of $P$, and the extension can be chosen consistently on the shared blocks because the values on $[v]$ are determined by $v$'s label.
This is the same argument as in the composition theorem of Lecture~\ref{lec:alm}, and the same one appears in the proof of the repetition theorem of Lecture~\ref{lec:raz}: a local view of a global object can be extended consistently when the object is global.


%% ===========================================================================
\section{The main theorem}
\label{sec:main}
%% ===========================================================================

\begin{theorem}[Main; \citet{Dinur07}, Theorem~4.1]\label{thm:main6}
For every constant alphabet $\Sigma$ there are constants $C>0$ and $0<\alpha<1$ such that every constraint graph $G=\langle(V,E),\Sigma,\mathcal C\rangle$ can be transformed in polynomial time into a constraint graph $G'$ with
\[
  \mathrm{size}(G')\le C\,\mathrm{size}(G),\qquad |\Sigma'|=O(1),
\]
\[
  \textsc{unsat}(G')\ \begin{cases}=\ 0&\text{if }\textsc{unsat}(G)=0,\\[2pt]
  \ge\ \min\big(2\,\textsc{unsat}(G),\ \alpha\big)&\text{otherwise.}\end{cases}
\]
\end{theorem}

\begin{proof}
Fix an assignment tester $P$ (Lemma~\ref{lem:at-exists}) and let $\beta_1,\beta_2,\beta_3$ be the constants of Lemma~\ref{lem:prep}, Lemma~\ref{lem:amp} and Lemma~\ref{lem:comp6}.  Put
\[
  G'=\big(\mathrm{prep}(G)\big)^{t}\circ P,\qquad
  t=\Big\lceil\Big(\frac{2}{\beta_1\beta_2\beta_3}\Big)^{2}\,\Big\rceil,
\]
a constant depending only on the alphabet and on the constants.  The three steps of the construction, with their effects on the $\textsc{unsat}$-value, are
\[
  G\ \xrightarrow[\text{Lemma~\ref{lem:prep}}]{\mathrm{prep}}\ H_1\ \xrightarrow[\text{Lemma~\ref{lem:amp}}]{\text{power }t}\ H_2=H_1^{t}\ \xrightarrow[\text{Lemma~\ref{lem:comp6}}]{\circ P}\ G'.
\]
Writing $\eps=\textsc{unsat}(G)$, we have $\textsc{unsat}(H_1)\ge\beta_1\eps$ and
\begin{align*}
  \textsc{unsat}(H_2)\ &\ge\ \beta_2\sqrt t\cdot\min\big(\beta_1\eps,\ \tfrac1t\big),\\
  \textsc{unsat}(G')\ &\ge\ \beta_3\,\textsc{unsat}(H_2)\ \ge\ \beta_1\beta_2\beta_3\sqrt t\cdot\min\big(\eps,\tfrac1{\beta_1 t}\big).
\end{align*}
Since $\sqrt t\ge 2/(\beta_1\beta_2\beta_3)$, the constant in front is at least $2$, so $\textsc{unsat}(G')\ge\min(2\eps,\ \alpha)$ with
\[
  \alpha=\beta_1\beta_2\beta_3\sqrt t\cdot\min\big(1,\tfrac1{\beta_1 t}\big)>0,
\]
a constant.
(For the last step we used $\min(\beta_1\eps,1/t)\ge\beta_1\min(\eps,1/(\beta_1 t))$: if $\eps\le1/(\beta_1t)$ the left side is $\beta_1\eps$ and the right is $\beta_1\eps$; if $\eps>1/(\beta_1t)$ the left side is $1/t$ and the right is $1/t$.)

\emph{Size.}  Preprocessing and composition cost a constant factor each (Lemmas~\ref{lem:prep} and~\ref{lem:comp6}); the power costs the factor $d^{t-1}$, a constant since $t$ is a constant (Remark~\ref{rem:size-power}).
So $\mathrm{size}(G')=C\,\mathrm{size}(G)$ with $C=d^{t-1}M$ where $M$ is the constant of Lemma~\ref{lem:comp6}.

\emph{Completeness.}  If $\textsc{unsat}(G)=0$ then $\textsc{unsat}(H_1)=0$ by Lemma~\ref{lem:prep}.  A satisfying assignment $\sigma$ of $H_1$ induces the labelling $\vec\sigma$ of $H_1^{t}$ that stores, at every $v$ and every $x\in B(v,h)$, the value $\sigma(x)$: for a walk $p$ the two tables agree at the middle and the merged colouring is $\sigma$ on the walk, so every edge of $p$ is satisfied.  Hence $\textsc{unsat}(H_2)=0$, and then $\textsc{unsat}(G')=0$ by the completeness of $P$.
\end{proof}

\begin{remark}[The three constants]\label{rem:three-constants}
$\beta_1$ is lost in preprocessing, $\beta_2$ is the amplification constant of Lemma~\ref{lem:amp}, $\beta_3$ is lost in composition.  The construction is designed so that a single round multiplies the gap by a factor that is \emph{independent of the instance}: the loss is a constant, the gain is $\sqrt t$, and $t$ is chosen so that the gain outweighs the loss twice over.
This is exactly what the algebra of Lectures~\ref{lec:as}--\ref{lec:alm} does \emph{not} do: there the gaps are preserved, not multiplied, and the constant at the end is that of the weakest module.
\end{remark}

\begin{pitfall}{the power is not applied to the composition, and the order cannot be swapped}
The order is $\mathrm{prep}$, then power, then compose.  Powering first would amplify a gap before regularizing, but the amplification lemma needs regularity and expansion; composing first would destroy the expansion with its gadget constraints and, worse, the alphabet after composition is fixed, so a subsequent powering would raise it again to $|\Sigma_0|^{d^{h}}$ --- the same place we started.
The single round is therefore a fixed pipeline, and it is the \emph{iteration} of the pipeline that produces the theorem.
\end{pitfall}


%% ===========================================================================
\section{The PCP theorem}
\label{sec:pcp6}
%% ===========================================================================

\begin{theorem}[PCP theorem, gap form; \citet{Dinur07}, Corollary~4.2]\label{thm:pcp6}
There is a constant $\alpha>0$ for which it is $\NP$-hard to decide, given a constraint graph $G$ with alphabet $\Sigma_0=O(1)$, whether $\textsc{unsat}(G)=0$ or $\textsc{unsat}(G)\ge\alpha$.
In particular $\Gap{0}{\alpha}{\textsc{CSP}}_{\Sigma_0}$ is $\NP$-hard, and so is $\Gap{1}{1-\alpha}{3SAT}$.
\end{theorem}
\begin{proof}
Start from $3$-colourability (Proposition~\ref{prop:np-cg}): given a graph on $n$ vertices, produce the constraint graph $G_0$ with $|\Sigma|=3$, so that $\textsc{unsat}(G_0)=0$ if the graph is $3$-colourable and $\textsc{unsat}(G_0)\ge1/(3m)$ otherwise, where $m$ is the number of edges.
Let $G_0,G_1,\dots$ be the iterates, $G_{i+1}$ the output of Theorem~\ref{thm:main6} on $G_i$, and let $\eps_i=\textsc{unsat}(G_i)$.  By that theorem,
\begin{equation}\label{eq:iter}
  \eps_{i+1}\ \ge\ \min\big(2\,\eps_i,\ \alpha\big),
\end{equation}
and $\mathrm{size}(G_i)\le C^i\,\mathrm{size}(G_0)$.

Since $\textsc{unsat}(G_0)\ge1/(3m)$ whenever it is positive, the sequence $\min(2^i\eps_0,\alpha/2)$ reaches $\alpha/2$ after at most
$k=\lceil\log_2(3m/\alpha)\rceil=O(\log n)$ steps, and $\eps_i\ge\min(2^i\eps_0,\alpha/2)$ for every $i$ by induction on~\eqref{eq:iter}.
Hence $\eps_k\ge\alpha/2$: after $O(\log n)$ rounds the gap is a constant.
$3$-colourability is $\NP$-hard, so distinguishing $\eps=0$ from $\eps\ge\alpha/2$ is $\NP$-hard.
Finally, a local gadget (Exercise~\ref{ex:gap3sat}) converts each constant-size constraint into a constant number of $3$-clauses while preserving the gap up to a constant, which gives $\Gap{1}{1-\alpha}{3SAT}$ and hence Theorem~\ref{thm:pcp2}.
\end{proof}

\noindent
Compare this with the proofs of Lectures~\ref{lec:as} and~\ref{lec:alm}.  There, the gap was created in one step (the PCP theorem was derived from algebraic tests whose soundness is a constant) and the rest was spent on the parameters: alphabet, queries, randomness.
Here the gap is \emph{grown}, from $1/n$ to a constant in $O(\log n)$ rounds, and no round is delicate; the parameters are handled by preprocessing (regularity) and composition (alphabet), both of which are lossy but only by constants.

\begin{bigpicture}{The arithmetic of one round, and of the iteration}
One round: $\textsc{unsat}\mapsto\min(2\cdot\textsc{unsat},\alpha)$ (Theorem~\ref{thm:main6}), size $\mapsto C\cdot$size.
Iteration: the size grows by $C$ per round, so after $k$ rounds it is $C^{k}\mathrm{size}(G_0)$.
Two regimes.  While $\textsc{unsat}<\alpha/2$ the gap doubles each round, so after $O(\log n)$ rounds it is a constant; the size paid is $C^{O(\log n)}=\poly(n)$, because $C$ is a constant.
Once the gap is a constant the round keeps it there (the minimum is $\alpha$), and it would cost nothing to stop.
The instance $G_0$ has size $m=O(n)$; the final instance has size $\poly(n)$, and its alphabet is the fixed $\Sigma_0$ of the assignment tester.
\end{bigpicture}


%% ===========================================================================
\section{The two proofs of the PCP theorem}
\label{sec:compare}
%% ===========================================================================

We now have both proofs.  It is worth putting them side by side, because they are good for different things, and neither supersedes the other.

\begin{table}[t]
\centering\small
\renewcommand{\arraystretch}{1.25}
\begin{tabular}{@{}l>{\raggedright\arraybackslash}p{5.0cm}>{\raggedright\arraybackslash}p{5.0cm}@{}}
\toprule
& Lectures~\ref{lec:as}--\ref{lec:alm}: algebraic & this lecture: combinatorial\\
\midrule
main object & a verifier reading few symbols of an encoded proof & a constraint graph and its power\\
source of the gap & a constant from local tests & amplification by $\sqrt t$ per round, iterated\\
number of rounds & $2$ compositions & $O(\log n)$ rounds\\
what is amplified & nothing: gaps are preserved, constants compound & the gap itself, by a fixed factor\\
algebra used & low-degree test, sum-check, Hadamard, curves & none (only an error-correcting code in composition)\\
queries at the end & $O(1)$ (parallelization) & $2$ throughout\\
proof length & $\poly(n)$, exponent depending on the constants & $\poly(n)$, with a known exponent $\log_2 C$ per round\\
good for & small constants, linear-length PCPs, hardness with explicit soundness & a self-contained proof, locality of codes, combinatorial intuition\\
\bottomrule
\end{tabular}
\caption{The two proofs of the PCP theorem.}
\label{tab:two-proofs}
\end{table}

\noindent
The two are connected by \emph{composition}, which both use and both need: in the algebraic proof it shrinks the alphabet of the inner verifier; here it shrinks the alphabet of the power.  In both, the inner object is applied only to an input of bounded size, and that is what makes composition affordable.

\begin{bigpicture}{What the iteration really does}
The gap doubles once per round.  Starting from $\textsc{unsat}\ge1/n$ ($n$ the input size) and $O(\log n)$ rounds, the gap becomes a constant.  The size grows by a constant factor at every round, so the total blow-up is $2^{O(\log n)}=\poly(n)$.
In other words: \emph{a constant number of rounds of amplification would only raise the gap from $1/n$ to $O(1/n)$, and it takes $\Theta(\log n)$ rounds to make it constant.}
This is why the theorem is about $O(\log n)$ random bits --- each round costs a constant number of random bits, and the logarithm comes from iterating until the gap stops improving.
\end{bigpicture}


%% ===========================================================================
\section{Amplification versus repetition}
\label{sec:amp-vs-rep}
%% ===========================================================================

Lecture~\ref{lec:raz} amplified a gap by repeating a game; this lecture amplifies it by powering a graph.  The two operations look similar --- both compare a vertex with a vertex at distance $t$ --- but they differ in the three ways that matter, and \citet[Section~9]{Dinur07} states the comparison explicitly.

\begin{table}[t]
\centering\small
\renewcommand{\arraystretch}{1.25}
\begin{tabular}{@{}l>{\raggedright\arraybackslash}p{4.4cm}>{\raggedright\arraybackslash}p{5.0cm}@{}}
\toprule
& parallel repetition (Lecture~\ref{lec:raz}) & powering (this lecture)\\
\midrule
questions & $k$-tuples of questions & the same question set\\
answer size & grows as $\sigma^{k}$ & grows as $|\Sigma|^{d^{h}}$, then bought back by composition\\
number of vertices & $|V|^{k}$ & $|V|$\\
size & exponential in $t$ & linear in $|E|$ (a factor $d^{t-1}$)\\
amplification per round & $\val\mapsto\val^{t}$ (roughly $1-t\eta$ for small $\eta$) & $\textsc{unsat}\mapsto\Theta(\sqrt t)\,\textsc{unsat}$ while $\textsc{unsat}\le1/t$\\
for constant $\textsc{unsat}$ & keeps improving, towards $1$ & stops improving \citep{Bogdanov05}\\
for tiny $\textsc{unsat}$ & $\Theta(t)\cdot\textsc{unsat}$ \citep[Section~9]{Dinur07} & $\Theta(\sqrt t)\cdot\textsc{unsat}$\\
\bottomrule
\end{tabular}
\caption{Two ways to amplify a gap.  Repetition is stronger per round for small gaps but blows up the instance; powering is weaker per round but keeps the size linear, which is what makes the iteration of Section~\ref{sec:pcp6} possible.}
\label{tab:amp-vs-rep}
\end{table}

\noindent
The decisive row is the size.  Repetition turns a game on $|V|$ questions into one on $|V|^k$, so a round of repetition of a $\poly(n)$-size instance has super-polynomial size as soon as $k$ is a constant bigger than $1$ (Lecture~\ref{lec:raz} only ever repeated a \emph{constant-size} game, which is why the Label Cover theorem survives).
Powering keeps the vertex set, so its cost is a constant factor, and a constant factor can be paid $O(\log n)$ times.  This is the whole reason the iterative proof works.

The price is the rate: $\sqrt t$ instead of $t$, and the stopping condition $\textsc{unsat}\le1/t$.  For a constant gap powering makes no progress at all \citep{Bogdanov05}, whereas repetition keeps improving.  The two operations are thus useful in opposite regimes.

\begin{bigpicture}{Why this lecture needs linear blow-up and Lecture~\ref{lec:raz} does not}
Lecture~\ref{lec:raz} used repetition once, on a game of constant size (the clause--variable game of a single constraint, in the composition step), to reach soundness $\delta$ at the cost of an alphabet $\poly(1/\delta)$.  The blow-up mattered there only through the alphabet.
This lecture uses amplification $O(\log n)$ times, on the whole instance, and the point of each round is that the \emph{value} improves while the \emph{size} does not get out of hand.
A $\poly(n)$-size growth per round, iterated $\log n$ times, would give an instance of size $n^{\Theta(\log n)}$ and the argument would collapse; a linear growth per round gives $\poly(n)$ overall, and the argument succeeds.
\end{bigpicture}

\begin{remark}[Radhakrishnan's modification and derandomized repetition]
The $\sqrt t$ in Lemma~\ref{lem:amp} can be improved to $t$ by a modification due to Jaikumar Radhakrishnan \citep[footnote in Section~9]{Dinur07}; the reason the basic proof stops at $\sqrt t$ is that it can only guarantee a hit inside a middle block of length $2\sqrt t$ (the set $I$ of the proof), and the gain is the size of that block.
One may also view the whole construction as a \emph{derandomized} repetition: instead of allowing all $t$-tuples of questions, we allow only the $t$-walks of an expander, and the expander makes the answers of the two ends almost independent.
\citet{FeigeKilian95} showed that no generic derandomization of the repetition theorem can work; as \citet{Dinur07} notes, their parameters do not cover this setting, and the question of how far the derandomization can be pushed is open.
\end{remark}


%% ===========================================================================
\section*{Exercises}
\addcontentsline{toc}{section}{Exercises}
%% ===========================================================================

\begin{exercise}[\stars{$\star$}]\label{ex:cg-game}
Write a constraint graph as a two-prover game, in the sense of Definition~\ref{def:game}: who are Alice and Bob, what are the question sets, what are the label sets, and what is the edge distribution?  Show that the game is a projection game iff every constraint has the form $\{(a,\pi(a))\}$.
Then show that the power $G^t$ of Definition~\ref{def:power} is a game whose two questions are both vertices of $V$ (not $t$-tuples), and identify its relation to the game of the clause--variable type analysed in Lecture~\ref{lec:games}.
\end{exercise}

\begin{exercise}[\stars{$\star$}]\label{ex:two-small-exact}
(a)~Show that the triangle with two colours and inequality has $\textsc{unsat}=\frac13$ and that every optimal assignment violates exactly one edge.
(b)~Show that adding the always-true self-loops makes it $\frac16$, and that this equals the fraction of the \emph{total} edge multiset.
(c)~For the triangle with self-loops, verify by hand that $\textsc{unsat}(G^2)=\frac{8}{27}$: the label set is $\{0,1\}^{3}$ per vertex, and the answer should be the minimum over the $512$ labellings.  (A program is fine, but describe the optimum.)
\end{exercise}

\begin{exercise}[\stars{$\star\star$}]\label{ex:parity}
Let $G$ be the triangle with two colours and inequality and \emph{no} self-loops.  Give a labelling of $G^2$ that satisfies every edge, so that $\textsc{unsat}(G^2)=0$ although $\textsc{unsat}(G)=\frac13$.
\emph{Hint: label each vertex $v$ by the table that assigns to every vertex $x$ of $B(v,1)$ the parity of $x$'s index relative to $v$.}  Where does the proof of Lemma~\ref{lem:amp} break for this graph?  (Check both the regularity hypothesis and the use of Lemma~\ref{lem:binom}.)
\end{exercise}

\begin{exercise}[\stars{$\star\star$}]\label{ex:fact25}
Prove that for every non-negative random variable $X$, $\Pr[X>0]\ge\frac{\E[X]^2}{\E[X^2]}$.  \emph{Hint: condition on $X>0$ and apply $\E[Y^2]\ge\E[Y]^2$ to $Y=[X\mid X>0]$.}
Where is this used in the proof of Lemma~\ref{lem:amp}, and why can it not be replaced by Markov's inequality?
\end{exercise}

\begin{exercise}[\stars{$\star\star$}]\label{ex:prep-degree}
Prove Lemma~\ref{lem:prep} (Step~1).  \emph{Hint: let $\sigma'$ be an assignment of the new graph and let $\sigma(v)$ be the plurality value of $\sigma'$ on the copies $[v]$; let $S$ be the set of copies whose value differs from the plurality.  Show that if $|S|$ is large then many expander edges inside some copy are violated, and otherwise many original edges are.}
Carry out the same count for Step~2 and compute the constants.
\end{exercise}

\begin{exercise}[\stars{$\star$}]\label{ex:prep-expand}
Show that if a $d$-regular graph $G$ is obtained from another graph by adding a $d_0$-regular expander with $h(X)\ge h_0$ and a self-loop at every vertex, then $h(G)\ge h_0$ and $\lambda(G)\le d_0+d+1-h_0^2/(d_0+d+1)$.
\end{exercise}

\begin{exercise}[\stars{$\star\star$}]\label{ex:robustization}
Let $e\colon\Sigma\to\{0,1\}^\ell$ be an encoding with relative distance $\rho$, and let $\tilde c$ be the Boolean constraint on $[u]\cup[v]$ that is satisfied iff both blocks are legal codewords of labels satisfying $c$.  Show that if a pair of blocks is at Hamming distance $<\rho/2$ from a satisfying input, then one of the two blocks is at distance $<\rho/2$ from a codeword; and deduce the $\rho/2$ used in the proof of Lemma~\ref{lem:comp6}.
Why does a na\"ive binary encoding of $\Sigma$ (no error-correcting code) fail?
\end{exercise}

\begin{exercise}[\stars{$\star\star\star$}]\label{ex:at-from-pcp}
Derive an assignment tester (Definition~\ref{def:at}) from the inner verifier system of Theorem~\ref{thm:alm17}: show that the soundness clause of Definition~\ref{def:inner} implies the proximity clause of Definition~\ref{def:at}, with $\epsilon$ a constant.  Where does the distance to $\mathrm{SAT}(\Phi)$ appear on the verifier side, and what has to be checked about the completeness of the encoded assignment?
\end{exercise}

\begin{exercise}[\stars{$\star$}]\label{ex:gap3sat}
Let $G$ be a constraint graph with alphabet $\Sigma_0$, $|\Sigma_0|=O(1)$, and $\textsc{unsat}$-value $\eps$.  Show that $G$ can be converted in polynomial time into a $3$-CNF whose optimal assignment violates at least $\eps/\mathrm{const}$ and at most $1$ of the clauses.  \emph{Hint: a constraint on two variables with alphabet size $k$ is a Boolean function of $2\log k$ bits; write each violated accepting-pattern as a clause and use one auxiliary variable per constraint.}
\end{exercise}

\begin{exercise}[\stars{$\star\star$}]\label{ex:iter-check}
Trace the constants of Theorem~\ref{thm:main6}.  Suppose $C=10^{6}$ and $\alpha=10^{-6}$.  How many rounds are needed for an instance with $m=2^{30}$ edges?  Compare with $\log_2(3m/\alpha)$.
\end{exercise}


%% ===========================================================================
\section*{Notes and further reading}
\addcontentsline{toc}{section}{Notes and further reading}
%% ===========================================================================

\paragraph{Where this lecture comes from.}
Everything is from \citet{Dinur07}: preprocessing is Section~3.1 (Lemmas~3.1--3.3 and Corollary~3.4), powering is Section~3.2 (Definition and Lemma~3.5), composition is Section~3.3 (Definition~3.1 and Lemma~3.6), the main theorem is Theorem~4.1 and the amplification lemma is proved in Section~5 (Lemmas~5.1--5.3 and Proposition~5.4).
The assignment-tester formalism and the composition theorem are due to \citet{DinurReingold04} and \citet{BGHSV06}; the expander-replacement transformation of Lemma~\ref{lem:prep} Step~1 is from \citet{PapadimitriouYannakakis91}.
The comparison of Section~\ref{sec:amp-vs-rep} is Section~9 of the paper, including the example of \citet{Bogdanov05} showing that amplification stalls at constant gaps and the improvement of Radhakrishnan to a rate $t$ instead of $\sqrt t$.

\paragraph{Presentational choices of these notes.}
(i)~\emph{The label of $G^t$ is a ball, not a walk.}  \citet{Dinur07} describes the label informally as ``values for all of its neighbors reachable in $t/2$ steps''; in a $d$-regular graph the number of walks grows like $d^{h}$ but the number of vertices is at most $d^{h}$, and we take the label to be a colouring of the ball $B(v,h)$, which is the smaller and more canonical object.  The constraints only ever read the label at the vertices of the given walk, so the two formulations give the same constraint.
(ii)~\emph{Self-loops are stipulated in the definitions.}  Dinur's Lemma~3.5 assumes them; Remark~\ref{rem:loops} shows that they cannot be dropped, with the parity example of Exercise~\ref{ex:parity}.
(iii)~\emph{The popular-opinion assignment is a separate definition} (Definition~\ref{def:popular}), since it is also the object that Corollary~5.5 of the paper needs for the short PCPs.
(iv)~\emph{Constants.}  We do not compute $\beta_1,\beta_2,\beta_3$; the paper does not either, and Theorem~\ref{thm:main6} only needs their existence.  The point of Theorem~\ref{thm:main6} is the \emph{shape} $2\cdot\textsc{unsat}$ rather than the number $2$.

\paragraph{On the numerical example.}
The values $\textsc{unsat}(G^2),\textsc{unsat}(G^4),\textsc{unsat}(G^6)$ of Example~\ref{ex:tri-power} were computed for these notes by exhaustive search over the $8^{3}$ labellings of the triangle power (the alphabet of $G^t$ has $2^{3}$ elements per vertex for every $t$, because the triangle has three vertices).
The claims about the spectra and about the $K_4$ example in the exercises were checked with the same programs.

\paragraph{What is not covered.}
We do not prove the soundness of the assignment tester outside the two sketches of Lemma~\ref{lem:at-exists}, we do not treat the extensions of the paper to assignment testers (its Section~8) or to short PCPs and locally testable codes (its Section~7), and we do not discuss the limits of amplification at constant gaps beyond the reference to \citet{Bogdanov05}.
The next lecture returns to the algebraic method and proves the optimal inapproximability constants, which the combinatorial proof of this lecture is not designed to give.

\chapter{Optimal Inapproximability: the Long Code and Fourier Analysis}
\label{lec:hastad}

{\sffamily\large\color{softgray}Guessing is optimal, and Fourier analysis proves it\par}
\vspace{0.5em}

\noindent
Lectures~\ref{lec:as}--\ref{lec:alm} built a verifier that reads a constant number of bits, and Lecture~\ref{lec:dinur} built one combinatorially.  Both end with a \emph{constant} gap whose size nobody computes: $\Gap{1}{1-\eps_0}{3SAT}$ is $\NP$-hard for \emph{some} $\eps_0>0$.  That is enough to prove that \sat{Max-3SAT} has no PTAS, and nothing more.  It does not determine \emph{which} constant,
and for the problems of Table~\ref{tab:scoreboard} the question is exactly that: the trivial algorithm satisfies $\tfrac78$ of the clauses of a $3$-CNF and $\tfrac12$ of the equations of a linear system,
and the theorem we want is that nothing does better.

This lecture proves it.  The result is that for \sat{Max-E3-SAT} and \sat{Max-E3-Lin}$(\F_2)$ the random assignment is optimal, up to $\eps$ in the approximation ratio:\footnote{Throughout this lecture $\eps>0$ is an arbitrary constant, fixed in advance and appearing in finitely many places.  It is not the $\eps$ of the UGC, and it is not the $\eps_0$ of the PCP theorem.}
\begin{itemize}
  \item for every $\eps>0$ it is $\NP$-hard to approximate \sat{Max-E3-Lin}$(\F_2)$ within $2-\eps$, \ie to do better than satisfying a $(\tfrac12+\eps)$-fraction when a $(1-\eps)$-fraction can be satisfied;
  \item the same with $2-\eps$ replaced by $\tfrac87-\eps$ for \sat{Max-E3-SAT}.
\end{itemize}
The proof is a PCP whose verifier reads \emph{three bits}, and whose analysis is Fourier analysis on a large Boolean cube.  The new ingredient is the \defn{long code}: instead of writing down an assignment $x\in\{\pm1\}^n$, the proof writes down the truth table of the map $f\mapsto f(x)$.
The table is astronomically redundant, and that redundancy is exactly what lets a verifier check consistency of a \emph{global} object while reading a constant number of bits.

Two features of the lecture are worth flagging in advance, because they are where the difficulty lies.

\emph{First, the Fourier analysis is on the wrong domain.}  A long code is a function of a function.  The characters are indexed by subsets of the domain of the inner function, and the orthogonality that makes everything work is orthogonality over a set of size $2^{2^n}$.
Section~\ref{sec:fourier-f} sets this up and proves the three identities that the rest of the lecture uses; there is nothing difficult in it, but it must be done carefully once.

\emph{Second, the verifier's soundness is proved by \emph{extracting} a strategy.}  If a table passes the test with probability $\tfrac12+\delta$, we do not show that the table is close to a codeword.  We show that its \emph{Fourier spectrum} contains a large coefficient, and we turn that coefficient into an answer for one of the two provers in the underlying game of Lecture~\ref{lec:raz}.
This is the ``two tasks integrated'' of H\aa stad's introduction, and it is the conceptual difference from Lectures~\ref{lec:as}--\ref{lec:alm}.

\begin{goalbox}{After this lecture you should be able to\ldots}
\begin{enumerate}[leftmargin=*,label=(\arabic*)]
  \item write down the long code, and compute its Fourier transform (Section~\ref{sec:longcode});
  \item state and use the orthonormality, inversion and Parseval identities on the domain $F_U$ of functions $\{\pm1\}^U\to\{\pm1\}$ (Section~\ref{sec:fourier-f});
  \item describe the test $\mathrm{L}_\eps(u)$ for linear equations, prove its completeness, and reduce its soundness to a statement about the parallel two-prover protocol (Section~\ref{sec:lin});
  \item carry out the Fourier proof of the key inequality ``acceptance $\tfrac{1+\delta}{2}$ forces a strategy of value $4\eps\delta^2$'', and say where the square comes from (Section~\ref{sec:soundness});
  \item deduce $\tfrac87-\eps$ for \sat{Max-E3-SAT} from $2-\eps$ for \sat{Max-E3-Lin}$(\F_2)$ by a four-clause gadget (Section~\ref{sec:e3sat});
  \item explain what ``non-approximable beyond the random assignment threshold'' means, and why this is where the course turns towards Unique Games (Sections~\ref{sec:optimal}--\ref{sec:towards-ugc}).
\end{enumerate}
\end{goalbox}


%% ===========================================================================
\section{Why the previous two proofs do not give constants}
\label{sec:why-constants}
%% ===========================================================================

It is worth being precise about what is missing.  The PCP theorem gives, for each $\NP$ language, a verifier that reads $q=O(1)$ bits and has soundness $\tfrac12$.  The reduction of Lecture~\ref{lec:why} turns it into a $3$-CNF by writing each accept predicate as at most $K=(q-2)2^q$ clauses, and an unsatisfiable input then forces a violated fraction of at least $1/(2K)$.
So $\eps_0$ is of order $2^{-q}/q$, and $q$ is whatever constant the construction happened to produce -- in Lecture~\ref{lec:alm} the Hadamard verifier alone reads $6k+12$ bits per run.

The obstacles to a sharp constant are different in the two proofs.
In the algebraic proof (Lectures~\ref{lec:as}--\ref{lec:alm}) the final constant comes out of a chain of modules -- the low-degree test, the Hadamard verifier with its soundness $\tfrac{19}{20}$, composition, repetition -- and each link loses a factor.  Any one of them can be improved, but the losses compound, and the resulting $\eps_0$ is not a number anybody tries to compute.
In the combinatorial proof (Lecture~\ref{lec:dinur}) the same is true in a cleaner form: the round multiplies the gap by $\sqrt t$, so the final constant is whatever $\min(2\,\textsc{unsat},\alpha)$ says at the end of the iteration, and $\alpha$ is a product of three unspecified constants $\beta_1\beta_2\beta_3\sqrt t$.

What is needed is a verifier that is \emph{designed together with the optimisation problem it is supposed to encode}, so that the acceptance predicate is exactly the predicate of the target problem and the soundness analysis does not lose a factor at each step.
H\aa stad's tests have this property: the acceptance criterion of the test for linear equations is the exclusive-or of three proof bits, and for satisfiability it is a clause on three proof bits.  Nothing is lost in the translation, and the soundness is analysed exactly, by Fourier analysis of the tables the prover writes down.

\begin{bigpicture}{The dictionary of this lecture}
A long code is an encoding of an assignment $x$ to a set $U$ of variables: the table $A_x\colon F_U\to\{\pm1\}$, $A_x(f)=f(x)$, where $F_U$ is the set of all maps $\{\pm1\}^U\to\{\pm1\}$.
The table has length $2^{2^{|U|}}$, so a constant-size $U$ gives a constant-size table.
A \emph{test} queries $A$ at a few inputs $f_1,f_2,\dots$ and accepts according to a predicate on the answers.
If $A$ is a correct long code, the test must accept with probability $1$; if the test accepts with probability $\tfrac{1+\delta}{2}$ for a table that is \emph{not} a long code, the Fourier coefficients of $A$ must be non-trivial, and those coefficients are the strategies for a two-prover game.
Good tests are those whose acceptance predicate is the target predicate -- an XOR of three bits for linear equations, a clause for satisfiability -- so that the PCP \emph{is} the CSP.
\end{bigpicture}


%% ===========================================================================
\section{Fourier analysis on $F_U$}
\label{sec:fourier-f}
%% ===========================================================================

Fix a finite set $U$ of variables, $|U|=u$, and let
\[
  F_U\;=\;\big\{\,f\colon \{\pm1\}^U\to\{\pm1\}\,\big\},\qquad |F_U|=2^{2^{u}} .
\]
We equip $F_U$ with the uniform measure, and for $A,B\colon F_U\to\R$ write (the second expression being the same thing written out)
\[
  \langle A,B\rangle\;=\;\E_{f\sim F_U}\big[A(f)B(f)\big]\;=\;2^{-2^{u}}\sum_{f\in F_U}A(f)B(f).
\]
The characters are indexed by \emph{subsets} of the domain of $f$, that is by $\alpha\subseteq\{\pm1\}^U$:
\begin{equation}\label{eq:char-f}
  \chi_\alpha(f)\;=\;\prod_{x\in\alpha}f(x),\qquad \chi_{\varnothing}\equiv1 .
\end{equation}
This is the only place where the double exponent shows up, and it is exactly parallel to the familiar case: a function on $\{\pm1\}^n$ has its characters indexed by subsets of $[n]$, and a function on $F_U$ has them indexed by subsets of $\{\pm1\}^U$.  If the reader keeps in mind that ``$[n]$'' has been replaced by ``$\{\pm1\}^U$'' nothing below is surprising.

\begin{lemma}[Orthonormality]\label{lem:orth-f}
$\E_f[\chi_\alpha(f)\chi_\beta(f)]=\mathbf 1[\alpha=\beta]$ for all $\alpha,\beta\subseteq\{\pm1\}^U$.
\end{lemma}
\begin{proof}
$\chi_\alpha\chi_\beta=\chi_{\alpha\triangle\beta}$, and a nonempty character is balanced: if $x_0\in\gamma\ne\varnothing$, the map that flips the single value $f(x_0)\mapsto-f(x_0)$ is a measure-preserving involution of $F_U$ that changes the sign of $\chi_\gamma(f)$, so $\E_f[\chi_\gamma(f)]=0$.
\end{proof}

\begin{definition}[Fourier transform]\label{def:fourier-f}
$\ \widehat A_\alpha=\langle A,\chi_\alpha\rangle=\E_f\big[A(f)\chi_\alpha(f)\big]$.
\end{definition}

\begin{theorem}[Inversion and Parseval]\label{thm:inv-f}
For every $A\colon F_U\to\R$: $\ A=\sum_\alpha\widehat A_\alpha\chi_\alpha$, and $\ \sum_\alpha\widehat A_\alpha^2=\E_f[A(f)^2]$.
In particular, if $A$ takes values in $\{\pm1\}$ then $\sum_\alpha\widehat A_\alpha^2=1$ and $\sum_\alpha\widehat A_\alpha^3\le\max_\alpha\widehat A_\alpha$.
\end{theorem}
\begin{proof}
By Lemma~\ref{lem:orth-f} the $2^{2^u}$ functions $\chi_\alpha$ are orthonormal in the $2^{2^u}$-dimensional space of real functions on $F_U$, hence an orthonormal basis; the coefficient of $\chi_\alpha$ in $A$ is $\langle A,\chi_\alpha\rangle$, which is the inversion formula.
Parseval is the same expansion applied to $\langle A,A\rangle$.
The last claim: $\sum_\alpha\widehat A_\alpha^3\le\max_\alpha\widehat A_\alpha\cdot\sum_\alpha\widehat A_\alpha^2=\max_\alpha\widehat A_\alpha$.
\end{proof}

\noindent
Two more identities, which are what the extraction argument uses.

\begin{lemma}[Products and restrictions]\label{lem:prod-f}
(a)~$\chi_\alpha(fg)=\chi_\alpha(f)\chi_\alpha(g)$, \ie the characters are multiplicative in the pointwise product $fg(x)=f(x)g(x)$.
(b)~$\chi_\alpha(f)\chi_\beta(f)=\chi_{\alpha\triangle\beta}(f)$.
(c)~\emph{Restriction.}  Let $U\subseteq W$, let $f\in F_U$, and extend $f$ to $\{\pm1\}^W$ by $f(y)=f(y|_U)$.  Then for every $\beta\subseteq\{\pm1\}^W$,
\[
  \chi_\beta(f)\;=\;\chi_{\pi_U(\beta)}(f),\qquad
  \pi_U(\beta)=\big\{x\in\{\pm1\}^U:\ x=y|_U\text{ for an \emph{odd} number of }y\in\beta\big\}.
\]
In particular, if every $x$ occurs an even number of times among the restrictions $y|_U$, $y\in\beta$, then $\chi_\beta\equiv1$.
\end{lemma}
\begin{proof}
(a) and (b) are immediate from~\eqref{eq:char-f}.  For (c), $\chi_\beta(f)=\prod_{y\in\beta}f(y|_U)=\prod_{x\in\{\pm1\}^U}f(x)^{m_x}$ where $m_x$ is the number of $y\in\beta$ with $y|_U=x$; since $f(x)^2=1$, only the parity of $m_x$ matters, which is exactly $\pi_U$.
\end{proof}

\begin{lemma}[Independence kills mixed terms]\label{lem:indep-f}
Suppose $f_1,\dots,f_k$ are random functions with ranges in possibly different $F_{U_i}$ and that some character $\chi_\gamma(f_i)$ depends on a coordinate $x_0$ that is uniform and independent of everything else in the product.  Then
$\E\big[\prod_i\chi_{\alpha_i}(f_i)\big]=0$.
In particular, if $f$ is uniform on $F_U$ and $\alpha\ne\varnothing$ then $\E_f[\chi_\alpha(f)]=0$.
\end{lemma}
\begin{proof}
Factor out $f_i(x_0)$, which is a uniform sign independent of the remaining factors; the expectation of a product containing an independent uniform sign is $0$.
\end{proof}

\begin{lemma}[The noise operator]\label{lem:noise}
Let $\mu$ be chosen by setting $\mu(y)=-1$ with probability $\eps$ and $\mu(y)=+1$ with probability $1-\eps$, independently over $y$.  Then for every $\beta\subseteq\{\pm1\}^W$,
\[
  \E_\mu\big[\chi_\beta(\mu)\big]\;=\;(1-2\eps)^{|\beta|}.
\]
\end{lemma}
\begin{proof}
$\chi_\beta(\mu)=\prod_{y\in\beta}\mu(y)$ is a product of independent random signs, each with mean $1-2\eps$.
\end{proof}

\begin{lemma}[$x^s\le e^{s(x-1)}$; H\aa stad's Lemma~5.3]\label{lem:exp}
For all $x>0$ and $s>0$: $x^s\le e^{s(x-1)}$.
\end{lemma}
\begin{proof}
It suffices to prove $x\le e^{x-1}$, since raising to the power $s$ gives $x^s\le e^{s(x-1)}$.  For $x\le1$ the left side is at most $1$ and the right side is at least $1$.  For $x>1$, the function $x-e^{x-1}$ has derivative $1-e^{x-1}<0$ and value $0$ at $x=1$, hence is $\le0$.
\end{proof}

\begin{pitfall}{do not confuse the two cubes}
The Boolean cube of an ordinary Fourier argument is $\{\pm1\}^n$, of size $2^n$; here the cube is $F_U=\{\pm1\}^{\{\pm1\}^U}$, of size $2^{2^u}$.  Everything is \emph{doubly} exponential, which is why the long code is only affordable for constant $u$; a table for $|U|=u$ has $2^{2^u}$ bits, and this is what makes the proof's $u$ a constant.
The character $\chi_\alpha$ is indexed by a subset of the \emph{inner} cube, and the size $|\alpha|$ that appears in Lemma~\ref{lem:noise} is a cardinality of a set of points, not of coordinates.
\end{pitfall}


%% ===========================================================================
\section{The long code}
\label{sec:longcode}
%% ===========================================================================

\begin{definition}[Long code; \citet{BGS98}, \citet{Hastad01}]\label{def:longcode}
The \defn{long code} of an assignment $x\colon U\to\{\pm1\}$ is the function
\[
  A_x\colon F_U\to\{\pm1\},\qquad A_x(f)=f(x).
\]
It is the truth table, of length $2^{2^{|U|}}$, of ``evaluation at $x$''.
\end{definition}

\noindent
The code is injective, and any two of its codewords differ on exactly half of $F_U$: if $x\ne x'$ then a uniform $f$ has $f(x)\ne f(x')$ with probability $\tfrac12$.  So the long code has relative distance $\tfrac12$, the largest possible, at the price of an \emph{astronomical} rate: $|U|$ bits are encoded by $2^{2^{|U|}}$ bits.
Only constant $|U|$ is affordable, and that is all we need, since the test is applied to sets of constant size.  The reason the code is used is the following computation.

\begin{lemma}[Long codes are exactly the monomials]\label{lem:longcode-fourier}
Let $x_0\in\{\pm1\}^U$ and $A=A_{x_0}$.  Then $\widehat A_{\{x_0\}}=1$ and $\widehat A_\alpha=0$ for every $\alpha\ne\{x_0\}$.  Conversely, if $A\colon F_U\to\{\pm1\}$ satisfies $\widehat A_{\{x_0\}}=1$ for some $x_0$, then $A=A_{x_0}$.
\end{lemma}
\begin{proof}
$A_{x_0}(f)=f(x_0)=\chi_{\{x_0\}}(f)$ by the definition of a character, and the characters are orthonormal (Lemma~\ref{lem:orth-f}), so $\widehat A_{\{x_0\}}=\langle\chi_{\{x_0\}},\chi_{\{x_0\}}\rangle=1$ and the other coefficients vanish.
Conversely, if $\widehat A_{\{x_0\}}=1$ then Parseval $\sum_\alpha\widehat A_\alpha^2=1$ forces $\widehat A_\alpha=0$ for $\alpha\ne\{x_0\}$, and the inversion formula gives $A=\chi_{\{x_0\}}=A_{x_0}$.
\end{proof}

\noindent
A \emph{single} Fourier coefficient of size $1$ pins the whole table down.  The tests below are designed so that a table passing them with probability noticeably above $\tfrac12$ must have a coefficient of size $\Omega(1)$ -- not necessarily $1$ -- and a large coefficient at a \emph{small} set $\alpha$ is then, for the purposes of the underlying game, as good as a long code: it yields an assignment of the variables in $\alpha$.

\begin{definition}[Folding and conditioning]\label{def:folding}
(a)~$A\colon F_U\to\{\pm1\}$ is \defn{folded over true} if $A(\bar f)=-A(f)$ for all $f$, where $\bar f=-f$ pointwise.  Any $A$ is turned into a folded table $\tilde A$ by choosing one function from each pair $\{f,\bar f\}$, say $f$, and setting $\tilde A(f)=A(f)$, $\tilde A(\bar f)=-A(f)$.  Then $\widehat{\tilde A}_\alpha=0$ for every even $|\alpha|$; in particular $\widehat{\tilde A}_\varnothing=0$.
(b)~Given $h\in F_U$ (read as a Boolean condition on assignments), the \defn{conditioning} of $A$ upon $h$ is $A_h(f)=A(f\wedge h)$, where $f\wedge h$ is the pointwise ``and''.  Since $f\wedge h$ depends only on the values of $f$ at points where $h$ is true, $\widehat{(A_h)}_\alpha=0$ whenever $\alpha$ contains a point at which $h$ is false.
\end{definition}

\noindent
Folding enforces a symmetry that every long code has anyway, $A_x(\bar f)=-f(x)=-A_x(f)$; it costs nothing on a correct proof and it removes all even coefficients, the empty set included.  Conditioning is used because the verifier knows something about $x$ -- namely that it satisfies the chosen clauses -- and it removes every coefficient $\alpha$ containing an assignment that violates them.  On a long code of an $x$ with $h(x)$ true, $A_h=A$.  Both are operations the verifier performs when it reads the table (it queries $A$ at $f$ or $\bar f$, at $f\wedge h$), and both facts are proved in \cite[Section~2.5]{Hastad01} by the pairing argument of Lemma~\ref{lem:indep-f} (Exercise~\ref{ex:folding}).


%% ===========================================================================
\section{The two-prover protocol and the standard written proof}
\label{sec:swp}
%% ===========================================================================

Everything rests on a game that is deliberately weak.

\begin{definition}[Basic protocol; \citet{Hastad01}, Section~3]\label{def:basic-protocol}
The input is a $3$-CNF $\varphi=C_1\wedge\dots\wedge C_m$ in which each variable occurs exactly five times.  The referee picks $j\in[m]$ and one of the three variables $x_k$ of $C_j$ uniformly, sends $j$ to $P_1$ and $k$ to $P_2$, receives from $P_1$ a satisfying assignment of $C_j$ and from $P_2$ a bit, and accepts iff the two agree on $x_k$.
\end{definition}

\begin{lemma}\label{lem:basic-sound}
If $\varphi$ is $c$-satisfiable then no strategy makes the referee accept with probability more than $\tfrac{2+c}{3}$.
\end{lemma}
\begin{proof}
$P_2$'s answers define an assignment $\sigma$ to all variables, and $\sigma$ satisfies at most a $c$-fraction of the clauses.  Fix a clause $C_j$ not satisfied by $\sigma$.  Either $P_1$ answers with an assignment not satisfying $C_j$, and the referee rejects outright, or $P_1$'s assignment satisfies $C_j$ and therefore disagrees with $\sigma$ on at least one of its three variables; the referee queries that variable with probability at least $\tfrac13$.
So for each such $j$ the referee rejects with probability at least $\tfrac13$, giving a rejection probability of at least $(1-c)/3$ overall.
\end{proof}

\noindent
Taking $c$ from the PCP theorem gives soundness $\tfrac{2+c}{3}<1$, a constant.  Parallel repetition drives it down.

\begin{lemma}[Parallel repetition of the basic protocol; \citet{Hastad01}, Lemma~3.2]\label{lem:parallel-protocol}
If $\varphi$ is $c$-satisfiable with $c<1$ then there is a constant $\bar c<1$ such that for every $u$ the $u$-parallel repetition of the protocol of Definition~\ref{def:basic-protocol} has value at most $\bar c^{\,u}$; if $\varphi$ is satisfiable its value is $1$.
\end{lemma}

\noindent
This is exactly the repetition of Lecture~\ref{lec:raz}: the referee draws $u$ independent clause--variable pairs $(j_i,k_i)$, sends $(j_i)_{i\le u}$ to $P_1$ and $(k_i)_{i\le u}$ to $P_2$, and demands that all $u$ coordinates be consistent and satisfying.
By Theorem~\ref{thm:raz} the value is $(1-\Omega(1))^{\Omega(u)}$, which is $\bar c^{\,u}$ for some $\bar c<1$; the answer sets of the basic protocol have constant size, so the logarithmic dependence on the alphabet is a constant.
The strategy of $P_2$ is a function of the set $U=\{x_{k_1},\dots,x_{k_u}\}$, of size at most $u$, and that of $P_1$ is a function of $W$, the set of all variables of the chosen clauses, of size at most $3u$.

Writing down the answers to \emph{all} questions means writing, for every set $V$ of at most $3u$ variables, a table over the assignments of $V$.  The long code turns this into a table over functions.

\begin{definition}[Standard written proof; \citet{Hastad01}, Definition~3.3]\label{def:swp}
A \defn{standard written proof with parameter $u$}, or $\mathrm{SWP}(u)$, consists, for every $V\subseteq[n]$ with $|V|\le3u$, of a string of length $2^{2^{|V|}}$, read as a function $A_V\colon F_V\to\{\pm1\}$.
It is \defn{correct} for $\varphi$ if there is an assignment $x$ satisfying $\varphi$ with $A_V=A_{x|_V}$ for every such $V$.
\end{definition}

\noindent
Its length is $n^{O(u)}2^{2^{3u}}$, polynomial for constant $u$.  Note that one table $A_V$ serves every clause that uses only variables of $V$; this is the sense in which the encoding is global.

\begin{bigpicture}{How the written proof is used}
A verifier chooses $U\subseteq W$, a set $h$ of clauses on $W$, and queries $A_U$ at one function $f\in F_U$ and $A_W$ at one or two functions $g_1,g_2\in F_W$, conditioned upon $h$ so that only assignments satisfying the clauses enter the Fourier expansion.
The whole design problem is to choose the joint distribution of $(f,g_1,g_2)$ so that (i)~a correct proof is accepted with the intended probability, and (ii)~the acceptance predicate is a fixed Boolean function of the three answers, so that the test \emph{is} a constraint of the target CSP.
Property~(ii) is what makes the final constant exact: there is no intermediate reduction in which a constant can be lost, and the \emph{random assignment threshold} of the target problem appears directly as the soundness.
\end{bigpicture}


%% ===========================================================================
\section{Testing a long code}
\label{sec:lctest}
%% ===========================================================================

The simplest test of this kind checks that a table over $F_U$ really is a long code.  It has no bearing on any optimisation problem; we present it because it isolates the argument.

\begin{definition}[Long-code test]\label{def:lctest}
Given $A\colon F_U\to\{\pm1\}$: choose $f_0,f_1\in F_U$ uniformly and independently, set $f_2=f_0f_1$ pointwise, and accept iff $A(f_0)A(f_1)A(f_2)=1$.
\end{definition}

\begin{lemma}[Completeness]\label{lem:lctest-complete}
If $A=A_{x_0}$ is a long code then the test accepts with probability $1$.
\end{lemma}
\begin{proof}
$A(f_0)A(f_1)A(f_2)=f_0(x_0)f_1(x_0)f_0(x_0)f_1(x_0)=1$.
\end{proof}

\begin{theorem}[Soundness of the long-code test; \citet{Hastad01}, Section~4]\label{thm:lctest}
Let $A\colon F_U\to\{\pm1\}$ be accepted with probability $\tfrac{1+\delta}{2}$.  Then $\sum_\alpha\widehat A_\alpha^3=\delta$; in particular $\max_\alpha\widehat A_\alpha\ge\delta$.
\end{theorem}
\begin{proof}
The test accepts exactly when the product is $+1$, so
\[
  \E_{f_0,f_1}\big[A(f_0)A(f_1)A(f_0f_1)\big]\;=\;2\Pr[\text{accept}]-1\;=\;\delta .
\]
Expand each factor by Theorem~\ref{thm:inv-f} and use Lemma~\ref{lem:prod-f}(a),(b):
\[
  \chi_{\alpha_0}(f_0)\chi_{\alpha_1}(f_1)\chi_{\alpha_2}(f_0f_1)=\chi_{\alpha_0\triangle\alpha_2}(f_0)\cdot\chi_{\alpha_1\triangle\alpha_2}(f_1).
\]
Substituting and using the independence of $f_0$ and $f_1$,
\[
  \E_{f_0,f_1}\big[A(f_0)A(f_1)A(f_0f_1)\big]
  =\sum_{\alpha_0,\alpha_1,\alpha_2}\widehat A_{\alpha_0}\widehat A_{\alpha_1}\widehat A_{\alpha_2}\ \E_{f_0}\big[\chi_{\alpha_0\triangle\alpha_2}(f_0)\big]\ \E_{f_1}\big[\chi_{\alpha_1\triangle\alpha_2}(f_1)\big].
\]
By Lemma~\ref{lem:indep-f} the two expectations are $\mathbf 1[\alpha_0=\alpha_2]$ and $\mathbf 1[\alpha_1=\alpha_2]$, so the sum collapses to $\sum_\alpha\widehat A_\alpha^3=\delta$.
The final claim is $\sum_\alpha\widehat A_\alpha^3\le\max_\alpha\widehat A_\alpha$ from Theorem~\ref{thm:inv-f}.
\end{proof}

\noindent
The identity $\E[A(f_0)A(f_1)A(f_0f_1)]=\sum_\alpha\widehat A_\alpha^3$ was checked by exhaustive computation for $|U|\le2$: it holds exactly for random tables and gives $1$ for a long code.
The moral is the sentence the whole method rests on: a table that passes the test with probability $\tfrac{1+\delta}{2}$ has a Fourier coefficient of size $\Omega(\delta)$, and that coefficient names a small set of variables.

Definition~\ref{def:lctest} is the prototype.  The tests that follow replace the product $f_0f_1$ by a \emph{noisy} version, which turns the acceptance predicate into the predicate of the target problem and simultaneously weakens the test, so that passing it does not require being a long code but only being close to a monomial.


%% ===========================================================================
\section{Linear equations: the test and its completeness}
\label{sec:lin}
%% ===========================================================================

Let $\varphi=C_1\wedge\dots\wedge C_m$ be a $3$-CNF with each variable in exactly five clauses.  We write $\{\pm1\}$ multiplicatively, so that a linear equation is a product: $x_ax_bx_c=b_0$ with $b_0\in\{\pm1\}$.

\begin{definition}[Test $\mathrm{L}_\eps(u)$; \citet{Hastad01}, Section~5]\label{def:testL}
The written proof is a $\mathrm{SWP}(u)$.  The verifier
\begin{enumerate}[label=\arabic*.,leftmargin=*]
  \item chooses $u$ clauses $(C_{j_i})_{i\le u}$ uniformly, and for each a uniformly random variable $x_{k_i}$ occurring in it; sets $U=\{x_{k_1},\dots,x_{k_u}\}$, $W=$ the set of all variables of the chosen clauses, and $h=\bigwedge_iC_{j_i}$;
  \item chooses $f\in F_U$ uniformly;
  \item chooses $g_1\in F_W$ uniformly;
  \item chooses $\mu\in F_W$ by $\mu(y)=-1$ with probability $\eps$ and $\mu(y)=+1$ with probability $1-\eps$, independently over $y$;
  \item sets $g_2=fg_1\mu$, \ie $g_2(y)=f(y|_U)\,g_1(y)\,\mu(y)$;
  \item accepts iff $A_{U;{\rm true}}(f)\cdot A_{W,h;{\rm true}}(g_1)\cdot A_{W,h;{\rm true}}(g_2)=1$.
\end{enumerate}
\end{definition}

\noindent
Here $A_{U;{\rm true}}$ is $A_U$ conditioned upon nothing (there is no $h$) and folded; $A_{W,h;{\rm true}}$ is conditioned upon $h$ and folded.  The three queried bits are the signs.  The point of the design is the last line: \emph{the acceptance criterion is the exclusive-or of three proof bits}, which is the predicate of \sat{Max-E3-Lin}$(\F_2)$.
So the CSP that the test defines is an instance of the target problem, with no reduction in between (Exercise~\ref{ex:testL-is-lin}).

The two queries $g_1,g_2$ to the same table are related: they differ only where $\mu=-1$, an $\eps$-fraction of the points.  This is what makes the test a test of consistency and the predicate a three-bit predicate at the same time.

\begin{lemma}[Completeness; \citet{Hastad01}, Lemma~5.1]\label{lem:L-complete}
For every $u$, $\ \mathrm{L}_\eps(u)$ accepts a correct proof with probability at least $1-\eps$.
\end{lemma}
\begin{proof}
Let $x$ satisfy $\varphi$ and let $A_V=A_{x|_V}$ for all $|V|\le3u$.  Folding and conditioning agree with the long code on a correct proof, so
\[
  A_{U;{\rm true}}(f)=f(x|_U),\qquad
  A_{W,h;{\rm true}}(g_1)=g_1(x|_W),
\]
\[
  A_{W,h;{\rm true}}(g_2)=g_2(x|_W)=f(x|_U)\,g_1(x|_W)\,\mu(x|_W).
\]
The product of the three is $\mu(x|_W)$, which is $-1$ with probability exactly $\eps$.  So the test fails only when $\mu(x|_W)=-1$.
\end{proof}

\noindent
Thus the verifier has completeness $1-\eps$ for every $\eps>0$, and the predicate is exactly the target predicate.  The remaining half is soundness, which is where the Fourier analysis of Section~\ref{sec:soundness} happens.


%% ===========================================================================
\section{Soundness: from acceptance to a strategy}
\label{sec:soundness}
%% ===========================================================================

The statement to prove is this.

\begin{lemma}[Soundness of $\mathrm{L}_\eps(u)$; \citet{Hastad01}, Lemma~5.2]\label{lem:L-sound}
Suppose $\mathrm{L}_\eps(u)$ accepts a proof with probability $\tfrac{1+\delta}{2}$.  Then the $u$-parallel two-prover protocol of Lemma~\ref{lem:parallel-protocol} has a strategy with value at least $4\eps\delta^2$.
\end{lemma}

\noindent
Before reading the proof, notice what it is not.  It does not say that the table is close to a long code.  It says that a \emph{large Fourier coefficient} of the table can be converted into a strategy for the game.  This is the difference from Lectures~\ref{lec:as}--\ref{lec:alm}, where soundness was proved by decoding, and it is what makes the analysis exact.

Fix $U,W,h$ for the moment and abbreviate $A=A_{U;{\rm true}}$ and $B=A_{W,h;{\rm true}}$.  By the assumption, the expectation of the accepted product is $\delta$:
\begin{equation}\label{eq:accept}
  \E_{U,W,h}\ \E_{f,g_1,\mu}\Big[A(f)\,B(g_1)\,B(g_2)\Big]\;=\;\delta,
  \qquad g_2=fg_1\mu.
\end{equation}

\begin{proof}[Proof of Lemma~\ref{lem:L-sound}]
\emph{Step 1: expand in Fourier series.}  Substitute the expansions of $A$ and of the two copies of $B$, and use $g_2=fg_1\mu$ together with part~(a) of Lemma~\ref{lem:prod-f}:
\[
  B(g_2)=\sum_{\beta_2}\widehat B_{\beta_2}\,\chi_{\beta_2}(f)\chi_{\beta_2}(g_1)\chi_{\beta_2}(\mu).
\]
Using part~(b) to combine the $g_1$-characters and part~(c) to push the $f$-characters down to $U$, and then the independence of $f,g_1,\mu$ (Lemma~\ref{lem:indep-f}),
\[
  \E_{f,g_1,\mu}\big[\chi_{\alpha}(f)\chi_{\beta_1}(g_1)\chi_{\beta_2}(f)\chi_{\beta_2}(g_1)\chi_{\beta_2}(\mu)\big]
  =\mathbf 1[\beta_1=\beta_2]\cdot\mathbf 1[\alpha=\pi_U(\beta_2)]\cdot\E_\mu[\chi_{\beta_2}(\mu)] .
\]
By Lemma~\ref{lem:noise} the last factor is $(1-2\eps)^{|\beta_2|}$.  So~\eqref{eq:accept} becomes
\begin{equation}\label{eq:dsquared}
  \E_{U,W,h}\left[\sum_{\beta}\widehat A_{\pi_U(\beta)}\,\widehat B_\beta^{\,2}\,(1-2\eps)^{|\beta|}\right]\;=\;\delta .
\end{equation}
(This is equation~(12)--(13) of the paper.  Note that the inner sum is over $\beta$, and only its projection to $U$ matters for $A$: the test ``sees'' the small table only through $\pi_U$.)

\emph{Step 2: strategies.}  Define two randomized strategies for the parallel protocol, in which $P_2$ receives $U$ and $P_1$ receives $(W,h)$.
\begin{itemize}
  \item $P_2$ on $U$: sample $\alpha$ with probability $\widehat A_\alpha^{\,2}$, then output a uniformly random $x\in\alpha$.
  \item $P_1$ on $(W,h)$: sample $\beta$ with probability $\widehat B_\beta^{\,2}$, then output a uniformly random $y\in\beta$.
\end{itemize}
Both sets are nonempty by the foldedness (Definition~\ref{def:folding}: only odd $|\alpha|$ survive, so $\varnothing$ never has a nonzero coefficient), and every $y$ that $P_1$ outputs satisfies the chosen clauses, by the conditioning (only $\beta$ whose points satisfy $h$ survive).  So the referee accepts whenever the two answers are consistent, $y|_U=x$.

For fixed $U,W,h$ and fixed sampled $\alpha,\beta$ with $\alpha=\pi_U(\beta)$, the answers are consistent with probability at least $|\beta|^{-1}$: whatever $x\in\alpha$ the prover $P_2$ outputs, some $y\in\beta$ restricts to it (an element of $\pi_U(\beta)$ is the restriction of at least one $y\in\beta$), and $P_1$ outputs that $y$ with probability at least $1/|\beta|$.
Each pair $(\alpha,\beta)$ is sampled with probability $\widehat A_\alpha^{\,2}\widehat B_\beta^{\,2}$, so the acceptance probability for fixed $U,W,h$ is at least
\begin{equation}\label{eq:strategy}
  \sum_{\beta}\widehat A_{\pi_U(\beta)}^{\,2}\,\widehat B_\beta^{\,2}\,\frac{1}{|\beta|}\;\ge\;\sum_\beta\widehat A_{\pi_U(\beta)}^{\,2}\,\widehat B_\beta^{\,2}\,4\eps\,(1-2\eps)^{|\beta|},
\end{equation}
where the inequality is the key estimate
\begin{equation}\label{eq:key-ineq}
  \frac{1}{k}\;\ge\;4\eps\,(1-2\eps)^{k}\qquad\text{for every integer }k\ge1,
\end{equation}
proved below.

\emph{Step 3: put the two together.}  Take expectations over $U,W,h$ in~\eqref{eq:strategy}, write $\mu_\beta=\widehat B_\beta^{\,2}(1-2\eps)^{|\beta|}\ge0$, and apply Cauchy--Schwarz to the inner sum:
\[
  \sum_\beta \widehat A_{\pi_U(\beta)}^{\,2}\mu_\beta\;\ge\;\Big(\sum_\beta\big|\widehat A_{\pi_U(\beta)}\big|\,\mu_\beta\Big)^2 \;\ge\;\Big(\sum_\beta\widehat A_{\pi_U(\beta)}\,\mu_\beta\Big)^2,
\]
using $\sum_\beta\mu_\beta\le\sum_\beta\widehat B_\beta^{\,2}=1$ in the first step and $\widehat A\le|\widehat A|$ in the second.
Taking expectations and applying $\E[X^2]\ge(\E X)^2$ to $X=\sum_\beta\widehat A_{\pi_U(\beta)}\mu_\beta$ gives
\[
  \E_{U,W,h}\Big[\sum_\beta\widehat A_{\pi_U(\beta)}^{\,2}\mu_\beta\Big]\;\ge\;\Big(\E_{U,W,h}\Big[\sum_\beta\widehat A_{\pi_U(\beta)}\mu_\beta\Big]\Big)^2\;=\;\delta^2
\]
by~\eqref{eq:dsquared}.  Substituting into~\eqref{eq:strategy} gives a strategy for the parallel protocol with value at least $4\eps\delta^2$.
\end{proof}

\begin{lemma}[The key inequality, \eqref{eq:key-ineq}]\label{lem:key-ineq}
For every $\eps\in(0,\tfrac12)$ and every integer $k\ge1$: $\ \dfrac{1}{k}\ \ge\ 4\eps\,(1-2\eps)^{k}$.
\end{lemma}
\begin{proof}
We must show $k\,(1-2\eps)^k\le\frac1{4\eps}$.  Write $w=1-2\eps\in(0,1)$.  The function $k\mapsto kw^k$ is maximal at $k^\ast=-\frac1{\ln w}$, where it equals
\[
  k^\ast w^{k^\ast}=\frac{1}{e\,(-\ln w)} .
\]
Now $-\ln w=-\ln(1-2\eps)\ge 2\eps$ (since $\ln(1-t)\le-t$), so the maximum is at most $\frac{1}{2e\eps}=\frac{0.1839\ldots}{\eps}<\frac{1}{4\eps}$, since $\frac1{2e}<\frac14$.
(The value of $k^\ast$ need not be an integer; the maximum over integers is no larger.)
\end{proof}

\noindent
It is worth separating the two losses in the proof.  The first is the factor $|\beta|^{-1}$: the extracted strategy outputs a single random point of $\beta$ rather than the whole set.  The noise is what pays for it -- the test only sees $\beta$ through the weight $(1-2\eps)^{|\beta|}$, so large sets contribute little to the acceptance probability, and Lemma~\ref{lem:key-ineq} converts this into the factor $4\eps$.  This is why the noise $\mu$ is in the test at all, and why the soundness bound degrades as $\eps\to0$.  The second loss is the square: the test measures a quantity \emph{linear} in the coefficients of $A$, the strategy succeeds with a quantity \emph{quadratic} in them, and Cauchy--Schwarz and Jensen in Step~3 turn $\delta$ into $\delta^2$.

\begin{pitfall}{what would fail without conditioning and folding}
Without folding, the empty set $\varnothing$ could carry a large coefficient $\widehat A_\varnothing$, and then the strategy of Step~2 would have to output a random point of the empty set.
Without conditioning, a set $\beta$ could have all its points violating the chosen clauses, and $P_1$'s answer would not satisfy them.
Both operations are therefore not cosmetic: they are what makes the extracted strategies well defined, and they are performed by the verifier on the table it is handed, at no cost to completeness.
\end{pitfall}

\begin{theorem}[Parallel repetition is enough to make $\delta$ arbitrary]\label{thm:L-sound-consequence}
Let $\eps>0$ be a constant.  If $\varphi$ is not satisfiable then $\mathrm{L}_\eps(u)$ accepts every proof with probability at most $\tfrac12+\delta(u)$, where $\delta(u)\to0$ as $u\to\infty$.
\end{theorem}
\begin{proof}
Suppose a proof is accepted with probability $\tfrac{1+\delta}{2}$.  By Lemma~\ref{lem:L-sound} the parallel protocol has value at least $4\eps\delta^2$, and by Lemma~\ref{lem:parallel-protocol} its value is at most $\bar c^{\,u}$.  Hence $\delta\le\bar c^{\,u/2}/(2\sqrt\eps)$.
\end{proof}

\noindent
Since $u$ is our choice, we may make $u$ large enough that $\bar c^{\,u/2}/(2\sqrt\eps)$ is below any prescribed $\delta_0$; $u=O(\log\frac1{\eps\delta_0})$ suffices, and then the proof has polynomial length.  The constants are now explicit in terms of $\bar c$ and $\eps$.


%% ===========================================================================
\section{Max-E3-Lin and Max-E3-SAT}
\label{sec:e3sat}
%% ===========================================================================

\begin{theorem}[\sat{Max-E3-Lin}$(\F_2)$; \citet{Hastad01}, Theorem~5.4]\label{thm:e3lin}
For every $\eps>0$ it is $\NP$-hard to approximate \sat{Max-E3-Lin}$(\F_2)$ within a factor $2-\eps$.  More precisely, for every $\eta>0$ it is $\NP$-hard to distinguish systems of three-variable equations over $\F_2$ in which a $(1-\eta)$-fraction of the equations can be satisfied from those in which at most a $(\tfrac12+\eta)$-fraction can.  (Completeness $1$ is impossible: a satisfiable linear system is solved by Gaussian elimination.)
\end{theorem}
\begin{proof}
Fix $\eps>0$ and choose $\delta=\delta(\eps)>0$ with
\begin{equation}\label{eq:choose-delta}
  \frac{1-\delta}{(1+\delta)/2}\ \ge\ 2-\eps ;
\end{equation}
$\delta\le\eps/4$ suffices, since $\frac{2(1-\delta)}{1+\delta}\ge2(1-\delta)^2\ge2-4\delta$.  We run the test with noise parameter $\delta$, \ie $\mathrm{L}_\delta(u)$, as \cite{Hastad01} does.
Let $L\in\NP$ and let $x$ be an input.  By Theorem~\ref{thm:pcp2} there is a $3$-CNF $\varphi$ with each variable in exactly five clauses, of size $\mathrm{poly}(|x|)$, with $\varphi$ satisfiable if $x\in L$ and at most $c$-satisfiable otherwise, $c<1$.
Choose $u$ so large that $\bar c^{\,u}<4\delta\cdot\delta^2$; this is a constant depending only on $\eps$ and $c$.
Consider the CSP given by $\mathrm{L}_\delta(u)$ on $\varphi$: its variables are the bits $b$ of the written proof, one for each query, and each accept condition is the equation $b_{U;f}\,b_{W,h;g_1}\,b_{W,h;g_2}=b_0$ with $b_0$ a fixed sign.
This is an instance of \sat{Max-E3-Lin}$(\F_2)$ of size $n^{O(u)}2^{2^{3u}}=\poly(|x|)$, and by Lemma~\ref{lem:L-complete} and Theorem~\ref{thm:L-sound-consequence}:
\[
  x\in L\ \Rightarrow\ \ge(1-\delta)\text{ of the equations satisfiable},\qquad
  x\notin L\ \Rightarrow\ <\tfrac{1+\delta}{2}\text{ satisfiable,}
\]
the second because acceptance probability $\tfrac{1+\delta}{2}$ would give a strategy of value $4\delta\cdot\delta^2>\bar c^{\,u}$ (Lemma~\ref{lem:L-sound}).
An algorithm with ratio better than $\frac{1-\delta}{(1+\delta)/2}\ge 2-\eps$ would therefore decide $L$; so none exists unless $\Ptime=\NP$.
\end{proof}

\noindent
Since $\eps$ is arbitrary, the ratio is arbitrarily close to $2$.  The threshold $\tfrac12$ is exactly the value of a random assignment, so the theorem says that \emph{no polynomial-time algorithm beats guessing}, even when almost all equations can be satisfied simultaneously.

\begin{theorem}[\sat{Max-E3-SAT}; \citet{Hastad01}, Theorem~6.1]\label{thm:e3sat}
For every $\eps>0$ it is $\NP$-hard to approximate \sat{Max-E3-SAT} within a factor $\tfrac87-\eps$.
\end{theorem}
\begin{proof}
The four-clause gadget of Lemma~\ref{lem:gadget}: replace each equation $x_ax_bx_c=b_0$ (an XOR of three literals) by the four clauses
\[
  (x_a\vee x_b\vee x_c),\quad (x_a\vee\bar x_b\vee\bar x_c),\quad (\bar x_a\vee x_b\vee\bar x_c),\quad (\bar x_a\vee\bar x_b\vee x_c),
\]
with the signs read according to $b_0$.
By Lemma~\ref{lem:gadget}, an assignment satisfying the equation satisfies all four clauses, and one violating it satisfies exactly three.  So if an assignment satisfies a fraction $s$ of the equations it satisfies a fraction $\tfrac{4s+3(1-s)}{4}=\tfrac34+\tfrac s4$ of the clauses, and conversely the best assignment for the clauses is the best for the equations.
Starting from the instances of Theorem~\ref{thm:e3lin} with completeness $1-\delta$ and soundness $\tfrac{1+\delta}{2}$, the clause instances have completeness $\ge1-\tfrac\delta4$ and soundness $\le\tfrac78+\tfrac\delta8$; the ratio $\frac{1-\delta/4}{7/8+\delta/8}$ tends to $\tfrac87$ as $\delta\to0$.
\end{proof}

\begin{lemma}[The four-clause gadget]\label{lem:gadget}
For the four clauses displayed in the proof of Theorem~\ref{thm:e3sat}: an assignment satisfies all four if and only if $x_ax_bx_c=b_0$, and exactly three of the four otherwise.
\end{lemma}
\begin{proof}
Direct enumeration of the eight assignments, checked by hand in Table~\ref{tab:gadget}.  It is also the \emph{only} set of four clauses with this property (Exercise~\ref{ex:gadget-unique}).
\end{proof}

\begin{table}[t]
\centering\small
\renewcommand{\arraystretch}{1.15}
\begin{tabular}{@{}ccc|cc@{}}
\toprule
$x_a$ & $x_b$ & $x_c$ & $x_ax_bx_c$ & clauses satisfied (out of $4$)\\
\midrule
$+$ & $+$ & $+$ & $+$ & $4$\\
$+$ & $+$ & $-$ & $-$ & $3$\\
$+$ & $-$ & $+$ & $-$ & $3$\\
$+$ & $-$ & $-$ & $+$ & $4$\\
$-$ & $+$ & $+$ & $-$ & $3$\\
$-$ & $+$ & $-$ & $+$ & $4$\\
$-$ & $-$ & $+$ & $+$ & $4$\\
$-$ & $-$ & $-$ & $-$ & $3$\\
\bottomrule
\end{tabular}
\caption{The four-clause gadget of Lemma~\ref{lem:gadget} (for $b_0=+1$; a literal is true when its value is $+$).  Satisfying the equation satisfies every clause; violating it costs exactly one clause.  The set is unique with this property.}
\label{tab:gadget}
\end{table}

\noindent
A variant of the same test \cite[Theorem~5.15]{Hastad01} shows that \emph{every} predicate on three bits that is implied by a linear equation is non-approximable beyond its own random threshold: if $P$ accepts whenever $x_ax_bx_c=1$, we may run the test of Definition~\ref{def:testL} with acceptance criterion $P$ instead of the XOR, which only makes the test more generous, and the completeness stays at $1-\eps$ while the soundness analysis goes through with the Fourier coefficients of $P$ in place of those of the XOR (Exercise~\ref{ex:implied-predicates}).

What the argument does \emph{not} give is the corresponding statement on \emph{satisfiable} instances: Theorem~\ref{thm:e3lin} is about $(\tfrac12+\eps)$ versus $1-\eps$, not about $(\tfrac78+\eps)$ versus $1$.  That sharper statement is true -- it is \cite[Theorem~6.5]{Hastad01} -- but its proof needs a substantially more delicate test, and we do not reproduce it.


%% ===========================================================================
\section{What ``optimal beyond the random threshold'' means}
\label{sec:optimal}
%% ===========================================================================

\begin{bigpicture}{Why these results are called optimal}
For \sat{Max-E3-Lin}$(\F_2)$ a uniformly random assignment satisfies exactly half of the equations, whatever the system is; there is no instance on which guessing does worse.  Theorem~\ref{thm:e3lin} says that for every $\eps>0$ no polynomial-time algorithm can guarantee more than a $(\tfrac12+\eps)$-fraction, even on instances where a $(1-\eps)$-fraction is achievable.
So the ratio $2$ is a bound on \emph{every} algorithm and the algorithm is the trivial one: the threshold of random assignment is exactly the threshold of approximability.
The same for \sat{Max-E3-SAT} with $\tfrac78$ in the place of $\tfrac12$.

The phrase ``non-approximable beyond the random assignment threshold'' is the name of this phenomenon.  Let $\mathrm{wt}(P)$ be the fraction of inputs accepted by the predicate $P$, \ie the expected fraction of constraints a random assignment satisfies.  The CSP of $P$ is \emph{non-approximable beyond the random assignment threshold} if, unless $\Ptime=\NP$, for every $\eps>0$ no polynomial-time algorithm has approximation ratio $\mathrm{wt}(P)+\eps$ \cite[Definition~2.16]{Hastad01}; it is so \emph{on satisfiable instances} if it is moreover $\NP$-hard to distinguish satisfiable instances from those in which at most a $(\mathrm{wt}(P)+\eps)$-fraction of the constraints can be satisfied \cite[Definition~2.17]{Hastad01}.
\sat{Max-E3-Lin}$(\F_2)$ has the first property but not the second (Gaussian elimination); \sat{Max-E3-SAT} has both, the second being his Theorem~6.5.
\end{bigpicture}

\noindent
For predicates on three bits, \cite{Hastad01} notes that the question is completely resolved by \citet{Zwick98}: the predicates implied by a linear equation are non-approximable beyond their threshold (as his Theorem~5.15 also shows), and the others are not.  For four bits \cite{Hastad01} leaves it open, saying so in his concluding remarks.  This is the frontier: the method is ``very strong'' but ``does not seem universal''.

The paper's table collects the constants that follow.  Table~\ref{tab:scoreboard} of Lecture~\ref{lec:why} quoted some of them; Table~\ref{tab:hastad-ratios} lists the main ones, all in the form ``$\NP$-hard to approximate within''.

\begin{table}[t]
\centering\small
\renewcommand{\arraystretch}{1.25}
\setlength{\tabcolsep}{4.5pt}
\begin{tabular}{@{}lccc>{\raggedright\arraybackslash}p{3.9cm}@{}}
\toprule
problem & ratio & previous & random & how it is obtained\\
\midrule
\sat{Max-E3-Lin}$(\F_2)$ & $2-\eps$ & $\tfrac87-\eps$ \citep{BGS98} & $\tfrac12$ & Test $\mathrm{L}_\eps(u)$, Theorem~\ref{thm:e3lin}\\
\sat{Max-E3-Lin}$(\F_p)$ & $p-\eps$ & $p^{c}$, some $c>0$ & $\tfrac1p$ & same test over $\F_p$\\
\sat{Max-E3-SAT} & $\tfrac87-\eps$ & $\tfrac{80}{77}-\eps$ \citep{BGS98} & $\tfrac78$ & four-clause gadget, Theorem~\ref{thm:e3sat}\\
\sat{Max-E2-SAT} & $\tfrac{22}{21}-\eps$ & $\tfrac{220}{217}-\eps$ & --- & gadget of \citet{TSSW00}\\
\sat{Max-Cut} & $\tfrac{17}{16}-\eps$ & $\tfrac{72}{71}-\eps$ & $\tfrac12$ & gadget of \citet{TSSW00}\\
\sat{Max-di-Cut} & $\tfrac{12}{11}-\eps$ & $\tfrac{72}{71}-\eps$ & --- & gadget of \citet{TSSW00}\\
\sat{Vertex Cover} & $\tfrac76-\eps$ & $\tfrac{233}{218}-\eps$ & --- & free bits, \eqref{eq:freebits}\\
\bottomrule
\end{tabular}
\caption{The ratios of \cite{Hastad01}.  ``Previous best'' is the earlier constant, in most cases from \citet{BGS98}; the \sat{Max-E2-SAT}, \sat{Max-Cut} and \sat{Max-di-Cut} rows come from the linear-equations test via gadgets of \citet{TSSW00}, the last row from a free-bit count; none of these four is proved here.}
\label{tab:hastad-ratios}
\end{table}

\noindent
Three points about the table.  First, the two rows we proved are the two where the \emph{predicate of the test} is the predicate of the problem; the others lose an unspecified but explicit constant in a gadget reduction.  Second, the improvement over the previous best is large in the first rows ($\tfrac87\to2$ for linear equations, $\tfrac{80}{77}\to\tfrac87$ for satisfiability) and smaller in the gadget rows, where part of the improvement is due to better gadgets.  Third, the last row uses a different mechanism, which is worth stating because it is the one place where the \emph{number of free bits} of the verifier -- not the predicate -- produces the constant.

\begin{pitfall}{free bits, and why they are not the same as queries}
For a verifier with perfect completeness, a bit it reads is \defn{non-free} if, at the moment of reading, the verifier will reject unless the bit has one prescribed value; otherwise it is \defn{free}.  Informally, the free bits count the accepting patterns: a verifier with $f$ free bits has at most $2^f$ accepting answer patterns per random string.  For a verifier with $f$ free bits, completeness $c$ and soundness $s$, the induced instance of \sat{Vertex Cover} has a gap of
\begin{equation}\label{eq:freebits}
  \frac{2^{f}-s}{2^{f}-c},
\end{equation}
and the test $\mathrm{L}_\eps(u)$ has $f=2$, $c=1-\eps$, $s=\tfrac12+\eps$, giving $\frac{4-\frac12-\eps}{4-1+\eps}=\frac{3.5-\eps}{3+\eps}\to\frac76$.
For \sat{Clique} the corresponding statement is the $N^{1-\eps}$ of \citet{Hastad99}, which is \emph{not} proved in this paper -- \cite{Hastad01} cites it -- and which likewise comes from counting the free bits of a PCP rather than from a sharp predicate (Exercise~\ref{ex:free-bits}).
\end{pitfall}


%% ===========================================================================
\section{Why this is where the course turns}
\label{sec:towards-ugc}
%% ===========================================================================

We now have the two sharpest theorems of the classical theory, and they are sharp in a strong sense: the algorithm is the trivial one and no other algorithm does better unless $\Ptime=\NP$.

Widening the picture, however, the results of this lecture have a shape that stops the method.  Each of them is proved by designing, for one predicate, a test whose acceptance predicate is that predicate, and then analysing that test.  The analysis relies on the Fourier spectrum of the predicate and on the long code, and it works for three-bit predicates and (with more difficulty) for the satisfiable case of \sat{Max-E3-SAT}.  What it does not give is a \emph{general} statement covering every CSP at once, and H\aa stad's own concluding remarks single out the four-bit case as open and say the method ``does not seem universal even limited to CSPs''.

Worse, for the two problems where the algorithmic side is genuinely clever, the results of this lecture do not close the gap.

\begin{bigpicture}{The two gaps that remain, and what they have in common}
\sat{Max-Cut} is approximable within $\alpha_{\mathrm{GW}}\approx0.8786$ by the semidefinite program of \citet{GW95}, and the best unconditional hardness is $\tfrac{16}{17}\approx0.941$ (\citealp{Hastad01,TSSW00}).  The gap is $0.8786$ against $0.941$.
\sat{Vertex Cover} is approximable within $2$ by a maximal matching, and the best unconditional hardness is $\sqrt2\approx1.414$ (\citealp{DKKMS18,KMS18}).  The gap is $1.414$ against $2$.
H\aa stad's own bounds in both rows ($\tfrac{17}{16}$ and $\tfrac76$) come from the linear-equations test through a gadget or a free-bit count, and each loses a constant that the algorithm does not; the later $\sqrt2$ for \sat{Vertex Cover} needed the $2$-to-$2$ theorem of the last part of these notes.
In both cases, assuming the Unique Games Conjecture, the gap closes exactly and the optimal ratio is the one the algorithm achieves (\citealp{KKMO07,MOO10} for \sat{Max-Cut}, \citealp{KhotRegev08} for \sat{Vertex Cover}).
\end{bigpicture}

\noindent
The pattern is this.  The results of this lecture are proved by a test whose predicate \emph{is} the target predicate, and they are optimal for the problems whose predicate is a linear equation or a clause.  For problems whose predicate is not of this kind, the reduction passes through a gadget or a cruder parameter, and the loss there is what the UGC, in these two examples, recovers.

There is also a more intrinsic reason for the turn.  Consider what the proof of Lemma~\ref{lem:L-sound} did: it extracted from a passing table a \emph{small set} of candidate values for a small set of variables, and used it as a strategy.  The natural generalisation is a CSP whose constraints are \emph{bijections} between label sets, so that a label of one side determines a label of the other: then a passing test yields a labeling, not just a set of candidate answers, and the correspondence between the test and the game is exact rather than approximate.  That object is the \emph{Unique Game}, and the conjecture that its gap problem is hard is the subject of the next part of these notes.

Lecture~8 introduces the conjecture; we shall see in Lecture~11 that it is strong enough to determine the optimal ratio for \emph{every} CSP, as the integrality gap of one semidefinite program \citep{Raghavendra08}; and we shall see in the last part of these notes that the $2$-to-$2$ games give it a relative that is now a theorem.


%% ===========================================================================
\section*{Exercises}
\addcontentsline{toc}{section}{Exercises}
%% ===========================================================================

\begin{exercise}[\stars{$\star$}]\label{ex:longcode-distance}
(a)~Show that any two distinct long codes differ on exactly half of $F_U$, so the relative distance is $\tfrac12$.
(b)~How many bits does a long code of $u$ bits have?  For which $u$ is the table shorter than $2^{1000}$?
(c)~Show that the long code of $x$ is the unique $A$ with $\widehat A_{\{x\}}=1$, and that if $\widehat A_{\{x\}}\ge1-\eta$ then $A$ agrees with $A_x$ on at least a $(\tfrac12+\tfrac{1-\eta}{2})$-fraction of $F_U$.
\end{exercise}

\begin{exercise}[\stars{$\star\star$}]\label{ex:folding}
Prove the two Fourier claims of Definition~\ref{def:folding}.  \emph{Hint for (a): for even $|\alpha|$, the substitution $f\mapsto\bar f$ pairs the terms of $\widehat{\tilde A}_\alpha$, and the two contributions cancel because $\chi_\alpha(\bar f)=(-1)^{|\alpha|}\chi_\alpha(f)$.}
For (b), pair $f$ with the function $f'$ that differs from it only at a point $x_0$ where $h$ is false.
\end{exercise}

\begin{exercise}[\stars{$\star$}]\label{ex:testL-is-lin}
Show that the CSP defined by the test $\mathrm{L}_\eps(u)$ is an instance of \sat{Max-E3-Lin}$(\F_2)$: write down its variables, its equations, and its weight distribution.  How many variables does it have (as a function of $u$ and $|\varphi|$)?  Why is this polynomial?
\end{exercise}

\begin{exercise}[\stars{$\star\star$}]\label{ex:key-ineq}
(a)~Prove $k\,(1-2\eps)^k\le\frac1{4\eps}$ for all integers $k\ge1$ and all $\eps\in(0,\tfrac12)$ by maximising over real $k$.  \emph{Hint: $\ln(1-t)\le-t$.}
(b)~Use Lemma~\ref{lem:exp} ($x^s\le e^{s(x-1)}$, with $x=4\eps|\beta|$ and $s=\frac12$) and $e^{-2\eps}\ge1-2\eps$ to prove the weaker bound $|\beta|^{-1/2}\ge(4\eps)^{1/2}(1-2\eps)^{|\beta|}$ that \cite{Hastad01} uses, and show that it also suffices for Lemma~\ref{lem:L-sound}.
(c)~What goes wrong in the proof of Lemma~\ref{lem:L-sound} if $\eps=0$, \ie if the test has no noise?  Exhibit a table that is not close to any long code but passes $\mathrm{L}_0(u)$ with probability $1$.  \emph{Hint: a character $\chi_\beta$ with $|\beta|$ large and $\pi_U(\beta)$ suitable.}
\end{exercise}

\begin{exercise}[\stars{$\star\star$}]\label{ex:gadget-unique}
Prove Lemma~\ref{lem:gadget} by enumerating the eight assignments, and show that the four clauses displayed are the \emph{only} set of four $3$-clauses (one variable each, all signs allowed) with the property ``satisfied equation $\Rightarrow4/4$, violated equation $\Rightarrow3/4$''.
\emph{Hint: check all $\binom{8}{4}=70$ possibilities, or argue that each of the four violated parity classes must be excluded by a different clause.}
\end{exercise}

\begin{exercise}[\stars{$\star\star$}]\label{ex:implied-predicates}
Let $P$ be a predicate on three bits that is satisfied whenever $x_ax_bx_c=1$ (in the $\{\pm1\}$ notation).  Show that the test of Definition~\ref{def:testL} with acceptance criterion $P$ has completeness $1-\eps$, and that its soundness analysis reduces to the one given here with the Fourier coefficients of $P$ in place of those of the XOR.
Check that $P=$ ``not exactly two of the three bits are true'' is implied by linearity, and compute its acceptance probability on a random input.
\end{exercise}

\begin{exercise}[\stars{$\star\star$}]\label{ex:free-bits}
(a)~Define free bits and verify \eqref{eq:freebits} on a PCP with $f$ free bits, completeness $c$ and soundness $s$: show that the induced \sat{Vertex Cover} instance has a gap of $(2^f-s)/(2^f-c)$.
(b)~Apply it to the test $\mathrm{L}_\eps(u)$ ($f=2$, $c=1-\eps$, $s=\tfrac12+\eps$) and compute the limit as $\eps\to0$.
(c)~Why does the ``$f=2$'' matter, and what would $f=3$ give?
\end{exercise}

\begin{exercise}[\stars{$\star\star\star$}]\label{ex:reading-hastad}
\emph{Reading.}  Read Section~6 of \cite{Hastad01} and describe the extra difficulty in proving the satisfiable case (Theorem~6.5) compared with the linear case: what does the test test, and why does conditioning alone not suffice?
Then read his Section~9, and state precisely which stronger hypothesis is needed to make $\eps$ a function of $n$, and what happens to the size of the PCP under that hypothesis.
\end{exercise}


%% ===========================================================================
\section*{Notes and further reading}
\addcontentsline{toc}{section}{Notes and further reading}
%% ===========================================================================

\paragraph{Where this lecture comes from.}
Everything is from \citet{Hastad01}: the Fourier toolkit on $F_U$ is Section~2.4 (our Lemmas~\ref{lem:orth-f}--\ref{lem:noise}), folding and conditioning are Section~2.5 (our Definition~\ref{def:folding}), the basic protocol and the standard written proof are Section~3 (our Definitions~\ref{def:basic-protocol} and~\ref{def:swp}, and Lemmas~\ref{lem:basic-sound}--\ref{lem:parallel-protocol}), the long-code test is Section~4 (our Theorem~\ref{thm:lctest}), the linear-equations test and its analysis are Section~5 (our Definition~\ref{def:testL}, Lemmas~\ref{lem:L-complete} and~\ref{lem:L-sound}) and the satisfiability results are Section~6 (our Theorems~\ref{thm:e3lin}--\ref{thm:e3sat} and Lemma~\ref{lem:gadget}).  The nonconstant-$\eps$ results are his Section~9 and the open problems are his Section~10.
The long code itself is due to \citet{BGS98}; the use of the full Fourier transform (rather than distance to a codeword) for extracting strategies is the innovation of \citet{Hastad99,Hastad01}, and \citet{Hastad01} says so in his introduction.

\paragraph{Presentational choices of these notes.}
(i)~\emph{Notation.}  Like the paper we work over $\{\pm1\}$ with exclusive-or written as multiplication; we let $+1$ denote true (sign conventions for ``true'' vary, and nothing depends on the choice), and we write $A_{V;h;{\rm true}}$ for the conditioned, folded table.  The reader who prefers $\F_2$ should replace $\{\pm1\}$ by $\{0,1\}$ and the product by the sum.
(ii)~\emph{Order.}  We prove the long-code test of Section~\ref{sec:lctest} before the linear-equations test, since it isolates the Fourier argument; the paper does the same and calls the first test ``pedagogical''.
(iii)~\emph{The key inequality.}  The paper's Lemma~5.3 is the estimate $x^s\le e^{s(x-1)}$; the paper applies it as $|\beta|^{-1/2}\ge(4\eps)^{1/2}(1-2\eps)^{|\beta|}$ inside a Cauchy--Schwarz bound.  We use instead the slightly stronger \eqref{eq:key-ineq}, isolated as Lemma~\ref{lem:key-ineq} and proved directly, and then apply Cauchy--Schwarz and Jensen in Step~3; the conclusion $4\eps\delta^2$ is the same.  The square on $\delta$ comes from the Cauchy--Schwarz--Jensen step, not from \eqref{eq:key-ineq}.
(iv)~\emph{The satisfiable case.}  We prove the gadget reduction of Theorem~\ref{thm:e3sat} but not the sharp satisfiable-case statement (Theorem~6.5 of the paper), and we say so.

\paragraph{Numerical checks made for these notes.}
The Fourier identities of Section~\ref{sec:fourier-f} (orthonormality, Parseval, the uniqueness of the coefficient of a long code, and the restriction formula of Lemma~\ref{lem:prod-f}(c)) were verified by exhaustive computation for $|U|\le2$; the identity of Theorem~\ref{thm:lctest} was verified on random tables; Lemma~\ref{lem:noise} and Lemma~\ref{lem:key-ineq} were verified by computation over a range of $\eps$; the four-clause gadget of Lemma~\ref{lem:gadget} was verified by enumeration of all eight assignments and of all $70$ four-subsets of the eight clauses; and the arithmetic of~\eqref{eq:choose-delta} and of Table~\ref{tab:hastad-ratios} was checked directly.

\paragraph{What is not covered.}
We do not prove Theorem~6.5 (the satisfiable case of \sat{Max-E3-SAT}), Set Splitting (Section~7 of the paper), the equations-over-$\F_p$ and Abelian-group extensions (Section~2.6 and Section~5.1), the gadget reductions to \sat{Max-2-SAT}, \sat{Max-Cut} and \sat{Vertex Cover} (Section~8), or the nonconstant-$\eps$ results (Section~9).  The classification of three-variable predicates combines Theorem~5.15 with \citet{Zwick98}; the four-variable case is open.
The next part of these notes turns to the Unique Games Conjecture, which is what remains when one asks for a sharp theory of \emph{every} CSP rather than one predicate at a time.


% ---------------------------------------------------------------- bibliography
\bibliographystyle{abbrvnat}
\bibliography{references}

\end{document}
